\documentclass[12pt,a4paper]{article}
\usepackage{graphicx}
\usepackage{slashed}
\usepackage{tikz}
\usepackage{colortbl}
\usepackage{calc}



% Macros-v2_181: active 31-paper xref binding for framework v3.028; mathematical definitions unchanged.
% Paper-series v2.381: new I-B torsion/modular results, bibliography and current source roots.
% Paper-series v2.380: new exceptional bibliography and synchronized paper references.
% Paper-series v2.379 current 31-paper macro/aux registry.
% Paper-series v2.378: current 31-paper AUX/bibliography registry updated for Paper I-B v1.2 literature additions.
% Framework v2.371: current-source reference registry refreshed; old theorem labels retained.
% Framework v2.370: current 30-paper source/AUX root rebinding; native mathematical macros unchanged.
% Framework v2.368: synchronized merged 30-root registry; mathematical notation unchanged.
% Framework v2.358: active 30-paper source registry synchronized; existing mathematical macros unchanged.
% Framework v2.357: selective new source-owned certificates; active reference roots updated for five revised papers.
% Framework v2.356: active 30-root registry synchronized to the 30 papers in this release.
% Framework v2.275: exact positive dense-twin word cubature, finite repeatability and eighth-moment classification, nested sampling frames, joint action and POVM-jet exports; preserves v2.274.
% Framework v2.274: filename and root-source synchronization.
% Framework v2.273: restored missing Problem 3 theorems from v2.268; exact exceptional frames, twin geometry, convolution repeatability and projector/occurrence firewalls; merged on v2.272 without overwriting independent branches.
% Framework v2.272: merged v2.270 covariant native-readout mathematics with v2.271 Problem 4 finite-action Hopf hypercharge/holonomy corrections; no universal bulk-gauge uniqueness claim.
% Independent branches: v2.270 readout/guard exports and v2.270 K16 action-consumed charge closure; all prior scoped assumptions retained.
% Framework v2.269: synchronized 30 active AUX roots; mathematics and bibliography preserved from Macros-v2_161.
% Framework v2.266: Paper II-C converse-derived outcome resources; II-E, III-C, IV-A/B, V-A and VI-B downstream mathematical closures; 30-root registry synchronized.
% Framework v2.265: II-C coordinate-stage constitution and exact first-symbol host inheritance; prior mathematical domains retained.
% Framework v2.264: paper II-B reverse/audit and IV-A observable-minimal exchange; synchronized external references to 30 active roots.
% Framework v2.263: Paper I-C adds bridge-controlled root-carry parity and source-relative old-cap cancellation; other root mathematical bodies retained; active 30-root references synchronized.
% Framework v2.262: integrates Paper II-F radius, Casimir and readout results; active 30-root references synchronized.
% Framework v2.261: merged Paper II-B GP2-level/geometric repairs; preserves v2.260 II-D repairs; active references synchronized.
% Framework v2.260: Paper II-D signed RS, finite-symbol typing, and conditional curvature/Jordan repairs; references synchronized.
% Framework v2.259: conditional native SU(2) atlas and explicit later discharge; v2.258 mathematical bodies retained elsewhere.
% Framework v2.258: preserves parallel II-A/dynamics update and adds complete readout/frequency closures.
% Framework v2.257: exact minimal complete position readout and arbitrary-preparation frequency law.
% Framework v2.256: exact native inertia and sharp bounded ordered-source closure; preserves concurrent v2.255 work.
% Framework v2.255: preserves parallel detector integration and adds exact finite readout closures.
% Framework v2.254: Pasterski bibliography and synchronized 30-paper external-reference registry.
% Framework v2.253: typed raw coordinate defect in I-A and synchronized VI-A import; active references synchronized.
% Framework v2.252: synchronized source registry; optional III-G shell-frame calculation.
% Framework v2.251: forward common-action execution, joint source jets and scoped coupled control.
% Framework v2.251: exact full-heat variations, fermion source jets and explicit physical completion frontier.
% Framework v2.251: certified interacting source jets, pole stability and scoped completion proofs.
% Framework v2.251: synchronized active source registry; mathematical notation unchanged.
% Framework v2.251: synchronized active source registry; mathematical notation unchanged.
% Framework v2.251: synchronized active source registry; mathematical notation unchanged.
% Framework v2.251: synchronized source registry; no mathematical notation changes.
% Framework v2.244: synchronized 30-paper AUX registry and rooted phase-source action-order tests.
% Framework v2.242: signed rooted phase lift, complete transverse return and exact family cubic.
% Framework v2.241: proved rooted conormal geometry, conditional cap realization and scoped action tests.
% Framework v2.239: all 30 active source roots synchronized; no mathematical macro changes.
% Framework v2.238: Paper IV A/B/C split and 30-root reference registry.
% Framework v2.237: Paper VII-D technical split and 28-root reference registry.
% Framework v2.236: Paper VI-C technical-volume reference added; all current root bindings advanced.
% Framework v2.235: 26-root AUX and bibliography bindings synchronized; mathematical macros unchanged.
% BEGIN INLINE SOURCE: Paper00-TeXBuildCompatibility-v1_0.tex
% Rendering-only repairs for inherited long appendices and bitmap text fonts.
% Integer counters through 26 retain the original labels; later values use AA,
% AB, ... . This file changes no action coefficient, state or theorem premise.
\RequirePackage{alphalph}
\makeatletter
\let\BHPrimitiveAlph\@Alph
\newalphalph{\BHExtendedAlph}{\BHPrimitiveAlph}{26}
\let\@Alph\BHExtendedAlph
\AtBeginDocument{%
  \@ifpackageloaded{microtype}{\microtypesetup{expansion=false}}{}%
}
\makeatother

% END INLINE SOURCE: Paper00-TeXBuildCompatibility-v1_0.tex

% Framework v2.164: current reference/bibliography profile; mathematical definitions retained.
% Framework v2.162: finite metric, functional Hilbert and scoped Ward integration.
% Framework v2.161: inherits Macros-v2_57.tex exactly except the current external-reference binding.
% Macros v2.57 / framework v2.156: current reference registry only; mathematical definitions unchanged.
% Framework v2.143 / Macros v2.55: current external-reference bindings; mathematical notation unchanged.
% Framework v2.110 / Macros v2.54: current external-reference registry and bibliography; mathematical notation inherited.
% Framework v2.67 / Macros-v2_53: adds Born-action-cell and joint-EH readout symbols; absolute unit realization is metrological; MATCH_{J/scal}^{EH} remains the sole G24 normalization bridge.
% Framework v2.35 / Macros-v2_52: grouped paper numbering, active AUX registry, and Born-v6_20 bibliography binding synchronized; notation unchanged.
% Framework v2.26: active-source AUX and bibliography bindings synchronized; notation unchanged.
% Framework v2.09: reviewed early-reciprocity integration; scoped physical imports.
% Framework v2.08: Born--Euler mathematics and conditional complete-action potential criterion.
% Framework v2.07: active AUX registry includes Paper I-B v1.2; shared bibliography Born-v6_18.
% Macros-v2_46 / framework v2.04: mathematical atlas and downstream realization correction.
% Macros-v2_46 / framework v2.03: scoped structure integration and current dependency bindings.
% Macros-v2_46 / framework v2.01: shared source bindings; prior revision comments are historical.
% v2.42 / framework v1.99: Synchronized RG/AS owner and bibliography bindings.
% v2.41 / framework v1.99: active manuscript AUX bindings and geometric-closure bibliography synchronized.
% Release-v1.80 centralized source synchronization for scalar-circle/comparison-band integration.
% Release v1.79: Hurwitz hierarchy and kinetic-normalized action response integration; physical mass matching remains conditional.
% Release v1.78: approved M8 closure-batch integration; physical and analytic domains retained.
% Current v1.78 registry: all active TeX roots are synchronized below.
% Earlier release comments are historical scope records, not current status.
% Release-v1.77 synchronization: native same-owner Noether closure, generator--flow/history typing, analytic-certification firewalls, scalaron/K16, LF/QCD, quantum-Ward, and RAR boundary-owner transport are registered; shared mathematical notation is unchanged.
% Release-v1.76 synchronization: registry points to the v1.76 active paper roots.
% Release-v1.76 synchronization: registers the local A+Phi+g Noether/constraint closure and dynamics-frontier propagation; shared mathematical notation is unchanged.
% Release-v1.75 synchronization: registry points to the v1.75 active paper roots.
% Release-v1.75 synchronization: integrates the joint-Dirac Einstein bulk real form and the local odd-scalar heat-coefficient action on explicit analytic first-jet/interior domains; v1.74 Yang--Mills closure is retained. Fermionic reality, pure-metric curvature-squared/exact-nonlocal terms, noncanonical boundaries, physical vacuum selection, and actual solution-history tokens remain open.
% Release-v1.74 synchronization: registry points to the v1.74 active paper roots and retains the v1.73 theorem labels.
% Release-v1.74 synchronization: canonical Macros-v2_34.tex; generated local nonlinear Yang--Mills action, constrained connection-field phase space, source-free Gauss momentum map, and Legendre closure are integrated on their proved domains. The remaining non-gauge ACT_cl sectors and actual solution-history selection remain open.
% Release-v1.71 synchronization: closes the former JOINT_x/JOINT_5
% connection-occurrence gap by canonical structure-group extension of the
% nonflat Born connection, faithful associated representation, and the local
% cotangent--Clifford operator constructor.  Actual fluctuation members,
% actions, vacua, solutions, calibrations, and analytic QFT domains remain
% separately typed.  The active-AUX registry is advanced once.
% Release-v1.71 canonical macro and active-AUX registry.
% ======================================================================
% Macros-v2_39.tex -- canonical publication macro file
% Release v1.71 retains the v1.70 null-ray/QCD boundaries and additionally
% removes the duplicate structural connection actuality premise.  It does not
% claim global bundle topology, select a gauge/Higgs member, or promote the
% finite construction to an interacting continuum theory.
% Release v1.68 constructs the finite light-front total-fraction operators and
% parton number/momentum operator measures, proves their simplex and Hausdorff-limit
% theorems, and separates positive cut measures from scheme-renormalized
% light-ray distributions.  It also separates finite real Schur poles, stable
% norm-resolvent/Riesz pole-residue limits, and continued resonance
% pole/residue limits.  The Paper-VII-C scattering package is claim-indexed:
% local QFT and spectral claims import no S-matrix or infrared premise, while
% asymptotic amplitudes, LSZ, S-matrix poles, completeness, and cross sections
% import only their named massive or massless/long-range clauses.
% Release v1.67 factorizes BND48_H into a same-action Hessian/Schur
% descendant and still-open typed boundary-trace, critical-lift, and
% refinement-limit conditions.  It proves the finite-collar Dirichlet-to-
% Neumann law sqrt(y)tanh(L sqrt(y)), its nonuniform critical limit, and the
% finite-return ultra-infrared turnover.  The Paper-I-A Hardy constant mode is
% an exact geometric near-hit, not yet the physical metric boundary mode.
% The current local f(R) Hessian remains massive and therefore still refutes
% a RAR prediction on the minimal action branch.
% Release v1.66 derives the transverse depth-six normalizer bridge from the
% already generated C6 Reynolds image and Paper-I-A binary-dihedral normalizer.
% The bridge is the unique least 2D3-stable scalar-containing realization of
% the faithful quotient 2D3/C6=C2, has trace defect six and hence K=48, but is
% transverse to the nested active carrier.  Exact K=44 and K=52 contraction,
% semigroup, and energy-midpoint identities are registered as structural
% readouts rather than projector events.  The active E48 carrier gap remains
% unchanged.  A same-action boundary-symbol/Calderon intertwiner is isolated
% as the necessary and still unproved owner bridge to the static DtN sort;
% representation equivalence alone neither derives RAR nor creates a K=48
% action junction, pole, occupation, or cosmological transition.
% Release v1.65 derives the native E_48-centered carrier mass gap from the
% GP1--GP3-generated active spectrum, distinguishes its carrier-free interval
% from a particle located at E_48, and withdraws the unsupported promotion to
% a positive continuum Yang--Mills mass gap.  The latter remains a valid
% conditional transfer theorem on an enlarged continuum branch, but it is not
% an active framework prediction or an unfinished native-gap proof.
% Release v1.64 synchronizes Derivation Batches 14--22.  It registers the
% GP1--GP3 heat-profile orbit, same-Hessian electroweak response and pole
% ownership, K16/K48 action-response corrections, perturbation state and
% non-Gaussianity typing, finite-vacuum/gap/QCD limit repairs, factorized
% position-area ownership, the actual static metric response, and the ordered
% Wilson/string-tension limits.  Strict continuum convergence, positive
% physical gap, surviving higher cumulants, and positive tension remain open.
% Release v1.61 closes TIMEWIDTH16 as an independent direct-production datum.
% JUNC_{2|7} continues to own the finite intrinsic microscopic profile/width,
% while the actual same-action background/production history determines its
% physical-time pullback, support, and duration.  The current Paper-I-A--VII-C AUX
% registry is advanced once; shared mathematical notation is unchanged.
% Release v1.60 derives the final-octet primitive-remnant depletion, canonical
% Hopf-current representative, exact Hopf descent to dOmega/(4pi), and
% closed-FLRW a^{-3} neutral-current shape.  It retypes K=48 as the onset
% of the final one-unit geometric transfer and records cDM/RAR as distinct
% physical projections of that common neutral parent while preserving
% occurrence, Hilbert-owner, stress-matching, and static-kernel premises.
% Release v1.59 transports the v1.58 exceptional-shell/K16 closure through
% Papers III-F and VI-B--VII-C.  It replaces Paper VI-B's obsolete mixed-Hessian
% output typing by the third-variation pair vertex and Bogoliubov occupation
% support, derives the post-seed Higgs-pair feedback vertex and its exact
% alpha--delta kernel factors, propagates the unique K16/K48 action-junction
% split, identifies the QCD delta count with the upstream 2O shell atom only
% as a common-scalar numerical compatibility, and narrows THERM16 to dynamical
% approach/beta selection because the finite Gibbs/KMS target is already
% unique on a supplied simple observable algebra once its Hamiltonian and
% inverse temperature are supplied; direct-sum central weights remain typed.
% Release v1.58 registers the exceptional nearest-shell closure, the derived
% transverse K=16 action event, the unique interior action junction, the
% type-correct ordered pair vertex, the thin-layer squeeze law, and the
% radiation-like/matter-like kinematic split.  REG_0, transfer magnitude,
% physical layer width, and thermalization remain explicitly conditional.
% Release v1.57 registers the active-depth-six obstruction, the native
% K=48 carrier desert, the typed 7|9 owner reduction, and the dynamical
% reheating handoff across Papers III-B--VI-B.
% Release v1.56 registers the least-support n_2 inflation dressing of the
% persistent n_7 Laurent response, the analytic global-convexity and positive
% endpoint-measure proofs, the derived inflation core and K=16 coefficient
% freeze, and the no-particle dark-sector split into static RAR plus the
% conditional homogeneous effective-curvature response at reflected K=48.
% There is no n_6 action, scalaron-DM mass, or second Einstein--Hilbert term.
% Release v1.55 registers the protected-octet resolved/residual reflection
% q_sp+Z_W=8, its j<->8-j and K<->64-K scale covariance, and the derived
% rank-one-exit quadratic grade C_1^2/U=C_2 in Papers III-C, III-E, and VIII.
% The reflection is a typed trace/scale theorem, not a time reversal or a
% physical event-occurrence theorem; reheating transfer remains conditional.
% Release v1.54 registers the exact positive-part Laurent indices, the
% active-domain/terminal-endpoint action split, the no-second-EH firewall, and
% the q_DM=6/K_DM=48 scalaron readout typing across Papers III-D/E/F, VI,
% VIII, and IX.
% Release v1.53 registers the common nonnegative Laurent obstacle projection
% for n_1^+ and n_7^+, discharges the three former exit/freeze assumptions,
% and advances the Paper-III-D/E/F, VIII, and IX consumers.  Shared notation
% and bibliography ownership are unchanged.
% Release v1.51 registers the PQTI derivation batch and its updated
% Paper-III-A/C/D/E/F, VIII, and IX consumers.  Shared notation and
% bibliography ownership are unchanged.
% Release v1.50 advances the Paper-III-B AUX target to v1.40 after the
% rigorous audit repair; shared notation and bibliography ownership are unchanged.
% Release v1.49 advances the Paper-II-E/F AUX target to v3.14 after the
% minimal audit repair; shared notation and bibliography ownership are unchanged.
% Release v1.46 synchronizes Paper II-A's minimal frozen constructor-origin
% theorem and its five dependent AUX targets.  Shared notation is unchanged.
% Release v1.45 updates the Paper-II-B AUX target to v2.23 after the
% GP3-rigidity and appendix-placement repair; shared notation is unchanged.
% Release v1.44 advances the centralized Paper-I-A--VII-C AUX registry after the
% Paper-I-A v5.16 protected-completion scope repair.  Every current
% paper imports this one file; no paper-level opt-in or \externaldocument list
% is permitted.  The bibliography remains centralized here through Born-v6_13.
% Shared bibliography, theorem infrastructure, colors, notation alignment,
% and current Paper-I-A--VII-C external-reference registration.
% This file centralizes common imports, bibliography sorting, colors, theorem
% infrastructure, and the notation shared by the complete series.
% v2.9 (release v1.43) made the Paper-I-A--VII-C AUX registry automatic and
% centralized the bibliography and all external-reference ownership.
% v2.8 release-registry refresh (v1.33) registers the rooted generated-parent
% cascade and the POL/REG/free-law/11+4/Born-group closure revisions.  No
% shared mathematical notation or Paper-II-A fundamental principle is changed.
% v2.8 release-registry refresh (v1.32) registers the universal generated-
% continuation revisions in Papers II-A, III-A, III-C, IV, and VIII.  No
% shared mathematical notation or Paper-II-A fundamental principle is changed.
% v2.8 release-registry refresh (v1.30) registers the CRP4/projective-state
% source revisions and their direct registry-only consumers.  No shared
% mathematical notation or Paper-II-A axiom is changed.
% v2.7 release-registry refresh (v1.27) updates only the current external
% Paper-I-A--VII-C basenames after the placement/assumption synchronization; no
% shared mathematical notation or Paper-II-A axiom is changed.
% v2.7 synchronizes the geometric defect-repair revision: Paper-II-A v3.42,
% Paper-III-C v1.26, Paper-III-D v1.27, Paper-III-E v1.32, and Paper-III-F
% v1.34.  It registers the common trace-defect/spectral-filtration theorem,
% its owner-typed response pullback, and the reduced GRED and TERM interfaces.
% No shared mathematical notation is changed.
% v2.6 synchronizes Paper-III-E v1.31's cone-type repair, minimal-response
% classification, and explicit classical action-role boundary, together with
% Paper-III-F v1.32's downstream response-atlas synchronization.  No shared
% mathematical notation is changed.
% v2.5 synchronizes the Paper-III-D v1.26 faithful-trace-duality proof,
% exact ordered-owner joint locus, and typed alternative-presentation
% bookkeeping repair, together with its Paper-III-E and closure-frontier
% references.  No shared mathematical notation is changed.
% v2.4 synchronizes the Paper-III-C v1.25 orbit/process/owner repair, its
% downstream Papers III-B--F, V, VI, VIII, and IX revisions, and Born-v6_13.
% v2.3 synchronizes the Paper-II-A v3.41 witness/support/principal inheritance
% layer, the Paper-III-A v1.24 joint-realizability/readout firewall, all dependent
% host-rate revisions, and the current Paper-I-A--VII-B source registry.
% v2.2 synchronizes the generated active-realization, determinant-depth,
% dimension-weighted hierarchy, and typed cosmological-event interfaces.
% v2.1 synchronizes the exceptional-support interface revisions: Paper I-A
% v5.9, Paper II-A v3.30, Paper II-E/F v3.7, and Paper IV v4.16.
% v2.0 introduces the stage-safe G0 occurrence notation, synchronizes the
% active Paper-I-A--IV source registry (current rigorous Paper-III closure:
% III-A v1.17, III-B v1.22, III-C--E v1.17, III-F v1.14, IV v4.14), and makes
% external-reference self-skip depend on logical paper role rather than a
% versioned job name.
% ======================================================================

% ======================================================================
% macros.tex — Shared macros for the Born Geometry series
% ======================================================================

% ======================================================================
% 1. PACKAGES, BIBLIOGRAPHY, COLORS, AND UNIT DEFINITIONS
% ======================================================================
% Core packages shared across the complete Paper-I-A--VII-C series.  Loading them here makes
% the macro file a genuine series preamble while remaining compatible with existing
% paper preambles that load the same packages before \input.
\usepackage{amsmath}
\usepackage{amssymb}
\usepackage{amsthm}
\usepackage{mathrsfs}
\usepackage{array}
\usepackage{booktabs}
\usepackage{xcolor}
\makeatletter
\@ifpackageloaded{geometry}{}{\usepackage{geometry}\geometry{margin=1in}}
\makeatother
\usepackage{mdframed}

% Shared support packages required by the macros below and by most papers.
\usepackage{ragged2e}
\usepackage{longtable}
\usepackage{tabularx}
\usepackage{enumitem}
\usepackage[most]{tcolorbox}

% Series-wide numeric bibliography convention using the portable BibTeX path.
% All Papers I-A--VII-C use the ordinary \cite interface.
\usepackage[numbers,sort&compress]{natbib}
\providecommand{\printbibliography}{%
  \bibliographystyle{plainnat}%
  \bibliography{Born-v6_59}%
}

% Shared colors.  \providecolor avoids clobbering a paper-specific override.
\providecolor{postulatecolor}{RGB}{230,242,255}
\providecolor{compare}{RGB}{230,245,230}
\providecolor{definitioncolor}{RGB}{230,242,255}

\providecolor{openriskcolor}{RGB}{255,235,235}
\providecolor{bridgecolor}{RGB}{240,240,255}
\providecolor{gold}{RGB}{255,215,0}

% ======================================================================
% 2. CUSTOM ENVIRONMENTS
% ======================================================================
% Common table column types used in all papers.  Guarded so that a paper may
% define its own versions before loading this file.
\makeatletter
\@ifundefined{NC@find@L}{%
  \newcolumntype{L}[1]{>{\RaggedRight\arraybackslash}p{#1}}%
}{}
\@ifundefined{NC@find@C}{%
  \newcolumntype{C}[1]{>{\Centering\arraybackslash}p{#1}}%
}{}
\@ifundefined{NC@find@Y}{%
  \newcolumntype{Y}{>{\RaggedRight\arraybackslash}X}%
}{}
\makeatother

% ======================================================================
% 3. THEOREM STYLES AND ENVIRONMENTS
% ======================================================================
% Custom underlined theorem style (useful for remarks or notes)
\newtheoremstyle{underlined}
  {3pt}   % Space above
  {3pt}   % Space below
  {}      % Body font (empty means normal/roman text, typical for remarks)
  {}      % Indent amount
  {}      % Theorem head font (can be \itshape, \bfseries, etc.)
  {.}     % Punctuation after theorem head
  { }     % Space after theorem head
  {\underline{\thmname{#1}\thmnumber{ #2}}\thmnote{ (#3)}} % Theorem head formatting

% Standard plain theorem style environments.  Each environment is guarded so
% the shared macro file can be loaded after a paper-specific theorem setup.
\theoremstyle{plain}
\makeatletter
\@ifundefined{theorem}{\newtheorem{theorem}{Theorem}[section]}{}
\@ifundefined{lemma}{\newtheorem{lemma}{Lemma}[section]}{}
\@ifundefined{proposition}{\newtheorem{proposition}{Proposition}[section]}{}
\@ifundefined{corollary}{\newtheorem{corollary}{Corollary}[section]}{}
\makeatother

% Definition style environments (typically non-italicized body)
\theoremstyle{definition}
\makeatletter
\@ifundefined{definition}{\newtheorem{definition}{Definition}[section]}{}
\@ifundefined{postulate}{\newtheorem{postulate}{Postulate}[section]}{}
\@ifundefined{closure}{\newtheorem{closure}{Closure}[section]}{}
\makeatother

% Applying the custom underlined style
\theoremstyle{underlined}
\makeatletter
\@ifundefined{remark}{\newtheorem{remark}{Remark}[section]}{}
\@ifundefined{convention}{\newtheorem{convention}{Convention}[section]}{}
\@ifundefined{principle}{\newtheorem{principle}{Principle}[section]}{}
\makeatother

% Framed proof environment.
% Every \begin{proof} ... \end{proof} is automatically framed/shaded.
% This avoids mdframed/trivlist nesting, which can drop the top border after
% theorem-like environments.
\makeatletter
\renewenvironment{proof}[1][\proofname]{%
  \par\addvspace{6pt}%
  \noindent
  \begin{tcolorbox}[%
    enhanced,
    breakable,
    colback=black!4,
    colframe=black!35,
    boxrule=0.35pt,
    arc=1pt,
    left=6pt,
    right=6pt,
    top=4pt,
    bottom=4pt,
    boxsep=0pt,
    before skip=0pt,
    after skip=6pt,
    parbox=false
  ]%
  \pushQED{\qed}%
  \normalfont\itshape #1\@addpunct{.}\enspace
  \normalfont\ignorespaces
}{%
  \unskip\nobreak\hfill\popQED
  \end{tcolorbox}%
  \par\addvspace{2pt}%
}
\makeatother

\mdfdefinestyle{postulatebox}{
  backgroundcolor=yellow!10,
  linecolor=black,
  linewidth=0.4pt,
  skipabove=\baselineskip,
  skipbelow=\baselineskip,
  nobreak=true
}

% ======================================================================
% PROOF-READABILITY HELPERS
% ======================================================================

\newcommand{\ProofPlan}[1]{%
  \par\smallskip
  \noindent\textbf{Proof plan. }#1%
  \par\smallskip
}

\newlist{proofsteps}{enumerate}{2}
\setlist[proofsteps,1]{
  label=\textbf{(\alph*)},
  ref=\alph*,
  leftmargin=2.2em,
  itemsep=0.45\baselineskip,
  topsep=0.35\baselineskip
}
\setlist[proofsteps,2]{
  label=\textbf{(\roman*)},
  ref=\roman*,
  leftmargin=2.0em,
  itemsep=0.25\baselineskip,
  topsep=0.2\baselineskip
}


% ======================================================================
% 4. NUMBER SYSTEMS AND ALPHABETS
% ======================================================================
\providecommand{\RR}{\mathbb{R}} % Real numbers
\providecommand{\PP}{\mathbb{P}} % Projective space / Probability
\providecommand{\CC}{\mathbb{C}} % Complex numbers
\providecommand{\ZZ}{\mathbb{Z}} % Integers
\providecommand{\NN}{\mathbb{N}} % Natural numbers
\providecommand{\HH}{\mathbb{H}} % Quaternions / Hyperbolic space
\providecommand{\LFRayGeom}{\mathsf{LF\!\mbox{-}RAY\!\mbox{-}GEOM}}
\providecommand{\LFRep}{\mathsf{LF\!\mbox{-}REP}}
\providecommand{\LFState}{\mathsf{LF\!\mbox{-}STATE}}
\providecommand{\LFLim}{\mathsf{LF\!\mbox{-}LIM}}
\providecommand{\LFScat}{\mathsf{LF\!\mbox{-}SCAT}}
% Number systems / projective spaces
\providecommand{\CP}{\mathbb{CP}}
\providecommand{\RP}{\mathbb{RP}}

% ======================================================================
% 5. GROUPS AND LIE THEORY
% ======================================================================
\providecommand{\SO}{\ensuremath{\mathrm{SO}}}     % Special Orthogonal group
\providecommand{\SU}{\ensuremath{\mathrm{SU}}}     % Special Unitary group
\providecommand{\Spin}{\ensuremath{\mathrm{Spin}}} % Spin group
\providecommand{\CS}{\ensuremath{\mathrm{CS}}}     % Chern-Simons (capitalized)
\providecommand{\cs}{\ensuremath{\mathrm{cs}}}     % Chern-Simons (lowercase)

% ======================================================================
% 6. COMMON MANIFOLDS AND TOPOLOGY
% ======================================================================
\providecommand{\Sone}{S^1}                 % Standard circle

% ======================================================================
% 7. MATH OPERATORS (Declared via \DeclareMathOperator)
% ======================================================================
\DeclareMathOperator{\End}{End}   % Endomorphisms
\DeclareMathOperator{\Aut}{Aut}   % Automorphisms
\DeclareMathOperator{\spec}{spec} % Lowercase spectrum
\DeclareMathOperator{\tr}{tr}     % Lowercase trace (if you already have \Tr)
\DeclareMathOperator{\supp}{supp} % Support
\DeclareMathOperator{\TrA}{Tr_{A}} % Partial trace over subsystem A
\DeclareMathOperator{\TrB}{Tr_{B}} % Partial trace over subsystem B
% Lie/algebra operators
\DeclareMathOperator{\Ad}{Ad}
\DeclareMathOperator{\Sym}{Sym}
\DeclareMathOperator{\su}{su}
% ======================================================================
% 8. STANDARD OPERATOR-LIKE SYMBOLS (Safe definitions via \providecommand)
% ======================================================================
\providecommand{\vol}{\operatorname{vol}}   % Volume form
\providecommand{\Vol}{\operatorname{Vol}}   % Total volume
\providecommand{\Ric}{\operatorname{Ric}}   % Ricci tensor
\providecommand{\Rm}{\operatorname{Rm}}     % Riemann tensor
\providecommand{\Tr}{\operatorname{Tr}}     % Uppercase trace
\providecommand{\Spec}{\operatorname{Spec}} % Uppercase spectrum
\providecommand{\diag}{\operatorname{diag}} % Diagonal matrix
\providecommand{\rank}{\operatorname{rank}} % Rank
\providecommand{\scal}{\mathcal R} % Scalar curvature
\providecommand{\Cl}{\operatorname{Cl}}     % Clifford algebra
\providecommand{\diam}{\operatorname{diam}} % Diameter

% ======================================================================
% 9. FUNCTIONS, TRIGONOMETRY, AND CALCULUS
% ======================================================================

% Codifferential naming conventions
% (Avoid \ddag: commonly defined by physics.sty and symbol packages)
\providecommand{\codiff}{\mathrm{d}^{\dagger}} % Codifferential operator

% Quantum mechanics bra-ket notation (needed for basis tables)
\providecommand{\ket}[1]{\lvert #1 \rangle}
\providecommand{\bra}[1]{\langle #1 \rvert}
\providecommand{\braket}[2]{\langle #1 \mid #2 \rangle}

% ======================================================================
% 10. PHYSICAL SCALES AND CONSTANTS (Paper 1c / Born Geometry)
% ======================================================================
\providecommand{\Rp}{\ensuremath{R_{\mathrm P}}}     % matched common EH curvature scale (not defined from observed G)
\providecommand{\hBorn}{\ensuremath{h_{\mathrm B}}} % physical representative 2\pi\hbar of one primitive Born scalar action cell
\providecommand{\UEHjoint}{\ensuremath{U_{J}}} % joint-Dirac EH curvature-coordinate scale 8 kappa_J
\providecommand{\Gread}{\ensuremath{G_{\rm scal}}} % scalar-action coordinate c_op^3/(h_B U); physical matching explicit
\providecommand{\GreadJoint}{\ensuremath{G_{J,\mathrm{raw}}}} % raw ordinary-trace coordinate = G_scal/96
\providecommand{\Rds}{R_{\mathrm{dS}}}    % de Sitter Curvature
\providecommand{\Mpl}{M_{\mathrm{Pl}}}    % Planck mass
\providecommand{\lp}{\ell_{\mathrm{Pl}}}  % Planck length
\providecommand{\Pmax}{P_{\mathrm{max}}}  % Maximum momentum
\providecommand{\Mplbar}{\bar{M}_{\mathrm{Pl}}} % Reduced Planck mass
\providecommand{\eff}{\mathrm{eff}}       % Subscript for effective quantities
\providecommand{\Cosmo}{\ensuremath{\Lambda}}                % Cosmological constant (use \Cosmo in source)
\providecommand{\LCDM}{\ensuremath{\Cosmo\mathrm{CDM}}}  % \LambdaCDM (typesets with \Cosmo)
\providecommand{\Lcut}{\ensuremath{\Lambda_{\mathrm{cut}}}}  % Cutoff scale (brace when used as script)
\providecommand{\Rcs}{\ensuremath{R_{\mathrm{cs}}}}      % Configuration-space radius
\providecommand{\CosmoEff}{\ensuremath{\Cosmo_{\eff}}}   % Effective cosmological constant

% Disambiguation: \Rmom is the momentum-space (S^3) radius. 
% Reserve plain R for the 4D curvature scalar in f(R).
\providecommand{\Rmom}{\ensuremath{R_{\mathrm{mom}}}} 
\providecommand{\Punit}{P_{0}} % Momentum unit set by \Rmom (= \hbar/\Rmom)
\providecommand{\Rvec}{\ensuremath{R_{\mathrm{vec}}}} % invariant vector-space radius in the reduced round sector
% Specific numeric values and expressions
\providecommand{\Vlogcap}{L_{\mathrm{Hopf}}^{(1)}}
\providecommand{\Hopf}{\operatorname{Hopf}}
% Scalar phase (S^1) sector + universal cover notation (Paper I-A)
\providecommand{\tper}{T_\star}      % Time period/scale
\providecommand{\Eper}{E_\star}      % Energy period/scale
\providecommand{\ttilde}{\tilde t}   % Universal cover time
\providecommand{\Etilde}{\tilde E}   % Universal cover energy

% ======================================================================
% 10A. SERIES-WIDE LOGICAL / NOTATIONAL DISAMBIGUATIONS
% ======================================================================
% Use these symbols to avoid collisions among: (i) the Paper--III
% no-insertion/geometric-closure class, (ii) Paper--I-A trimer diagonals,
% (iii) Paper--V-A/V-C family-space diagonals, (iv) scalar-index running,
% and (v) protected shell-clock labels.
\providecommand{\Rprim}{\ensuremath{\mathsf R_{\mathrm{prim}}^{\Sigma}}} % primitive finite structural readouts
\providecommand{\Rloop}{\ensuremath{\mathsf R_{\circ}^{\Sigma}}} % closed-loop readout role
\providecommand{\Rbdry}{\ensuremath{\mathsf R_{\partial}^{\Sigma}}} % boundary-conditioned readout role
\providecommand{\Rgen}{\ensuremath{\mathsf R_{\mathrm{fin}}^{\Sigma}}} % finite generated closure of primitive roles
\providecommand{\PhysOcc}{\operatorname{PhysOcc}}
\providecommand{\Actual}{\operatorname{Actual}}
\providecommand{\GeomLic}{\operatorname{GeomLic}}
\providecommand{\InvDist}{\operatorname{InvDist}}
\providecommand{\Dcl}{\ensuremath{\mathsf{D}_1}}      % unmarked positive Dirac-even closed spectral-readout form class
\providecommand{\Dheat}{\ensuremath{\mathsf D_{\mathrm{heat}}}} % one-parameter heat-semigroup subclass
\providecommand{\DtriI}{\ensuremath{D^{\mathrm{tri}}_1}} % Paper--I-A cyclic-trimer diagonal
\providecommand{\DtriII}{\ensuremath{D^{\mathrm{tri}}_2}}
\providecommand{\DfamI}{\ensuremath{D^{\mathrm{fam}}_1}} % Family-space diagonal selector
\providecommand{\DfamII}{\ensuremath{D^{\mathrm{fam}}_2}}
\providecommand{\alpharun}{\ensuremath{\alpha_{\mathrm{run}}}} % scalar spectral-index running
\providecommand{\Ksh}{\ensuremath{K_{\mathrm{sh}}}} % protected shell-clock coordinate
\providecommand{\ddef}{\ensuremath{\delta_{\mathrm{def}}}} % universal defect scalar after Paper--III realization

% ======================================================================
% 10B. SERIES-WIDE LOGICAL STATUS TAGS
% ======================================================================
% Upright small-caps wrapper for audit/status labels.  It prevents LaTeX from
% trying to use unavailable small-caps italic font shapes inside theorem bodies
% such as TU/lmr/m/scit.
\providecommand{\AuditLabel}[1]{}
\renewcommand{\AuditLabel}[1]{{\normalfont\scshape #1}}

% These tags are bookkeeping labels for consistency audits. They do not
% create new mathematical hypotheses; they record how a datum is being used
% in the proof chain.
\providecommand{\StatusPrimitive}{}
\renewcommand{\StatusPrimitive}{\AuditLabel{Primitive}}
\providecommand{\StatusDerived}{}
\renewcommand{\StatusDerived}{\AuditLabel{Derived}}
\providecommand{\StatusImported}{}
\renewcommand{\StatusImported}{\AuditLabel{Imported}}
\providecommand{\StatusConditional}{}
\renewcommand{\StatusConditional}{\AuditLabel{Conditional}}
\providecommand{\StatusBridge}{}
\renewcommand{\StatusBridge}{\AuditLabel{Bridge-derived}}
\providecommand{\StatusComparison}{}
\renewcommand{\StatusComparison}{\AuditLabel{Comparison-only}}

% ======================================================================
% 10C. SERIES-WIDE CLAIM-WORDING AUDIT LABELS
% ======================================================================
% These are textual audit labels used to keep modal language consistent
% across the series.  They introduce no new mathematical assumptions.
\providecommand{\ClaimForced}{}
\renewcommand{\ClaimForced}{\AuditLabel{forced}}
\providecommand{\ClaimUnique}{}
\renewcommand{\ClaimUnique}{\AuditLabel{unique}}
\providecommand{\ClaimCanonical}{}
\renewcommand{\ClaimCanonical}{\AuditLabel{canonical}}
\providecommand{\ClaimMinimal}{}
\renewcommand{\ClaimMinimal}{\AuditLabel{minimal}}
\providecommand{\ClaimClosed}{}
\renewcommand{\ClaimClosed}{\AuditLabel{closed}}
\providecommand{\ClaimPrediction}{}
\renewcommand{\ClaimPrediction}{\AuditLabel{prediction}}

% ======================================================================
% 10D. SERIES-WIDE THEOREM-CHAIN AUDIT LABELS
% ======================================================================
% These labels mark the direction of load-bearing theorem chains. They are
% audit annotations only; they introduce no new mathematical assumptions.
\providecommand{\ChainInput}{}
\renewcommand{\ChainInput}{\AuditLabel{Chain input}}
\providecommand{\ChainOutput}{}
\renewcommand{\ChainOutput}{\AuditLabel{Chain output}}
\providecommand{\ChainGuard}{}
\renewcommand{\ChainGuard}{\AuditLabel{Chain guard}}
\providecommand{\ChainNoBackflow}{}
\renewcommand{\ChainNoBackflow}{\AuditLabel{No backflow}}

% ======================================================================
% 10E. SERIES-WIDE LOCAL THEOREM-AUDIT LABELS
% ======================================================================
% These labels mark a statement-level dependency audit around individual
% theorem/definition nodes. They are bookkeeping labels only.
\providecommand{\LocalAudit}{}
\renewcommand{\LocalAudit}{\AuditLabel{Local audit}}
\providecommand{\UsesOnly}{}
\renewcommand{\UsesOnly}{\AuditLabel{Uses only}}
\providecommand{\Exports}{}
\renewcommand{\Exports}{\AuditLabel{Exports}}
\providecommand{\DoesNotUse}{}
\renewcommand{\DoesNotUse}{\AuditLabel{Does not use}}

% ======================================================================
% 10F. SERIES-WIDE NUMERICAL-ROW AUDIT LABELS
% ======================================================================
% These labels classify individual rows in numerical, benchmark, status,
% and comparison tables.  They are audit annotations only and do not change
% any theorem, value, formula, or comparison number.
\providecommand{\RowAudit}{}
\renewcommand{\RowAudit}{\AuditLabel{Numerical-row audit}}
\providecommand{\RowTheorem}{}
\renewcommand{\RowTheorem}{\AuditLabel{Theorem-fixed row}}
\providecommand{\RowBridge}{}
\renewcommand{\RowBridge}{\AuditLabel{Bridge-fixed row}}
\providecommand{\RowRepresentative}{}
\renewcommand{\RowRepresentative}{\AuditLabel{Representative row}}
\providecommand{\RowComparison}{}
\renewcommand{\RowComparison}{\AuditLabel{Comparison row}}
\providecommand{\RowCalibrated}{}
\renewcommand{\RowCalibrated}{\AuditLabel{Calibrated row}}
\providecommand{\RowEstimate}{}
\renewcommand{\RowEstimate}{\AuditLabel{Estimate row}}
\providecommand{\RowOpen}{}
\renewcommand{\RowOpen}{\AuditLabel{Open row}}
\providecommand{\RowNoSelector}{}
\renewcommand{\RowNoSelector}{\AuditLabel{No selector}}

% ======================================================================
% 10G. TABLE HELPERS
% ======================================================================
% Table helpers used by Paper V-C.  The unused readout-label preset has been
% removed; only active table helpers remain.
\providecommand{\gpsection}[1]{%
  \addlinespace[0.55em]
  \multicolumn{3}{@{}l}{\textsc{#1}}\\[-0.25em]
  \midrule
}
\providecommand{\mathcell}[1]{%
  \(\begin{gathered}#1\end{gathered}\)%
}

\newcommand{\DerBy}[1]{%
  \begin{tabular}[t]{@{}l@{}}#1\end{tabular}%
}

\newcommand{\TableNotes}[1]{%
  \par\medskip
  \begin{minipage}{0.94\linewidth}
  \footnotesize\emph{Notes.} #1
  \end{minipage}\par
}
% ======================================================================
% 10H. PAPER VI-A / QCD NOTATION
% ======================================================================
% Shared symbols used by Paper VI-A and downstream QCD/Yang--Mills handoffs.
\providecommand{\Pqcd}{\mathfrak P_{\rm VIII}^{\rm QCD}}
\providecommand{\Iqcd}{\mathfrak I_{\rm QCD}^{(7)}}
\providecommand{\muQCD}{\mu_{\rm QCD}^{\rm int}}
\providecommand{\LambdaQCD}{\Lambda_{\rm QCD}}
\providecommand{\LambdaBH}{\Lambda_{\rm BH}^{\rm QCD}}
\providecommand{\HQCD}{\mathcal H_{\rm QCD}^{\rm BH}}
\providecommand{\HQCDfin}{H_{\rm QCD}^{\rm fin}}
\providecommand{\Hraw}{\mathcal H_{\rm raw}}
\providecommand{\HphysQCD}{\mathcal H_{\rm phys}^{\rm QCD}}
\providecommand{\Vcol}{V_C}
\providecommand{\Ps}{P_{\rm sing}}
\providecommand{\PiOct}{\Pi_8}
\providecommand{\MSbar}{\overline{\rm MS}}
\providecommand{\YM}{\mathrm{YM}}
\providecommand{\Ktri}{K_{\triangle}}
\providecommand{\KYM}{K_{\YM}^{\rm BH}}
\providecommand{\aBH}{a_{\rm BH}}
\providecommand{\epsBH}{\varepsilon_{\rm BH}}
\newcommand{\RowConstructive}{\textnormal{\textsc{Constructive}}}
\newcommand{\RowObligation}{\textnormal{\textsc{Bridge obligation}}}
\newcommand{\RowTarget}{\textnormal{\textsc{Target row}}}
\newcommand{\RowValidation}{\textnormal{\textsc{Validation row}}}
\providecommand{\slashed}[1]{\not{#1}}

% ======================================================================
% 10I. PAPER VI-B / CONTINUOUS-QFT NOTATION
% ======================================================================
% Shared symbols used by Paper VI-B and any downstream continuous-QFT or
% Yang--Mills handoff paper.  All definitions are guarded so individual
% papers may override them locally before loading this file if needed.
\providecommand{\A}{\mathfrak A}
\providecommand{\Hphys}{\mathcal H_{\rm phys}}
\providecommand{\QCD}{\mathrm{QCD}}
\providecommand{\BH}{\mathrm{BH}}
\providecommand{\loc}{\mathrm{loc}}

\providecommand{\RBH}{\mathcal R_{\rm BH}}
\providecommand{\DG}{\mathfrak D_G}
\providecommand{\PW}{\mathcal{P\!W}} % Peter--Weyl spectral-symbol subspace
\providecommand{\Toe}{\operatorname{Toe}} % Toeplitz compression/representative
\providecommand{\Hol}{\operatorname{Hol}} % Holonomy

% ======================================================================
% 11. LATE-LOADED PACKAGES (Must be at the end of the preamble)
% ======================================================================
% Load late packages only if a paper has not already loaded them, avoiding
% option clashes in paper-specific preambles.
\makeatletter
\@ifpackageloaded{hyperref}{}{\usepackage{hyperref}}
\makeatother
% --- Series-wide disambiguations (Pass 3) ---
\providecommand{\diracdelta}{\delta_{\rm D}} % Dirac delta distribution (delta reserved for hierarchy)
\providecommand{\Hg}{H_{\mathrm g}} % Hodge harmonic projector (keep H for Hubble)
\providecommand{\Gr}{G_{\mathrm r}} % Green operator (keep G for Newton)
\providecommand{\FD}{\mathrm{D}} % functional/variational derivative symbol (delta reserved)
\providecommand{\Id}{\mathbf{1}} % identity/Kronecker tensor (avoid \delta_{ij})


% BEGIN INLINE SOURCE: Paper00-ActiveExternalReferences-v2_130.tex
% Framework v2.191: current reference binding; existing definitions retained.
% Framework v2.190: current22-root registry with native clock-evolution proofs.
% Framework v2.189: exact current22-root external registry, including retained v2.188 hardening owners.
% Framework v2.187: current registry synchronization only; source release v2.186 preserved.
% Framework v2.186: authoritative current 22-paper registry.
% Framework v2.182: all current 22 logical paper AUX bindings.
% Framework v2.170: current I-B v1.43 external AUX binding.
% Framework v2.167: exact current 22-root registry after CP--BL2 integration.
% Framework v2.166: exact current 22-root compact-release AUX registry.
% Framework v2.165: authoritative current root bindings.
% Framework v2.164: exact current 22 manuscript AUX bindings.
% Framework v2.162: finite metric, functional Hilbert and scoped Ward integration.
% Framework v2.161: current 22 manuscript AUX bindings; source only.
% Framework v2.156 / registry v2.113: synchronized current 22 manuscript roots; source-only registry.
% Framework v2.143 / registry v2.111: active manuscript AUX basenames; no generated AUX or PDF included.
% ======================================================================
% Central external-reference imports, numbers only, no PDF jumping
% ======================================================================

% xr-hyper is loaded centrally so that papers may retain their existing
% \usepackage{xr-hyper} lines without redefining \externaldocument after this
% registry is installed.  If plain xr was loaded earlier, retain it.
\makeatletter
\@ifpackageloaded{xr-hyper}{}{%
  \@ifpackageloaded{xr}{}{\usepackage{xr-hyper}}%
}
\makeatother

% Make external-document registration idempotent.  The registry identifies a
% logical paper by its Paper<number><optional letter> prefix.  Prefix-qualified
% imports remain independent.  Current paper roots contain no local
% \externaldocument calls; this guard also makes accidental duplicates harmless.
\ExplSyntaxOn
\prop_new:N \g_bh_external_refs_prop
\cs_new_eq:NN \BHOriginalExternalDocument \externaldocument

\cs_new_protected:Npn \bh_external_document:nnnn #1#2#3#4
 {
  \tl_set:Nn \l_tmpa_tl {#3}
  \regex_replace_once:nnN
   { \A (?: xref/ )? (Paper[0-9]+[A-F]?) .* \Z }
   { \1 }
   \l_tmpa_tl
  \tl_set:Nx \l_tmpb_tl { #1 | \tl_use:N \l_tmpa_tl }
  \prop_if_in:NVTF \g_bh_external_refs_prop \l_tmpb_tl
   {
    \PackageInfo{BornHopf}
     {Skipping~duplicate~external~refs~for~#3}
   }
   {
    \prop_gput:NVn \g_bh_external_refs_prop \l_tmpb_tl {#3}
    \IfNoValueTF{#4}
     { \BHOriginalExternalDocument[#1][#2]{#3} }
     { \BHOriginalExternalDocument[#1][#2]{#3}[#4] }
   }
 }

\RenewDocumentCommand{\externaldocument}{O{} O{} m o}
 { \bh_external_document:nnnn {#1}{#2}{#3}{#4} }

\NewDocumentCommand{\BHExternalRefs}{m}
 {
  % Do not import the aux file of the logical paper currently being compiled,
  % even when the job and registry use different version suffixes.
  \tl_set:Nx \l_tmpa_tl { \c_sys_jobname_str }
  \regex_replace_once:nnN
   { \A (Paper[0-9]+[A-F]?) .* \Z }
   { \1 }
   \l_tmpa_tl
  \tl_set:Nn \l_tmpb_tl { #1 }
  \regex_replace_once:nnN
   { \A (Paper[0-9]+[A-F]?) .* \Z }
   { \1 }
   \l_tmpb_tl
  \str_if_eq:eeTF { \tl_use:N \l_tmpa_tl } { \tl_use:N \l_tmpb_tl }
   {
    \PackageInfo{BornHopf}
     {Skipping~external~refs~for~#1~because~this~is~the~current~document}
   }
   {
    \IfFileExists{xref/#1.aux}
     { \externaldocument[][nocite]{xref/#1}[#1.pdf] }
     {
      \IfFileExists{#1.aux}
       { \externaldocument[][nocite]{#1} }
       {
        \PackageInfo{BornHopf}
         {External~aux~file~#1.aux~not~found;~no~labels~imported}
       }
     }
   }
 }
\ExplSyntaxOff

% External references are automatic.  This file is the sole owner of both the
% bibliography and the current Paper-I-A--VII-D AUX registry (30 roots including I-C).  Missing auxiliaries
% are informational on the first pass and are populated by the source-only
% two-pass series build.  Basenames deliberately omit temporary download/copy
% suffixes such as ``(1)'' or ``-1''.
  % Paper I-A: intrinsic S3 and exceptional group introduction.
  \BHExternalRefs{Paper1A-Geometry-v5_178}

  % Paper I-B: complete exceptional geometry and finite spectral structures.
  \BHExternalRefs{Paper1B-ExceptionalGeometry-v1_16}
  % Paper I-C: reciprocal product M8 geometry.
  \BHExternalRefs{Paper1D-TechnicalMaterial-v1_67}
  % Paper I-D: terminal technical proofs and calculations.
  \BHExternalRefs{Paper1C-M8Geometry-v1_142}

\ifdefined\BHMathOnly
  % Papers I-A and I-B import only the mathematical atlas labels.
\else
  % Paper II: five physical-mathematical main papers plus one technical proof volume.
  \BHExternalRefs{Paper2A-GeometricConstitution-v3_224}
  \BHExternalRefs{Paper2B-MomentumCarrier-v2_145}
  \BHExternalRefs{Paper2C-ReciprocalLocalization-v2_170}
  \BHExternalRefs{Paper2D-SymmetrySpectra-v2_143}
  % The former II-F export labels remain in II-E; current II-F is the technical volume.
  \BHExternalRefs{Paper2E-KinematicExports-v3_138}
  % Paper II-F: complete technical proof and calculation owner for II-A--II-E.
  \BHExternalRefs{Paper2F-TechnicalMaterial-v1_67}

  % Paper III: Dynamics and Intrinsic Geometric Constitution (six components).
  \BHExternalRefs{Paper3A-GeneratedDynamics-v1_177}
  \BHExternalRefs{Paper3B-SpectralAction-v1_189}
  \BHExternalRefs{Paper3C-ProtectedHierarchy-v1_191}
  \BHExternalRefs{Paper3D-RunningResponse-v1_166}
  \BHExternalRefs{Paper3E-GravityAction-v1_198}
  \BHExternalRefs{Paper3F-CompleteDynamics-v1_265}
  % Paper III-G: complete preserved technical proof owner.
  \BHExternalRefs{Paper3G-TechnicalMaterial-v1_54}

  % Paper IV-A: generated finite quantum mechanics.
  \BHExternalRefs{Paper4A-FiniteQuantumMechanics-v4_215}
  % Paper IV-B: complete finite quantum processes and physical response.
  \BHExternalRefs{Paper4B-QuantumResponse-v1_53}
  % Paper IV-C: extended proofs and finite technical calculations.
  \BHExternalRefs{Paper4C-TechnicalMaterial-v1_49}

  % Paper V-A: The Standard Model.
  \BHExternalRefs{Paper5A-FiniteMatterGeometry-v3_212}

  % Paper V-B: Finite Spectrum and Flavor Closure.
  \BHExternalRefs{Paper5B-TerminalSpectrum-v5_201}

  % Paper V-C: Finite Interaction Boundaries of the Standard Model.
  \BHExternalRefs{Paper5C-InteractionResponse-v2_179}
  % Paper V-D: preserved technical proof and calculation owner.
  \BHExternalRefs{Paper5D-TechnicalMaterial-v1_55}

  % Paper VI-A: Gravity and Closed Cosmology.
  \BHExternalRefs{Paper6A-GravityCosmology-v2_262}

  % Paper VI-B: Geometric Dark Sector and Cosmological Readouts.
  \BHExternalRefs{Paper6B-CosmologyReadouts-v1_209}

  % Paper VI-C: complete technical proof and comparison owner.
  \BHExternalRefs{Paper6C-TechnicalMaterial-v1_51}

  % Paper VII-A: Finite QCD Kernel and Nonperturbative Closure.
\BHExternalRefs{Paper7A-QCDKernel-v1_177}

  % Paper VII-B: Hadronic, Thermal, and Nuclear QCD Readouts.
  \BHExternalRefs{Paper7B-QCDReadouts-v1_163}

  % Paper VII-C: Continuous Quantum Field Theory.
  \BHExternalRefs{Paper7C-ContinuousQFT-v2_206}

  % Paper VII-D: complete technical material for VII-A through VII-C.
  \BHExternalRefs{Paper7D-TechnicalMaterial-v1_56}

\fi

% END INLINE SOURCE: Paper00-ActiveExternalReferences-v2_130.tex


\hypersetup{hidelinks}
\makeatletter
\renewcommand*\l@subsection{\@dottedtocline{2}{1.5em}{3.25em}}
\makeatother





\usetikzlibrary{positioning, calc}

\title{\textbf{Born--Hopf Geometry II-A: Geometric Constitution and Selection}\\[0.5em]
\large G0, GP1--GP3, least realization, and licensed application interfaces}

\author{Peter Kugelmann\thanks{\texttt{pkugelmann@proton.me}}}
\date{Version 3.227 --- 20 November 2026}

\setlength{\emergencystretch}{3em}
\begin{document}

\maketitle

\begin{abstract}
Paper~II-A states the constitutional grammar of Born--Hopf Geometry:
G0 gives the occurrence semantics of invariantly distinguishable geometric
comparisons, while GP1 (finite/bounded realization), GP2 (generated,
exact and host-faithful realization), and GP3 (irreducible/least
realization) act on complete candidate presentations of an already constituted
geometrically primary role.  The primitive role predicate
$\operatorname{Geo}_R$ includes its comparison-bearing occurrence,
identity-rooted reapplication and intrinsic comparison diagrams;
G0 alone neither provides that strong role nor selects an occurring witness.
GP1 imposes stagewise typed capacity and no numerical radius or harmonic level.
GP2 retains exactly the outputs of previously licensed typed constructors,
including their relations and intrinsic defects, without making ambient
mathematical maps into physical constructors. GP3 compares whole admissible
candidate realizations; if no unique least class exists, the unresolved
family is retained rather than arbitrarily resolved.

On the stated primitive-role and comparison domains, the resulting least
operation-parent skeletons are $U(1)$ for the scalar role and
$SU(2)\simeq S^3$ for the noncommutative transfer role, with the
additional $SO(3)$ quotient relation excluded by source-exactness.
The overall momentum-sphere scale and possible unmarked faithful carrier
lifts remain separately typed. The finite momentum-seed support and
$m_p=1$ are established in Paper~II-B under their own constructor and
whole-candidate comparison requirements, not posited here.
For licensed rooted finite words, generated structural content and its
co-realized presentation are canonically equivalent in their declared
sorts.  The same-owner invariance and reciprocal-branch transport rules
preserve complete input marks, role boundaries and constructor licenses.

Constitutional descent is not physical occurrence: a finite exceptional
incidence design, spectral diagnostic, action response, state/effect role,
or conditional complete-$2O$ particle carrier does not by itself
actualize a particular source, preparation, detector, history, mass, or
readout.  The paper records the precise hypothesis and ownership
boundaries for these applications, including the conditions under which
least comparisons, finite-band identification and action-source
inheritance are valid.  No additional constitutive principle,
independently imposed physical mechanism, or numerical calibration is
introduced.
\end{abstract}

\tableofcontents
\newpage

\section{Constitutional scope, inputs, and series contract}
\label{paper2a:sec:organized-scope}

\subsection{Introduction, series allocation, and notation}
\label{p2a:sec:introduction}

\paragraph{Full technical derivations.}
The active auxiliary classifications, exact calculations, and negative
results supporting this paper are collected without abridgment in
Paper~II-F, Section~\ref{p2ftech:sec:constitution}.
That companion is a proof owner, not an additional constructor,
physical occurrence, state, or calibration assumption.

\paragraph{Typed co-realization and inheritance.}
The universal statement is proved in
Theorem~\ref{p2a:thm:universal-corealization-bridge-normal-form}.
Theorem~\ref{p2a:thm:evaluation-dual-corealization} applies it to operations
and dynamics, and
Corollary~\ref{p2a:cor:born-invertible-corealization} identifies Born
reciprocity as its invertible self-dual case, and
Theorem~\ref{p2a:thm:generated-born-inheritance} propagates that generated
involution through every licensed descendant word for which it remains
unmarked.  Proposition~\ref{p2a:prop:typed-born-branch-transport} handles the
complementary case in which exchange transports the complete typed
presentation to a distinct reciprocal branch.  The rules preserve the G0
token, constructor-license, target-sort, family-selection, calibration,
analytic, and actual-history firewalls.


Paper~II-A is the constitutional component of the five-component Paper~II
sequence, whose terminal component is the merged Paper~II-E.  Its native
task is to allocate the compact parent operation geometries and to state the
common finite-level/discharge contract before any coordinate, spectral,
gravity-sector, or internal-symmetry consequence is used.  Papers~II-B--II-E
are coordinated consumers of these exports.  Their results do not enter the
parent-allocation proof below; conversely, they own
the constructor-relative finite-carrier consequences and, for roles whose
signature explicitly constitutes a finite tower, the corresponding tower,
trajectory, and recovery consequences rather than deferring those consequences
to dynamics.

Paper~I-A~\cite{Kugelmann2026a} supplies the mathematical source layer used after
parent allocation and during finite-carrier constitution: the round Hopf-fibered sphere \(S^3(\Rmom)\), its derived
\(SO(4)\) isometry action and \(\Spin(4)\) lift, and the declared finite
\(\Spin(4)\)-equivariant module packages
\[
\mathcal H_m=\bigoplus_{j_L+j_R\le m}V_{j_L}\otimes V_{j_R},
\]
and the associated representation, Clifford, Hurwitz-link, contact, and trimer
structures.  Paper~II does not assume at the outset that the primary transfer parent is
\(S^3\).  It first derives the abstract A-side transfer group and
then solves the narrower geometric realization problem required
for kinematics.

The derivation uses an explicit cumulative dependency calculus organized by
Geometric Primacy and Minimal Constitution.  G0 validates the physical
occurrence represented by the initial existence assertion; GP1--GP3 then close
its structured realization:
\[
\boxed{\begin{gathered}
\mathrm{G0}[\mathsf E_R]:\ \exists\mathfrak M\,\operatorname{Geo}_R(\mathfrak M)\\
\longrightarrow\
\operatorname{Adm}_R
=\{\mathfrak M:\operatorname{Geo}_R(\mathfrak M),
\ \mathrm{GP1}(\mathfrak M),\ \mathrm{GP2}_R(\mathfrak M)\}\\
\longrightarrow\
\mathcal M_R
:=\operatorname{Close}_{\mathrm{GP3}}\!\left(\operatorname{Adm}_R\right)\\
\longrightarrow\
P_R.
\end{gathered}}
\]
If \(\mathcal M_R=\{[\mathfrak M_R]\}\) is uniquely least, then
\(P_R=P_{<R}\cup\operatorname{Pres}_{\mathfrak M_R}(\mathfrak M_R)\).
If instead \(\mathcal M_R\) is an unresolved irreducible family, \(P_R\)
contains only \(P_{<R}\) and candidate-generated items proved common to every
member.
GP1 constrains the carrier and represented operation layer together and its
comparison specialization GP1-TB supplies total boundedness where a
constituted comparison is present.  GP2 determines the exact typed content of
each candidate: only outputs of its geometrically licensed physical
constructors and operations enter, and every such output is retained in its
declared sort.  Hence
\[
 \forall\mathfrak M\in\operatorname{Adm}_R:\qquad
 \operatorname{Pres}_{\mathfrak M}(\mathfrak M)
 =\operatorname{Gen}_R(\mathfrak M).
\]
This candidate-internal equality does not impose the generated self-structure
of one rejected candidate on every other candidate.  GP3 acts only after each
candidate has been exactified and compares whole complete realizations through
role-faithful structural comparisons.  It retains a unique least class when
one exists; otherwise it keeps the irreducible family unresolved.
Candidate-specific dimension, rank, center, quotient relations,
representation labels, metrics, and operation coefficients may occur inside a
candidate but cannot witness that candidate's own selection.

Once a stage is closed, its constituted finite presentation is frozen and may
not be discarded by a later GP3 comparison.  Theorem consequences remain
available for inference without automatically becoming new signature data.
The distinction between a theorem consequence and a generated item is
load-bearing: physical generation follows the constituted constructor grammar,
not the entire mathematical closure of the representation category.

The primitive input is the assertion
\(\mathsf E_R:\exists\mathfrak M\,\operatorname{Geo}_R(\mathfrak M)\).
Its physical-occurrence status is the G0 witness; G0 does not choose its
candidate realization or prove that the witness occurs.
The predicate \(\operatorname{Geo}_R\) is strong: it requires an actual
effective comparison-bearing occurrence, a full identity-rooted local
reapplication constructor, and the intrinsic role diagrams that compare
successive reapplications.  It does not choose a completed global parent or
freeze one candidate's multiplication table.  GP1 and GP2 derive, in every
admissible candidate, a nontrivial connected compact \emph{topological}
parent and exactify the intrinsic order-swap diagram; a compact-Lie upgrade
of the entire parent is conditional on a faithful continuous finite port or
an independently certified Lie structure.  That diagram is, exhaustively, commuting or
order-sensitive; the resulting branch is geometric rather than a preassigned
role-type bit.  Peter--Weyl separation supplies finite Lie \emph{image germs} of arbitrary
compact parents; their admissible GP2 exactifications and compact root
subroles expose complete \(U(1)\) and \(SU(2)\) endpoint candidates
before any candidate's self-data are frozen.  GP3 then compares the complete exact candidates and
derives the unique least local/global parent-role skeleton.  Relation exactness
retains the compact source form \(SU(2)\) rather than the quotient \(SO(3)\).
The endpoint theorem is a comparison between complete role realizations, not a
claim that the selected endpoint contains every self-operation of each rejected
candidate.

To make the closure comparable in scope to the complete \(S^3\) catalogue of
Paper~I-A, this paper records all standard consequences used by the series in
three tiers.  The tag \([{\rm CL}]\) denotes consequences common to compact
connected Lie parents.  The tag \([{\rm MR}]\) denotes consequences of the
GP3 irreducible/least scalar/noncommutative role allocation.  The tag
\([{\rm SU2}]\) denotes consequences of GP2 relation-exact globalization.
The labels (FD), (LR), (OR1), and center-oddness denote
regularity or representation-theoretic diagnostics; none is a separate
premise of the primary allocation theorem.

Here \emph{complete consequence catalogue} means closure under the standard
compact-Lie, homogeneous-metric, topology, generated-test, representation, and
elliptic-analysis constructions used by this series.  It does not mean the
literally infinite set of all formal corollaries of the postulates.

At the bare level, A denotes active comparison-preserving transformations and
B denotes generated localizations and tests.  A acts on B by pullback; map and
relation presentations may overlap.  Identity and sequential composition are
ordinary function operations.  Initially they form only an
identity-containing forward monoid.  The GP1-TB comparison specialization
--- total boundedness of the constituted comparison, with the
finite-operational uniform bridge of Corollary~\ref{p2:cor:gp1-tb} ---
provides the canonical compact completion.  Each completed exact isometry is
surjective, and the closure of the forward monoid inside the compact isometry
group is automatically a group
(Theorem~\ref{p2:thm:gp1-bounded-isometry-closure}); reversals are therefore
derived.  The full metric-profile family canonically separates points and operations
mathematically.  It belongs to the physical B-test
family only under the metric-partial-evaluation closure declared below.

This compact group is a \emph{candidate realization}, not yet the selected
physical role.  At the allocation stage Geometric Primacy asserts
\(\mathsf E_R\) and fixes the role-geometric condition
\(\operatorname{Geo}_R\), not the complete self-description of any candidate.
Applying GP2 candidatewise derives connectedness and exactifies the intrinsic
abelian/noncommutative order-swap diagram.  No designated operation pair or
its generated dense subgroup belongs to the previously frozen inventory.  The
compact group just constructed is one complete exact candidate; its full
profiles and operation table are internally retained but are not common prior
data for all candidates.
Peter--Weyl exposes a finite-dimensional noncommutative candidate, and a
compact root construction exposes a complete type-\(A_1\) candidate.  Because
GP2 first exactifies every candidate and GP3 then compares those whole
candidates before one candidate's self-data are frozen, the least class is the
structurally irreducible compact connected realization of the geometrically
primary role.  This gives \(U(1)\) for the scalar role and the \(A_1\) local
type for the order-sensitive role.

The quotient \(SU(2)\to SO(3)\) cannot be decided by a local metric argument,
but it is decided inside the prior GP2 admissibility filter.  For the
\(\mathfrak{su}(2)\) candidate, \(SU(2)\) is the compact
source-simply-connected integration.  Passing to \(SO(3)\) adds \(-1=1\), a
global equality not generated by the role geometry.  GP2 relation exactness
therefore removes the quotient before GP3 compares the remaining complete
relation-exact structured realizations.
The resulting selected transfer parent is
\[
M_p\cong SU(2)\cong S^3_A(\Rmom).
\]
The negative Killing form supplies its canonical automorphism-invariant
round conformal structure; the overall scale \(\Rmom>0\) is not fixed by that
form, since every positive multiple of \(-B\) is equally Ad-invariant.  The initial FRC-b comparison presentation
\(D_\infty\) need not be this metric.  Conditions (FD) and (LR) below retain
their valid status as sufficient tests for differentiability of
\(D_\infty\), but the round downstream geometry is the Killing geometry and
does not depend on those tests.


\begin{corollary}[Momentum-sphere curvature coordinate and homothety frontier]
\label{p2:cor:v398-momentum-curvature-homothety}
On the selected nonabelian momentum branch
\[
 M_p\cong SU(2)\cong S^3_A(\Rmom),
\]
the inherited Paper-I-A round metric has intrinsic scalar curvature
\[
 \boxed{\mathcal K_p:=\mathcal R_p^{(3)}:=\operatorname{Scal}(M_p)=\frac6{\Rmom^2}}.
\]
The pair \((\Rmom,\mathcal K_p)\) gives two coordinates on the same
one-dimensional positive homothety modulus.  GP1--GP3 do not choose its
numerical representative.  A later generated dimensionful curvature standard
may actualize this modulus without creating a second momentum parent, but it
must retain every dimensionless normalization generated by the branch geometry.
\end{corollary}

\begin{remark}[Downstream actualization and normalization firewall]
\label{p2:rem:v398-momentum-curvature-firewall}
No equality between \(\mathcal K_p\) and a later Einstein/response scale is a
Paper-II structural theorem.  One-dimensionality of the positive homothety
modulus fixes the dimensionful ray but does not force the coefficient of a
cross-owner relation to be one.  If a later constituted constructor consumes
a generated dimensionless normalization datum together with this curvature
coordinate, Theorem~\ref{p2:thm:gp2-generated-owner-normalization-inheritance}
requires that datum to be transported unless the frozen grammar independently
generates a deformation owner.  Such a downstream actualization does not
retroactively become a GP1--GP3 numerical scale selector.
\end{remark}

Paper~I-A is the mathematical compendium for the round
\(S^3(\Rmom)\), its scalar, spinor, vector-field, and bracket
structures and its declared finite module family. Paper~I-B
owns the exceptional finite-subgroup and incidence constructions.  Under the revised constructor-relative GP2
semantics, Paper~II-B may use those already generated round-host constructions
to determine the initial momentum carrier without first selecting a spectral
level.  Its assigned chronology is
\[
   \text{round-host generated content}
   \Longrightarrow \mathcal G^{\rm car}_{\Sigma_p}\cong\mathcal H_1
   \Longrightarrow m_p=1,
\]
where the final implication is only the intrinsic word-length classification
of the already generated six-type carrier support.  Full endpoint multiplicity-one
is a separately proof-gated comparison conclusion, not an antecedent of
\(m_p=1\).  The optional \(\mathsf P_{\rm sli}\), tower,
and generated-boundary interfaces remain available for later open-ended roles
but are not momentum-seed premises.  A particular ambient coordinate action,
its omitted Wigner--Eckart data, absolute protected radius, raw coordinate
normalization, and coordinate--coordinate commutator magnitude remain relative
to a chosen point \(\mathsf C\in\mathfrak M_{\rm coord}\), with
\(\mathsf C_{\rm PI}\) denoting Paper~I-A's representative; the carrier fixes
only the completion skeleton.  This completion dependence does not extend to
every normalized Clifford statement: current Paper~II-C proves that on the
nondegenerate protected stratum \(\varrho_{\rm sp}>0\), normalization followed
by passage to the protected even products removes both the nonzero scale and
the relative chiral phase, yielding a represented real algebra canonically
isomorphic to \(\Cl(0,3)\).  Hurwitz/trimer theorems use the unmarked
core-support output \(\mathsf{Supp}_{2T}\).  An occurring point readout gives
the generic finite-resolution role by
Theorem~\ref{p2:thm:derived-finite-point-support-role}; it gives
\(\mathsf{Supp}_{2T}\) only under the additional domain bridge or direct
core-constructor occurrence of
Corollary~\ref{p2:cor:core-support-inheritance-no-duplicate-witness}.  The
integer equation \(m_p=1\) creates neither a coordinate-completion point nor
either occurrence.

\subparagraph{Notation boundary.}
Bare \(m\) denotes a local truncation or representation index unless explicitly
promoted to a realized role level.  The II-B theorem determining \(m_p\) from generated six-type support, an
actual finite position-frontier history \((\mathrm{HIST}_x^{\rm fin})\), Born
normalization, Lorentzian orientation, Standard-Model couplings, and comparison
representatives are not assumed here; each requires whichever later
structural, dynamical, occurrence, or identification theorem actually proves
or supplies it.  Once such a position frontier occurs, its
\(m_x(\tilde t)\) is a derived resolution statistic rather than an independent
trajectory selector.  Neutral trimer constructions based on the round-host geometry of
Paper~I-A and the exceptional support atlas of Paper~I-B remain
mathematical geometry until an identification theorem assigns further roles.

\subparagraph{Optional finite-role interfaces.}
The symbols \(\mathsf P_{\rm sli}\), \(P_{\rm ceil}\), \(s_\rho\), and
\(m_{\rm adm}\) are retained only for later roles that explicitly invoke the
legacy finite-level interface.  Their exact capacity and floor-function
relations are collected in Paper~II-F, Section~\ref{p2:app:selector-limitations}.  They
are not premises of the initial momentum-seed derivation and do not determine
\(m_p\).

\subsubsection{Notation.}
Throughout this paper:
\begin{itemize}
    \item For roles that explicitly invoke the legacy finite-level interface,
    \(P_{\rm ceil}\in\mathbb R_{\ge0}\) denotes an auxiliary, possibly
    unattained component-spectrum admissibility ceiling, with associated labels
    \(s_\rho\) and \(m_{\rm adm}\).  Their exact formulas and non-selection
    consequences are deferred to Paper~II-F, Section~\ref{p2:app:selector-limitations}.
    They are not momentum-seed selectors.  The symbol \(\Pmax\) is reserved
    for the sharp actual component maximum of an independently constituted
    realized carrier.

    \item \(\Rmom\) is the round compact momentum-sphere radius, equivalently
    the unique positive homothety modulus of the compact parent.  Its intrinsic
    scalar-curvature owner and scalar-curvature length are
    \[
        \mathcal K_p:=\frac6{\Rmom^2},\qquad
        \ell_{\mathcal K}:=\mathcal K_p^{-1/2}=\frac{\Rmom}{\sqrt6}.
    \]
    Thus \(\Rmom\) itself is not called the scalar-curvature length.  A later
    action owner such as \(U\) may differ from \(\mathcal K_p\) by a generated
    dimensionless normalization and must not be identified coefficient-one by
    dimensionality alone.  With \(\Rmom>0\), the momentum-unit normalization is
    \[
        \Punit := \hbar/\Rmom\in\mathbb R_{>0}.
    \]

    \item \(S^3(\Rmom)\) denotes the smooth round parent operation
    geometry of radius \(\Rmom\), not an infinite realized carrier.

    \item \(S^3_m\) is shorthand for the paired object consisting of that
    round parent and the declared finite level-\(m\)
    \(\Spin(4)\)-equivariant module package, with \(j_L+j_R\leq m\).  It is
    neither the group \(\Spin(4)\cong S^3\times S^3\) nor, without further
    qualification, a scalar Peter--Weyl truncation.  The realized momentum-side physical label \(m_p\) is the derived
    half-word-length label of the already constituted momentum carrier and is
    available after the II-B host theorem; it is not an independently
    registered selector.

    \item \(\mathcal H_m\) is the finite module space attached to \(S^3_m\), with
    \[
        \dim(\mathcal H_1)=15.
    \]

    \item Once a realized carrier has been independently constituted, the
    phenomenological comparison \(\lp\sim\hbar/\Pmax\) may be made using
    its sharp \(\Pmax\); this relation is not used as a Paper-II selector.

    \item \(\mathfrak h,q\in\ZZ\) denote integer geometric Hopf charges,
    normalized by
    \[
        \partial_\psi\phi=i\mathfrak h\,\phi/2.
    \]
    The finite diagonal spin weight
    \[
        \mathfrak s=m_L+m_R\in\frac12\ZZ
    \]
    is a separate Cartan grading on the finite \(\Spin(4)\)-module; it is not
    identified with the geometric Hopf weight \(q/2\) unless an additional
    convention is explicitly imposed.

\item The truncation labels are as follows:
\par\smallskip\noindent
\begin{tabularx}{\linewidth}{@{}
  >{$}l<{$}
  @{\quad}
  >{\RaggedRight\arraybackslash}X
@{}}
  m
  & generic truncation-family index, with \(j_L+j_R\le m\), \\[2pt]

  m_{\mathrm{adm}}
  & Paper-II maximal integer Toeplitz-compressed module member below the component
    ceiling, \\[2pt]

  m_p
  & derived compact momentum half-word-length label of the host-selected
    carrier, read off after the II-B theorem \(\mathcal H_p\cong\mathcal H_1\), \\[2pt]

  m_x(\tilde t)
  & derived least scalar harmonic-resolution statistic
    \(m_x^{\rm res}(F_x(\tilde t))\) of an occurring finite generated
    position-test frontier history \((\mathrm{HIST}_x^{\rm fin})\); it is
    not an independently selected integer trajectory.
\end{tabularx}
\par\smallskip

    \item \(c_{\rm op}\) denotes the possibly unattained supremum of the
    invariant rate image of a constituted full timed propagation role.  A
    later symbol \(c_{\rm field}\) or \(c_{\rm light}\) is type-distinct until
    its support and principal ancestry are proved.  On an undeformed
    same-principal-owner branch,
    Theorem~\ref{p2a:thm:undeformed-host-principal-inheritance} then forces
    the corresponding front speed to equal \(c_{\rm op}\); a medium or
    independently generated constitutive deformation may define a different
    effective speed.

    \item \(k=2j\) is the Laplacian eigenvalue index for diagonal
    representations.
\end{itemize}


\subsubsection{Geometric inputs from Papers I-A and I-B}

Paper~I-A~\cite{Kugelmann2026a} supplies the round Hopf/contact geometry, the
derived \(SO(4)\) isometry group and \(\Spin(4)\) lift, and the declared
\(\Spin(4)\)-equivariant module packages \(\mathcal H_m\), their dimension
and minimal-sector results, normalized Toeplitz compression, the
Hurwitz--Walsh--Clifford carrier, and the neutral cyclic trimer grammar.  Paper~II
uses that source layer after
Theorem~\ref{p2:thm:finite-realization-momentum-selection} has derived the
relation-exact \(S^3\) transfer parent.  Under the revised constitutional semantics, the primary momentum handoff is
not a level conditional.  Paper~II-B must use the already constituted round-host
scalar, spinor, coordinate, and canonical finite bracket-image content to prove
\(\mathcal H_p\cong\mathcal H_1\), and only then records
\((\mathrm{LVL}_p^1)\), equivalently \(m_p=1\), as a derived label.  The
legacy \(\mathsf P_{\rm sli}\) interface remains available only for later
roles that explicitly constitute it; details are in
Paper~II-F, Section~\ref{p2:app:selector-limitations}.  Raw ambient coordinate
actions, their omitted Wigner--Eckart data, absolute coordinate/protected
normalizations, and coordinate--coordinate commutator magnitudes are
completion-relative, and Paper~I-A's quoted representative values use
\(\mathsf C_{\rm PI}\in\mathfrak M_{\rm coord}\).  By contrast, on the
nondegenerate protected stratum \(\varrho_{\rm sp}>0\), current Paper~II-C
shows that the normalized protected even products are independent of both the
nonzero completion scale and the relative chiral phase and generate the
represented real algebra \(\Cl(0,3)\).  Finite Hurwitz/trimer
statements use \(\mathsf{Supp}_{2T}\), one role output of the exceptional
support atlas.  A labelled Hurwitz vertex, orientation, or neutral-trimer
representative is calculation notation inside a symmetry orbit, not a new
physical principle and not something made unique by \(m_p=1\).

Paper~II adds the G0 semantic root and three constitutional layers.  G0 types
actual occurrence and distinguishable readout without supplying their
witnesses.  First, GP1-TB makes the exact geometric operation closure compact.  In the
constituted state-affine exact-operation sector, the finite state/test span
provides a faithful finite-dimensional affine action and hence an automatic
compact Lie parent; Haar averaging unitarizes that representation.  A
faithful linear lift on the original Hilbert carrier is not automatic and
requires its own generated operation port.  For a separately constituted
comparison presentation, GP1-TB gives a canonical compact completion; compact
recurrence makes its forward isometry closure reversible
(Theorem~\ref{p2:thm:gp1-bounded-isometry-closure}).  Profiles generated
canonically from each comparison metric detect its geometric action kernel,
and metric-curry closure is the G0 certificate that realizes the full profile
family as physical separating B-tests.  GP2 then
forms the origin-admissible, candidate-internally exact-and-host-faithful model
class, and GP3 orders its complete candidates by role-faithful comparison before candidate
self-data are frozen.  The compact-Lie comparison lemmas and the
admissible-globalization operator give \(U(1)\) as the unique least connected
abelian parent-role skeleton and source-global \(SU(2)\), not \(SO(3)\), as
the unique least connected noncommutative parent-role skeleton.  Direct
candidatewise endpoint comparison shows that every GP3-irreducible complete
primitive realization has the corresponding endpoint skeleton, while distinct
unmarked carrier presentations over it remain unresolved.  The actual generated
multiplication is retained after that allocation.
Operational rank one and center-oddness are consequences or diagnostics,
whereas (FD)+(LR) remain optional primitive-comparison regularity conditions;
none is a selection input.
Second, it places the Paper-I-A finite structures into a kinematic dictionary whose entries are classified by logical status.  Third, it derives explicit operator representatives and finite spectral consequences.  The phrase \emph{integer Toeplitz-compressed
module family} denotes the integer Paper-I-A module subfamily together with its
declared coordinate-sector compression rule.  It does not mean that the
single momentum parent has Lie algebra
\(\mathfrak{su}(2)_L\oplus\mathfrak{su}(2)_R\); that is the Killing algebra of
the derived round metric.

Later papers may assign additional reciprocal, scalar, dynamical, quantum, or
field-theoretic meanings to the exported structures.  None of those later
assignments is part of Paper~II and none is used to select the momentum parent.





\label{p2:part:kin-core}
\label{p2:part:carrier-constitution}

This part contains the native constitutional derivation of the compact
momentum parent.  FRC-a makes the faithfully represented exact reversible
operation closure a compact Lie group on a finite ontic carrier.  GP1-TB
separately compactifies any constituted comparison presentation, and compact
recurrence derives reversal for its forward isometry closure.  Full metric
profiles give the canonical geometric effective quotient; their operational
B-test use is conditional on metric-curry closure.  GP2 makes each
prior-geometrically admissible paired model candidate-internally exact and
relation-exact; GP3 then performs irreducible/least whole-candidate comparison
before candidate-specific group data are frozen.  Compact-Lie comparison
followed by admissible globalization
gives \(S^1\cong U(1)\) for the connected abelian role and source-global
\(SU(2)\), rather than \(SO(3)\), for the
connected noncommutative role.  Lemma~
\ref{p2:lem:no-su2-u1-character} proves that the scalar role is neither a
character nor a canonically selected circle inside the transfer group.  Only
after the parent allocation is the Paper-I-A round/Hopf/\(\Spin(4)\) theorem
bank applied to the finite kinematic construction.  Paper~II exports the
B-side localization and scalar-readout obligations but does not realize them
as position and energy or orient the circle as causal time.

\section{Geometric primacy and the constitutional principles}
\label{paper2a:sec:organized-principles}

\subsection{Derivation of Kinematic Structure}
\label{p2:sec:derivation}

We derive the geometry of the typed directed-transfer role without
postulating a vector manifold or global group.  The derivation separates total
boundedness, prior-geometric origin, candidate-internal exactness, and
irreducible/least whole-candidate comparison.  The
invariant bounded momentum scale is discussed in
Section~\ref{p2:app:axiom-motivation}; the derivation itself begins with the
three geometric principles and the cumulative allocation order.

\subsubsection{Geometric principles and primary parent constitution}
\label{p2:subsec:constitution-clauses}

Ideas involving compact momentum space have appeared in lower-dimensional
gravity and curved momentum-space models~\cite{tHooft1996}.  Here compact
momentum is treated as a consequence of a dependency-separated construction
rather than as an assumed
kinematic background.


\begin{definition}[Finite presentation convention]
\label{p2:def:finite-presentation-convention}
A \emph{finitely presented constituted signature} is a tuple
\[
 \Sigma=(\mathrm{Ob},\mathrm{Op},\mathrm{Con},\mathrm{Rel},\mathrm{Test},\mathrm{Mark})
\]
given by finitely many generated constants together with finitely many object,
operation, construction, relation, test, and marking schemata over already
constituted content.  A family canonically indexed by an already constituted
parameter object counts as one generated schema, not as one signature item for
each parameter value.  Thus finite presentation concerns the generating
description and does not assert finite cardinality of all values produced by
that description.
\end{definition}

\begin{definition}[Structured role realization and retained presentation]
\label{p2:def:structured-role-realization}
A candidate realization of a finitely presented constituted signature
\(\Sigma\) is a structured model
\[
 \mathfrak M
 =\bigl(\mathcal H_{\mathfrak M},\mathcal A_{\mathfrak M},
          \mathcal C_{\mathfrak M},\mathcal R_{\mathfrak M},
          \mathcal T_{\mathfrak M},\mathcal M_{\mathfrak M}\bigr),
\]
where \(\mathcal H_{\mathfrak M}\) is the finite ontic carrier,
\(\mathcal A_{\mathfrak M}\subseteq\operatorname{End}(\mathcal H_{\mathfrak M})\)
is a finite-dimensional unital \(*\)-algebra containing the represented
operation envelope, \(S(\mathcal A_{\mathfrak M})\) is the full compact state
space fixed in Definition~\ref{p2:def:standard-finite-state-space},
\(\mathcal C_{\mathfrak M}\) is the typed constructor envelope,
and the remaining slots record the generated relations, tests, and markings
required in the candidate.  Write
\(\operatorname{Pres}_{\mathfrak M}(\mathfrak M)\) for this complete
candidate-internal presentation, including the type and provenance of every
item generated in \(\mathfrak M\).  This must not be confused with the
chronologically frozen inventory \(P_n\): before GP3, an internally retained
candidate item is not yet a constituted physical item common to the candidate
class.  Operation, constructor, and carrier retention are distinct typed
statements.
A mathematical module occurring as the parameter or image type of an
operation does not thereby become a carrier sector unless the constituted
host typing objectifies that image.  Conversely, an internal object image
cannot be represented merely as an operation without changing its sort.

A parent operation group may be continuous while its represented envelope is
finite-dimensional or finitely presented.  Exactness, host faithfulness, irreducibility, and leastness are
properties of the whole structured model, not of the carrier or operation
parent separately.  The following presentation guards are part of structured
realization:
\begin{enumerate}[label=(SR\arabic*),leftmargin=2.8em]
\item \label{p2:sr:no-extra-marked-datum}
No marked object, relation, operation, constructor, test, coefficient, scale,
basis choice, or selector is retained unless it is required by the previously
frozen inventory, intrinsic to the geometrically primary role condition, or
generated internally by the candidate's constituted role constructors.
\item A theorem consequence of
\(\operatorname{Pres}_{\mathfrak M}(\mathfrak M)\) remains
available for deduction but is not thereby promoted to a new retained
presentation item.
\end{enumerate}
At this initial stage \(\mathcal H_{\mathfrak M}\) denotes a finite-dimensional
complex carrier; no particular Hermitian form is primitive data.  GP1 supplies
the finite linear operation envelope, while GP2 preserves candidatewise the
generated composition law, action topology, and all operation distinctions.
In the state-affine exact-operation sector, the closed parent has a
continuous faithful complex-linear representation on the homogeneous
finite-state span, not automatically on \(\mathcal H_{\mathfrak M}\).
The invariant Hermitian form on that span is derived in
Theorem~\ref{p2:thm:frc-automatic-compact-lie-closure}; a Hilbert-carrier
lift remains a separate typed operation certificate.  A proof-witness inner
product used there is not retained as a new marked datum.
\end{definition}

\begin{definition}[Standard finite state space, effects, and evaluations]
\label{p2:def:standard-finite-state-space}
Every finite represented \(*\)-algebra in a structured realization is used in
its finite-dimensional \(C^*\)-algebra form.  Thus, for
\[
 \mathcal A\cong\bigoplus_{\lambda=1}^{r}M_{d_\lambda}(\CC),
\]
the notation already used above means the full state space
\[
 S(\mathcal A)
 :=\{\omega:\mathcal A\to\CC:
       \omega\text{ is linear},\quad
       \omega(a^*a)\ge0,\quad
       \omega(\mathbf1)=1\}.
\]
Its effect interval is
\[
 \operatorname{Eff}(\mathcal A)
 :=\{E\in\mathcal A_{\rm sa}:0\le E\le\mathbf1\},
\]
and each \(a=a^*\in\mathcal A\) has the canonical evaluation
\[
 \widehat a:S(\mathcal A)\longrightarrow\RR,
 \qquad \widehat a(\omega):=\omega(a).
\]
This convention fixes the admissible state role of \(\mathcal A\).  When
\(\mathcal A\) is generated in a physical test slot, its evaluation role is
generated with it.  The convention does not assert that every state is
prepared, every test is performed, or any outcome is registered.
\end{definition}

\begin{remark}[Scope of the complex state grammar]
\label{p2a:rem:complex-state-grammar-scope}
The complex-linear carrier in the structured-realization definition and the
full complex $C^*$ state space in
Definition~\ref{p2:def:standard-finite-state-space} are the domain of the
following reconstruction.  Within this domain, the density-matrix
evaluation is a theorem, rather than a separately adjustable probability
rule.  The argument does not compare that domain with real or quaternionic
carrier grammars or with general Jordan state spaces, and therefore does not
exclude those alternatives from GP1--GP3 alone.  Such an exclusion would
require a generated comparison or composition argument in the enlarged
candidate class.  This scope statement changes neither the current
structured candidate class nor its canonical evaluation.
\end{remark}

\begin{theorem}[Canonical finite state--test evaluation]
\label{p2:thm:canonical-finite-state-test-evaluation}
Let \(\mathcal A=\bigoplus_\lambda M_{d_\lambda}(\CC)\) be a finite-dimensional
unital \(C^*\)-algebra.  Then:
\begin{enumerate}[label=(\roman*),leftmargin=2.2em]
\item every \(\omega\in S(\mathcal A)\) has a unique block density
\[
 \rho=\bigoplus_\lambda\rho_\lambda,
 \qquad \rho_\lambda\ge0,
 \qquad \sum_\lambda\Tr(\rho_\lambda)=1,
\]
such that
\[
 \omega(a)=\Tr_{\mathcal A}(\rho a)
 :=\sum_\lambda\Tr(\rho_\lambda a_\lambda);
\]
\item \(S(\mathcal A)\) is compact and convex;
\item the evaluation map
\[
 \Gamma:\mathcal A_{\rm sa}\longrightarrow
 \operatorname{Aff}(S(\mathcal A),\RR),
 \qquad \Gamma(a)=\widehat a,
\]
is a real-linear unital order isomorphism onto the full finite-dimensional
affine-function space, and is order-unit isometric:
\[
 \|a\|=\sup_{\omega\in S(\mathcal A)}|\omega(a)|;
\]
\item \(E\in\operatorname{Eff}(\mathcal A)\) exactly when
\(0\le\widehat E(\omega)\le1\) for every \(\omega\in S(\mathcal A)\);
\item for every unital \(*\)-homomorphism
\(\phi:\mathcal A\to\mathcal B\), evaluation is natural:
\[
 \widehat{\phi(a)}(\eta)=\widehat a(\eta\circ\phi)
 \qquad(\eta\in S(\mathcal B)).
\]
\end{enumerate}
Consequently, once a structured role contains its finite algebra and full
state space, the state--test evaluation is canonical generated content rather
than an independent linear-evaluation postulate.
\end{theorem}

\begin{proof}
The trace pairing on each matrix summand is nondegenerate, so every linear
functional has a unique representing block matrix \(\rho\).  Positivity on
rank-one projections is equivalent to \(\rho_\lambda\ge0\) in every block, and
normalization at \(\mathbf1\) is exactly
\(\sum_\lambda\Tr(\rho_\lambda)=1\).  Conversely every such block density
defines a positive normalized functional.  These densities form a closed and
bounded subset of a finite-dimensional real vector space, hence a compact
convex set.

For \(a=a^*\), the order inequality
\(-\|a\|\mathbf1\le a\le\|a\|\mathbf1\) gives
\(|\omega(a)|\le\|a\|\) for every state.  Choose a unit eigenvector in a
block on which an eigenvalue of \(a\) has modulus \(\|a\|\); its vector state
attains equality.  This proves the norm formula and injectivity.  To prove
surjectivity, let \(f:S(\mathcal A)\to\RR\) be affine.  Every nonzero positive
functional \(\varphi\in\mathcal A_{\rm sa}^*\) has the unique form
\(\varphi=\varphi(\mathbf1)\omega\) with \(\omega\in S(\mathcal A)\).
Therefore
\[
 \widetilde f(\varphi):=\varphi(\mathbf1)
 f\!\left(\frac{\varphi}{\varphi(\mathbf1)}\right),
 \qquad \widetilde f(0):=0,
\]
is additive and positively homogeneous on the positive dual cone.  Extending
by differences gives a well-defined real-linear functional on
\(\mathcal A_{\rm sa}^*\).  Finite-dimensional biduality supplies a unique
\(a=a^*\in\mathcal A\) with \(\widetilde f(\varphi)=\varphi(a)\), and hence
\(f=\widehat a\) on states.  The same vector-state argument shows that
\(0\le E\le\mathbf1\) exactly when every state evaluates \(E\) in
\([0,1]\).  Finally,
\((\eta\circ\phi)(a)=\eta(\phi(a))\) proves naturality.

In a structured realization in which \(\mathcal A\) is already generated in
the physical test slot, these maps are precisely the canonical partial
evaluations closed by GP2.  The self-adjoint, projection, and effect roles and
their normalized state-evaluation rule are then generated physical role
content, not separately assigned semantics.  This does not prepare a named
state, perform a named test, produce an outcome, or identify a finite sample
frequency: each such token remains an ordinary G0 occurrence.
\end{proof}

\begin{theorem}[State-dual reconstruction of finite test operations]
\label{p2:thm:state-dual-affine-operation-reconstruction}
Let \(\mathcal A\) and \(\mathcal B\) be finite-dimensional unital
\(C^*\)-algebras.  Every affine map
\[
 F:S(\mathcal A)\longrightarrow S(\mathcal B)
\]
has a unique complex-linear unital positive dual map
\[
 F^\sharp:\mathcal B\longrightarrow\mathcal A
\]
characterized by
\[
 \omega\!\left(F^\sharp(b)\right)=F(\omega)(b),
 \qquad
 \omega\in S(\mathcal A),\quad b\in\mathcal B.
\]
Conversely, every unital positive map
\(T:\mathcal B\to\mathcal A\) defines the affine state transformation
\(F_T(\omega)=\omega\circ T\), and these assignments are mutually inverse.
If \(F\) is an affine isomorphism, then \(F^\sharp\) is an order-unit
isomorphism and a Jordan \(*\)-isomorphism:
\[
 F^\sharp(b\circ c)=F^\sharp(b)\circ F^\sharp(c),
 \qquad b\circ c:=\tfrac12(bc+cb).
\]
Thus an occurring affine state transformation canonically
generates its Heisenberg test operation; no second operation-selection datum
is required.

The state space and its affine structure alone recover the ordered
self-adjoint test space, effect interval, projection lattice, spectral
squaring, and Jordan product.  They do not recover an orientation of the
associative product: \(\mathcal A\) and \(\mathcal A^{\rm op}\) have the same
state/effect/Jordan data.  Nor do they supply complete positivity, a marked
tensor-composition law, or an actual time flow.  Those require the
corresponding generated sequential/composite structure or occurrence
antecedents.
\end{theorem}

\begin{proof}
For \(b=b^*\in\mathcal B\), the function
\(\omega\mapsto F(\omega)(b)\) is real affine on \(S(\mathcal A)\).
Surjectivity of \(\Gamma\) in
Theorem~\ref{p2:thm:canonical-finite-state-test-evaluation} gives a unique
\(F^\sharp(b)\in\mathcal A_{\rm sa}\) representing it.  Real linearity,
unitality, and positivity follow by evaluating against all states and using
state separation; complexification gives the unique complex-linear positive
map.  Conversely, composition with a unital positive map preserves positivity
and normalization and is affine.  The two constructions are inverse by their
defining evaluation identity and state separation.  Applying the same
argument to \(F^{-1}\) proves the order-isomorphism statement.

For completeness, the extreme points of the effect interval
\([0,\mathbf1]\) of a finite-dimensional \(C^*\)-algebra are exactly its
projections: a nonprojection has a spectral value strictly between zero and
one and can be perturbed in the associated spectral subspace, whereas
positivity on a projection and its complement proves extremality.  An
order-unit isomorphism therefore maps projections bijectively to projections.
It also preserves the order relation \(p+q\le\mathbf1\); for projections this
is equivalent to orthogonality.  If
\(b=\sum_j\lambda_jp_j\) is a spectral resolution, then
\[
 F^\sharp(b)^2
 =\sum_j\lambda_j^2F^\sharp(p_j)
 =F^\sharp(b^2).
\]
Polarization gives preservation of \(b\circ c\), and complexification gives
the Jordan \(*\)-statement.  Reversing every product leaves these data
unchanged, while amplification, tensor factorization, and a distinguished
one-parameter family have not entered the proof.  This establishes both the
additional reconstruction and its scope boundary.
\end{proof}

\begin{corollary}[Finite state--test dual dynamics is one co-realized process]
\label{p2a:cor:finite-state-test-dual-dynamics}
Let \(M\) be a monoid and let
\[
 F_m:S(\mathcal A)\longrightarrow S(\mathcal A),
 \qquad F_{mn}=F_m\circ F_n,
\]
be an affine operation family that occurs or is generated in a structured
finite role together with its family law.  Then the standard test slot of
Definition~\ref{p2:def:standard-finite-state-space} generates the unique dual
family
\[
 F_m^\sharp:\mathcal A\longrightarrow\mathcal A,
 \qquad
 \omega(F_m^\sharp a)=F_m(\omega)(a),
\]
and
\[
 F_{mn}^\sharp=F_n^\sharp\circ F_m^\sharp.
\]
Thus the state and test descriptions are the covariant and contravariant
faces of one operation occurrence.  No second operation or dynamics
occurrence is required for the dual description.

For a one-parameter additive family, commutativity of the parameter law makes
the dual family a semigroup on the same parameter domain.  If the original
domain is \(t\ge0\), this conclusion supplies \(F_t^\sharp\) for \(t\ge0\)
only: it supplies neither \(F_{-t}\) nor physical time reversal.  Complete
positivity, a tensor/composite interpretation, a selected generator, an
initial state, an actual member history, a time scale, and detailed balance
remain available only when their corresponding typed structure is generated
or occurs.
\end{corollary}

\begin{proof}
Apply Theorem~\ref{p2:thm:state-dual-affine-operation-reconstruction} to each
\(F_m\).  For every state \(\omega\) and test \(a\),
\[
 \omega(F_{mn}^\sharp a)
 =F_m(F_n(\omega))(a)
 =\omega\!\left(F_n^\sharp(F_m^\sharp a)\right),
\]
with the displayed order adjusted exactly by the convention
\(F_{mn}=F_m\circ F_n\); state separation gives the asserted opposite-law
identity.  The maps are unique generated outputs of the canonical evaluation
constructor, so retaining a second occurrence would duplicate the same
operation in its dual test presentation.  An additive one-parameter monoid is
commutative, whereas inversion, token selection, calibration, and the other
listed structures were not used.
\end{proof}

\begin{corollary}[Generated finite state--test role needs no semantic witness]
\label{p2:cor:generated-state-test-role-no-duplicate-witness}
If a G0-occurring structured role contains a finite algebra \(\mathcal A\) in
its generated physical test slot, then
\[
 \bigl(\mathcal A,S(\mathcal A),\operatorname{Eff}(\mathcal A),
       (\omega,E)\mapsto\omega(E)\bigr)
\]
is generated role content.  The terms state, real test, sharp event, effect,
and normalized outcome weight introduce no further role-level occurrence
assumption.  A named preparation, performed test, registered outcome,
empirical frequency, or selected affine process is not thereby actualized and
still satisfies G0 at the token level.
\end{corollary}

\begin{proof}
The algebra is already test-typed by hypothesis.  Its full positive-normalized
dual, effect interval, and evaluation are canonical and functorial by the two
preceding theorems, so GP2 no-insertion and no-omission retain them without a
new selector.  The final statement is the distinction between the generated
role of admissible tokens and the contingent occurrence of one such token.
\end{proof}

\begin{definition}[Parent geometric object and induced operation occurrences]
\label{p2:def:parent-geometric-object-occurrence}
Let \(R\) be a constituted smooth host role with frozen parent geometry
\(M_R\).  A section of a natural or associated finite-rank bundle whose bundle
structure is already generated by that frozen geometry is a \emph{parent
geometric object occurrence} of \(R\).  This is a presentation-side object
sort and does not, by itself, place the full section space in the finite ontic
carrier.  If such a section canonically induces an action on another host
object, that action remains a distinct operation-typed output and does not
retype the section.  When the frozen grammar licenses the action constructor
and declares its operation target, the section occurrence and the induced
operation are co-realized by the same rooted witness; the type distinction
does not require a duplicate G0 occurrence.  When that constructor or target
is absent, mathematical inducibility alone does not promote the action to a
physical operation.

In particular, for the tangent bundle write
\[
 \Gamma(TM_R)_{\rm sec,obj}
\]
for the vector-field section/object occurrence.  Its canonical derivation
occurrence is
\[
 \delta:\Gamma(TM_R)_{\rm sec,obj}
 \longrightarrow\operatorname{Der}(C^\infty(M_R))_{\rm op},
 \qquad \delta_X(f):=X(f).
\]
The two typed nodes are linked by \(\delta\).  If \(\delta\) is licensed as
displayed, they are co-realized outputs of one root rather than independent
occurrences.
Only a GP1-finite generated subimage of a parent object sort can create a
finite-carrier obligation through GP2; the ambient infinite-dimensional
section space is not thereby objectified as a finite carrier.
\end{definition}

\begin{mdframed}[style=postulatebox,nobreak=false]
\begin{principle}[Geometric Primacy and Minimal Constitution]
\label{p2:prin:geometric-primacy-minimal-constitution}
Geometric Primacy has a semantic root, G0, followed by the constitutional
closure GP1--GP3.  Every primitive physical role begins
with an actual geometric occurrence specified by a strong role predicate
\(\operatorname{Geo}_R\), rather than with a named manifold or group.  This
predicate constitutes the comparison-bearing role, its local reapplication
constructor, and its intrinsic role diagrams.  It does not freeze the complete
self-description of a candidate realization.

Throughout this paper, G0 and GP1--GP3 are the constitutional axioms of the
framework.  They are not theorem obligations and are not claimed to follow
from weaker mathematical premises.  Theorems below derive consequences inside
the theory they define and identify exactly which axiom performs each step.

\medskip
\noindent\textbf{G0: geometric actualization of distinguishability.}
Let \(P_n\) be the constituted presentation frozen before the claimed
occurrence.  For a role-typed datum \(q\), let \(\Delta_q\) be the part of its
comparison-bearing signature not already present in \(P_n\), with every
dependency and interface fixed independently of the eventual token and
result.  Define the stage at which the occurrence is tested by
\[
 P_n[q]:=
 \begin{cases}
   P_n,
     & \Delta_q=\varnothing,\\[2pt]
   \operatorname{Fr}_{\mathrm{pres}}
   \!\left(\operatorname{Cl}^{P_n}_{\mathrm{GP1:3}}
   (P_n\cup\Delta_q)\right),
     & \Delta_q\ne\varnothing,
 \end{cases}
\]
where relative closure preserves the complete frozen content of \(P_n\).
The physical-occurrence predicate is
\[
\boxed{
 \PhysOcc_{P_n}(q)
 \ \Longleftrightarrow\
 \exists\mathcal E\,
 \bigl[
   \Actual_{\mathcal E}(q)
   \wedge
   \GeomLic_{P_n\rightsquigarrow P_n[q]}(q;\mathcal E)
   \wedge
   \InvDist_{P_n[q]}(q;\mathcal E)
 \bigr].}
\tag{\ensuremath{\mathrm{G0}}}
\]
Here \(\Actual_{\mathcal E}(q)\) is the primitive contingent-token predicate;
it asserts that this admissible token occurs and carries no structural or
selection content.  The predicate \(\GeomLic\) requires the role type of
\(q\), its comparison operations, and every retained source, preparation,
test, sink, field, projector, grading, or channel mark to be intrinsic to,
generated by, or explicitly constituted relative to the indicated stage and
faithfully represented in its declared sort.

The predicate \(\InvDist\) is claim-relative: a licensed invariant comparison
or test must distinguish the particular alternatives asserted by the claim,
modulo unmarked presentation isomorphisms.  It does not require a family that
separates every mathematically available carrier state unless that stronger
tomographic role is separately constituted.  Thus G0 does not import the
optional carrier-resolving event extension of Paper~IV.

In a reported experiment
\[
       O_{\mathcal E}=\ell_{\mathcal E}(O_{\rm geom}),
\]
where the comparison, apparatus, boundary provenance, calibration standard,
and readout map \(\ell_{\mathcal E}\) are fixed before the reported value is
known.  Actuality may select a contingent token of an admissible role type; it
may not use the token or its result to select, modify, or enlarge upstream
geometry.  An already constituted standard and map give an instantiational
calibration.  A new standard, comparison law, or interface is constitutive and
must enter through \(\Delta_q\); neither case creates or changes the carrier.
G0 certifies what counts as physical actualization.  It does not derive which
admissible state, process, history, apparatus, outcome, coupling, or
calibration actually occurs.  The schema is universal in role-type scope but
one-way in actuality.  Writing \(\operatorname{Adm}_{P_n}(q)\) only for
membership in the admissible role family at the frozen stage,
\[
 \boxed{\begin{gathered}
  \text{G0 applies to every role-typed physical occurrence},\\[2pt]
  \operatorname{Adm}_{P_n}(q)\not\Longrightarrow \PhysOcc_{P_n}(q).
 \end{gathered}}
\]
Thus ``whatever physically occurs is G0-licensed'' is a scope statement, not
an assertion that every admissible or mathematically available token occurs.

Every further physical datum, relation, operation, constructor, test, or
distinction must be intrinsic to, or generated from, this geometry.  GP1--GP3
therefore select the geometry and
its operation content jointly: GP1 restricts complete candidates, GP2 makes
each candidate internally exact and faithful, and GP3 selects the least whole
candidate before its complete self-content is frozen.  Accordingly:
\[
 \begin{gathered}
  \boxed{\mathrm{G0}:\ \text{Actualized/Distinguishable}}\\[4pt]
  \boxed{\begin{gathered}
   \mathrm{GP1}:\ \text{Finite/Bounded}\\
   \mathrm{GP2}:\ \text{Exact/Faithful}\\
   \mathrm{GP3}:\ \text{Irreducible/Least}.
  \end{gathered}}
 \end{gathered}
\]
``Irreducible'' means that no strict role-faithful comparison remains, while
``least'' means that the realization lies below every admissible realization
in the structural preorder.  Neither term refers to a numerical complexity
functional.  A realization with fewer carrier states is not preferred if it
fails to realize the same role geometry or omits previously frozen structure;
inside every candidate, GP2 already forbids omission of that candidate's own
generated structure.  If a needed datum is neither intrinsic nor generated,
it remains an explicit constitutive
residue rather than being supplied by an unmarked choice.  In particular,
``nontrivial'', ``connected'', and ``abelian/noncommutative'' are not free
labels in the primitive existence assertion used below.  The strong
\(\operatorname{Geo}_R\) predicate instead constitutes effectiveness, connected
local generation, and the order-swap diagram.  GP1--GP2 realize that content as
a compact exact parent, after which nontriviality, connectedness, and the
abelian/noncommutative status are theorem classifications of the realized
constituted data.  GP3 compares only candidates realizing that same role and
selects the irreducible/least whole realization.  Only after that selection is
its full operation law frozen.  The provenance chain is therefore
\([\mathrm{CONST}]\to[\mathrm{REAL}]\to[\mathrm{CLASS}]\to[\mathrm{MIN}]\),
not constitution by a neutral downstream rule.
\end{principle}
\end{mdframed}

\begin{definition}[Instantiational and constitutive G0 occurrences]
\label{p2:def:g0-occurrence-types}
A G0 occurrence is \emph{instantiational} relative to \(P_n\) when it selects
an actual token of a role type already frozen in \(P_n\), so that
\(P_n[q]=P_n\).  It is \emph{constitutive} when its result-independent
signature contains a new comparison-bearing role, relation, operation, or
interface.  Only the latter opens a dependency frontier \(\Delta_q\) and
requires relative GP1--GP3 closure.  Both preserve the complete frozen content
of \(P_n\); this is cumulative, frozen-base preservation and is not an
unqualified model-theoretic conservativity claim.
\end{definition}

\begin{corollary}[Common geometric origin of G0 witnesses]
\label{p2:cor:g0-witness-normal-form}
Every physical occurrence, preparation, test, process, trajectory, cross-role
bridge, readout, and calibration used in the series is a role-specific
instance of G0: an actual token of a geometrically licensed, claim-relative
invariant distinction with result-independent provenance.  A role-specific
tag states what is instantiated; the tag alone neither asserts occurrence nor
supplies a separate actualization principle.

For a frozen occurring composite, every nonredundant generated constituent is
co-realized in its native sort by
Theorem~\ref{p2:thm:generated-constituent-corealization}; no further G0
certificate is required merely for that constituent.  A different admissible
state, apparatus, process, or diagram does not thereby occur.  Finite
structural scalar readouts specialize, by
Theorem~\ref{p2:thm:loop-boundary-readout-exhaustion}, to closed-loop or
boundary-conditioned provenance.  This two-branch theorem refines G0 for
finite scalar readouts and is not its general definition.
\end{corollary}

\begin{remark}[Momentum has no independent role-level occurrence witness]
\label{p2:rem:g0-mom-open}
The primitive order-sensitive role is an actual comparison/reapplication
operation occurrence under G0.  GP1--GP3 select its relation-exact operation
parent, and Paper~II-B generates the finite carrier from the frozen
round-host constructor presentation.  Therefore the parent, its momentum
role, and the generated carrier are descendants of the same occurring root.
Theorems~\ref{p2:thm:generated-role-corealization} and
\ref{p2:cor:operational-momentum-role-occurrence} show that no independent
``G0-MOM'' certificate exists or is needed.  That expression may be used only
as a descriptive tag for the original G0 specialization to the occurring
transfer operations.  A particular momentum measurement, numerical value,
unit standard, or scale calibration remains a token-level G0 occurrence.
\end{remark}

\begin{definition}[Intrinsic locally finite atomic generation grammar]
\label{p2:def:intrinsic-generation-grammar}
Let \(q\) be a typed object/carrier or retained operation-presentation slot of
a frozen signature.  An \emph{intrinsic locally finite atomic generation
grammar} \(\Gamma_q\) consists of a prior generated atomic constructor family,
the typed constructor/composition/decomposition schemata actually constituted
for that role, and the minimal construction
length
\[
 \ell_q(Y):=\min\{r:\ Y\text{ is generated using }r
 \text{ atomic constructor occurrences}\}.
\]
For every finite \(N\), the sublevel
\[
 \Gamma_q^{\le N}:=\{Y:\ell_q(Y)\le N\}
\]
must have finite type-support/finitely presented generated content appropriate
to the slot.  The atomic family, rules, and \(\ell_q\) must be invariant under
every unmarked automorphism of the frozen signature.  The grammar is
\emph{intrinsic} only when these data are generated before the present
selection and are independent of the candidate, bound, Casimir, scale,
coordinate choice, or other datum whose selection they could otherwise
certify.  This definition supplies terminology only; it asserts neither that
every role has such a grammar nor that a physical realization is a sublevel of
one.  In particular, a tensor product or decomposition operation belongs to
physical generation only when its constructor schema belongs to the frozen
grammar; mathematical availability alone is insufficient.
\end{definition}

\begin{definition}[Primitive constructor atom]
\label{p2:def:primitive-constructor-atom}
Let a frozen host grammar carry a generated commutative construction monoid
\(M_\Sigma\) of constructor types.  A nontrivial type \(a\in M_\Sigma\) is a
\emph{primitive constructor} exactly when it is an atom of that monoid:
\[
 a=b+c,\qquad b,c\in M_\Sigma
 \quad\Longrightarrow\quad b=0\ \text{or}\ c=0,
\]
with the analogous indecomposability statement for a nonadditive composition
grammar.  Thus ``primitive'' is a structural term of the frozen construction
grammar, not a synonym for smallest dimension, lowest Casimir, or a preferred
representation chosen after the fact.  A decomposable generated type may occur
as a constructor output but cannot be relabelled primitive without adding an
extra constitutive mark.
\end{definition}

\begin{mdframed}[style=postulatebox,nobreak=false]
\begin{postulate}[GP1: finite and bounded realization]
\label{p2:post:GP1}
At every realized stage and region GP1 imposes two role-typed capacity
conditions on the complete structured realization.
\begin{enumerate}[label=(\roman*),leftmargin=2em]
\item \emph{Finite realized ontic/compositional capacity.}  Every realized
ontic item and represented exact-operation envelope is finite, and the total
retained composition/occupation envelope of the stage is finite.  In
particular
\[
 0<\dim_{\CC}\mathcal H_{R,s}(\mathcal D)<\infty,
 \qquad
 \mathcal A_{R,s}(\mathcal D)\subseteq
 \operatorname{End}(\mathcal H_{R,s}(\mathcal D)).
\]
A continuous compact operation group may act through this finite represented
envelope; GP1 does not require finitely many operation values.  The historical
names (FRC-a) and (FRC-c) refer respectively to the finite ontic/represented
operation instance and the finite retained composition/occupation instance of
this one condition.  The latter is load-bearing: finite one-particle dimension
does not by itself close
\(\bigoplus_{n\ge0}\operatorname{Sym}^n\mathcal K\).

\item \emph{Totally bounded constituted families.}  Every constituted
comparison, scalar-readout, real-valued operation-parameter, and complete
timed-rate family has totally bounded image in its already generated
comparison/uniform structure.  The historical name (FRC-b) refers to this
condition.  The word \emph{family} is load-bearing: a fixed finite carrier or
compact endpoint group may coexist with the unbounded families
\(H_n=nH\) or \(\Phi_t^{(n)}(z)=e^{int}z\), which GP1 excludes only when the
whole family is the constituted physical datum.
\end{enumerate}
The bounds are stagewise and role-typed; no stage-independent numerical upper
bound is asserted.  Smooth manifolds, continuum function spaces, field
algebras, and infinite-refinement limits may occur as parent, symbolic,
filtered, fixed-band, or asymptotic presentations without being infinite
physical stages.  GP1 selects no dimension, level, radius, matrix basis,
Hamiltonian, grading, operation coefficient, or numerical ceiling.
\end{postulate}
\end{mdframed}

\begin{remark}[GP1 obstruction is not a successor constructor]
\label{p2a:rem:gp1-obstruction-no-successor-constructor}
If a claimed realized continuation, complete candidate family, or declared
physical member violates GP1 on its stated domain, that realization is
inadmissible on that domain.  This obstruction does not itself constitute a
retained process output, successor parent, boundary member, action, law,
Hamiltonian, radius, endpoint geometry, or history.  Any successor
realization still requires the already existing GP2 generated-output/owner
conditions and, for dynamics, the corresponding Paper~III-A origin and
continuation certificate.  GP3 may compare already generated admissible
candidates but does not manufacture a missing successor candidate.
\end{remark}

\begin{definition}[Registered generated-boundary hypothesis]
\label{p2:def:generated-boundary-hypothesis}
Suppose a frozen role slot \(q\) carries an intrinsic grammar \(\Gamma_q\) in
the sense of Definition~\ref{p2:def:intrinsic-generation-grammar}.  The
registered hypothesis
\[
 \mathsf P_{\rm gbd}(q,N)
\]
asserts that the realized generated presentation-support of that slot lies in
the intrinsic sublevel \(\Gamma_q^{\le N}\).  It is a boundary declaration,
not a clause of GP1, and it supplies no numerical minimization.  When the
boundary itself is generated or explicitly assumed, it becomes part of the
host signature and GP2 exactness applies inside that declared scope.  GP3 may
then compare the resulting structured realizations only through role-faithful
structural comparisons.  A theorem that derives a specific
intrinsic grammar or length does not thereby derive \(\mathsf P_{\rm gbd}\).
\end{definition}

\begin{remark}[Capacity versus generated-boundary semantics]
\label{p2:rem:frcd-boundary-semantics}
The former name ``FRC-d'' is retired.  Its capacity part was redundant once a
finite retained presentation with finite item lengths was fixed; its genuine
content was the assertion that the physical boundary is an intrinsic generated
sublevel rather than an arbitrary finite subset or a candidate-specific
Casimir/valuation cut.  That content is now explicit in the registered
hypothesis \(\mathsf P_{\rm gbd}\), outside GP1.  This is a status change, not
merely a rename: every consumer must cite the registered boundary hypothesis
or a theorem that independently discharges it.
\end{remark}

\begin{corollary}[GP1-TB: comparison specialization of FRC-b]
\label{p2:cor:gp1-tb}
Let \(X=X_{R,s}(\mathcal D)\) be a constituted physical comparison carrier,
let
\[
 \iota:X\longrightarrow S\bigl(\mathcal A_{R,s}(\mathcal D)\bigr)
\]
be its faithful realized-state map, and suppose its comparison \(d\) has a
modulus \(\omega\), with \(\omega(r)\to0\) as \(r\to0\), such that
\[
 d(x,y)\le
 \omega\!\left(\|\iota(x)-\iota(y)\|_*\right).
\]
Then \((X,d)\) is totally bounded.  In particular,
\[
 N_X(\varepsilon)
 :=\sup\{|E|:E\subseteq X\text{ is }\varepsilon\text{-separated}\}
 <\infty
 \qquad(\varepsilon>0).
\]
The same conclusion holds for a comparison generated by a uniformly
norm-bounded test family
\(d(x,y)=\sup_{a\in\mathcal T}|\iota(x)(a)-\iota(y)(a)|\).
The corollary asserts neither finite cardinality nor a positive minimum
separation, polynomial packing exponent, finite parent dimension, manifold
structure, or numerical propagation ceiling.
\end{corollary}

\begin{proof}
The state space of the finite-dimensional algebra
\(\mathcal A_{R,s}(\mathcal D)\) is compact in any finite-dimensional state
norm and hence totally bounded.  The displayed modulus makes the identity
from the realized state image to \((X,d)\) uniformly continuous.  A uniformly
continuous image of a totally bounded space is totally bounded.  In the
test-family form one may take
\(\omega(r)=\bigl(\sup_{a\in\mathcal T}\|a\|\bigr)r\).
No bound on an independently varying scalar or timed-rate family is inferred;
those bounds are exactly FRC-b.
\end{proof}

\begin{theorem}[Affine state-span criterion for automatic compact Lie closure]
\label{p2:thm:frc-automatic-compact-lie-closure}
At a fixed role, region and stage, let \(X\) be the GP1-totally-bounded
constituted comparison carrier and let
\(G=\overline{S_0(X)}\leq\operatorname{Isom}(\widehat X)\)
be its compact reversible exact-operation closure from
Theorem~\ref{p2:thm:gp1-bounded-isometry-closure}.
Assume the frozen finite-state/test presentation supplies a faithful
realization \(j:\widehat X\hookrightarrow S(\mathcal A)\), where
\(\mathcal A\subseteq\operatorname{End}(\mathcal H)\) is a
finite-dimensional unital \(C^*\)-algebra, and that \(j\) is continuous.
Write \(V=\operatorname{aff}_{\mathbb R}j(\widehat X)\),
\(r=\dim_{\mathbb R}V\), and assume each forward exact operation
\(s\in S_0(X)\) has a \emph{state-affine, role-faithful} action on this
finite state span: there is an affine automorphism \(A_s:V\to V\) with
\(A_s j(x)=j(sx)\), and the represented action preserves the generated
operation composition law.  (This is automatic for a constituted
state-affine exact-operation/test-dual port; it does not follow from
being an arbitrary metric isometry.)

Then the action extends uniquely to a continuous faithful homomorphism
\[
 \rho_{\rm aff}:G\longrightarrow
 \operatorname{GL}(\mathbb R^{r+1})
 \subset\operatorname{GL}(\mathbb C^{r+1}),
 \qquad r+1\leq\dim_{\mathbb R}\mathcal A_{\rm sa}
 \leq(\dim_{\mathbb C}\mathcal H)^2,
\]
using homogeneous affine coordinates.  Its image is compact, hence
\(G\) is a compact Lie group.  Haar averaging derives a positive
\(G\)-invariant form on this finite faithful representation and a
unitary complex presentation of dimension at most
\((\dim_{\mathbb C}\mathcal H)^2\).

A unitary action on the \emph{original} carrier \(\mathcal H\) follows
by Haar averaging only if a continuous faithful complex-linear
\emph{carrier lift}
\(\rho_{\mathcal H}:G\to\operatorname{GL}(\mathcal H)\)
was independently supplied by the frozen exact-operation grammar.
Finite ontic capacity and faithful affine action on state/test data do
\emph{not} imply such a lift.  No chosen Hermitian form is frozen by the
averaging construction, and the Lie algebra consists of the actual
one-parameter subgroups of the closed faithful image, not arbitrary
logarithms of individual operations.  The theorem does not make a
general completely positive semigroup into a group.
\end{theorem}

\begin{proof}
Theorem~\ref{p2:thm:gp1-bounded-isometry-closure} supplies compactness
and the group structure of \(G\).  A finite affinely independent frame
\(j(x_0),\ldots,j(x_r)\) of \(V\) uniquely determines each affine
map \(A_s\) by its values on that frame.  If a net \(s_i\to g\)
uniformly on \(\widehat X\), then continuity of \(j\) on the compact
completion implies
\(j(s_i x_k)\to j(gx_k)\) at every frame point.  Thus
the coefficients of \(A_{s_i}\) converge to those of a unique affine
map \(A_g\), and for every \(x\in\widehat X\),
\(A_gj(x)=j(gx)\).  Since \(g^{-1}\in G\), the corresponding map
\(A_{g^{-1}}\) is its inverse.  Composition, continuity and
faithfulness follow respectively from uniqueness on the frame, the
frame-coordinate formula, and injectivity of \(j\).

The homogeneous lift of \(A_g(v)=L_gv+b_g\) is
\(\rho_{\rm aff}(g)=\left(\begin{smallmatrix}L_g&b_g\\0&1\end{smallmatrix}\right)\),
a continuous faithful group homomorphism.  The real linear span of
the state vectors in \(\mathcal A_{\rm sa}^*\) has dimension
\(r+1\), because each state has value one on the unit.  Hence the
stated dimension bound.  Compactness makes the image closed in the
finite matrix group, so the closed subgroup theorem makes \(G\) Lie.
For any proof-witness positive Hermitian form \(h_0\) on the complexified
homogeneous space, normalized Haar measure gives
\[
 h_G(u,v)=\int_G h_0(\rho_{\rm aff}(g)u,
                    \rho_{\rm aff}(g)v)\,d\mu(g).
\]
It is positive and invariant.  The same argument applies to a
\emph{separately given} \(\rho_{\mathcal H}\), but it does not
construct one.  The proof never retains \(h_0\) as an output mark.
\end{proof}

\begin{proposition}[Sharp carrier-lift obstruction from a qubit state action]
\label{p2a:prop:v281-qubit-affine-lift-obstruction}
Finite state-test faithfulness does not imply a faithful carrier-linear
lift of the exact-operation group of the same Hilbert dimension.  In
particular, on \(\mathcal A=M_2(\mathbb C)\), with \(\mathcal H=\mathbb C^2\),
the exact conjugation operations on density matrices have the compact
faithful affine group
\(\operatorname{PU}(2)\cong SO(3)\), acting on the Bloch ball, but there
is no faithful continuous complex two-dimensional linear representation
of \(SO(3)\).  Its faithful affine/test-space representation is real
four-dimensional (or three-dimensional on traceless test directions).
The linear spinor representation lives on the double cover \(SU(2)\)
and is not a representation of the quotient operation group.
\end{proposition}
\begin{proof}
Every qubit state is
\(\rho=\frac12(I+\mathbf r\cdot\boldsymbol\sigma)\),
\(\mathbf r\in\overline B^3\).  Conjugation by \(U\in SU(2)\)
acts by a rotation of \(\mathbf r\), and its exact kernel is
\(\{I,-I\}\), so the effective operation group is \(SO(3)\).
The standard complex irreducible representations of \(SU(2)\) of
dimension at most two are trivial or the spin-half representation; the
latter sends the central element \(-I\) to \(-I_2\) and therefore
cannot descend to \(SO(3)\).  Every complex two-dimensional
representation of \(SO(3)\) is consequently trivial and not faithful.
The action \(\mathbf r\mapsto R\mathbf r\) is faithful on the
traceless Hermitian test space, and its homogeneous affine extension is
faithful on \(\mathbb R^4\).  The finite state/test operation is exact
and affine; only the unlicensed two-dimensional quotient lift fails.
\end{proof}

\begin{remark}[Firewall between isometry compactness and affine Lie regularity]
\label{p2a:rem:v281-isometry-affine-lie-firewall}
GP1-TB and compact recurrence alone make \(G\) a compact \emph{topological}
group; compactness does not generally imply that it is Lie.  The finite
affine-state action is the additional \emph{typed operation port} needed for the finite-dimensional Lie representation.  Merely embedding
a compact metric comparison carrier into the state space of a finite
algebra does not make all its metric isometries affine.  Any downstream
Lie assertion must use the generated affine operation port, a separate
faithful finite-dimensional linear representation, or an independent
Lie-regularity certificate.  The latter may be mathematical and need
not select new physically marked states or bases.
\end{remark}

\section{Scalar and invariant readout exhaustion}
\label{paper2a:sec:organized-readouts}

\subsection{Exhaustion of finite structural scalar readouts}
\label{p2:sec:finite-structural-readout-exhaustion}

The endpoint classifications \(U(1)\cong S^1\) and
\(SU(2)\cong S^3\) exhaust the primitive scalar and order-sensitive parent
operation roles.  The following results perform the analogous exhaustion for
finite scalar \emph{readout} roles.  They classify constructor-relative normal
forms, not all set-theoretic functions and not which admissible readout occurs.

\begin{theorem}[Exhaustion of the ambient scalar target]
\label{p2:thm:scalar-target-exhaustion}
Let \(I=\mathbb C\) be the tensor unit of a finite complex-linear carrier
grammar.  Evaluation at \(1\in I\) gives a canonical isomorphism
\[
 \operatorname{End}(I)
 =\operatorname{End}_{\mathbb C}(\mathbb C)
 \cong\mathbb C.
\]
Consequently every primitive complex scalar output is one complex number, and
the standard dagger/ordered predicates used below select the slices
\[
\begin{array}{c|c}
 \text{condition on }z&\text{target slice}\\ \hline
 z^*=z&\mathbb R\\
 z=w^*w&\mathbb R_{\ge0}\\
 0\le z\le1&[0,1]\\
 z^*z=1&U(1).
\end{array}
\]
A separately generated real form restricts the target to \(\mathbb R\); it
does not create a second unconstrained scalar universe.  The exhaustive claim
is the ambient target \(\operatorname{End}(I)\cong\mathbb C\), not that the
four displayed predicates are the only generated subsets of \(\mathbb C\).
For example, adjoining the idempotence predicate \(z^2=z\) selects
\(\{0,1\}\); every such further restriction must retain the provenance of its
additional predicate.
\end{theorem}

\begin{proof}
Every complex-linear \(a:I\to I\) satisfies
\(a(\lambda)=\lambda a(1)\), and multiplication by an arbitrary
\(z\in\mathbb C\) supplies the converse.  The dagger is complex conjugation,
so self-adjointness gives \(z=\bar z\).  Positivity gives
\(z=\bar w w=|w|^2\); the effect order adds \(z\le1\); and unitarity gives
\(\bar z z=1\).  These alternatives exhaust the displayed structural slices
of \(\operatorname{End}(I)\).  No occurrence or empirical interpretation is
asserted.
\end{proof}

\begin{corollary}[Cardinality of finite readout families]
\label{p2:cor:finite-readout-cardinality}
Write \(\mathfrak c:=|\mathbb C|=2^{\aleph_0}\).  Then:
\begin{enumerate}[label=(\roman*),leftmargin=2em]
\item the exhaustive primitive scalar target has cardinality
\(\mathfrak c\);
\item for every nonzero finite-dimensional \(H\), the vector space of all
complex-linear maps
\(\operatorname{End}(H)\to\mathbb C\) has cardinality \(\mathfrak c\);
\item for a fixed countable constructor grammar whose diagrams carry only
finitely many finite-dimensional complex parameters, the union of all such
finite parametrized readout families has cardinality at most
\(\mathfrak c\), and has cardinality \(\mathfrak c\) whenever arbitrary
constant scalars are licensed;
\item if \(X\) is a nonempty separable parameter/orbit space, then
\(|C(X,\mathbb C)|=\mathfrak c\), whereas if \(|X|=\mathfrak c\), the set of
all unrestricted functions \(X\to\mathbb C\) has cardinality
\(2^{\mathfrak c}\).
\end{enumerate}
The later two-role theorem therefore exhausts two primitive
\emph{provenance/topology types}; it does not reduce the continuum of
possible values or parametrized process laws to two elements.
\end{corollary}

\begin{proof}
The first claim is
\(|\operatorname{End}_{\mathbb C}(\mathbb C)|=|\mathbb C|=\mathfrak c\).
If \(n=\dim_{\mathbb C}H>0\), then
\(\operatorname{End}(H)^*\cong\mathbb C^{n^2}\), whose cardinality is
\(\mathfrak c\), proving (ii).  A finite diagram with finitely many
finite-dimensional complex parameters has at most \(\mathfrak c\)
instantiations; a countable union of such sets still has cardinality at most
\(\aleph_0\mathfrak c=\mathfrak c\), proving (iii).  For (iv), a continuous
map on a separable space is determined by its values on a countable dense
set, so
\[
 |C(X,\mathbb C)|\le\mathfrak c^{\aleph_0}
 =(2^{\aleph_0})^{\aleph_0}=\mathfrak c.
\]
Constant maps give the reverse inequality.  Finally,
\[
 |\mathbb C^X|=\mathfrak c^{\mathfrak c}
 =(2^{\aleph_0})^{\mathfrak c}=2^{\mathfrak c}.
\]
\end{proof}

\begin{lemma}[Canonical finite duality and the trace]
\label{p2:lem:canonical-finite-duality-trace}
For a finite-dimensional complex vector space \(H\), put
\(H^*=\operatorname{Hom}_{\mathbb C}(H,\mathbb C)\).  For a basis
\((e_i)_{i=1}^n\) and its dual basis \((e^i)_{i=1}^n\), define
\[
 \operatorname{coev}_H(1)=\sum_{i=1}^n e^i\otimes e_i,
 \qquad
 \operatorname{ev}_H(\epsilon\otimes\psi)=\epsilon(\psi).
\]
The coevaluation tensor is basis-independent, the two maps obey the triangle
identities, and every \(T\in\operatorname{End}(H)\) satisfies
\[
 \operatorname{ev}_H\circ(1_{H^*}\otimes T)
 \circ\operatorname{coev}_H
 =\sum_i e^i(Te_i)=\operatorname{Tr}_H T.
\]
\end{lemma}

\begin{proof}
Under the canonical isomorphism
\(H^*\otimes H\cong\operatorname{End}(H)\), the tensor
\(\epsilon\otimes\psi\) is the rank-one map
\(x\mapsto\psi\epsilon(x)\).  Thus
\(\sum_i e^i\otimes e_i\) corresponds to \(1_H\), proving basis
independence.  The identities
\(\sum_i e_i e^i(x)=x\) and
\(\sum_i\epsilon(e_i)e^i=\epsilon\) are the two triangle identities.
Closing \(T\) gives the displayed diagonal sum, hence the ordinary algebraic
trace.  No basis, state, projector, or normalization weight survives in this
closure.
\end{proof}

\begin{lemma}[Algebraic exhaustion of linear scalar readouts]
\label{p2:lem:linear-readout-exhaustion}
Every complex-linear functional
\(\ell:\operatorname{End}(H)\to\mathbb C\) has a unique
\(B\in\operatorname{End}(H)\) such that
\[
 \ell(T)=\operatorname{Tr}(BT)\qquad(T\in\operatorname{End}(H)).
\]
If \(B=\sum_{a=1}^r\psi_a\otimes\epsilon_a\) is a rank-one
decomposition, then
\[
 \ell(T)=\sum_{a=1}^r\epsilon_a(T\psi_a).
\]
Thus every finite linear scalar readout is a finite sum of boundary matrix
coefficients; the state-free closing tensor \(B=1_H\) gives the trace.
\end{lemma}

\begin{proof}
The pairing \((B,T)\mapsto\operatorname{Tr}(BT)\) is nondegenerate.  Indeed,
for matrix units,
\(\operatorname{Tr}(E_{ij}E_{k\ell})=\delta_{i\ell}\delta_{jk}\),
so the values of \(\ell\) on the matrix units determine a unique \(B\).
Every finite matrix is a finite sum of rank-one maps, and
\[
 \operatorname{Tr}\bigl((\psi\otimes\epsilon)T\bigr)
 =\epsilon(T\psi),
\]
which proves the normal form.
\end{proof}

\begin{theorem}[Uniqueness of the canonical unmarked duality closure]
\label{p2:thm:unmarked-natural-closure}
Let $H$ be a generated finite-dimensional complex carrier in a frozen
constructor grammar whose licensed duality is the canonical finite duality of
Lemma~\ref{p2:lem:canonical-finite-duality-trace}. Suppose the closing word
applied to $T\in\operatorname{End}(H)$ is precisely
\[
 \ell_H(T):=
 \operatorname{ev}_H\circ(1_{H^*}\otimes T)\circ\operatorname{coev}_H,
\]
with the licensed symmetry of tensor factors understood. Then
\[
 \boxed{\ell_H(T)=\operatorname{Tr}_H(T).}
\]
The closing functional is independent of the duality presentation and is
unitary-natural, complex-linear, and positive on positive endomorphisms.
It is additive on every licensed direct sum, and multiplicative under every
licensed tensor aggregation. No decomposition into physically generated
one-dimensional blocks is required. A weighting, retained insertion, extra
loop, or scalar postprocessor belongs to a different declared word or process.
\end{theorem}
\begin{proof}
Under $H^*\otimes H\cong\operatorname{End}(H)$, a triangle identity forces
$\operatorname{coev}_H(1)$ to be the identity endomorphism. In a basis it is
$\sum_i e^i\otimes e_i$, so the licensed closure evaluates to
$\sum_i e^i(Te_i)=\operatorname{Tr}_H T$. This is a host calculation of the
already licensed word and does not constitute a new direct-sum operation.
Uniqueness follows from uniqueness of the tensor corresponding to $1_H$.
Changing the presentation changes evaluation and coevaluation compatibly and
leaves this contraction unchanged. Positivity follows from
$\operatorname{Tr}(B^*B)=\sum_{ij}|B_{ij}|^2$; naturality, additivity and
tensor multiplicativity follow from the ordinary finite index sums.
\end{proof}

For arbitrary natural linear functional families, unit normalization fixes
only the coefficient classes connected to the unit by the actually licensed
sum constraints. Lemma~\ref{p2a:lem:repair-natural-functional-coefficients}
classifies that freedom. The present theorem uses the specified canonical
duality word and therefore needs no extra sum license.


\begin{definition}[Finite structural readout predicate]
\label{p2:def:finite-structural-readout}
Let \(\Sigma\) be a frozen finite-linear rigid symmetric-monoidal physical
constructor grammar.  Its licensed structure includes the tensor unit,
associative/symmetric tensor aggregation of already licensed ports, sequential
composition, duals with evaluation and coevaluation, and exactly those
biproduct/linear-alternative operations explicitly listed in \(\Sigma\).
No other monoidal, duality, or trace operation is imported from the ambient
category.  The statement \(\operatorname{FSR}_{\Sigma}(r)\) means:
\begin{enumerate}[label=(\roman*),leftmargin=2em]
\item \(r\) is the scalar output of a finite process diagram built only from
\(\Sigma\)-generated objects, arrows, and licensed composition or
partial-evaluation constructors;
\item its apparatus map is fixed independently of the result, as required by
G0, and is natural under unmarked presentation isomorphisms;
\item every primitive arrow is linear, generated alternatives are combined
linearly, and sequential composition uses the licensed categorical
composition; every nonlinear estimator is separately typed;
\item every retained boundary state, effect, field, source, grading,
projector, or channel label has generated provenance;
\item the primitive diagram is connected; disconnected diagrams are generated
composites.
\end{enumerate}
This predicate neither licenses every mathematically writable scalar function
nor asserts that a particular admissible readout occurs.
\end{definition}

\begin{theorem}[Loop--boundary exhaustion of primitive readouts]
\label{p2:thm:loop-boundary-readout-exhaustion}
Every \(r\) satisfying \(\operatorname{FSR}_{\Sigma}(r)\) lies in exactly
one of the two provenance roles
\[
 \boxed{\Rprim=\Rloop\sqcup\Rbdry.}
\]
After aggregating the already licensed ports, its normal form is
\[
\boxed{
\begin{aligned}
 r\in\Rloop
 &\Longrightarrow r=\operatorname{Tr}_XW,
 &&W:X\to X,\\
 r\in\Rbdry
 &\Longrightarrow r=\epsilon T\psi,
 &&\psi:I\to X,
 \quad T:X\to Y,
 \quad\epsilon:Y\to I.
\end{aligned}}
\]
Boundary provenance remains present even if a boundary coefficient is
algebraically rewritten as a trace with a rank-one insertion.
\end{theorem}

\begin{proof}
Because the diagram is finite, tensor all \emph{licensed} incoming ports into
one object \(X\), all licensed outgoing ports into one object \(Y\), and
compose its internal coupons into one effective arrow \(T:X\to Y\).  This is
only aggregation of constructors already in \(\Sigma\).

If no external boundary remains, every port is paired by generated duality.
Cut one closure wire.  The cut produces an endomorphism \(W:X\to X\), and
re-closing it gives the categorical trace, which is the algebraic trace by
Lemma~\ref{p2:lem:canonical-finite-duality-trace}.  Naturality and cyclicity
make the result independent of the cut.

If an external boundary remains, aggregate the incoming preparation or source
ports into \(\psi:I\to X\), and the outgoing test or sink ports into
\(\epsilon:Y\to I\).  Sequential composition gives \(\epsilon T\psi\).
The two cases exhaust the possibilities because the external boundary is
either empty or nonempty.  They are disjoint as provenance roles because that
boundary inventory is frozen before numerical evaluation.  Rewriting
\(\epsilon T\psi\) as
\(\operatorname{Tr}((\psi\otimes\epsilon)T)\) records the rank-one boundary
insertion and therefore does not turn it into an unmarked loop.
\end{proof}

\begin{corollary}[Generated closure and absence of a third primitive role]
\label{p2:cor:generated-readout-closure}
Every finite structural scalar readout has the constructor-relative form
\[
 \Rgen
 =\operatorname{Gen}_{\Sigma}
 \bigl(\Rloop\sqcup\Rbdry\bigr).
\]
For disconnected diagrams, licensed tensor composition multiplies component
scalars, licensed biproduct alternatives add them, and a further estimator is
an explicit generated map \(F_{\rm gen}\).  Decorations such as
\[
 \operatorname{Tr}(\Gamma W),\quad
 \operatorname{Tr}(PW),\quad
 (\operatorname{Tr}W)^2,\quad
 \det D,\quad \operatorname{Pf}A,\quad
 -\operatorname{Tr}(\rho\log\rho),\quad
 \operatorname{Tr}[E\Phi(\rho)]
\]
are respectively marked-loop subclasses, multiloop composites, exterior or
oriented closures, and boundary-conditioned evaluations; they are not new
primitive outer topologies.  Their admissibility requires the corresponding
generated grading, projector, exterior power, orientation, state, effect,
process, or postprocessing constructor.  Hence GP3 has no third readout-role
candidate to delete or minimize: it may compare carrier realizations of a
fixed role, but it cannot erase a retained boundary mark or select a process
law.
\end{corollary}

\begin{definition}[Finite generated-constituent presentation]
\label{p2:def:finite-generated-constituent-presentation}
Let \(P\) be a G0-occurring composite in a frozen constructor grammar
\(\Sigma\).  A presentation
\[
 P=G_{\Sigma}(r_1,\ldots,r_N)
\]
is a \emph{finite generated-constituent presentation} when:
\begin{enumerate}[label=(\roman*),leftmargin=2em]
\item \(G_{\Sigma}\) is built only from licensed constructors and its
      constituent ports are fixed before evaluation;
\item every \(r_i\) is generated from the occurring input data of \(P\), has
      a frozen declared object, operation, boundary, or scalar output sort,
      and is not merely a mathematically available alternative;
\item the list is nonredundant modulo generated presentation isomorphisms;
      equal numerical values in different frozen sorts are not identified;
\item deleting any \(r_i\), or replacing its target sort by that of another
      constituent, changes the generated presentation of \(P\).
\end{enumerate}
This definition concerns the internal generated content of one occurring
composite.  A declared target sort is a requirement of the frozen presentation,
\emph{not} an assertion that a candidate has already realized the constituent;
realization of every required constituent is the conclusion of
Theorem~\ref{p2:thm:generated-constituent-corealization}.  The definition does
not turn the class of all diagrams satisfying
\(\operatorname{FSR}_{\Sigma}\) into an occurring family.
\end{definition}

\begin{theorem}[GP2 co-realization of generated constituents]
\label{p2:thm:generated-constituent-corealization}
Let \(P\) have a finite generated-constituent presentation
\[
 P=G_{\Sigma}(r_1,\ldots,r_N).
\]
If the G0 occurrence of \(P\) and this presentation are frozen before the
result is evaluated, then GP2 exactness, typed host faithfulness, and
no-omission require the simultaneous realization of every \(r_i\) in its
native generated sort.  No additional selection or occurrence certificate is
required for an \(r_i\) merely to be a constituent of \(P\).

This conclusion does not identify distinct sorts, does not add the values of
distinct constituents, and does not assert the occurrence of a different
apparatus, preparation, estimator, or process diagram that is only admissible
in the same grammar.
\end{theorem}

\begin{proof}
Because \(P\) occurs in the G0 sense of Principle~\ref{p2:prin:geometric-primacy-minimal-constitution} with the displayed presentation frozen, each \(r_i\)
belongs to the constructor-relative generated content
\(\operatorname{Gen}_{\Sigma}(P)\).  If some \(r_i\) were absent from the
realization of \(P\), clause~(iv) of
Definition~\ref{p2:def:finite-generated-constituent-presentation} would make
the represented composite a proper omission of its generated presentation,
contrary to Postulate~\ref{p2:post:GP2} (exactness and no-omission).
Replacing \(r_i\) by an isomorphic copy in another sort would contradict the
same postulate's typed host faithfulness.  Hence all
constituents are represented simultaneously in their native sorts.

The argument begins with the occurrence of \(P\); it therefore has no force
for a different diagram that is merely a member of the admissible readout
class.  Finally, GP2 preserves licensed composition and typing but does not
insert a new sum, tensor product, empirical performance, or numerical
identification.  Those operations would require their own generated
constructor.  This proves the scope clauses.
\end{proof}

\begin{definition}[Finitely presented generated-role descendant]
\label{p2:def:finitely-presented-generated-role-descendant}
Let \(P\) be a G0-occurring structured composite with frozen constructor
grammar \(\Sigma\).  A role \(R\) is a \emph{finitely presented
generated-role descendant} of \(P\), written \(R\preccurlyeq_\Sigma P\),
when:
\begin{enumerate}[label=(\roman*),leftmargin=2em]
\item \(R\) is the image or exact closure of constructors and operations
      already generated from the occurring data of \(P\);
\item its equivalence relation, topology, operation law, and target sort are
      fixed by \(\Sigma\), not selected from the eventual result;
\item it is invariant under every unmarked automorphism of the frozen
      presentation and nonredundant for typed host faithfulness; and
\item it has a GP1-finite exact presentation in the structured-realization
      sense.  Its underlying parent set may be continuous, provided its
      represented operation envelope and retained presentation satisfy GP1.
\end{enumerate}
The notation concerns a role occurring as content of \(P\), not the assertion
that every mathematically admissible element of that role occurs as a token.
\end{definition}

\begin{theorem}[Generated-role co-realization; no duplicate actualization]
\label{p2:thm:generated-role-corealization}
Let \(P\) be G0-occurring and let
\(R\preccurlyeq_\Sigma P\).  Then the same occurrence of \(P\) co-realizes
\(R\) in its native sort.  No independent actuality witness is required
merely to pass from the occurring root to this generated role:
\[
 \boxed{
 \PhysOcc_{P_n}(P)\ \wedge\ R\preccurlyeq_\Sigma P
 \quad\Longrightarrow\quad
 \operatorname{CoReal}_{P_n[P]}(R\mid P).}
\]
Here \(\operatorname{CoReal}(R\mid P)\) abbreviates constituent occurrence
inside the same G0 witness for \(P\); it is not a new actualization predicate
or principle.  The conclusion does not imply
\(\PhysOcc(r)\) for every \(r\in R\), nor the occurrence of a different
preparation, apparatus, process, history, or calibration whose type merely
belongs to \(R\).
\end{theorem}

\begin{proof}
By clause (i), the exact presentation of \(R\) belongs to
\(\operatorname{Gen}_\Sigma(P)\).  Omitting it from the realization of the
occurring \(P\) would violate GP2 no-omission; changing its closure law,
equivalence, topology, or target sort would violate exactness or typed host
faithfulness.  Clause (iii) prevents GP3 from deleting it as an unmarked
implementation redundancy, while clause (iv) places it in the admissible
finite structured presentation.  Thus the original witness for \(P\) realizes
\(P\) together with \(R\).  No second contingent token has entered the proof.
The last sentence follows because membership in a generated role is an
admissibility statement, whereas token occurrence still contains the
primitive predicate \(\Actual_{\mathcal E}\).
\end{proof}

\begin{theorem}[Typed cross-relation firewall]
\label{p2a:thm:typed-cross-relation-firewall}
Let a generated relation be constituted in one of the typed sorts
\[
 \mathsf{Prov},\ \mathsf{Comp},\ \mathsf{Eval},\ \mathsf{Symp},\
 \mathsf{Act},\ \mathsf{Polar},\ \mathsf{Cal},\ \mathsf{DynGraph}.
\]
GP2 exactness preserves both its content and its sort.  In particular, a
relation generated in sort \(\sigma\) does not by itself constitute a relation
in a distinct sort \(\sigma'\neq\sigma\).  Promotion between sorts requires an
already licensed constructor whose typed input is the first relation and whose
required output is the second.
\end{theorem}

\begin{proof}
A candidate that replaces a generated relation by one in another target sort
changes the frozen presentation and therefore has nonzero typing realization
defect.  Conversely, if a licensed constructor with the stated source and
target sorts is present, omission of its required nonredundant output is a GP2
closure defect by the no-idle rule.  This is exactly the distinction between
constructor-relative generation and mere mathematical availability.
\end{proof}

\begin{corollary}[Zero-member test for nonzero occurrence]
\label{p2a:cor:zero-member-occurrence-test}
Suppose a later role \(\mathcal R\) has already been generated, but its zero
member (or the absence of a nonzero token) preserves every frozen relation,
type, intrinsic defect signature, and licensed-constructor closure.  Then GP2
does not force a nonzero member of \(\mathcal R\) to occur.  A nonzero field,
connection, interaction coefficient, state history, or calibration therefore
remains a G0 token unless its nonzero value is itself the required output of a
licensed constructor on an occurring root.
\end{corollary}

\begin{corollary}[Declared-output actualization boundary]
\label{p2a:cor:declared-output-actualization-boundary}
Let \(C\) be a licensed constructor in the frozen grammar and let all of its
required generated or occurring inputs be present.  GP2 no-idle closure forces
only the nonredundant output sorts explicitly declared by the constructor
schema.  A mathematically canonical composition, partial evaluation, solution
map, orbit, or spectral subspace does not acquire an additional physical target
sort merely because it can be formed.  In particular, a canonical orbit
\[
  t\longmapsto \Phi_t(x_0)
\]
need not be an actual-history token unless the licensed schema has an
actual-history output, and a canonical positive-frequency mode space need not
have one-particle semantics unless a licensed field/species constructor assigns
that role.  Promotion between these sorts is governed by
Theorem~\ref{p2a:thm:typed-cross-relation-firewall}.
\end{corollary}

\begin{proof}
The no-idle clause retains required outputs of licensed schemata, while the
mathematical-availability firewall forbids adding undeclared outputs.  Changing
only the target sort of a canonical mathematical consequence changes the
frozen typed presentation and hence has nonzero typing realization defect.
\end{proof}

\begin{definition}[Frozen constructor-origin certificate]
\label{p2a:def:frozen-constructor-origin-certificate}
Let \(P\) be a frozen typed presentation with licensed physical constructor
grammar \(\Sigma\), let \(D\) be a jointly occurring/generated input tuple,
and let \(\mathsf Q\) be a target sort.  The predicate
\(\operatorname{Orig}_{\Sigma}(w;D,\mathsf Q)\) means that, before any output
is evaluated, \(P\) contains a finite typed word
\[
 x_0=D\xrightarrow{C_1}x_1\xrightarrow{C_2}\cdots
 \xrightarrow{C_m}x_m,
 \qquad q:=\operatorname{ReqOut}_{\Sigma}(w,D)\in\mathsf Q,
\]
with the following properties.
\begin{enumerate}[label=(\roman*),leftmargin=2em]
\item Every \(C_j\), including every partial evaluation, composition,
      dualization, contraction, quotient, completion, or functional-calculus
      step actually used, is licensed in \(\Sigma\) at this chronological
      stage with its displayed domain and codomain.
\item Every input port is filled by an entry of \(D\) or by a previously
      retained output in its native sort.  Each intermediate output needed
      later is retained and licensed for reuse; no object, operation, carrier,
      relation, or token is silently retyped.
\item The target \(\mathsf Q\), the required nonredundant output \(q\), and any
      finite exact closure retained with \(q\) are declared by the constructor
      schemas before evaluation.  The word respects the licensed iteration and
      closure boundary.
\item The word is frozen independently of the result and imports no
      candidate-specific selector, later-frontier constructor, new scale, or
      undeclared parameter.
\end{enumerate}
A one-step formula is not an origin certificate unless its constructor was
already licensed.  This definition records provenance; it does not add a
constructor to \(\Sigma\).
\end{definition}

\begin{theorem}[Rooted retention and co-realization of a generated image]
\label{p2a:thm:rooted-generated-image-corealization}
\label{p2a:thm:g0-witness-transport-derived-constructor}
Let \(P,\Sigma,D,\mathsf Q,w\) satisfy
\(\operatorname{Orig}_{\Sigma}(w;D,\mathsf Q)\), let
\(q=\operatorname{ReqOut}_{\Sigma}(w,D)\), and write
\(\operatorname{RootReal}_P(D)\) for the conjunction that every token root in
\(D\) occurs and every generated entry of \(D\) is co-realized with its
complete licensed root.  Then
\[
 \boxed{
 \operatorname{RootReal}_P(D)
 \Longrightarrow
 \operatorname{Obs}_{\rm FRC}(w,D)
 \ \vee\ 
 \operatorname{CoReal}_{\mathsf Q}(q\mid D;w).}
\]
Here \(\operatorname{Obs}_{\rm FRC}(w,D)\) means that the complete word
presentation has no GP1-finite-and-bounded, GP2-exact-and-faithful structured
realization.  On the nonobstructed branch, every such realization retains
\(q\) in \(\mathsf Q\); no duplicate G0 occurrence is needed.

The conclusion is role/constituent co-realization.  It neither actualizes every
member of an output role nor promotes an output to an undeclared sort.  GP3 is
not used: it may subsequently retain a unique least realization or leave the
irreducible family unresolved, but it does not generate \(q\).
\end{theorem}

\begin{proof}
The origin certificate places every step of \(w\), every retained intermediate
output, and \(q\) in the constructor-relative closure
\(\operatorname{Gen}_{\Sigma}(D)\).  If the complete presentation has no
finite-and-bounded exact faithful realization, the first alternative holds by
definition.  Otherwise choose any candidate in that class.  Omitting \(q\)
violates GP2 no-omission/closure; changing its target sort violates typed host
faithfulness; adding content not declared by the word violates no-insertion.
Thus \(q\) is present in every nonobstructed candidate as a generated
constituent or generated-role descendant of the same occurring root tuple.
That universal retained presence in the required native sort is precisely
\(\operatorname{CoReal}_{\mathsf Q}(q\mid D;w)\), giving the second
alternative directly from GP2.  If GP3 is applied later, its admissible
reductions remain GP2-exact and therefore cannot omit or retype \(q\).
No transport, naturality, leastness, or zero-defect premise was used to license
the constructors: license entered only through
\(\operatorname{Orig}_{\Sigma}\), while zero GP2 defect is the retained
candidate conclusion.
\end{proof}

\section{Faithful host realization and constructor closure}
\label{paper2a:sec:organized-hosting}

\subsection{Hosting, construction and retained incidence}

\begin{definition}[Host, constructor and incidence contract]
\label{p2a:def:host-constructor-incidence-contract}
Fix a constituted geometric host \(H\), its enriched subroles, a typed
constructor grammar and the retained input diagram.  Write
\(X\preceq_H H\) for geometric hosting: \(X\) is a constituted subrole
or a bundle, section, connection, representation, mode, state or other
licensed geometric construction supported by those roles.  This is not
literal inclusion of every output as a subset of the underlying manifold.
Write
\[
 (X_1,\ldots,X_k)\xrightarrow{D_f}Y
\]
for the specific licensed constructor \(f\), including its input marks,
codomain, domain predicates and proof of transport.  Write
\(R(X_1,\ldots,X_k;Y_1,\ldots,Y_j)\) for a retained generated incidence.
An incidence can have more than two arguments and need not be an
isomorphism.  Its automorphisms must transport the whole diagram.
For an occurrence of a constructor word, its consumed host support is
the union of supports of the inputs actually used by that word.
A minimal host is minimal for that certified word and support
ordering; this asserts neither a unique constructor nor a universal
smallest host among inequivalent realizations.
\end{definition}

\begin{corollary}[Recursive coverage of a licensed construction]
\label{p2a:cor:recursive-host-coverage}
If the initial constituted inputs are hosted by \(H\), every output
of a finite composition of host-compatible licensed constructors is
hosted by \(H\).  The retained incidence diagram is covered at the
same time.  Shared hosting does not imply equality of outputs,
equivalence of their actions, or occurrence of every possible member.
\end{corollary}
\begin{proof}
At depth zero this is the input hosting certificate.  Suppose the
claim holds through depth \(d\).  Each input to a constructor at depth
\(d+1\) has a host certificate by induction.  The constructor's
host-compatibility certificate supplies the output certificate and
transports its specified incidences.  This proves the induction.
Equality or contraction requires a separate same-sort identity theorem;
neither is a clause of a hosting certificate.
\end{proof}

The frozen-origin and co-realization theorems specify when such a
word is licensed in the occurring parent.  G0 refers to the actual
global parent occurrence; its consequences do not require a new
independent actuality premise at every readout.  GP1--GP3 are unchanged.
They do not turn availability of a carrier into selection of a state,
connection member, vacuum, solution or physical interpretation.
The later compact \(M_8\) presentation is an overview of the completed
host.  It does not reverse the proof order by making the completed
product an antecedent of the factors from which it is constructed.
Thus a package can have a host path through an internal factor without
being the output of one constructor applied to the whole product.



\begin{theorem}[GP2 same-owner equivariance under unmarked automorphisms]
\label{p2a:thm:v368-same-owner-unmarked-equivariance}
Let \(P\) be a G0-occurring frozen typed presentation with intrinsic physical
constructor grammar \(\Sigma\).  Let \(D\) be a root-realized input tuple and
let
\[
  \operatorname{Orig}_{\Sigma}(w;D,\mathsf Q)
\]
be a frozen constructor-origin certificate with required output
\(q=\operatorname{ReqOut}_{\Sigma}(w,D)\).  Let
\(a\in\operatorname{Aut}_0(P;w)\) be an unmarked automorphism of the
\emph{complete frozen presentation entering the word}: it preserves every
retained mark, relation, source, boundary/interface datum, constructor schema,
and typing used by \(w\), and it induces the declared target action
\(a_{\mathsf Q}\) on \(\mathsf Q\).  Then on the non-FRC-obstructed branch,
\[
 \boxed{
 \operatorname{ReqOut}_{\Sigma}(w,aD)
 =a_{\mathsf Q}\,
  \operatorname{ReqOut}_{\Sigma}(w,D).}
 \tag{SO-EQ}
\]
Equivalently, every retained same-owner descendant is equivariant in its
native target sort under every unmarked automorphism of the complete
presentation from which it is generated.

If \(\mathsf Q\) carries the trivial action, then the descendant is invariant:
\[
 \boxed{q(aD)=q(D).}
 \tag{SO-INV}
\]
In particular, a licensed scalar action descendant
\(S=\mathscr A_{\Sigma}(D)\) satisfies
\[
 \boxed{S(aD)=S(D)}
 \tag{ACT-INV}
\]
for every such unmarked automorphism.

The theorem has no force when \(a\) changes a generated source, boundary
condition, coupling, orientation, grading, state-dependent mark, or any other
part of the complete action-owning presentation.  Such a transformation is
not an unmarked automorphism of the antecedent presentation and may represent
an explicitly broken symmetry or a different physical input.  Likewise, the
theorem does not assert that every generated operation is an action symmetry.
\end{theorem}

\begin{proof}
Write the certified word as
\[
 x_0=D\xrightarrow{C_1}x_1\xrightarrow{C_2}\cdots
 \xrightarrow{C_m}x_m=q.
\]
Because \(a\) is an automorphism of the complete frozen presentation entering
\(w\), it preserves every typed relation and constructor schema used by the
word.  Inductively, if \(x_{j-1}\) is transported to \(a x_{j-1}\), then the
same frozen constructor schema gives
\[
 C_j(a x_{j-1})=a C_j(x_{j-1})
\]
in the declared target sort.  Hence the transported terminal output is
\(a_{\mathsf Q}q\), proving \((\mathrm{SO\text{-}EQ})\) at the level of the
constructor presentation.  Theorem~\ref{p2a:thm:rooted-generated-image-corealization}
and GP2 no-omission retain that transported output physically, while typed
host faithfulness forbids replacing it by an inequivalent output in another
sort.  A failure of equivariance would therefore make the frozen constructor
word distinguish two presentations related only by an unmarked automorphism;
that would insert a presentation-dependent mark not generated by the complete
host, contrary to GP2 no-insertion and the claim-relative unmarked-isomorphism
clause of G0.  Thus the equivariance is forced in every GP2-exact
nonobstructed realization.

When the target action is trivial, \(a_{\mathsf Q}=\mathrm{id}\), so
\((\mathrm{SO\text{-}INV})\) follows.  The action specialization is exactly
this trivial-target case.  The final firewall follows because a transformation
that changes any retained antecedent mark is not an element of
\(\operatorname{Aut}_0(P;w)\).
\end{proof}

\begin{lemma}[Automatic regularity of a finite linear one-parameter law]
\label{p2a:lem:repair-finite-generator-regularity}
Let $T_t$ be a norm-continuous one-parameter group, or a norm-continuous
semigroup for $t\ge0$, on a finite-dimensional real or complex vector space,
with $T_0=I$ and $T_{s+t}=T_sT_t$. There is a unique matrix $A$ such that
$T_t=e^{tA}$. If the law is unitary, $A^*=-A$.
\end{lemma}
\begin{proof}
For small $a>0$, continuity makes
$B=\int_0^aT_t\,dt$ invertible because $\|B/a-I\|<1$.
The law gives $T_sB=\int_s^{s+a}T_t\,dt$. Differentiating the integral yields
\[
 T_s'=(T_{s+a}-T_s)B^{-1}
     =T_s(T_a-I)B^{-1}=T_sA.
\]
At zero in the semigroup case use the right derivative. Thus $A=T'_0$ is
unique, and differentiation of $T_se^{-sA}$ proves the exponential formula.
Differentiation of $T_s^*T_s=I$ proves skew-adjointness. A mathematical
negative-time extension of a semigroup is not thereby a licensed physical
inverse in its original process class.
\end{proof}

\begin{corollary}[Native infinitesimal action invariance on its derivative domain]
\label{p2a:cor:v368-native-noether-antecedent}
Under Theorem~\ref{p2a:thm:v368-same-owner-unmarked-equivariance}, let
$s\mapsto g_s$ be a constituted family of unmarked automorphisms of the
complete action-owning presentation, with $g_0=1$. Let $S$ be its licensed
scalar action descendant. At a configuration $D$, require that the orbit
$s\mapsto g_sD$ is differentiable at zero in its stated tangent topology and
that $S$ has the corresponding chain rule there; Fr\'echet differentiability
in a normed configuration chart is one sufficient condition. Define
\[
 R_\xi(D)=\left.\frac{d}{ds}g_sD\right|_{s=0}.
\]
Then
\[
 S(g_sD)=S(D),\qquad
 \boxed{dS_D(R_\xi(D))=0.}
\]
This is the native invariant variational antecedent. A continuous
finite-dimensional linear one-parameter group automatically has the required
smooth orbit by Lemma~\ref{p2a:lem:repair-finite-generator-regularity}.
Continuity of an arbitrary family alone is insufficient.

A local Noether current or identity still requires the licensed variational
and integration-by-parts calculus. A global charge requires its boundary and
global hypotheses; a quantum Ward identity requires its measure, regulator
and anomaly conditions. No occurring history or numerical charge is selected.
\end{corollary}
\begin{proof}
The same-owner theorem gives the invariant finite identity. On the stated
domain, write $g_sD=D+sR_\xi(D)+o(s)$. Differentiability of $S$ gives
\[
 0=S(g_sD)-S(D)=s\,dS_D(R_\xi(D))+o(s).
\]
Divide by $s$ and pass to zero. More generally the admitted curve chain rule
gives the same result. The final conditions identify the distinct downstream
operations not contained in this local identity.
\end{proof}

\begin{theorem}[Evaluation-dual co-realization of operations and dynamics]
\label{p2a:thm:evaluation-dual-corealization}
Let a root-realized tuple \(D\) generate typed roles \(X,E,S\), a total
evaluation
\[
 \varepsilon:X\times E\longrightarrow S,
\]
and an operation \(T:X\to X\).  Assume:
\begin{enumerate}[label=(\roman*),leftmargin=2em]
\item \(E\) is separated by evaluation: if
      \(\varepsilon(x,e_1)=\varepsilon(x,e_2)\) for every \(x\in X\), then
      \(e_1=e_2\);
\item the frozen grammar licenses evaluation precomposition by \(T\), retains
      its output in the operation sort of \(E\), and declares \(E\) closed
      under that constructor; and
\item this constructor word, including its target and family law when
      present, has a frozen origin certificate in the sense of
      Definition~\ref{p2a:def:frozen-constructor-origin-certificate}.
\end{enumerate}
Then there is a unique generated dual operation \(T^\vee:E\to E\) satisfying
\[
 \boxed{\varepsilon(Tx,e)=\varepsilon(x,T^\vee e)}
 \qquad(x\in X,\ e\in E).
\]
The complete tuple
\[
 (X,E,S,\varepsilon,T,T^\vee)
\]
together with every operation graph and law actually declared by the grammar
is co-realized from the same root \(D\).  In particular, \(T^\vee\) requires
no duplicate operation occurrence.

Dualization is contravariant:
\[
 \operatorname{id}_X^\vee=\operatorname{id}_E,
 \qquad
 (T\circ U)^\vee=U^\vee\circ T^\vee.
\]
Consequently a generated monoid action on \(X\) yields an opposite-monoid
action on \(E\).  A generated one-parameter group or semigroup yields its
dual state/test, observable/test, pullback/pushforward, or
Schrödinger/Heisenberg description on the same parameter domain.  This is one
co-realized dynamics in two evaluation-dual presentations, not two physical
processes.

The theorem does not create the precomposition constructor, its target sort,
or closure of \(E\).  It does not turn a semigroup into a group, produce
negative-time evolution, select a Hamiltonian or generator, establish
complete positivity or detailed balance, choose an initial state, actualize a
trajectory/history, or calibrate the time parameter.  Those conclusions
require their own licensed outputs or occurring roots.
\end{theorem}

\begin{proof}
Closure gives, for every \(e\in E\), an element \(T^\vee e\in E\) with the
displayed evaluation identity.  Evaluation separation makes it unique.  It
also gives the identity law.  For composable \(U,T:X\to X\),
\[
 \varepsilon((T\circ U)x,e)
 =\varepsilon(Ux,T^\vee e)
 =\varepsilon(x,U^\vee(T^\vee e)),
\]
so separation yields
\((T\circ U)^\vee=U^\vee\circ T^\vee\).
The same calculation memberwise proves the opposite action law for a
generated family.

By hypothesis the evaluation, operation, precomposition word, output graphs,
and any asserted family law are required typed outputs of one frozen origin
certificate.  Theorem~\ref{p2a:thm:rooted-generated-image-corealization}
therefore retains them from \(D\) on every nonobstructed branch without a
second G0 witness.  None of the structures excluded in the final paragraph
is used in the construction; in particular, contravariance reverses
composition order but does not invert an operation or reverse physical time.
\end{proof}


\begin{theorem}[Exact stabilizer criterion for evaluation-dual operations]
\label{p2a:thm:v280-evaluation-dual-stabilizer}
Let $X,E,S$ be sets, let $\varepsilon:X\times E\to S$ be a total
pairing and suppose its evaluation representation
\[
 j:E\longrightarrow S^X,\qquad j(e)(x)=\varepsilon(x,e)
\]
is injective.  For any map $T:X\to X$ write
$P_T(F)=F\circ T$ on $S^X$.  Then a map $T^\vee:E\to E$ satisfying
\[
 \varepsilon(Tx,e)=\varepsilon(x,T^\vee e)
 \quad\text{for all }x,e
\]
exists \emph{if and only if}
\[
 \boxed{P_T(j(E))\subseteq j(E).} \tag{ED-STAB}
\]
When it exists it is unique and is given by
$T^\vee=j^{-1}P_Tj$, where $j^{-1}$ is restricted to $j(E)$.
The set of maps satisfying \textup{(ED-STAB)} is a submonoid of
$\operatorname{End}(X)$, and dualization is an antihomomorphism:
$(T\circ U)^\vee=U^\vee\circ T^\vee$.

Thus closure of the declared test/effect family under precomposition is an
\emph{exact mathematical existence certificate}, not an extra dual-operation
postulate.  The certificate does not itself license this precomposition
constructor as a \emph{physical} output.  Its frozen origin, typed target,
operation graph and any claimed family law must independently satisfy the
conditions of Theorem~\ref{p2a:thm:evaluation-dual-corealization}.
\end{theorem}
\begin{proof}
If $T^\vee$ satisfies the pairing identity, then
$P_T(j(e))=j(T^\vee e)$ for every $e$, proving necessity.  Conversely,
\textup{(ED-STAB)} allows a unique definition
$T^\vee e=j^{-1}(P_Tj(e))$; injectivity of $j$ proves both the identity
and uniqueness.  Because $P_{T\circ U}=P_UP_T$, stability is preserved by
composition and the dual operation reverses the order.  The identity map
also preserves $j(E)$.  The final boundary is exactly the separation
between mathematical closure and the frozen constructor-origin hypothesis
in Theorem~\ref{p2a:thm:evaluation-dual-corealization}.
\end{proof}

\begin{corollary}[Multi-root and role/token boundary]
\label{p2a:cor:g0-witness-transport-multiroot-boundary}
In Theorem~\ref{p2a:thm:evaluation-dual-corealization} the origin certificate is for the complete tuple
\(D\), not an arbitrarily chosen proper subtuple.  If a word consumes a
calibration, provenance arrow, action, support map, or output of an unresolved
family, that datum must occur or be generated in \(D\) in its native sort.  A
member-dependent word applies only on the complete surviving joint locus.  A
generated history doctrine may therefore be co-realized while a particular
member history remains a token-level G0 occurrence.
\end{corollary}

\begin{remark}[Constructor-origin firewall and necessity]
\label{p2a:rem:g0-witness-transport-firewall}
Intrinsicness, functoriality, naturality, conservativity, faithfulness, a
minimal essential image, and even a zero-defect proposed graph do not create
physical license.  To see necessity, take two presentations with identical
root data and host mathematics, \(\Sigma_0\) and
\(\Sigma_1=\Sigma_0\cup\{F:X\to\mathsf Q\}\), where \(F\) is natural and
faithful but is licensed only in \(\Sigma_1\).  Any rule admitting \(F\) from
those host properties alone generates the same output in \(\Sigma_0\),
contradicting constructor-relative closure.  The origin certificate separates
the two grammars.

The remaining clauses are sharp in the same local sense.  Without
clause~(ii), a word may consume a missing root or an intermediate that the
presentation did not retain, so its output is not generated from \(D\).
Without clause~(iii), the same mathematical value can be promoted after the
fact to an undeclared physical sort or closure.  Without clause~(iv), one may
choose a result-dependent constant word \(w_q\) after observing \(q\), making
every desired output vacuously “generated” and violating G0 chronology.
Finally, deleting the nonobstruction alternative is unsound: a licensed typed
word with no GP1-finite, GP2-exact faithful realization has no output that can
be co-realized.  Thus the theorem is clausewise premise-minimal relative to
the stated frozen-grammar semantics; no claim is made that syntactically
different foundations must encode provenance with this particular predicate.

Consequently an ambient tensor algebra, \(C^*\)-completion,
spectrum, solution space, quotient, orbit, or functional calculus contributes
physical output only through an already licensed typed word with retained and
reusable outputs.
\end{remark}

\begin{corollary}[Root--descendant normal form for occurrence assumptions]
\label{p2:cor:root-descendant-occurrence-normal-form}
Consider a finite chronological dependency presentation and collapse every
same-stage mutually generated block to one node.  Let
\(Q_1,\ldots,Q_s\) be the remaining nodes with no generated predecessor.
Then every physical node in the presentation is either one of these G0 roots
or a typed generated descendant co-realized from them.  Consequently the
premise-minimal occurrence ledger is
\[
 \boxed{
   \{\text{G0 occurrences of the independent roots }Q_i\}_{i=1}^{s}
   +\operatorname{Cl}_{\rm GP1:3}^{\rm typed}(Q_1,\ldots,Q_s).}
\]
An ungenerated cross-role bridge, a particular state or process token, and a
new calibration standard are new roots.  A canonical role, carrier, dual,
state/effect space, or representation generated from an existing root is not.
Equivalently, no second G0 root is introduced merely to supply a licensed
constructor output whose absence would give the existing host a GP2 closure
realization defect; Corollary~\ref{p2:cor:no-idle-licensed-constructor} and
Theorem~\ref{p2:thm:generated-role-corealization} already retain that output at
its generated role level.
\end{corollary}

\begin{proof}
After same-stage generated blocks are collapsed, chronological dependence is
acyclic.  Induct on its topological order.  Root nodes require G0 by
definition.  At each later node, GP2 no-insertion says that it is either
generated from earlier nodes or is an additional root; in the generated case
Theorem~\ref{p2:thm:generated-role-corealization} or
Theorem~\ref{p2:thm:generated-constituent-corealization} supplies
co-realization.  This proves the displayed normal form.
\end{proof}

\begin{corollary}[Operational momentum role and carrier]
\label{p2:cor:operational-momentum-role-occurrence}
Let \(\mathscr O_{\rm tr}\) be the transfer-operation content of the
G0-occurring primitive order-sensitive role, and let
\[
 \mathcal P
 :=\operatorname{Fr}_{\rm pres}\!\left(
   \operatorname{Cl}^{\rm typed}_{\rm GP1:3}
   (\mathscr O_{\rm tr}/\!\sim)\right)
\]
use exactly the generated relation, operation law, and topology of that role.
Then \(\mathcal P\preccurlyeq\mathscr O_{\rm tr}\) and
\[
 \PhysOcc(\mathscr O_{\rm tr})
 \Longrightarrow
 \operatorname{CoReal}(\mathcal P\mid\mathscr O_{\rm tr}).
\]
For the order-sensitive primitive branch, the allocation theorems identify
this role with the relation-exact parent
\(SU(2)\cong S^3_A(\Rmom)\).  Paper~II-B's
\(\mathcal H_p\cong\mathcal H_1\) is a further canonical finite generated-role
descendant and is therefore co-realized as the physical momentum carrier.
Neither step needs an independent G0-MOM witness.  A measured numerical
momentum and an identification such as
\(P_0=\hbar/\Rmom\) with a laboratory unit remain G0 readout/calibration
tokens.
\end{corollary}

\begin{proof}
The primitive G0 occurrence already contains the effective comparison and
local reapplication operations.  GP1--GP3 derive their exact least parent
without adding a result-dependent relation, and hence the parent satisfies
Definition~\ref{p2:def:finitely-presented-generated-role-descendant}.
Theorem~\ref{p2:thm:generated-role-corealization} gives its co-realization.
Paper~II-B proves that the finite carrier is the canonical nonredundant
generated image of the frozen round-host constructor frontier, so the same
theorem applies once more.  Numerical readout and calibration are not parts
of either closure, proving the final clause.
\end{proof}

\begin{corollary}[Canonical identification by a separating signature]
\label{p2:cor:canonical-constituent-identification}
Let \(\mathcal C(P)=\{r_1,\ldots,r_N\}\) be the co-realized constituent
family of Theorem~\ref{p2:thm:generated-constituent-corealization}, and let
\(\mathcal T(P)=\{T_1,\ldots,T_N\}\) be a generated family of terms in a
downstream presentation of the same composite.  Suppose a frozen invariant
signature
\[
 \operatorname{sig}
 =(
   \text{role},\text{support},\text{outer provenance},
   \text{Born parity},\text{tensor degree},\text{scale dependence})
\]
is injective on each family and has the same image on
\(\mathcal C(P)\) and \(\mathcal T(P)\).  Then there is a unique
type-preserving identification
\[
 \iota_P:\mathcal C(P)\xrightarrow{\ \cong\ }\mathcal T(P),
 \qquad
 \operatorname{sig}(\iota_P(r))=\operatorname{sig}(r).
\]
Thus no routing or selection datum remains.  If the signature is not
injective, exhaustion and co-realization still hold but do not choose among
the degenerate identifications.
\end{corollary}

\begin{proof}
For each \(r\in\mathcal C(P)\), equality of signature images gives at least
one \(T\in\mathcal T(P)\) with the same signature.  Injectivity on
\(\mathcal T(P)\) makes it unique.  Applying the same argument in the reverse
direction proves bijectivity.  The final statement is immediate: two terms
with the same complete signature are not separated by the stated invariant
data.
\end{proof}

\begin{remark}[State/test cross-link]
The finite state--test theorem later in this paper is the positive normalized
specialization of \(\Rbdry\): preparation supplies a state boundary, an
effect supplies a test boundary, and an occurring completely positive process
lies between them.  By contrast, \(\operatorname{Tr}1_H=\dim H\) is an
extensive loop normalization.  Even when
\(\rho=1_H/\dim H\), its preparation provenance and
\(\operatorname{Tr}\rho=1\) prevent a type identification with the unmarked
loop.
\end{remark}

\begin{remark}[Downstream consumption convention and parent-side re-typing]
\label{p2:rem:gp1-tb-consumption}
Whenever an argument below uses compactness of a physical comparison image,
it consumes Corollary~\ref{p2:cor:gp1-tb} with its explicit finite-operational
comparison typing.  Whenever it uses bounded scalar or timed-rate data, it
consumes FRC-b directly.  Whenever it uses a closed reversible operation
parent on a finite carrier, it may instead consume
Theorem~\ref{p2:thm:frc-automatic-compact-lie-closure}.  None of these routes
identifies an arbitrary primitive metric with the canonical Killing metric.
The continuum objects produced by the allocation chain---canonical compact
presentation completions and the least parent-role skeletons
\(U(1)\cong S^1\) and \(SU(2)\cong S^3\)---are parent operation geometries.
Finite physical realizations are determined by the constructor-relative host
signature.  A paired tower is used only for a role that explicitly constitutes
that grammar; the initial momentum host is selected directly and classified in
the mathematical tower only afterward.  The allocation theorems are therefore
re-typed, not weakened.
\end{remark}

\begin{definition}[Auxiliary primitive-comparison regularity conditions]
\label{p2:def:lie-regularity-conditions}
After the compact operation closure \(G\) has been derived, the notation
\(({\rm FD})+({\rm LR})\) means:
\begin{enumerate}[label=(\roman*), leftmargin=2em]
\item \(({\rm FD})\): \(G\), with its uniform metric, has finite topological
dimension;
\item \(({\rm LR})\): there are \(\kappa\in(0,1]\) and \(C_R\ge1\) such that
every \(g\in G\) with \(D_\infty(g,e)<\kappa D_X\) is joined to \(e\) by a
path in \(G\) of \(D_\infty\)-length at most
\(C_R D_\infty(g,e)\).
\end{enumerate}
These are optional regularity hypotheses on the primitive uniform comparison,
not clauses hidden inside GP1 and not inputs to the global role-allocation
theorem below.
A finite polynomial packing bound on \(G\) is one stronger sufficient
criterion for (FD); it is not required for compact completion.
\end{definition}

\begin{definition}[Intrinsic generated defect signature]
\label{p2:def:intrinsic-generated-defect}
Fix a previously frozen inventory \(P_{<R}\), a geometrically primary role
\(R\), and its frozen typed constructor grammar \(\Sigma_R\).  Let \(D\) be a
typed comparison diagram whose objects and arrows are generated by
\(P_{<R}\), \(\operatorname{Geo}_R\), and \(\Sigma_R\).  Suppose the same
frozen grammar contains a typed diagnostic constructor \(d_D\) that is fixed
before candidate selection and is invariant under unmarked role-isomorphisms.
The output
\[
  \delta_D:=d_D(D)
\]
is the \emph{intrinsic generated defect signature} of \(D\).  Its target sort
is part of the constructor data.  Depending on the role, \(\delta_D\) may be a
zero/nonzero order-swap distinction, kernel or cokernel, rank residual, parity
residual, support complement, stabilizer, boundary/cut residual, or another
typed invariant produced by a licensed diagnostic schema.

A nonzero \(\delta_D\) is not a GP2 failure.  It is generated content of the
role and must be represented faithfully.  If \(D\) is a candidate-independent
diagram of \(\operatorname{Geo}_R\), its defect status is fixed before the
candidate comparison.  If \(D\) is candidate-internal, its defect signature is
retained inside that candidate but is not promoted into \(P_{<R}\) and may not
serve as a result-dependent selection mark.  Mathematical residuals obtained
from diagnostic constructions not licensed by \(\Sigma_R\) create no physical
defect obligation.
\end{definition}

\begin{definition}[GP2 realization defect]
\label{p2:def:gp2-realization-defect}
Let \(\mathfrak M\) be GP1-admissible for the current role and let
\(\operatorname{Gen}_R(\mathfrak M)\) denote its complete typed closure under
the frozen constructor grammar, as in Postulate~\ref{p2:post:GP2}.  Its
\emph{GP2 realization defect} is the diagnostic tuple
\[
 \boxed{
 \mathfrak D_R^{\rm GP2}(\mathfrak M)
 :=
 \bigl(
   \mathfrak D_{\rm ins},
   \mathfrak D_{\rm omit},
   \mathfrak D_{\rm rel},
   \mathfrak D_{\rm type},
   \mathfrak D_{\rm clos}
 \bigr).}
\]
The components have the following meanings.
\begin{enumerate}[label=(\roman*),leftmargin=2em]
\item \(\mathfrak D_{\rm ins}\) records physically retained or selection-bearing
content not generated by the previously frozen inventory, the geometrically
primary role, or the frozen grammar, including candidate-only self-data
promoted chronologically so as to witness that candidate's own selection.
\item \(\mathfrak D_{\rm omit}\) records generated objects, operations, tests,
relations, sectors, constructors, or nonredundant generated images absent from
the candidate after all prior or internally generated identifications have
been applied.
\item \(\mathfrak D_{\rm rel}\) records failure to preserve a generated
relation or intrinsic defect signature, or the insertion of an ungenerated
relation.  Thus collapsing a generated distinction and imposing an ungenerated
quotient equality are both relation defects.
\item \(\mathfrak D_{\rm type}\) records generated content represented in the
wrong target sort or discharged only by an occurrence in an inequivalent
operation/object/state/effect slot.
\item \(\mathfrak D_{\rm clos}\) records a licensed constructor or operation
schema whose generated prerequisites are present but whose required
nonredundant typed output, partial evaluation, canonical composition, or output
image is absent or mistyped.  This component is diagnostic and may overlap the
omission or typing components.
\end{enumerate}
The candidate is \emph{zero-realization-defect} when all five components
vanish.  This definition introduces no second obligation set beside
\(\operatorname{Gen}_R(\mathfrak M)\); the latter remains the authoritative
candidate-internal generated closure.
\end{definition}

\begin{mdframed}[style=postulatebox,nobreak=false]
\begin{postulate}[GP2: exact and host-faithful generated realization]
\label{p2:post:GP2}
For every realized causal region \(\mathcal D\), GP2 is evaluated relative to
the previously frozen inventory and, during an open allocation, relative to
the geometrically primary role predicate.  Each candidate must contain exactly
the content generated in that candidate and represent it faithfully in its
generated typing.  Equivalently, with
Definition~\ref{p2:def:gp2-realization-defect}, GP2 admits exactly the
zero-realization-defect candidates:
\[
 \boxed{
 \mathrm{GP2}_R(\mathfrak M)
 \quad\Longleftrightarrow\quad
 \mathfrak D_R^{\rm GP2}(\mathfrak M)=0.}
\]
The following four coupled clauses are the authoritative expansion of that
vanishing condition.

\begin{enumerate}[label=(\roman*),leftmargin=2em]
\item \emph{No insertion.}  Let \(P_{<R}\) denote the content frozen before
the current role allocation.  No carrier datum, relation, operation,
constructor, test, law, valuation, scale, quotient, basis, projector, grading,
channel, cross-role identification, or other marked distinction may enter the
frozen inventory or the candidate-comparison criteria unless it is intrinsic
to the geometrically primary role or generated from \(P_{<R}\) by the frozen
host grammar.  A candidate's own internally generated carrier presentation,
operation law, relations, and tests are required inside that candidate, but
they are not promoted into \(P_{<R}\) and may not witness the candidate's
selection.

\item \emph{No omission.}  No generated object, state, relation, test, sector,
operation, constructor, or nonredundant generated image may be deleted from
the candidate in which it is generated.  Content may be identified at the GP2
exactification stage only when prior or candidate-internal generated data
already identify the apparently separate item as the same generated content.
GP3 does not delete an internally generated operation from a candidate; it may
select another complete exact candidate when a role-faithful comparison places
that candidate below the first.

\item \emph{Typed host faithfulness.}  Every item generated in a candidate is
represented in its generated sort and every generated distinction and relation
is preserved inside that candidate.  In particular, the comparison topology
and continuity of a generated action are part of its geometric relation data
and are preserved by an exact realization.  If a generated datum is typed as an internal object/carrier role,
an isomorphic occurrence only in the operation slot does not discharge that
object obligation; if it is operation-typed, it does not automatically become
a carrier state.  A generated subimage of an already constituted parent
geometric object sort inherits that object sort under its genuine inclusion;
retaining only an induced operation occurrence would be a retyping rather than
a restriction.  A structured subrealization or invariant subcarrier inherits every applicable
host item by genuine restriction.  Thus a retained operation \(A:H\to H\) restricts to a
subcarrier \(K\subseteq H\) only when \(A(K)\subseteq K\).  Replacing it by
\(P_KAP_K\) is a new compression/quantization construction and is admissible
only when that construction is itself generated and constituted.

\item \emph{Constructor-relative closure.}  Candidate generation is closure
under the typed constructors and operations licensed by the frozen host
signature and the geometrically primary role in that candidate, together with
the canonical compositions, partial evaluations, and output images required
by those schemata.  Mathematical
availability of a direct sum, tensor product, symmetric or exterior power,
isotypic decomposition, polynomial function, quotient, or theorem consequence
does not by itself constitute a physical constructor.  Such a construction
enters the generated closure only when its schema has prior geometric origin.
\end{enumerate}

Writing \(\operatorname{Gen}_R(\mathfrak M)\) for this complete typed
candidate closure, the exactness clause is
\[
 \forall\mathfrak M\in\operatorname{Adm}_R:\qquad
 \operatorname{Pres}_{\mathfrak M}(\mathfrak M)
 =\operatorname{Gen}_R(\mathfrak M).
 \tag{GP2-exact}
\]
This is candidate-internal exactness.  It does not assert the cross-candidate
monotonicity
\[
 \mathfrak M_1\preceq_R\mathfrak M_2
 \quad\Longrightarrow\quad
 \operatorname{Gen}_R(\mathfrak M_2)
 \subseteq\operatorname{Pres}_{\mathfrak M_1}(\mathfrak M_1),
 \tag{not GP2}
\]
which would freeze the complete self-structure of a rejected candidate before
GP3 and make leastness vacuous.

The word ``prior'' fixes the selection chronology: a proposed realization is
tested before any datum belonging only to that candidate may serve as a
selection mark or enter the chronological inventory.  Internally, the proposed
realization is nevertheless complete by \((\mathrm{GP2\text{-}exact})\).  For
globalization relations, the prior role geometry determines the generated
equality congruence.  A quotient
relation is not admissible merely because it becomes canonical after the
candidate has been constructed.  In particular, if the source integration
already satisfies FRC, an ungenerated quotient relation is excluded.  If the
source fails the closed compact-parent requirement, FRC first restricts the
admissible globalizations; a quotient class is forced only when that
restricted unmarked class is unique, and no member is selected when several
nonisomorphic classes remain.  Within such a compact globalization, however,
a central kernel lying wholly in the semisimple source factor does not
participate in compactifying the abelian covering directions.  When removing
that pure semisimple kernel leaves the same projected full lattice and a
compact paired realization, GP2 no-insertion excludes the additional
ungenerated semisimple quotient relation.  Lemma~\ref{p2:lem:no-gratuitous-semisimple-kernel}
proves exactly this narrower statement.  It does not compare or select among
genuinely different mixed lattice--center compactifications that survive the
same test.  The source integration and relation-exact globalization are made
precise in Proposition~\ref{p2:prop:admissible-globalization-operator} and
Theorem~\ref{p2:thm:relation-minimal-globalization}.

GP2 contains the following universal typed-generation clauses.
\begin{enumerate}[label=(\alph*),leftmargin=2em]
\item \emph{Canonical-image objectification.}  Let \(F:X\to Y\) be a
constituted, candidate-independent, host-natural constructor, with \(X\)
already constituted and \(Y\) an admissible host sort.  Suppose the canonical
image \(\operatorname{im}F\) admits a GP1-finite realization in its declared
target sort, is intrinsically characterized and nonredundant, and carries an
autonomous transformation law induced by already constituted host structure.
Its typing is then preserved as follows.
\begin{enumerate}[label=(\alph*),leftmargin=2em]
\item If the target is a parent geometric object sort already constituted
inside the \emph{same physical role}, then the finite image inherits that
object sort.  Subject to the absence of a generated autonomous role split,
Lemma~\ref{p2:lem:gp2-same-role-no-splitting} places its finite object copy in
that role's carrier.
\item If the target is an autonomous object sort distinct from the generating
physical role, then the image constitutes that separate object role and its
own carrier obligation; it is not silently inserted into the original role's
carrier.
\item If the target is operation-only, the image remains operation-typed.
\end{enumerate}
If a generated isomorphism identifies the image with already retained object
content of the same sort, the isomorphism is retained and no second independent
constitutive copy is added.  A mathematical module that also represents
infinitesimal operations therefore does not identify those operations with an
object copy.  The rule objectifies the canonical finite image, not the entire
ambient codomain \(Y\).  Any application that enlarges a finite carrier must
establish its object target typing either directly or by inheritance from an
already constituted same-role parent object sort; representation isomorphism
alone is insufficient.

\item \emph{Generated-role objectification.}  Whenever a generated separating
test spectrum carries its own faithful operation role and autonomous variance,
that role is constituted as an object of the physical category rather than
retained only as a parameter set attached to another object.

\item \emph{Relation-free composition.}  When two generated roles have no
prior generated cross-relation, their joint realization is their universal
relation-free composition.  Here a cross-relation means a prior generated
item in the relation slot that restricts the jointly admissible values---for
example a diagonal identification, graph constraint, cross-role quotient, or
selected tensor relation.  A total generated test, comparison map, duality
constructor, or role isomorphism does not by itself impose such a support
restriction; replacing it by the graph equation that its two arguments must
satisfy would be new inserted structure.  In the compact-Hausdorff parent
presentations used below, the relation-free joint carrier is the Cartesian
product.
\end{enumerate}

The primitive input to the present closure stage is the geometrically
constituted assertion \(\mathsf E_R\) of
Definition~\ref{p2:def:geometric-role-existence}.  It fixes the strong
role-geometric condition \(\operatorname{Geo}_R\), but no completed candidate
group, dimension, rank, global form, representation, scale, or multiplication
table.  GP2 faithfully realizes the constituted effectiveness and connected
local generation, while the exact candidate classifies the common commutator
branch of the constituted order-swap diagram, as stated in
Theorem~\ref{p2:thm:geometric-existence-exactification}.  Once GP3 has selected
a realization, all relations, typings, and distinctions in its exact generated
presentation are chronological facts and may not be deleted or retyped later.
\end{postulate}
\end{mdframed}

\begin{corollary}[No idle licensed constructor in a GP2-exact realization]
\label{p2:cor:no-idle-licensed-constructor}
Let \(\mathfrak M\) be GP2-exact for a frozen role grammar \(\Sigma_R\).  Let
\(C\in\Sigma_R\) be a typed constructor or operation schema, and suppose all of
its required inputs
\[
 x_1,\ldots,x_n\in\operatorname{Gen}_R(\mathfrak M)
\]
are generated in the declared source sorts.  If the schema requires a
nonredundant typed output
\[
 y=C(x_1,\ldots,x_n),
\]
then
\[
 \boxed{y\in\operatorname{Pres}_{\mathfrak M}(\mathfrak M)}.
\]
The same conclusion holds for required canonical compositions, partial
evaluations, and output images belonging to the licensed schema.  No conclusion
follows for a mathematically definable operation whose schema is absent from
\(\Sigma_R\).
\end{corollary}

\begin{proof}
By constructor-relative closure, the required typed output belongs to
\(\operatorname{Gen}_R(\mathfrak M)\).  If it were absent or retained only in
the wrong sort, then
\(\mathfrak D_{\rm clos}\ne0\), and respectively also
\(\mathfrak D_{\rm omit}\ne0\) or
\(\mathfrak D_{\rm type}\ne0\), contradicting GP2 exactness.  The final
sentence is exactly the mathematical-availability firewall in clause~(iv) of
Postulate~\ref{p2:post:GP2}.
\end{proof}

\begin{remark}[Intrinsic defect preservation versus realization defect]
\label{p2:rem:defect-preservation-vs-realization-defect}
An intrinsic generated defect signature \(\delta_D\) of
Definition~\ref{p2:def:intrinsic-generated-defect} need not vanish.  A
noncommutative order-swap distinction, a kernel, stabilizer, support
complement, or other generated residual may be the physical content that must
be preserved.  What GP2 forbids is a nonzero \emph{realization} defect: loss or
mistyping of that signature, insertion of ungenerated structure, or failure of
licensed closure.  GP2 therefore rejects a defective candidate; it does not
modify that candidate until the defect disappears.  If no GP1-finite,
GP2-exact candidate exists, the existing FRC obstruction logic applies, and
any repair category must itself have independent generated provenance.
\end{remark}

\begin{theorem}[GP2 generated reciprocal-role closure]
\label{p2:thm:gp2-generated-reciprocal-role-closure}
Let \(A\) be an already constituted physical role and let \(B_A\) be a
reciprocal separating test spectrum generated from the frozen physical test
or comparison grammar of \(A\).  Suppose:
\begin{enumerate}[label=(\roman*),leftmargin=2em]
\item the generated spectrum carries a faithful operation role with the
contravariant/opposite variance induced by the constituted pullback action;
\item that operation law and variance distinguish \(B_A\) as an autonomous
object sort rather than merely as a parameterization of the original carrier;
and
\item the frozen relation slot contains no prior generated joint-support
constraint restricting the simultaneously admissible pairs \((a,b)\in
A\times B_A\).
\end{enumerate}
Then GP2 generated-role objectification constitutes \(B_A\) as an autonomous
physical reciprocal role, and GP2 relation-free composition constitutes the
joint role as the universal relation-free product
\[
             A\times B_A.
\]
A generated comparison isomorphism
\(\beta_A:A\stackrel{\sim}{\longrightarrow}B_A\) is compatible with this
product and does not imply the graph restriction \(b=\beta_A(a)\).  Whenever
such a generated isomorphism is present, the exchange
\[
 \mathcal B_{\beta_A}(a,b)
 :=\bigl(\beta_A^{-1}(b),\beta_A(a)\bigr)
\]
is an involution:
\[
             \mathcal B_{\beta_A}^{\,2}=\operatorname{id}_{A\times B_A}.
\]
Consequently the legacy labels \((\mathrm{BR}_{\rm ind})\) and
\((\mathrm{IND}_{xp})\) are antecedent-certification statuses for this GP2
theorem, not additional constitutive principles.  Once their stated
antecedents have been proved by a reciprocal-spectrum construction, no further
physical postulate is required to objectify or independently compose the two
roles.  This theorem does not supply a dimensionful reciprocal normalization
or a local position--momentum interpretation; those remain separately typed.
\end{theorem}

\begin{proof}
Clauses (i)--(ii) are exactly the generated-role objectification antecedent of
GP2, so GP2 requires the generated spectrum to be represented in its own
object role rather than left as a parameter set.  Clause (iii) is exactly the
relation-free-composition antecedent.  GP2 no-insertion therefore forbids a
diagonal, graph, quotient, or other joint-support relation, while GP2
no-omission requires simultaneous realization of all generated values; in the
compact-Hausdorff presentation this is the Cartesian product.  A total
comparison isomorphism is a constructor relating the two role types, not a
relation restricting admissible pairs, so its retention does not alter the
product support.  Finally,
\[
 \mathcal B_{\beta_A}^{\,2}(a,b)
 =\mathcal B_{\beta_A}\bigl(\beta_A^{-1}(b),\beta_A(a)\bigr)
 =(a,b),
\]
which proves the involution.  The final status statement is therefore a
restatement of GP2 chronology: the labels certify proved antecedents and add
no new datum.
\end{proof}

\begin{lemma}[GP2 finite generated-image inheritance]
\label{p2:lem:gp2-finite-generated-image-inheritance}
Let \(R\) be an already constituted physical role with frozen parent geometry,
and let \(Y\) be a parent geometric object sort of \(R\) in the sense of
Definition~\ref{p2:def:parent-geometric-object-occurrence}.  Let
\[
 F:X\longrightarrow Y
\]
be a constituted, candidate-independent, host-natural constructor whose typed
inputs are already generated in \(R\).  Suppose
\(Z:=\operatorname{im}F\) is canonical, invariant, nonredundant, and admits a
GP1-finite realization.  Then GP2 generates \(Z\) in the inherited object sort
of \(Y\).  Any canonically induced operation on the same mathematical image
is a distinct operation-typed node and cannot discharge the object node.
When that operation is a required output of a licensed constructor, both
nodes are co-realized from the same root and no duplicate G0 occurrence is
needed.  The conclusion concerns the finite image \(Z\), not the entire
ambient parent object space \(Y\).
\end{lemma}

\begin{proof}
By GP2 constructor-relative closure, the output image of a constituted
constructor belongs to the generated closure in its declared target sort.
Since \(Y\) is already a parent geometric object sort of \(R\), the genuine
inclusion \(Z\hookrightarrow Y\) inherits that object typing; no new target
choice is made.  GP2 no-omission forbids deleting the nonredundant generated
image, and typed host faithfulness forbids replacing its object occurrence by
an isomorphic or induced operation occurrence.  GP1 finiteness makes the image
eligible for finite realization but does not promote the ambient section
space.  Hence \(Z\) is generated as a finite object image of the same parent
role.  Whether that object image is realized in \(R\)'s carrier or in a
separately generated autonomous role is then decided by the generated role
typing; Lemma~\ref{p2:lem:gp2-same-role-no-splitting} supplies the same-role
case.
\end{proof}

\begin{lemma}[GP2 same-role no-splitting for canonical internal object images]
\label{p2:lem:gp2-same-role-no-splitting}
Let \(R\) be an already constituted physical role with finite realized carrier
\(\mathcal H_R\), and let
\[
 F:X\longrightarrow Y
\]
be a constituted, candidate-independent, host-natural constructor whose inputs
are internal to \(R\).  Suppose the finite nonredundant image \(Z:=\operatorname{im}F\) has
already been typed as an object rather than merely as an operation parameter,
either directly by the frozen target typing or by
Lemma~\ref{p2:lem:gp2-finite-generated-image-inheritance}.  If the frozen prior
signature generates no autonomous role label, separate carrier, cross-role
variance, or separating test that distinguishes \(Z\) as a role different
from \(R\), then GP2 no-insertion forbids introducing such a role split.
Consequently the object copy of \(Z\) is internal object content of \(R\) and
must be represented in \(\mathcal H_R\).  Any canonically corresponding
operation copy remains a distinct operation-typed node; when licensed, it is
co-realized rather than independently actualized.
\end{lemma}

\begin{proof}
The object typing of \(Z\) is already established before this lemma is
invoked.  A new role \(R_Z\neq R\) would require at least one additional
generated distinction:
a role label, independent carrier assignment, cross-role variance, or a test
that separates the new role from \(R\).  By hypothesis none is present in the
frozen signature, so adding such a distinction violates GP2 no-insertion.
Hence \(Z\) remains object content of the already constituted role \(R\).
Within a structured realization, internal object/carrier content of one role
is realized in that role's carrier slot.  Typed host faithfulness therefore
places the object copy of \(Z\) in \(\mathcal H_R\).  The conclusion concerns
the object copy only: if the same mathematical module also parameterizes or
realizes infinitesimal operations, GP2 sort faithfulness retains that
operation node in the operation slot rather than identifying states with
operators.  The universal bridge normal form,
Theorem~\ref{p2a:thm:universal-corealization-bridge-normal-form}, retains the
constructor incidence between the two typed nodes without requiring a second
root.
\end{proof}

\begin{remark}[Scope of the same-role lemma]
\label{p2:rem:gp2-same-role-no-splitting-scope}
The lemma does not objectify arbitrary theorem consequences and does not turn
every parent operation into carrier content.  Its hypotheses begin only after
a constructor has entered the frozen physical grammar and object typing has
been established independently, including by finite generated-image
inheritance from an already constituted parent object sort.  In
particular, a mathematically available parent multiplication, tensor product,
or polynomial map creates no carrier obligation unless its constructor schema
has been constituted at that stage.
\end{remark}


\begin{theorem}[GP2 generated-owner normalization inheritance]
\label{p2:thm:gp2-generated-owner-normalization-inheritance}
Let \(R\) be an already constituted physical role and let a generated parent
relation \(\mathcal R_P\) carry a normalization datum \(\nu_P\) in a declared
typed owner slot \(N_P\).  Let
\[
 F:P\longrightarrow C
\]
be a constituted, candidate-independent, host-natural constructor whose child
\(C\) has already been proved to be a generated descendant in the same
relation/owner sort.  Assume that the constructor transports
\(\mathcal R_P\) to \(\mathcal R_C\), transports the owner slot to \(N_C\),
and that the frozen generated grammar contains no independent scale,
constitutive tensor, clock, grading weight, calibration, central weight, or
other nontrivial endomorphism that deforms that inherited normalization slot.
Then the child normalization is the transported parent normalization,
\[
 \boxed{\nu_C=F_\nu(\nu_P).}
\]
In particular, when the inherited owner line is canonically the same line and
\(F_\nu=\mathrm{id}\),
\[
 \boxed{\nu_C=\nu_P.}
\]
A proposed replacement \(\nu_C=\lambda\nu_P\) with \(\lambda\neq1\) is
admissible only if the frozen grammar independently generates the corresponding
nontrivial deformation of the normalization owner.
\end{theorem}

\begin{proof}
The same-role and target typing are antecedents, not conclusions, of this
theorem.  By GP2 constructor-relative closure, the transported relation and its
owner slot belong to the generated closure in their declared sorts.  Typed host
faithfulness therefore requires a genuine restriction/pullback/pushforward of
the parent relation rather than replacement by an isomorphic but independently
normalized copy.  If \(F_\nu=\mathrm{id}\), changing the inherited coefficient
by \(\lambda\neq1\) either changes an already generated normalized relation or
introduces the endomorphism \(\lambda\mathbf1_{N_P}\).  The former violates
exact host faithfulness and the latter violates GP2 no-insertion unless that
deformation is itself generated.  Hence the inherited value is forced.
\end{proof}

\begin{theorem}[Unique neutral support forces support ancestry]
\label{p2a:thm:unique-neutral-support-ancestry}
Let \(H\) be an already constituted host role and let
\(\mathcal N_H\) be its unique nonzero irreducible object in a frozen neutral
support sort \(\mathsf{Supp}^0_H\), in the precise sense that every nonzero
generated \(H\)-invariant subobject of that sort equals \(\mathcal N_H\).
Let \(Q\) be a generated nonzero neutral descendant with a typed, nonzero,
host-equivariant support map
\[
 j_Q:\operatorname{Supp}(Q)\longrightarrow\mathcal N_H,
\]
and suppose the frozen grammar supplies no independent carrier, autonomous
support owner, multiplicity label, cross-role variance, or separating test that
distinguishes a second neutral host support.  Then
\[
 \boxed{\operatorname{im}j_Q=\mathcal N_H,}
\]
and \(Q\) is a same-support descendant of \(H\).  Introducing a second copy of
the neutral host support has nonzero GP2 insertion/type defect.

This conclusion establishes support ancestry only.  It does not identify a
principal differential relation, kinetic normalization, clock, or
characteristic speed.
\end{theorem}

\begin{proof}
Host equivariance makes \(\operatorname{im}j_Q\) an \(H\)-invariant object in
\(\mathsf{Supp}^0_H\), and nonzeroness excludes the zero object.  The defining
uniqueness of \(\mathcal N_H\) therefore gives
\(\operatorname{im}j_Q=\mathcal N_H\).  A separately labelled neutral support
would add an autonomous owner, multiplicity, variance, or distinguishing test
absent from the frozen presentation and hence violate GP2 no-insertion or
typed host faithfulness.  Nothing in this argument refers to a differential
operator or its principal symbol, which proves the scope sentence.
\end{proof}

\begin{corollary}[Unique unlabelled neutral tensor-unit support]
\label{p2a:cor:unique-neutral-tensor-unit-support}
In the current finite complex-linear carrier grammar, define the unlabelled
neutral support sort to contain irreducible objects on which the constituted
host operation acts trivially and which carry no charge, grading, variance,
multiplicity, or species label.  Its unique nonzero structured isomorphism
class is the tensor unit \(I=\mathbb C\); on the dagger-fixed real scalar
slice it is \(I_{\mathbb R}=\mathbb R\).  Consequently any nonzero generated
primitive unlabelled neutral scalar support with the canonical tensor-unit
map satisfies the support-ancestry theorem above.  A second neutral copy is
not another member of this unlabelled irreducible sort: it requires a
multiplicity/species owner and belongs to an enlarged typed branch.
\end{corollary}

\begin{proof}
An irreducible complex representation with trivial host action is
one-dimensional and hence isomorphic to \(\mathbb C\) with its tensor-unit
action.  The dagger-fixed self-adjoint scalar line is \(\mathbb R\).  A direct
sum of two copies is reducible and its distinction requires a multiplicity or
species label.  Thus there is one nonzero irreducible structured isomorphism
class in the stated unlabelled sort, and
Theorem~\ref{p2a:thm:unique-neutral-support-ancestry} applies to its nonzero
generated descendants.
\end{proof}

\begin{proposition}[Same support alone does not determine speed]
\label{p2a:prop:neutral-support-speed-countermodel}
On one and the same neutral real line support over a stationary host, the
operators
\[
 P_c=\partial_t^2-c^2\Delta_h,
 \qquad c>0,
\]
have identical carrier and support typing but distinct characteristic cones
when their positive constants differ.  Therefore uniqueness of neutral support
cannot by itself imply \(c=c_{\rm op}\).
\end{proposition}

\begin{proof}
The principal symbol is
\(\tau^2-c^2h^{-1}(\xi,\xi)\).  For \(c_1\ne c_2\), choose
\(\xi\ne0\) and \(\tau=c_1\|\xi\|_{h^{-1}}\); the symbol for \(c_1\)
vanishes while that for \(c_2\) does not.  Thus the characteristic sets differ
although the support is unchanged.
\end{proof}

\begin{theorem}[Undeformed host-principal inheritance]
\label{p2a:thm:undeformed-host-principal-inheritance}
Let the operational host relation have frozen quadratic principal form
\[
 q_{\rm op}(\tau,\xi)
 =\tau^2-c_{\rm op}^2h^{-1}(\xi,\xi),
 \qquad c_{\rm op}>0,
\]
in a declared principal-relation owner slot.  Let \(Q\) be a generated
descendant whose support ancestry has been established independently, for
example by Theorem~\ref{p2a:thm:unique-neutral-support-ancestry}.  Suppose a
licensed typed principal constructor proves that the second-order kinetic
relation of
\(Q\) is a same-owner descendant of \(q_{\rm op}\), and suppose the complete
frozen grammar generates no independent clock, spatial cometric, constitutive
tensor, medium response, anisotropy, disformal map, or other principal-order
deformation.  Additional generated mass, potential, mixing, and interaction
terms may be present only at differential order below \(2\).  On the claimed
propagation domain assume additionally that \(P_Q\) is a hyperbolic finite-
propagation system for which the wave-front boundary is governed by its real
characteristic cone.

Then there is an everywhere invertible zeroth-order owner endomorphism
\(A_Q\), independent of \((\tau,\xi)\), such that
\[
 \boxed{
 \sigma_2(P_Q)(\tau,\xi)
 =q_{\rm op}(\tau,\xi)A_Q.}
\]
For a single unlabelled real neutral line with positive kinetic
normalization, \(A_Q=a_QI_Q\) with \(a_Q>0\).  More generally the invertible
factor may change normalization or mix an already generated internal
multiplicity, but it cannot change the characteristic set.
Consequently
\[
 \boxed{
 \operatorname{Char}(P_Q)=\{q_{\rm op}=0\},
 \qquad c_{{\rm front},Q}=c_{\rm op}.}
\]
The overall invertible factor is normalization of the whole principal
equation and does not change its characteristic set.  A different relative
temporal/spatial coefficient is admissible only on a branch carrying an
independently generated principal deformation.
\end{theorem}

\begin{proof}
The typed principal constructor, not support isomorphism alone, establishes the
same principal owner.  Apply
Theorem~\ref{p2:thm:gp2-generated-owner-normalization-inheritance} to the
relative temporal/spatial normalization of that owner.  In the absence of a
generated deformation, replacing \(c_{\rm op}\) by another positive speed
would insert a nontrivial endomorphism of the normalized principal relation.
Hence the temporal and spatial terms retain their common relative coefficient;
only a \((\tau,\xi)\)-independent invertible owner factor \(A_Q\) remains.
On an unlabelled real line, positivity reduces it to \(a_QI_Q\) with
\(a_Q>0\).  Lower-order terms do not enter \(\sigma_2(P_Q)\), and an
invertible factor does not alter the vanishing of the principal determinant.
The displayed characteristic identity follows algebraically.  The stated
hyperbolic finite-propagation hypothesis then identifies that cone with the
wave-front boundary and gives the front-speed identity.
\end{proof}

\begin{corollary}[Joint saturation of undeformed neutral descendants]
\label{p2a:cor:joint-neutral-principal-saturation}
Let \(Q_1,\ldots,Q_N\) be simultaneous nonzero neutral descendants of the
same host.  Assume for each \(Q_i\) the support and principal ancestries of the
preceding theorems.  In the absence of a generated principal deformation, the
joint GP2-exact locus is contained in
\[
  \boxed{
  \left\{(Q_i,c_i)_{i=1}^{N}:
   \operatorname{Supp}_H(Q_i):=\operatorname{im}j_{Q_i}
                  =\mathcal N_H,
   \ c_i=c_{\rm op}\ \text{for all }i\right\}.}
\]
Here \(\operatorname{Supp}_H(Q_i)\) denotes the \emph{retained
host-visible neutral support}, not an assertion that the entire source
domain \(\operatorname{Supp}(Q_i)\) is injectively identified with it.
If a claim requires the stronger source-domain equality, it must also
establish \(\ker j_{Q_i}=0\) (for example from an independently
certified irreducible unlabelled source-support domain); any genuinely
generated kernel or multiplicity remains retained until then.
Thus individually writable assignments \(c_i\ne c_{\rm op}\) are removed by
their simultaneous owner/normalization defect.  A medium or effective branch
with an independently generated constitutive tensor is a different typed
joint component and is not excluded.
\end{corollary}

\begin{corollary}[Later photon and primitive-scalar rates]
\label{p2a:cor:later-light-scalar-speed-inheritance}
If a later vacuum photon channel or primitive unlabelled neutral scalar is
proved to be an undeformed principal descendant on \(\mathcal N_H\), then its
front speed is \(c_{\rm op}\).  In particular, after a later theorem identifies
the physical vacuum light channel \(\gamma\),
\[
 c_{\rm light}^{\rm th}:=c_{{\rm front},\gamma}=c_{\rm op}.
\]
An empirical identification of an observed signal with that channel is a G0
measurement bridge in the already fixed spatial/time calibration; it does not
introduce a second structural speed.  Finite-frequency group or phase
velocities and medium-modified effective speeds are not covered.
\end{corollary}

\begin{remark}[Owner-typing firewall]
\label{p2:rem:generated-owner-normalization-firewall}
The theorem does not prove ancestry.  Numerical equality, equal dimension,
representation isomorphism, or the mathematical availability of a canonical
functor does not identify two independently typed owner slots.  The same owner
or relation sort must first be established by the frozen constructor grammar,
including by the generated-image and same-role results above.  Nor does the
theorem determine an external unit calibration or a genuinely independent
cross-role bridge.  Its content is only that an already inherited normalized
relation cannot acquire a second free normalization without a generated
deformation owner.  Theorem~\ref{p2a:thm:unique-neutral-support-ancestry}
proves support ancestry on its stated domain but, by
Proposition~\ref{p2a:prop:neutral-support-speed-countermodel}, still does not
prove principal ancestry.  That second typing obligation is exactly the
hypothesis of
Theorem~\ref{p2a:thm:undeformed-host-principal-inheritance}.
\end{remark}

\begin{lemma}[GP2 canonicity / no-new-choice criterion]
\label{p2:lem:gp2-canonical-construction}
Let \(C\) be an already licensed typed constructor in the frozen grammar
\(\Sigma\), with all required inputs present.  If all admissible outputs of
\(C\) lie in one orbit of the
already generated equivalence relation and no generated test distinguishes an
additional choice of output, then retaining its required output introduces no new datum for
GP2(i).  Conversely, if inequivalent outputs remain and selecting one requires
data not generated by \(\Sigma\), that selection is an inserted datum.
Nontrivial automorphisms of a canonical output are allowed; canonicity means
absence of a new marked choice, not rigidity of the output groupoid.
\end{lemma}

\begin{proof}
In the first case all outputs represent the same constituted construction
relative to the equivalence already present in the frozen signature, so no
new distinction is added by retaining the licensed required output.  This
criterion does not license \(C\); that premise is explicit in the statement.
In the second case
a choice of orbit separates alternatives that the prior generated data do not
separate, which is exactly GP2(i) insertion.
\end{proof}

\begin{theorem}[Geometric defect repair on retained data]
\label{p2a:thm:geometric-defect-repair-admission}
Let $\Sigma$ be a frozen constituted presentation and let $D$ be a retained
jointly occurring/generated diagram. The following are sufficient
no-new-physical-premise repair rules.
\begin{enumerate}
\item The exact object $X$, its physical sort and its retained relations are
unchanged, and the replacement invariant predicate is proved true of that
same $X$ from the retained data. This is mathematical reclassification.
If a single strongest predicate in a predeclared preorder is additionally
claimed, its true set must have a greatest equivalence class. A unique
maximal class suffices when every true class lies below a maximal class,
as in a finite true poset, or when the true set is upward directed.
Without such a condition the conclusion is only unique maximality.
\item A target sort $\mathsf Q$ is predeclared and the proposed repair is the
required output $q$ of a licensed finite word with
$\operatorname{Orig}_\Sigma(w;D,\mathsf Q)$. Every step has the required
typing, output retention and licensed reuse, and the resulting presentation
is GP1-admissible and GP2-exact. Naturality, uniqueness and universal
properties identify this output but do not replace its constructor license.
\end{enumerate}
In the second case, the rooted generated-image theorem gives the existing
obstruction/co-realization alternative. A word whose output is an unresolved
family retains that family without selecting a member. These sufficient
rules do not assert a classification of all conceivable repair operations
outside the displayed domain.
\end{theorem}
\begin{proof}
In the first case the proved predicate already holds in every model of the
retained data, so no new physical object or distinction is introduced.
For the order assertion, if each true class lies below a maximal class and
that class is unique, every true class lies below it. If the true set is
directed, a common upper bound of a maximal class and any other true class
must equal the maximal class; hence it is greatest.
In the second case the output is the value of the already licensed word.
A uniquely specified value can be named by definitional extension. If the
output is specified only up to an already generated equivalence, retain it at
that equivalence level, without a marked representative or a claim of unique
isomorphism; nontrivial automorphisms remain allowed.
Induction over the word and GP2 output retention give the generated
descendant, with its original type and all incidences. A family-valued output
is not a section of that family. Neither a uniqueness proof nor permission
to use an ambient object supplies a missing word.
\end{proof}

\begin{corollary}[No bare-permission relaxation]
\label{p2a:cor:no-bare-permission-relaxation}
Permission to use a weaker ambient object, or a proof that such an object is
mathematically unique, does not supply a frozen constructor-origin
certificate. A same-object reclassification must be proved true of that
object; a new physical output requires its licensed generated word or an
explicit additional occurrence/constitution premise on the stated branch.
\end{corollary}

\begin{theorem}[GP2 exact-fidelity endpoint criterion]
\label{p2:thm:gp2-exact-fidelity-endpoint}
Let \(F_r\), \(0\le r\le r_{\max}\), be a family of realizations of one
already generated transfer, embedding, or fidelity relation, and let
\(\mathcal R_{\rm gen}\) be that relation with its generated normalization.
Assume
\[
 \mathcal R(F_r)=\mathcal R_{\rm gen}
 \quad\Longleftrightarrow\quad
 r=r_{\max},
\]
and assume that every \(r<r_{\max}\) either omits part of
\(\mathcal R_{\rm gen}\) or requires an additional loss/defect channel not
generated in the frozen signature.  Then GP2 exactness and typed host
faithfulness admit only
\[
 r=r_{\max}.
\]
No separate endpoint-selection rule is used.

The conclusion does not follow from boundedness, injectivity, positivity, or
the absence of a preferred interior number alone.  In particular, if
\(F_r=rF\) is injective for every \(r>0\), faithfulness of the underlying map
does not distinguish the interior values.  Nor does the theorem apply when
\(r\) is a contingent state value, response amplitude, coupling, mass ratio,
or trajectory variable rather than a constitutive fidelity coefficient.
\end{theorem}

\begin{proof}
For \(r<r_{\max}\), the first alternative violates GP2 no-omission and
relation exactness, while the second violates GP2 no-insertion.  The endpoint
realizes the generated normalized relation without either defect and is
therefore the only GP2-admissible member of the stated family.  Bare
injectivity preserves only the kernel distinction and supplies no equality
with \(\mathcal R_{\rm gen}\), proving the converse boundary.
\end{proof}

\begin{proposition}[GP1--GP2 exact host closure]
\label{p2:prop:gp1-gp2-bounded-generated-exactness}
Fix a frozen constituted signature \(\Sigma\), a GP1-admissible structured
realization \(\mathfrak M\), and a typed slot \(q\).  Let
\(\mathcal G_{\Sigma,q}\) be the \(q\)-slot of the constructor-relative GP2
generated closure and let
\[
 \operatorname{Adm}_{\Sigma,q}
 :=\{Y\in\mathcal G_{\Sigma,q}:
 Y\text{ is not identified as redundant by prior generated data}\}.
\]
If a GP1-finite/bounded, GP2-faithful realization exists, GP2 exactness gives,
at generated presentation-item/type-support level,
\[
 \operatorname{Pres}^{\rm gen}_q(\mathfrak M)
 =\operatorname{Adm}_{\Sigma,q}.
\]
The equality is typed: it neither identifies multiplicity spaces nor concrete
operator coefficients, and it does not turn arbitrary theorem consequences or
mathematically available composites into new frozen items.  If the required
typed generated closure has no GP1-admissible faithful realization, the
finite/bounded realization class is empty and the role is FRC-obstructed.

If a registered generated boundary \(\mathsf P_{\rm gbd}(q,N_q)\) is itself
part of the frozen signature, the same statement applies after replacing
\(\mathcal G_{\Sigma,q}\) by
\(\mathcal G_{\Sigma,q}\cap\Gamma_q^{\le N_q}\).
\end{proposition}

\begin{proof}
GP2 no-insertion excludes retained items without prior generated origin.  GP2
no-omission requires every nonredundant generated item, and typed host
faithfulness preserves its constituted sort and relations.  These give the
two inclusions in the displayed equality.  The final sentence is simply the
failure of GP1-admissibility for the exact faithful closure.  When a generated
boundary belongs to the signature, it restricts the constructor-relative
closure before the same argument is applied.
\end{proof}

\begin{corollary}[Finite/bounded exact-and-faithful admissible family]
\label{p2:cor:gp1-gp2-admissible-family}
For a fixed frozen signature, GP1 followed by GP2 leaves only structured
realizations that are finite or totally bounded, generated-exact, and
host-faithful.  GP1+GP2 do not numerically minimize a generated boundary or
choose among residual presentation moduli.  If one structured isomorphism
class remains, GP3 has nothing further to compare.  If several survive, GP3
retains a unique least class only when role-faithful structural comparisons
place one below all the others; otherwise the irreducible family remains
unresolved.
\end{corollary}

\begin{proof}
The first sentence is Proposition~\ref{p2:prop:gp1-gp2-bounded-generated-exactness}
combined with the stagewise GP1 restrictions.  The remaining alternatives
are exactly the GP3 least-class rule in
Postulate~\ref{p2:post:GP3}; neither GP1 nor GP2 supplies a
numerical minimizer or a selector on residual unmarked moduli.
\end{proof}

\begin{remark}[Same-stage divisibility reading of connectedness]
\label{p2:rem:divisibility-signature}
After GP1--GP2 faithfully realize the constituted local-generation clause, the
classified status ``non-isolated connected exact-operation role'' may
equivalently be expressed algebraically as
same-stage divisibility:
\[
 \forall s,\mathcal D\ \forall g\in G_{R,s}(\mathcal D)\
 \forall n\ge2\ \exists h\in G_{R,s}(\mathcal D):\ h^n=g.
\]
The root must lie in the same closed exact-operation parent at the same
realized stage; deferring it to a later enlargement would reopen the finite
group loophole.  On compact groups this is equivalent to connectedness by
Lemma~\ref{p2:lem:compact-divisible-connected}.  No nontrivial finite group
can satisfy it.  For example, the binary tetrahedral group \(2T\) contains elements of
order \(4\) and none of order \(8\), so an order-\(4\) element has no square
root in the same group.
\end{remark}

\begin{lemma}[Compact divisible equals compact connected]
\label{p2:lem:compact-divisible-connected}
A compact Hausdorff group \(G\) is connected if and only if it is divisible.
\end{lemma}

\begin{proof}
Suppose first that \(G\) is compact and connected.  Write it as an inverse
limit of compact Lie quotients \(G_i\), each connected as a continuous image
of the connected \(G\).  For fixed \(g\in G\) and
\(n\ge2\), let
\[
 R_i=\{h_i\in G_i:h_i^n=g_i\}.
\]
Each \(R_i\) is nonempty and compact because the power map is surjective on a
compact connected Lie group: by the maximal-torus theorem every \(g_i\) is
\(\exp X\) for some \(X\), whence \(h_i=\exp(X/n)\) satisfies \(h_i^n=g_i\).  Pulling these conditions back to \(G\) gives a
family of nonempty closed subsets with the finite-intersection property: any
finite set of quotient conditions is dominated by one common quotient, where
a root exists.  Compactness of \(G\) therefore gives an \(h\) lying in every
member of the family, so \(h^ng^{-1}\in\bigcap_iN_i\), where \(N_i\) is the
kernel of the \(i\)th quotient.  The compact Lie quotients separate the points
of \(G\), so that intersection is trivial and \(h^n=g\).  Hence
\(G\) is divisible.

Conversely, quotients of divisible groups are divisible, while
\(G/G_0\) is profinite, since in a compact group the identity component is the
intersection of the open subgroups, so the quotient is a compact totally
disconnected group and hence an inverse limit of finite groups.  Every nontrivial finite group \(F\) fails
divisibility: for \(g\ne e\) and \(n=|F|\), one has \(h^n=e\) for all
\(h\in F\).  Thus every finite quotient of \(G/G_0\) is trivial, so the
profinite group \(G/G_0\) is trivial and \(G=G_0\) is connected.
\end{proof}

\begin{mdframed}[style=postulatebox,nobreak=false]
\begin{postulate}[GP3: irreducible and least structured realization]
\label{p2:post:GP3}
GP3 compares the complete GP1-admissible, candidate-internally
GP2-exact-and-host-faithful structured realizations of one geometrically
primary role.  Each candidate retains its whole generated carrier--operation--
constructor--relation--test structure.  The comparison is between those whole
candidates; carrier and operations are never minimized independently.

A \emph{role-faithful structural comparison} from \(\mathfrak M_2\) to
\(\mathfrak M_1\) is a generated comparison which:
\begin{enumerate}[label=(\alph*),leftmargin=2em]
\item preserves every item already frozen before the current allocation;
\item preserves the effective role occurrence and every intrinsic comparison,
composition, and order-swap diagram required by
\(\operatorname{Geo}_R\); and
\item has complete GP1-admissible, GP2-exact-and-faithful realizations at both
endpoints.
\end{enumerate}
Write \(\mathfrak M_1\preceq_R\mathfrak M_2\) when such a comparison, or a
finite chain of them, exhibits \(\mathfrak M_1\) as containing no more
independent role structure than \(\mathfrak M_2\), modulo structured
isomorphism.  The comparison need not realize
\(\operatorname{Gen}_R(\mathfrak M_2)\) as a subpresentation of
\(\mathfrak M_1\); those are the internally exact self-operations of the
source candidate, not frozen target obligations.  If a source item realizes a
role diagram required by \(\operatorname{Geo}_R\), however, the target must
realize that diagram exactly.  Thus genuinely richer geometrically primary
directions, brackets, relations, or tests cannot be discarded.

Depending on the typed domain, witnesses of role-faithful comparison include
passage to an effective image, construction of a complete exact candidate from
a closed connected local subrole, quotient by an already generated
congruence, deletion of a factor unused by the geometrically primary role, and
replacement by a canonically equivalent presentation.  An internal
subrealization claim still requires genuine restriction of every operation it
claims to retain; compression is a new construction and is admissible only
when generated.  At first allocation, however, a comparison may instead
construct a different complete exact candidate of the same role.  This is
candidate selection, not deletion from the source candidate.  A concrete
domain may prove additional comparison lemmas without introducing a new
minimization principle.

A realization \(\mathfrak M\) is \emph{irreducible} when no proper
role-faithful structural comparison from \(\mathfrak M\) remains inside the
GP1/GP2-admissible class.  It is \emph{least} when
\[
 \mathfrak M\preceq_R\mathfrak N
 \qquad\text{for every GP1/GP2-admissible }\mathfrak N,
\]
up to structured isomorphism.  If a unique least isomorphism class exists, GP3
retains it.  If no unique least class exists---whether because distinct
irreducibles are incomparable or because nonisomorphic minima are
preorder-equivalent---GP3 retains the unresolved irreducible family and
selects none.  On such a family, only the previously frozen inventory, the
common role-geometric content, and candidate-generated items proved common to
every member are frozen.  If the finite-and-bounded exact
faithful realization class is empty, the role is FRC-obstructed and any repair
must belong to a separately generated repair category.

If GP3 selects a least candidate \(\mathfrak M_R\), closure updates the
chronological inventory by
\[
 P_R=P_{<R}\cup
 \operatorname{Pres}_{\mathfrak M_R}(\mathfrak M_R)
 =P_{<R}\cup\operatorname{Gen}_R(\mathfrak M_R).
\]
Every later GP application must preserve all of \(P_R\); candidate-internal
exactness is no longer a license to reopen the selected multiplication,
topology, relation, test, or typing.  Thus the first-allocation comparison of
whole candidates and the later fixed-signature comparison are two chronological
uses of the same leastness principle, not two different GPs.

No Lie dimension, rank, root datum, named group, carrier dimension, operation
count, truncation label, Casimir, or numerical complexity functional is a GP3
selector unless it has already entered the role geometry or frozen inventory
as part of a generated structural relation.  GP3 neither changes the
GP2-exactified intrinsic commutator branch nor fixes a dimensionful scale.
After a least realization is closed, its complete candidate-internal generated
presentation becomes constituted and frozen.  For an unresolved irreducible
family, only the common content specified above is frozen; theorem
consequences remain available for inference without automatically becoming new
signature symbols.
\end{postulate}
\end{mdframed}


\begin{theorem}[Generated least-carrier closure]
\label{p2a:thm:generated-least-carrier}
Let $X$ be an already constituted typed ambient object and let $S$ be the
complete datum produced in $X$ by a licensed frozen constructor from already
occurring roots.  Let $\mathcal C_X(S)$ be the GP1/GP2-admissible class of
complete, host-faithful realizations of the downstream carrier role for $S$.
Assume that the frozen grammar generates a candidate
$\operatorname{Car}_X(S)\in\mathcal C_X(S)$ with the following universal
property: for every $Y\in\mathcal C_X(S)$ there is a generated role-faithful
structural comparison
\[
  \operatorname{Car}_X(S)\preceq_R Y,
\]
and any two generated candidates with this property are structurally
isomorphic.  Then $\operatorname{Car}_X(S)$ is the unique GP3-least carrier of
the complete datum $S$.  It is retained without an additional carrier
selector.

The conclusion is only a carrier statement.  It applies, on the domains where
the corresponding constructor and universal property have actually been
generated, to familiar cases such as subgroup or $*$-algebra closure,
exact kernel/fibre carriers, generated quotient images, fixed-point or
stabilizer roles, and primitive symmetry-eigenspace carriers.  If no such
least admissible carrier exists, if several incomparable faithful minima
remain, or if the downstream role consumes extra member/state/history data
not contained in $S$, this theorem abstains and GP3 retains the appropriate
unresolved family.
\end{theorem}

\begin{proof}
GP2 supplies $\operatorname{Car}_X(S)$ as generated typed content and makes
its claimed operations, relations and host incidence exact.  By hypothesis it
role-faithfully compares to every other member of $\mathcal C_X(S)$, so it is
least in the preorder of Postulate~\ref{p2:post:GP3}.  The uniqueness clause
makes its least isomorphism class unique, and GP3 therefore retains it.  An
independent carrier choice would add structure not generated by the frozen
constructor grammar.  The abstention clauses are exactly the unresolved-family
clauses of GP3.
\end{proof}


\section{Least tests and same-owner dynamical closure}
\label{paper2a:sec:organized-dynamics}

\subsection{Least tests closed under the complete owned dynamics}
\label{p2a:sec:least-dynamic-tests}
\begin{corollary}[Generated least invariant test algebra]
\label{p2a:cor:least-dynamic-test-algebra}
Let $\mathcal A$ be an already licensed finite-dimensional unital star
algebra, $\mathcal T$ its declared initial tests, and
$\{\mathcal L_a\}$ the complete owned family of linear star-preserving
generators on $\mathcal A$. Assume their application and unital star
closure are licensed constructors in this role. Define
\[
 B_0=\operatorname{alg}^{*}(I,\mathcal T),\qquad
 B_{n+1}=\operatorname{alg}^{*}
       \left(B_n\cup\bigcup_a\mathcal L_a(B_n)\right)
       \subseteq\mathcal A.
\]
The sequence stabilizes after at most
$\dim_{\mathbb C}\mathcal A-\dim_{\mathbb C}B_0$ strict increases.
Its limit $B_*$ is the unique least invariant unital star subalgebra
containing the tests. If this inclusion is the generated role-faithful
comparison, Theorem~\ref{p2a:thm:generated-least-carrier} retains this
carrier without an additional accessibility selector.
\end{corollary}
\begin{proof}
The $B_n$ are nested linear subspaces of a finite-dimensional space.
Every strict inclusion increases dimension. Once $B_{n+1}=B_n$, every
generator preserves $B_n$, and later iterations agree. Any invariant unital
star algebra containing the tests contains $B_0$ and, inductively, every
$B_n$. It thus contains the stabilized algebra. This proves leastness and
uniqueness, and supplies the stated carrier universal property.
\end{proof}
The domain includes all owned generators, also at all times if the family
varies with time. It does not enlarge a declared host to its full
presentation envelope or select an unresolved Hamiltonian or state.
For a consumer using only linear expectations, the least invariant linear
operator space can suffice in place of an algebra. A stress-response
consumer seeds its complete energy and source-jet tests: a central
$c(g)I$ is invisible to commutators but can have nonzero metric derivatives.
Paper~IV gives the common-flow and finite source realizations on these
declared domains.

\begin{proposition}[Surjective state restriction for invariant finite tests]
\label{p2a:prop:v257-observable-process-target}
Let $\mathcal A$ be an already licensed finite-dimensional unital complex
$C^*$-algebra. Let $T_t=e^{t\mathcal L}$, $t\ge0$, be an already owned
positive-unital Heisenberg semigroup, and let $B_*$ be the least invariant
unital star subalgebra generated by the retained tests as in
Corollary~\ref{p2a:cor:least-dynamic-test-algebra}.
Assume that closure, evaluation restriction, target retention and output
reuse are licensed in this role. Then
\[
 q:S(\mathcal A)\longrightarrow S(B_*),\qquad
 q(\omega)=\omega|_{B_*}
\]
is surjective. For the state process $F_t\omega=\omega\circ T_t$,
there is a unique continuous affine descended semigroup on all of $S(B_*)$,
specified by
\[
 [\bar F_t(q\omega)](b)=\omega(T_tb),\qquad b\in B_*.
\]
It does not choose an extension state or assert that every state is prepared.
\end{proposition}
\begin{proof}
The restriction image is compact convex. A state $\eta$ outside it could
be strictly separated by a self-adjoint $b\in B_*$, giving
$\eta(b)>\sup_{\omega\in S(\mathcal A)}\omega(b)=\lambda_{\max}(b)$.
Since $\lambda_{\max}(b)1-b\ge0$ belongs to $B_*$, positivity and
normalization of $\eta$ contradict this inequality. Thus restriction is
onto. Invariance implies $T_tB_*\subseteq B_*$, so the displayed formula
is independent of the extension $\omega$ and is a state. Surjectivity
gives uniqueness on the entire target. The semigroup law follows by
composition. Finite-dimensional linear evaluation makes this affine
semigroup continuous. This proves the conclusion for this constructed target,
not for an arbitrary larger readout codomain.
\end{proof}

\begin{corollary}[Minimal linear observations and finite evolution]
\label{p2a:cor:v257-minimal-linear-observations}
For a one-generator role retaining only linear expectations of the
declared self-adjoint tests $E_i$, assume the licensed closure
\[
 W_0=\operatorname{span}_{\mathbb C}\{1,E_i\},\qquad
 W_{n+1}=W_n+\mathcal L(W_n)
\]
has retained output $W_*$. It stabilizes after at most
$\dim\mathcal A-\dim W_0$ strict increases and equals
$\operatorname{span}_{\mathbb C}\{1,T_tE_i:t\ge0,i\}$.
Restriction to $W_*$ is surjective onto the normalized positive Hermitian
functionals on its inherited ordered space. Two ambient states restrict
identically exactly when all retained future expectations agree. Every
other readout determining those expectations factors to this restriction
on its actual image. If
$\mathcal Lb_j=\sum_kM_{kj}b_k$, then the expectation vector satisfies
$\dot y=M^Ty$.
\end{corollary}
\begin{proof}
Finite-dimensional stabilization proves invariance and leastness. The
time-orbit span lies in $W_*$; taking its derivatives at zero proves the
opposite inclusion. State extension follows from the same separating
functional proof as above, which needs only a unital self-adjoint space.
The equality of readouts is then exactly the equality of restrictions,
and the finite equation follows by evaluating $\mathcal Lb_j$.
The operator-system target is physical only where that target sort and
restriction word are licensed. Product closure is a separate consumer role.
\end{proof}
For a role owning several generators, apply the algebra closure to their
complete owned family and restrict every actual positive-unital process.
The one-generator minimal linear statement does not discard another owned
process. All time-dependent generators and required source or metric tests
are retained when the consumer requires them.





\begin{corollary}[Carrier--member firewall]
\label{p2a:cor:carrier-member-firewall}
A unique generated least carrier does not by itself select an occurring
physical member, source, state, vacuum, history, action, measure,
normalization, calibration, analytic continuation or trajectory inside that
carrier.  Such data descend without a new G0 or selector only when a separately
licensed constructor maps the already occurring roots to them uniquely, or
when the relevant readout is proved constant on the surviving unresolved
family.
\end{corollary}

\begin{proof}
Theorem~\ref{p2a:thm:generated-least-carrier} compares complete carrier
realizations of the already generated datum $S$; it does not enlarge the
output sorts of the frozen constructor grammar.  The declared-output and G0
boundaries therefore remain in force.  The only exceptions are the existing
GP2 cases of a uniquely generated descendant or a readout proved common to all
surviving members.
\end{proof}

\begin{corollary}[Generated-constructor chronology beyond carriers]
\label{p2a:cor:generated-constructor-chronology}
Let a frozen constituted presentation $P$ admit a licensed GP2 constructor
\[
 \mathcal C:P\longrightarrow Y
\]
into any already declared downstream sort, including an operator, process,
action owner, measure/state ray, response, or coarse-comparison object.  If
$Y=\mathcal C(P)$ is generated without using candidate-specific downstream
consequences as witnesses, then $Y$ is descendant content of $P$ rather than
a new independent premise.  GP3 may compare faithful realizations of this
generated object, remove redundant presentations, or retain an unresolved
family when several incomparable faithful minima survive; it does not supply
an independent simplicity ordering on arbitrary candidate actions,
functionals, measures, states, or trajectories.
\end{corollary}

\begin{proof}
This is the chronological GP2--GP3 rule of
Postulate~\ref{p2:post:GP3} applied to a downstream declared output sort rather
than specifically to a carrier.  The constructor must be available before the
candidate-specific output is evaluated, so a later prediction cannot provide
its own GP2 origin.  If the constructor has one faithful least realization,
GP3 retains it; if not, the unresolved-family clause applies.
\end{proof}

\begin{corollary}[Convenience compatible-family and observable-quotient closure]
\label{p2a:thm:compatible-family-observable-closure}
Let $P$ be an already constituted occurring owner.  Let
$\mathfrak F(P)$ be the complete GP1/GP2-admissible family of downstream
realizations of one declared role, modulo a result-independent generated
congruence $\sim_P$.  Assume the family constructor and congruence are frozen
before any candidate-specific observable is evaluated.

If the quotient $\mathfrak F(P)/{\sim_P}$ has exactly one structured
isomorphism class, that class is generated descendant content of $P$ and no
second structural selector is licensed.  This conclusion is structural only:
it does not create a new occurrence token, state, history, or member when the
claim separately consumes occurrence.

More generally, write $Q_P=\mathfrak F(P)/{\sim_P}$ for the complete
licensed generated quotient, and let $S_P=\operatorname{Surv}_{\rm GP3}(Q_P)$
be the complete structured family retained \emph{after} all licensed,
currently relevant GP1/GP2 frontier obligations and GP3 comparisons have
been consumed at the declared frozen stage.  The family and quotient must be
predeclared generated outputs; in an application of
Theorem~\ref{p2a:thm:generated-exact-solution-moduli} they retain its
certified $\operatorname{Orig}_{\Sigma}$ word and typed GP1/GP2 exactness
hypotheses.  Let $O:S_P\to\mathsf Y$ be a licensed,
congruence-invariant generated readout into a predeclared target sort, fixed
without using its desired value as its constructor certificate.

There is exactly one common determined generated readout value $O_*$ at this
frozen stage if and only if
\[
  \boxed{S_P\neq\varnothing\quad\text{and}\quad O(S_P)=\{O_*\}.}
\]
A nonempty retained family can therefore determine a common observable
without selecting a source class or physical member.  If the retained
family is empty, no positive common value follows (an empty complete exact
solution locus is FRC-obstructed).  If two \emph{GP3-surviving complete
structured classes} $s,t\in S_P$ have $O(s)\neq O(t)$, then no unique common
value is determined by this currently retained family; GP3 supplies no
member selector merely to force a preferred prediction.  Later independently
generated frontier data are processed under GP1--GP3 at a new stage, not
used retroactively to change this frozen comparison.
\end{corollary}

\begin{proof}
The singleton structured-isomorphism-class clause retains its own
nonempty witness and follows from the GP2 generated-family and GP3
whole-candidate rules.  For the readout clause, the complete licensed source
and quotient are formed before the GP3 survivor test.  Theorem~\ref{p2a:thm:generated-exact-solution-moduli}
separates the nonempty exact family from FRC obstruction and evaluates a
licensed generated readout only on the complete retained classes.  A nonempty
image is a singleton $\{O_*\}$ exactly when every survivor has the common
value $O_*$ and no different common value is possible.  With two different
surviving values no such singleton exists.  Pre-GP3 differences do not prove
post-GP3 nondetermination: a different complete candidate may be removed by
an already licensed GP3 comparison.  Neither clause adds a member occurrence
or a selector to the frozen physical presentation.
\end{proof}

\begin{corollary}[Generated exact-interface closure]
\label{p2a:cor:generated-exact-interface-closure}
Let an occurring owner generate a finite piecewise-regular same-owner object
\(A\) and a codimension-one interface \(\Sigma\) at which the exact traces
required by its declared variational role agree.  Assume the one-sided local
jets used by the owner are finite and that the frozen constructor grammar
contains no additional smoothing kernel, profile, surface source, width,
or interpolation parameter.

Then GP1--GP3 retain the exact interface germ already generated by \(A\):
GP1 admits its finite one-sided data; GP2 no-omission retains the matching
traces and finite one-sided jets, while GP2 no-insertion forbids adding an
ungenerated smoothing/profile/source datum; and GP3 may remove only redundant
presentations of that same exact germ.  In particular, a finite-width
microscopic profile is not new structural data merely because it would make
an interface visually smooth.

This structural closure does not actualize a trajectory through \(\Sigma\).
If a physical claim uses a crossing, its occurrence/history remains a G0
member.  Any additional junction condition imposed by the same-owner Euler
or field equation is a derived dynamical condition on that history, not a new
GP postulate and not a license to insert a profile.
\end{corollary}

\begin{proof}
Finite one-sided jets and the matched exact traces place the generated
interface in the GP1-admissible class.  Removing a required trace or one-sided
jet violates GP2 no-omission; adding a width, interpolation, surface source,
or smoothing kernel absent from the frozen constructor grammar violates GP2
no-insertion.  Thus every admissible faithful realization contains the same
interface germ.  GP3 can identify redundant presentations of that germ but
cannot manufacture an extra process parameter or history member.  The final
sentence is the law--flow--history and carrier--member firewall: same-owner
Euler equations act on an already generated structure, whereas actual
traversal is occurrence content.
\end{proof}

\begin{corollary}[Local observable-insensitivity test]
\label{p2a:cor:local-observable-insensitivity}
Let a smooth connected stratum of the compatible family near $q$ be defined
by generated constraints $\mathcal C(q;P)=0$, modulo generated gauge
relations. Let $V_q$ be the actual tangent of this smooth gauge-reduced
stratum, consisting of velocities of admitted curves. It equals the kernel
of the linearized constraints modulo gauge only when that equality is
certified, for example on a regular constraint and regular quotient
presentation. If a
$C^1$ readout $O$ satisfies
\[
  D O_q(v)=0\qquad\text{for every }q\text{ on the stratum and every }v\in V_q,
\]
then $O$ is constant on each path-connected component of that stratum and is
therefore closed there by
Corollary~\ref{p2a:thm:compatible-family-observable-closure}.  Failure of this
tangent condition on an actual tangent at one point proves local member
sensitivity. A nonintegrable linearized direction alone does not.
\end{corollary}

\begin{proof}
For any piecewise $C^1$ path $q(s)$ in the stratum,
$\dot q(s)\in V_{q(s)}$, so the chain rule gives
$\frac{d}{ds}O(q(s))=0$.  Hence $O$ is constant along every path in the
component.  A nonzero derivative along an allowed tangent gives two nearby
compatible realizations with different readout values.
\end{proof}

\begin{corollary}[Generated orbit and unresolved-family inheritance]
\label{p2:cor:generated-orbit-family-inheritance}
Let \(P\) be an already constituted occurring parent at a frozen stage.  Let
\(\mathcal F(P)\) be a finite or totally bounded candidate family obtained
functorially from the frozen parent by constructors admitted to the current
GP2 closure, and let \(\sim_P\) be a result-independent congruence already
generated from the frozen inventory or from candidate-internal exact data.
Assume that passage to the quotient preserves the constituted host, all
generated relation and comparison data, and role typing.  If $\mathcal F(P)$
is infinite and only totally bounded, require in addition that the retained
quotient comparison is a uniformly continuous generated image of its
comparison (or separately certify GP1 total boundedness of the entire
quotient family).  Mere set-theoretic quotienting need not preserve GP1.
Then
\[
 \mathcal Q(P):=\mathcal F(P)/{\sim_P}
\]
is generated role-level descendant content of \(P\).

If \(\mathcal Q(P)\) has one structured isomorphism class, that class is the
inherited output and no representative inside it is new physical data.  If the complete admissible quotient candidates have no unique GP3-least
structured class and several inequivalent irreducible classes survive, the
inherited output is the unresolved quotient family itself; GP3 selects no
member, and only the frozen content common to every surviving class may be
used without an additional generated separator.  Multiple quotient classes
\emph{before} GP3 do not by themselves preclude a unique least class.  In neither case does role-level co-realization imply
that every mathematical member of \(\mathcal F(P)\) occurs as an independent
physical token.
\end{corollary}

\begin{proof}
GP2 no-insertion requires both the family constructor and the congruence used
for the quotient to have generated provenance, while exactness and host
faithfulness preserve the role typing of the quotient.  Quotient by an
already generated congruence is an admissible role-faithful structural
comparison in Postulate~\ref{p2:post:GP3}; the stated comparison-control
condition ensures that an infinite quotient stays inside its GP1 domain.
A unique surviving quotient class is therefore retained without a further
choice.  If no unique least class exists, the last paragraph of
Postulate~\ref{p2:post:GP3} retains the whole irreducible family and
freezes only its common generated content.  This conclusion does not infer
absence of a least class from the presence of several initial candidates.  Generated
role-level occurrence is inherited from the occurring parent by the
co-realization theorem; token-level occurrence remains governed separately by
G0.  Thus no symmetry representative or unresolved family member is
actualized merely by choosing notation for it.
\end{proof}

\paragraph{Typed rooted structural-realization equivalence}
\label{p2a:subsec:typed-structural-geometrization}

\begin{definition}[Structural projection, generated identity, and the two rooted presentations]
\label{p2a:def:generated-realized-structural-presentations}
Let \(P_0\) be a frozen base presentation, let
\(\mathbf Q=(Q_1,\ldots,Q_s)\) be the complete tuple of independent G0 roots
of a finite chronological dependency presentation \(P\), and let \(\Sigma\)
be its frozen typed constructor grammar.  For each root, \(\Delta_{Q_i}\) is
the result-independent signature frontier of G0.

The \emph{structural-sort projection} \(\pi_{\rm str}\) is the full typed
subpresentation on nodes whose frozen tags are structural.  It retains
carriers, objects, operations, constructors, laws, relations, test/state/effect
roles, and structural scalar or readout roles, together with every generated
incidence among retained nodes.  It does not retain a particular preparation,
performance, history, outcome, empirical readout value, nonzero member, or
calibration token, and it does not retain an analytic theorem-domain
hypothesis.  A structural interface generated from an occurring token may be
retained even though the token itself is removed by \(\pi_{\rm str}\); the
decision is made by the frozen output sort, not by mathematical isomorphism.

Write \(\equiv_{\Sigma}\) for the smallest typed congruence generated by
identity or equality relations that the frozen grammar itself licenses and
retains.  A bare isomorphism, equal dimension, equal invariant, or
representation equivalence is not an element of \(\equiv_{\Sigma}\).
Moreover, the congruence may not erase a generated provenance, owner,
multiplicity, operation law, or relation slot.

Define the \emph{rooted generated structural presentation} by
\[
 \operatorname{GenStr}_{\Sigma,P_0}(\mathbf Q)
 :=
 \pi_{\rm str}\!\left(
  \operatorname{Fr}_{\rm pres}
  \operatorname{Cl}^{P_0,\Sigma,\rm typed}_{\rm GP1:3}
  \!\left(P_0\cup\{Q_1,\ldots,Q_s\}
                 \cup\bigcup_{i=1}^{s}\Delta_{Q_i}\right)
 \right)\Big/\!\equiv_{\Sigma}.
\]
The root tokens in the displayed closure are occurring inputs; they survive
the structural projection only through their structural signatures and
licensed structural outputs.

If GP3 retains an unresolved family \(\mathcal U\), this presentation retains
\(\mathcal U\) itself as a family-valued role node and retains the frozen
subpresentation proved common to every member.  It retains no chosen member
and no member-specific datum.

Define \(\operatorname{RealStr}_{P,\Sigma}(\mathbf Q)\) independently as the
full exactified typed subpresentation of \(P\) consisting of structural items
co-realized from the same occurring roots, with only the generated identities
\(\equiv_{\Sigma}\) imposed.  Constructor evaluation gives a canonical map
\[
 \rho_{\mathbf Q}:
 \operatorname{GenStr}_{\Sigma,P_0}(\mathbf Q)
 \longrightarrow
 \operatorname{RealStr}_{P,\Sigma}(\mathbf Q)
\]
which sends each generated node and incidence to its retained realization.
For concise geometric language set
\[
 \operatorname{GeoEff}_{\Sigma,P_0}(\mathbf Q)
 :=\operatorname{GenStr}_{\Sigma,P_0}(\mathbf Q),
 \qquad
 \operatorname{PhysStr}_{P,\Sigma}(\mathbf Q)
 :=\operatorname{RealStr}_{P,\Sigma}(\mathbf Q).
\]
These are aliases for the independently defined source and target of
\(\rho_{\mathbf Q}\), not definitions of one in terms of the other.
\end{definition}

\begin{theorem}[Typed rooted structural-realization equivalence]
\label{p2a:thm:typed-structural-geometrization}
For every finite chronological presentation satisfying
Definition~\ref{p2a:def:generated-realized-structural-presentations}, the
canonical realization map is an isomorphism of typed presentations:
\[
 \boxed{
 \rho_{\mathbf Q}:
 \operatorname{GeoEff}_{\Sigma,P_0}(\mathbf Q)
 \xrightarrow{\ \cong_{\rm typed}\ }
 \operatorname{PhysStr}_{P,\Sigma}(\mathbf Q).}
\]
It preserves target sort, source--target incidence, provenance, owner,
multiplicity, operation law, relation slots, and every frozen intrinsic defect
signature.  Thus rooted generated structural content has no second
structural-actuality requirement in its native generated sort, and the
co-realized structural presentation contains no ungenerated structural item.
\end{theorem}

\begin{proof}
The structural frontier \(\pi_{\rm str}(\Delta_{Q_i})\) of each independent
root is realized with its G0 witness.  For every later generated structural
node, GP2 no-omission and
Corollary~\ref{p2:cor:no-idle-licensed-constructor} retain the required output
in its declared sort.  Theorems~\ref{p2:thm:generated-constituent-corealization}
and~\ref{p2:thm:generated-role-corealization} give the corresponding
co-realization statements for finite constituents and finitely presented
roles; the general object, operation, constructor, law, relation, test, state,
effect, and structural-readout cases are already included in GP2 exactness.
Hence \(\rho_{\mathbf Q}\) is defined on the whole generated structural
presentation.

Conversely, Corollary~\ref{p2:cor:root-descendant-occurrence-normal-form}
and GP2 no-insertion say that every co-realized structural node is either in a
root structural frontier or is a typed generated descendant of earlier
nodes.  Since the tuple \(\mathbf Q\) contains every independent root, every
target node is therefore in the image of \(\rho_{\mathbf Q}\).

Suppose two source nodes have the same realized image.  Typed host
faithfulness forbids this collapse when the nodes differ in sort, provenance,
owner, multiplicity, operation law, relation incidence, or intrinsic defect.
If none differs, exactification can identify them only through a generated
identity belonging to \(\equiv_{\Sigma}\), by which they were already
quotiented in the source.  Thus \(\rho_{\mathbf Q}\) is injective on the
quotient and preserves all typed presentation structure.

Finally, on an unresolved GP3 family, the family-valued role node is itself a
generated descendant and is co-realized as such, while GP3 freezes only the
common structural subpresentation and selects no member.  The same argument
therefore applies without inserting a representative.  The inverse node and
incidence map is consequently well defined, proving the typed presentation
isomorphism.
\end{proof}


\begin{corollary}[Born-series rooted structural-descendant principle]
\label{p2a:cor:born-series-rooted-structural-descendant}
Let a later Born-series structural object \(X\) be a nonredundant output of a
finite licensed constructor word rooted in the G0-actualized geometric
presentation, with every input and typing hypothesis of that word already
satisfied.  Then \(X\) is co-realized in its declared structural sort under
the typed equivalence of
Theorem~\ref{p2a:thm:typed-structural-geometrization}.  A mathematical theorem
used to compute, classify, or certify that output does not create a second
physical root or a second structural-actuality premise.

The conclusion is structural only.  It does not select a member of an
unresolved GP3 family, an actual state or history, a numerical calibration,
or a hypothesis belonging to a separately stated analytic theorem domain.
\end{corollary}
\begin{proof}
Adjoin the licensed constructor word for \(X\) to the finite chronological
presentation.  GP2 no-omission retains its typed nonredundant output, while
Theorem~\ref{p2a:thm:typed-structural-geometrization} identifies the generated
structural presentation with its co-realized structural presentation.  The
listed exclusions are exactly the sorts not converted into selected members,
calibrations, or analytic certificates by that theorem.
\end{proof}

\begin{theorem}[Universal rooted co-realization bridge normal form]
\label{p2a:thm:universal-corealization-bridge-normal-form}
For a rooted generated presentation in
Theorem~\ref{p2a:thm:typed-structural-geometrization}, form its typed incidence
graph \(\mathsf I_{\rm gen}\) by taking retained root, role, object, operation,
relation, test, state, scalar, family, and output nodes as vertices and every
constructor input/output port or retained relation incidence as an edge.
Define \(\mathsf I_{\rm real}\) in the same way for the co-realized
presentation.  Then \(\rho_{\mathbf Q}\) induces an isomorphism
\[
 \boxed{
 \mathsf I_{\rm gen}\xrightarrow{\;\cong_{\rm typed}\;}
 \mathsf I_{\rm real}.}
\]
Hence every genuinely co-realized descendant is connected in
\(\mathsf I_{\rm real}\) to each root on which its frozen constructor word
actually depends.  Any two co-realized nodes in the same rooted constructor
component are connected by a finite typed incidence zigzag.  Their common
ultimate origin is that rooted component, although their immediate
constructors and bridge sorts need not be the same.

This universal bridge is constructor incidence and provenance.  It does not
assert a direct map or an isomorphism between every pair of co-realized nodes.
The stronger forms are exactly typed specializations:
\begin{enumerate}[label=(\alph*),leftmargin=2em]
\item a retained total evaluation with the closure hypotheses of
      Theorem~\ref{p2a:thm:evaluation-dual-corealization} gives a
      contravariant operation/dynamics bridge;
\item a retained comparison isomorphism between distinct generated reciprocal roles,
      together with relation-free joint support, gives the internal
      involutive Born bridge;
\item an inclusion, projection, quotient, readout, partial evaluation, Schur
      complement, process graph, or other retained non-isomorphism relation
      gives only its declared directed, relational, or noninvertible bridge
      unless an inverse is separately generated; and
\item an unresolved GP3 family gives a bridge through the family-valued node
      and its common subpresentation, but selects no member and identifies no
      member-specific nodes.
\end{enumerate}
If two independent roots lie in disconnected generated components, their
simultaneous presence alone generates no bridge between them.  Adding one
would violate GP2 no-insertion.  Thus \emph{co-realized} always means connected
to its complete licensed origin; pairwise exchangeability holds only on the
invertible self-dual locus.
\end{theorem}

\begin{proof}
Theorem~\ref{p2a:thm:typed-structural-geometrization} states that
\(\rho_{\mathbf Q}\) preserves every node, source--target incidence,
provenance, owner, operation law, and relation slot, and that its inverse does
the same.  Applying the incidence-graph construction therefore gives the
displayed typed graph isomorphism.  A frozen constructor word is a finite
input/output path from every input it actually consumes to its required
output, so its realized image supplies the asserted rooted path.  Concatenating
two such paths through their common component gives the finite zigzag.

No-insertion proves the disconnected-component converse.  Case (a) is
Theorem~\ref{p2a:thm:evaluation-dual-corealization}.  Case (b) is
Theorem~\ref{p2:thm:gp2-generated-reciprocal-role-closure}; its involution
uses the generated comparison and its inverse.  In case (c), promoting a
one-way typed constructor to an isomorphism inserts an undeclared inverse.  In
case (d), choosing a member inserts a selector forbidden by the GP3
unresolved-family clause.  These cases exhaust the bridge forms because every
edge of the exactified realized presentation is, by no-insertion, the image
of a retained generated constructor port or relation incidence.
\end{proof}


\begin{corollary}[Derived three-sphere ancestry is inherited, not recited]
\label{p2a:cor:v282-transitive-s3-ancestry}
Fix a primitive role \(R\) in the noncommutative endpoint-comparison domain
of Theorem~\ref{p2:thm:master-compact-role-allocation}.  Assume its
GP1/GP2-admissible candidate class is nonempty and its constituted exact
branch is noncommutative.  Then its round operation parent
\(H_R\cong SU(2)\cong S^3(\Rmom)\) is a \emph{derived} structural parent,
not an additional geometric postulate or an independently actualized
three-sphere.

Let \(X\) be any later structural output for which a finite frozen,
GP2-licensed constructor word is supplied and genuinely consumes this
parent (possibly together with other independently occurring roots).
Assume all input-role, typing, and generated-relation conditions of the
word.  Then the co-realized structural incidence presentation contains
a directed constructor path from the derived \(S^3\) parent to \(X\),
with any other consumed roots recorded separately.  This ancestry is
\emph{transitive}: the proof of a later result about \(X\) need not
repeat a citation to Paper~II-A, and the mere absence of such a citation
is not an independent physical occurrence obligation.

The conclusion concerns the declared generated structural output, not
all mathematical lemmas about its geometry.  It neither creates a
constructor when one has not been supplied nor proves a contingent
member selection, state/history occurrence, calibration, or analytical
hypothesis.  A source theorem's citation graph and the generated
constructor-incidence graph are therefore distinct dependency relations.
\end{corollary}
\begin{proof}
Theorem~\ref{p2:thm:master-compact-role-allocation}, in its stated
admissible noncommutative domain, supplies the relation-exact
\(SU(2)\) globalization and round metric up to scale.  Appending the
finite licensed word for \(X\) to its rooted structural presentation
and applying Theorem~\ref{p2a:thm:typed-structural-geometrization}
co-realizes the exact generated structural output.  Theorem~
\ref{p2a:thm:universal-corealization-bridge-normal-form} preserves
all declared constructor incidences and gives a finite rooted path
from every input actually consumed by the word to its output.
Transitivity is concatenation of those paths.  No direct citation of
the original constitution theorem is required at a later node,
because its parent is already part of the generated input presentation.
The final exclusions are the explicit structural-only boundary of
Corollary~\ref{p2a:cor:born-series-rooted-structural-descendant}.
\end{proof}

\begin{theorem}[Conjunctive finite-word closure and essential-root criterion]
\label{p2a:thm:v283-conjunctive-constructor-closure}
Fix a finite acyclic slice of the frozen typed constructor grammar
$\Sigma$.  Write $D$ for its required root tokens, and let each
retained constructor occurrence $w$ have the complete finite input set
$I(w)$, output $o(w)$, and a separately checked frozen-license and
admissibility predicate $C(w)$.  A true $C(w)$ certifies the source
sort, origin, operation/relation slot and any additional antecedents; it
cannot be inferred from a graph arrow.  For $R\subseteq D$ define
\[
 \begin{aligned}
  S_0(R)&=R,\\
  S_{n+1}(R)&=S_n(R)\cup\{o(w): C(w),\ I(w)\subseteq S_n(R)\},\\
  \operatorname{Cl}_\Sigma(R)&=\bigcup_{n\ge0}S_n(R).
 \end{aligned}
\]
Then $x\in\operatorname{Cl}_\Sigma(R)$ if and only if $x$ admits a
finite rooted \emph{AND-derivation tree}: each leaf belongs to $R$,
and every internal constructor has all its required input children
with $C(w)$ verified.  Alternative licensed constructors for the same
output give OR-alternatives among complete AND-trees; this does not
select a physical member of an unresolved GP3 realization family.
For a root $r\in R$, indispensability within this \emph{fixed}
licensed grammar is characterized exactly by
\[
 x\in\operatorname{Cl}_\Sigma(R)
 \quad\text{and}\quad
 x\notin\operatorname{Cl}_\Sigma(R\setminus\{r\}).
\]
An ordinary directed path $r\leadsto x$ obtained by drawing separate
arrows $i\to o(w)$ for all $i\in I(w)$ establishes neither closure of
$x$ nor indispensability of $r$.  Consequently the transitive
$S^3$-ancestry of
Corollary~\ref{p2a:cor:v282-transitive-s3-ancestry} records a
\emph{consumed} origin of a complete licensed word, not that the derived
sphere by itself supplies all of the word's inputs.  The theorem
creates no physical constructor: independent mathematical lemmas
verifying $C(w)$ remain proof inputs, and G0-occurring calibrations,
other roots and member-specific certificates remain in the input tuple.
\end{theorem}
\begin{proof}
Induct on the least stage $n$ containing $x$.  A root in $S_0(R)$
is a one-vertex tree.  If $x=o(w)$ enters $S_{n+1}(R)$, attach the
previously constructed trees of \emph{all} inputs $i\in I(w)$ under
$w$, whose predicate $C(w)$ is true.  Conversely induction on the
height of a finite AND-tree places its leaves in $S_0(R)$ and each
internal output in the next stage.  The finite slice stabilizes after
finitely many newly added tokens.  The indispensability criterion is
exactly the closure-membership test after removing $r$ without changing
$\Sigma$ or its predicates.  To disprove sufficiency of a directed
path, take $R=\{a\}$ and one licensed rule $\{a,b\}\to x$, with
$b$ neither a root nor generated: $a\to x$ exists but $x$ is absent
from the closure.  For indispensability, take $R=\{a,c\}$ and two
licensed rules $\{a\}\to x$ and $\{c\}\to x$: the path $a\to x$
exists but $a$ is dispensable.  On a complete licensed word,
Theorem~\ref{p2a:thm:universal-corealization-bridge-normal-form}
preserves all constructor input incidences and
Corollary~\ref{p2a:cor:v282-transitive-s3-ancestry} transports
actually consumed $S^3$ ancestry.  Neither theorem waives the AND-test.
\end{proof}

\begin{corollary}[Born reciprocity is the invertible self-dual co-realization case]
\label{p2a:cor:born-invertible-corealization}
Under the hypotheses of
Theorem~\ref{p2:thm:gp2-generated-reciprocal-role-closure}, suppose the frozen
grammar also generates the comparison isomorphism
\(\beta_A:A\overset{\sim}{\to}B_A\).  Then the same rooted occurrence
co-realizes
\[
 \bigl(A,B_A,A\times B_A,\beta_A,\beta_A^{-1},
       \mathcal B_{\beta_A}\bigr),
 \qquad
 \mathcal B_{\beta_A}(a,b)
 =\bigl(\beta_A^{-1}b,\beta_Aa\bigr),
\]
and \(\mathcal B_{\beta_A}^2=\operatorname{id}\).  Born reciprocity is
therefore not an additional bridge principle placed over co-realization.  It
is the involutive internal form taken by the universal rooted bridge when the
two roles are autonomous, relation-free, and joined by a generated
isomorphism.
\end{corollary}

\begin{proof}
The reciprocal-role theorem generates and co-realizes the autonomous role and
relation-free product.  The isomorphism retains its inverse in the declared
comparison sort, and GP2 no-idle closure retains the displayed exchange
operation.  Direct substitution gives
\(\mathcal B_{\beta_A}^2=\operatorname{id}\).  The universal bridge theorem
then identifies this as case~(b), while its other cases show why the same
involution cannot be inferred for arbitrary co-realized nodes.
\end{proof}

\begin{theorem}[Generated Born inheritance on licensed descendants]
\label{p2a:thm:generated-born-inheritance}
Let \(P\) be a G0-occurring frozen typed presentation with intrinsic physical
constructor grammar \(\Sigma\).  Suppose a rooted reciprocal-role closure and
a generated comparison isomorphism have produced an involutive Born
operation
\[
 \mathcal B:P\longrightarrow P,
 \qquad \mathcal B^2=\operatorname{id},
\]
as in Corollary~\ref{p2a:cor:born-invertible-corealization}.  Let
\[
 \operatorname{Orig}_{\Sigma}(w;D,\mathsf Q),
 \qquad
 q=\operatorname{ReqOut}_{\Sigma}(w,D),
\]
be a frozen constructor-origin certificate.  If
\[
 \boxed{\mathcal B\in\operatorname{Aut}_0(P;w)},
\]
meaning that Born exchange preserves every retained mark, relation, source,
boundary/interface datum, calibration token, constructor schema and typing
actually consumed by \(w\), and if \(\mathcal B_{\mathsf Q}\) denotes the
induced involution in the declared target sort, then on the non-FRC-obstructed
branch
\[
 \boxed{
 \operatorname{ReqOut}_{\Sigma}(w,\mathcal B D)
 =\mathcal B_{\mathsf Q}\,
  \operatorname{ReqOut}_{\Sigma}(w,D).}
 \tag{BR-INH}
\]
Thus generated Born reciprocity is hereditary through the licensed
constructor tree wherever it remains unmarked.  The general conclusion is
equivariance, not pointwise invariance: the target action may fix a scalar,
exchange components, reverse a parity, conjugate an operator, or transport an
action graph within its declared target sort.

This theorem does not generate the reciprocal role, the comparison
isomorphism, a constructor word, its target action, a calibration, a state,
a boundary condition, an actual history, an analytic limit, or an unresolved
family member.  If Born exchange changes any datum consumed by \(w\), then it
is not an element of \(\operatorname{Aut}_0(P;w)\) and this theorem does not
assert an internal symmetry.
\end{theorem}

\begin{proof}
Corollary~\ref{p2a:cor:born-invertible-corealization} supplies the generated
involutive automorphism.  Under the displayed unmarkedness hypothesis it is
an admissible choice \(a=\mathcal B\) in
Theorem~\ref{p2a:thm:v368-same-owner-unmarked-equivariance}.  Substitution in
\((\mathrm{SO\text{-}EQ})\) gives \((\mathrm{BR\text{-}INH})\).  Every
exclusion in the final paragraph is exactly an antecedent or firewall of the
same-owner theorem or of the reciprocal-role construction and therefore is
not created by this specialization.
\end{proof}

\begin{corollary}[Born-fixed descendants and parity elimination]
\label{p2a:cor:born-fixed-descendant-parity}
Under Theorem~\ref{p2a:thm:generated-born-inheritance}, if the complete input
tuple is Born fixed, \(\mathcal B D=D\), then
\[
 \boxed{q=\mathcal B_{\mathsf Q}q.}
 \tag{BR-FIX}
\]
If moreover \(\mathsf Q\) is a real or complex linear target and
\(\mathcal B_{\mathsf Q}^2=1\), write
\[
 \mathsf Q_\pm:=\ker(\mathcal B_{\mathsf Q}\mp1).
\]
Then the Born-odd component of the generated output vanishes:
\[
 q=q_++q_-\quad\Longrightarrow\quad\boxed{q_-=0}.
\]
For a family-valued GP3 output this statement applies to the generated family
object or to content common to every surviving member; it does not select an
unresolved member merely because a Born-fixed member exists.
\end{corollary}

\begin{proof}
Set \(\mathcal B D=D\) in \((\mathrm{BR\text{-}INH})\).  For a linear
involution, applying the projectors
\(\frac12(1\pm\mathcal B_{\mathsf Q})\) gives the parity conclusion.  The
family firewall is the GP3 abstention clause already retained by the frozen
origin calculus.
\end{proof}

\begin{proposition}[Typed reciprocal-branch transport]
\label{p2a:prop:typed-born-branch-transport}
Let \(\mathcal B:P\overset{\sim}{\to}P^{\mathcal B}\) be a generated typed
Born isomorphism between two complete frozen presentations rather than an
unmarked automorphism of one presentation.  Let
\[
 w=C_m\circ\cdots\circ C_1:D\to\mathsf Q
\]
have a frozen origin certificate in \(P\).  Suppose the reciprocal
presentation independently contains the transported licensed word
\[
 w^{\mathcal B}=C_m^{\mathcal B}\circ\cdots\circ C_1^{\mathcal B}:
 \mathcal B D\to\mathsf Q^{\mathcal B}
\]
with its own frozen origin certificate, and that every declared constructor
square commutes with the typed Born maps:
\[
 C_j^{\mathcal B}\circ\mathcal B_{j-1}
 =\mathcal B_j\circ C_j.
\]
Then
\[
 \boxed{
 \operatorname{ReqOut}_{\Sigma^{\mathcal B}}
   (w^{\mathcal B},\mathcal B D)
 =\mathcal B_{\mathsf Q}
   \operatorname{ReqOut}_{\Sigma}(w,D).}
 \tag{BR-TR}
\]
This is reciprocal branch transport, not invariance of the original marked
presentation.  In particular, unequal finite levels may be exchanged without
requiring equality of those levels.
\end{proposition}

\begin{proof}
Induct on the transported constructor word.  The commuting-square hypothesis
gives
\(C_1^{\mathcal B}(\mathcal BD)=\mathcal B_1C_1(D)\), and the same identity
at each successive stage gives
\(x_j^{\mathcal B}=\mathcal B_jx_j\).  At the terminal stage this is
\((\mathrm{BR\text{-}TR})\).  The two frozen origin certificates and GP2
no-omission retain the two outputs in their respective presentations.  No
claim that \(P=P^{\mathcal B}\) has been made.
\end{proof}

\begin{remark}[Reciprocity hierarchy and scope]
\label{p2a:rem:reciprocity-hierarchy-and-scope}
The following is a hierarchy of \emph{separately guarded} licensed
construction steps, not an unconditional chain of implications from
separating evaluation alone.  In particular, a dual operation requires
precomposition stability and a frozen physical license; an autonomous
reciprocal role additionally requires the GP2 objectification and
relation-free-product conditions, and self-duality requires an independently
generated comparison isomorphism.  Under those guards the hierarchy is
\[
 \boxed{\begin{aligned}
  &\text{separating evaluation}
    \Rightarrow\text{evaluation-dual operation}\\
  &\Rightarrow\text{GP2-generated reciprocal role}\\
  &\Rightarrow\text{generated self-duality}\\
  &\Rightarrow\text{Born involution}\\
  &\Rightarrow\text{hereditary covariance on unmarked licensed descendants}.
 \end{aligned}}
\]
The last arrow does not identify reciprocal quantities numerically and does
not erase marks.  Dimensionful position--momentum normalization, clock units,
\(2\pi\hbar\), Newton normalization and analogous interfaces remain
calibrations unless separately generated.  A state, source, orientation,
boundary condition, cutoff, regulator or other retained mark that is not Born
fixed either breaks the internal symmetry or moves the claim to the typed
transport setting of Proposition~\ref{p2a:prop:typed-born-branch-transport}.
\end{remark}

\begin{corollary}[No duplicate structural-actuality premise]
\label{p2a:cor:generated-structural-manifestation}
Let \(R\) be a nonredundant node of
\(\operatorname{GeoEff}_{\Sigma,P_0}(\mathbf Q)\).  Then
\(\rho_{\mathbf Q}(R)\) is co-realized in the same generated structural sort.
A further premise whose only content is that this same rooted descendant is
``physically real'' in that same sort contributes no structural datum and is
not an independent occurrence assumption.

This does not imply the occurrence of an arbitrary member \(r\in R\), a
representative of an unresolved family, or a different preparation, apparatus,
process, history, outcome, nonzero field/connection, empirical value, or
calibration token.
\end{corollary}

\begin{corollary}[Typed generated-identity deduplication]
\label{p2a:cor:typed-generated-identity-deduplication}
Suppose \(X\) and \(Y\) are generated in the same declared structural sort and
the frozen grammar generates and retains an identity congruence
\(X\equiv_{\Sigma}Y\).  Equivalently, in the canonical-image case the exact
GP2 canonical-image clause identifies the generated image with already
retained object content in that same object sort.  If no generated
provenance, owner, multiplicity, operation law, intrinsic defect, or relation
slot distinguishes the two, the exactified presentation contains one node,
and a second constitutive or structural-physicality premise is redundant.

A generated isomorphism that is not itself an identity congruence does not
satisfy the hypothesis.  Neither do abstract isomorphism, equality of
dimension or numerical invariants, representation equivalence, or a canonical
mathematical functor between differently typed or independently owned roles.
\end{corollary}

\begin{remark}[Structural-geometrization boundary]
\label{p2a:rem:typed-structural-geometrization-boundary}
The name ``effective rooted geometry'' refers only to the structural
projection of the licensed typed closure.  Ambient atlas objects, theorem
consequences, operators, quotients, tensor constructions, or solution families
do not enter it unless the frozen physical grammar generates them in the
asserted target sort.  In particular, a canonical geometric operator is not
thereby a physical kinetic Hamiltonian, and a mathematical separating family
is not thereby an actually performed test family.

The theorem supplies no token plenitude.  The typed cross-relation firewall,
zero-member test, declared-output boundary, and G0 actuality predicate remain
in force.  A missing cross-role bridge, sort-changing promotion, selected
family member, particular preparation, process, history, outcome, nonzero
field/connection, empirical value, or calibration still requires its own
licensed generated route or G0 occurrence.
\end{remark}

\begin{remark}[Dependency-separated convention]
\label{p2:rem:cumulative-gp-domain}
The logical cores of G0 and GP1--GP3 remain distinct.  G0 types the physical
actualization of a datum or relation and forbids outcome-dependent
retro-selection; it neither generates an occurrence witness nor replaces an
analytic theorem-domain hypothesis.  GP1 supplies typed
finite/totally-bounded admissibility.  GP2 exactifies each candidate's generated
host content and requires its faithful sort-preserving realization, using only
geometrically licensed physical constructors rather than the unrestricted
mathematical closure of a representation category.  GP3 then asks whether the
resulting class of complete exact candidates has an irreducible least
structured realization under role-faithful comparisons.  The principles are
applied cumulatively and in that order; the selected candidate's complete
self-content is frozen only after the final step.

A generated-depth or spectral boundary is not part of GP1.  When a later role
really requires such a boundary it enters through an explicit generated or
registered item such as \(\mathsf P_{\rm gbd}\), after which GP2 and GP3 apply
to the richer host signature.  The existence of a bounded filtration does not
itself make every mathematical tensor word physical, nor does GP3 minimize the
filtration label numerically.  Conversely, no separate level-specific
\emph{principle} is required when a concrete role-faithful structural
comparison is already proved in that domain.

Higher tori and higher-rank compact Lie groups such as \(SU(3)\) are genuine
FRC-admissible primitive-role candidates before GP3 and are exact to their own
generated geometries.  They are not least for the primitive scalar or
order-sensitive role.  If \(\operatorname{Geo}_R\) itself requires their extra
directions, brackets, relations, or tests, they belong to a richer role domain
and the primitive endpoint comparison is unavailable.  The quotient \(SO(3)\)
is removed earlier by GP2 relation exactness when its kernel relation is
ungenerated.  GP3 cannot turn an untyped role into a noncommutative one,
invent an operation or test missing from the host signature, or supply a
physical radius.  Ordinary analytic assumptions used later---for example regularity of a
primitive parent comparison or ellipticity of an operator---remain displayed
separately.  GP2 no-insertion governs constituted physical content; a
candidate may still contain internally exact self-operations or redundant
factors that realize no additional geometrically primary role item.  Those
self-data are not frozen as physical content unless GP3 retains the candidate.
\end{remark}



\begin{definition}[Recursive cumulative GP application]
\label{p2:def:recursive-gp-application}
The cumulative order
\[
   \mathrm{G0}\ \longrightarrow\ \mathrm{GP1}\ \longrightarrow\
   \mathrm{GP2}\ \longrightarrow\ \mathrm{GP3}
\]
is not restricted to the first constitution of a role.  It is reapplied
whenever later structural, kinematic, dynamical, or operational reasoning
generates new physical content not present in the signature of the preceding
closed stage.  G0 supplies the occurrence/bridge/test/readout certificate by
which the new datum enters the physical ledger, after which GP1--GP3 close the
enlarged signature.  This rule licenses a later GP application; it does not
itself generate the new content or its occurrence witness.

Let \(\Sigma_n\) be the finitely presented constituted signature at stage
\(n\).  The ordered GP1--GP3 application first checks whether its
complete GP1-admissible, GP2-exact candidate class is nonempty.  If it is
empty, the already constituted FRC obstruction is recorded; no selected
class, unresolved surviving family, or new frozen physical presentation
is inferred from that obstructed attempt.  On the non-FRC-obstructed branch
write
\[
   C_n=(P_n,\mathcal M_n)
   :=\operatorname{Cl}_{\mathrm{GP1:3}}(\Sigma_n),
\]
where \(P_n\) is the frozen constituted presentation established by closure
and \(\mathcal M_n\) is either the selected structured isomorphism class
or the nonempty unresolved irreducible family.  If
\(\mathcal M_n\) is unresolved, \(P_n\) contains only the prior signature and
those generated presentation items proved common to every irreducible member; it
is not the full candidate-specific presentation of an arbitrarily chosen
member.  The symbol \(\operatorname{Cl}_{\mathrm{GP1:3}}\) is notation for
the ordered constitutional application only; G0 validates its physical input
and is not an additional mathematical closure operator.  Define
\[
   \operatorname{Fr}_{\mathrm{pres}}(C_n):=P_n.
\]
Thus the frozen presentation remains well defined even when GP3 returns an
unresolved family.  Theorem consequences of \(P_n\) remain available for
deduction, but do not become signature members merely by being theorems.

The recursive chronology is a dependency chronology rather than a publication
or paper-number chronology.  Let \(D_n\) be the declared dependency ledger
available at stage \(n\).  A \emph{dependency frontier}
\(\mathcal F_{n+1}\) is a recorded collection of new generated data such that
the prerequisites of every \(\delta\in\mathcal F_{n+1}\) lie in already
frozen stages and no member of the frontier depends on another member of the
same frontier.  Put
\[
   \Delta[\mathcal F_{n+1}]
   :=\bigcup_{\delta\in\mathcal F_{n+1}}\delta,
\]
and form
\[
   \Sigma_{n+1}
   :=P_n\cup\Delta[\mathcal F_{n+1}].
\]

The next closure is relative to the frozen stage.  Its candidate class consists
only of structured realizations of \(\Sigma_{n+1}\) in which \(P_n\) remains
faithfully represented as prior constituted content.  Denote this ordered
relative application by
\[
   C_{n+1}
   :=\operatorname{Cl}^{P_n}_{\mathrm{GP1:3}}(\Sigma_{n+1})
\]
\emph{only on its non-FRC-obstructed branch}.  If the relative candidate
class is empty, this attempted extension is FRC-obstructed: the preceding
\(P_n\) remains frozen, but no \(C_{n+1}\) or \(P_{n+1}\) is constituted
by that attempt.  Otherwise GP1 has rechecked finite/bounded admissibility
of the enlarged carrier--operation structure, GP2 exactifies the enlarged
generated content, and GP3 performs only role-faithful comparisons preserving
the frozen base.  On this nonempty branch the frozen presentations are
cumulative:
\[
   P_n\hookrightarrow P_{n+1}.
\]
A later GP3 comparison may therefore remove neither a previously frozen carrier
distinction nor a previously frozen operation, relation, test, or marking.

Generation at the later stage remains subject to GP2 chronology.  No
\(\delta\in\mathcal F_{n+1}\) may presuppose the branch, parameter value,
quotient, representation, carrier, operation coefficient, or other distinction
whose later selection it is used to establish.  In particular, a specialized
conclusion cannot be used to generate a new datum and then be recovered from
that datum by a later GP3 comparison.

A proposed discharge relative to \(D_n\) is licensed only with a recorded
frontier containing every independently generated datum in \(D_n\) whose
prerequisites lie in the same frozen base and whose generated tests, relations,
or operation distinctions bear on the retained GP3 family being resolved.
Data generated from different subsets of the same frozen stages enter the same
comparison whenever they bear on that family; identical prerequisite sets are
not required.  If two frontier data resolve the family incompatibly, no
discharge is licensed until the incompatibility is resolved at the level of
their derivations.  Frontier completeness is relative to the declared current
dependency ledger: a genuinely new datum proved in a later extension may enter
a later frontier and trigger another GP application, but it does not
retroactively change what was derivable from the earlier ledger.

Accordingly, an unresolved structured family at stage \(n\) may become
resolved later only when independently generated frontier data distinguish its
members.  Such resolution is another application of the same GP1--GP3
calculus, not a new selection principle.  If the enlarged signature has a nonempty exact candidate class but still
no unique least structured isomorphism class, GP3 retains the unresolved
irreducible family exactly as before.  If that class is empty, the FRC
obstruction branch above applies instead.
\end{definition}


\begin{definition}[Internal theorem closure and analytic certification]
\label{p2a:def:v368-theorem-closure-status}
Let a proved mathematical theorem \(\mathfrak T\) be applied at a frozen stage
\(P_n\).  Its use is an \emph{internal theorem closure} when every physical
owner, relation, operation, and target sort appearing in its hypotheses is
already generated or constituted at frozen stages, the theorem introduces no
new physical selector or token, and its conclusion is obtained only by
ordinary mathematical implication from those frozen data.  A same-owner
variational identity derived from
Corollary~\ref{p2a:cor:v368-native-noether-antecedent} is of this type once the
required variational constructor is already licensed.

Its use is an \emph{analytic certification} when the theorem additionally
requires hypotheses such as regularity, coercivity, hyperbolicity, domain or
closedness conditions, compactness/precompactness, convergence, resolvent
bounds, existence, uniqueness, or continuous dependence that are not
consequences of G0 or GP1--GP3 merely from correct provenance.  Those
hypotheses must be proved on the stated domain before the conclusion may enter
the physical dependency ledger.

The distinction is logical rather than bibliographic: either theorem may be
standard mathematics.  Neither status turns the theorem into a GP, and neither
licenses an undeclared physical target sort.
\end{definition}

\begin{corollary}[Licensed theorem-generated dependency frontier]
\label{p2a:cor:v176-licensed-theorem-generated-frontier}
Let \(P_n\) be a frozen constituted presentation and let \(\mathfrak T\) be
a proved theorem whose physical prerequisites lie entirely in already frozen
stages.  Suppose its declared outputs \(\Delta_{\mathfrak T}\) land in
predeclared carrier, operation, relation, test, readout, constraint, or
comparison sorts and do not presuppose the branch or member that a later GP
comparison is intended to resolve.  If the use of \(\mathfrak T\) is an
internal theorem closure in the sense of
Definition~\ref{p2a:def:v368-theorem-closure-status}, no additional physical
premise is required.  If it is an analytic certification, every additional
analytic hypothesis must first be independently proved on the quantified
domain.  Once the resulting outputs are explicitly registered in the current
dependency ledger, they form a licensed dependency frontier and may be
submitted to the recursive closure.  On its non-FRC-obstructed branch,
this yields
\[
 P_n\xrightarrow{\ \mathfrak T\ }
 P_n\cup\Delta_{\mathfrak T}
 \xrightarrow{\ \operatorname{Cl}^{P_n}_{\mathrm{GP1:3}}\ }
 P_{n+1}.
\]
If instead the enlarged GP1/GP2 exact candidate class is empty,
the extension reports FRC obstruction, not an occurring \(P_{n+1}\).
The theorem does not become a new GP and GP repetition does not generate
\(\Delta_{\mathfrak T}\).  Conversely, a mathematical consequence of
\(P_n\) does not become physical signature content merely because it is true:
its theorem output must have a licensed physical target sort and be recorded
as a new dependency frontier.
\end{corollary}

\begin{proof}
This is the dependency-frontier clause of
Definition~\ref{p2:def:recursive-gp-application} specialized to proved theorem
outputs.  In the internal-closure case, all physical hypotheses and target
sorts are already frozen and the theorem contributes only mathematical
implication.  In the analytic-certification case, the extra analytic
hypotheses are part of the theorem's prerequisite ledger and must be certified
before the output exists on the claimed domain.  Frozen-prerequisite,
chronology, target-typing, and frontier-completeness conditions are otherwise
exactly those already required by the recursive definition; the conclusion
therefore introduces no new selection principle.
\end{proof}

\begin{remark}[Native Noether closure is not a fourth GP]
\label{p2a:rem:v176-noether-not-gp4}
On a native same-owner branch satisfying
Theorem~\ref{p2a:thm:v368-same-owner-unmarked-equivariance}, action invariance
is not an independently inserted physical premise:
Corollary~\ref{p2a:cor:v368-native-noether-antecedent} derives the invariant
variational relation \(dS(R_\xi)=0\).  Once the already licensed local
variational calculus converts that relation into an Euler--Lagrange identity,
current, moment map, constraint relation, or Ward identity, the result is an
\emph{internal theorem closure} and hence a canonical instance of
Corollary~\ref{p2a:cor:v176-licensed-theorem-generated-frontier}.  Noether is
therefore not a fourth GP and, on this native branch, not an independent
physical conservation principle.

For an externally supplied action or symmetry, or whenever the proposed
transformation is not an unmarked automorphism of the complete action-owning
presentation, invariance must still be separately proved.  An integrated
global charge still requires its boundary/global hypotheses; a quantum Ward
identity requires its anomaly/measure hypotheses; and a numerical charge value
still requires a state or history.  These types are not collapsed by the
recursive rule.
\end{remark}


\begin{corollary}[Recursive fixed-point property]
\label{p2:cor:recursive-gp-fixed-point}
If no new dependency frontier is present, recursive application changes
nothing:
\[
   \Delta[\mathcal F_{n+1}]=\varnothing
   \quad\Longrightarrow\quad
   C_{n+1}\cong C_n.
\]
In particular, repeated GP application alone creates no new carrier, operation,
role, scale, coefficient, or relation.
\end{corollary}

\begin{proof}
With an empty frontier one has \(\Sigma_{n+1}=P_n\).  The relative candidate
class is exactly the already closed structured class extending \(P_n\), and
GP1, GP2, and GP3 receive no new admissibility, exactness, or faithfulness
datum on which to act.  The selected class or unresolved family is therefore
unchanged up to structured isomorphism.
\end{proof}


\begin{corollary}[Joint exhaustion of simultaneous unresolved families]
\label{p2:cor:joint-unresolved-family-exhaustion}
Let a frozen stage retain unresolved role families
\(
 \mathfrak A_1,\ldots,\mathfrak A_N
\)
that are consumed simultaneously by one later physical claim.  Their formal
product
\[
 \mathfrak A_{\rm atlas}:=\prod_{i=1}^{N}\mathfrak A_i
\]
is an enumeration atlas, not a claim that every Cartesian tuple is jointly
admissible.  Let \(\mathcal F\) contain every independently generated frontier
datum, relation, test, and licensed constructor output whose prerequisites lie
in the same frozen base and which bears on any of these roles.  For a tuple
\(a\in\mathfrak A_{\rm atlas}\), form the complete enlarged presentation using
only those licensed data, and define the joint zero-defect locus
\[
 \boxed{
 \mathfrak Z_{\rm GP}
 :=\left\{a\in\mathfrak A_{\rm atlas}:
      {\rm GP1}(a),\ 
      \mathfrak D_{\rm joint}^{\rm GP2}(a)=0
   \right\}.}
\]
GP3 compares only the complete structured realizations in
\(\mathfrak Z_{\rm GP}\).  Thus several families that survive separately may
lose Cartesian combinations through insertion, omission, relation, type, or
closure defect once they are realized together; an empty joint locus cannot be
repaired by GP3.

If no generated cross-role relation and no newly enabled common constructor is
present, the relation-free composition theorem remains controlling and no
cross-family equality may be inserted merely to force a selector.  Likewise, when later data are a generated readout \(B=f(A)\), their raw
joint incidence is the graph \(\{(A,f(A))\}\).  This graph alone is not a
member selector.  It preserves the previous GP3 preorder \emph{if} its
addition is conservative for joint GP1/GP2 admissibility and both preserves
and reflects whole-candidate GP3 comparisons.  Without those certificates,
the graph's set-theoretic bijection with \(A\) does not rule out new
obstructions or changes to the comparison order.  Genuine reduction can
arise from independently generated cross-role incidence \emph{or} from a
newly consumed, previously licensed constructor that renders some completed
candidate inadmissible.  This is a
multi-family application of the existing recursive GP1--GP3 calculus and is
not an additional selection principle.
\end{corollary}

\begin{proof}
For \(a\in\mathfrak A_{\rm atlas}\), let
\(\mathcal R_0(a;\mathcal F)\) denote the class of complete structured
realizations of the enlarged presentation based on \(a\), the frozen base,
and the closure of all licensed frontier data in \(\mathcal F\), restricted
to GP1-admissible realizations with zero joint GP2 realization defect.
Candidate-internal GP2 exactness and the definition of
\(\mathfrak Z_{\rm GP}\) therefore give the next GP3 candidate class exactly
as
\[
 \operatorname{Adm}_{\rm joint}
 =
 \left(
   \bigcup_{a\in\mathfrak Z_{\rm GP}}
      \mathcal R_0(a;\mathcal F)
 \right)\!\big/\!\cong .
\]
This proves both the exhaustion statement and the possibility that tuples
which survive every coordinatewise check fail after joint relations, common
constructors, or typed closure obligations are imposed.  If
\(\mathfrak Z_{\rm GP}=\varnothing\), the GP3 preorder has no object on which
to attain a least class; GP3 cannot create a candidate and hence cannot repair
the empty locus.

Suppose next that the frozen signature contains neither a generated
cross-role relation nor a newly enabled common constructor.  The
relation-free-composition clause of Postulate~\ref{p2:post:GP2} then supplies
the universal product joint realization and forbids inserting a diagonal,
graph, equality, quotient, or other proper support relation.  Thus no
cross-family restriction beyond the already frozen coordinatewise
admissibility can be used as a selector.

Finally, let \(B=f(A)\) be only a generated readout of one unresolved family.
Its joint locus is the graph
\[
 \Gamma_f:=\{(A,f(A)):A\in\mathfrak A\}.
\]
The projection \(\pi_A:\Gamma_f\to\mathfrak A\) is a bijection with inverse
\(i_f(A)=(A,f(A))\).  The maps are inverse on \emph{sets}, but this fact alone does not
identify GP1/GP2-complete candidates or their GP3 comparisons.  On the
specified conservative branch, every candidate and every whole-structured
comparison survives the extension and restriction, in both directions.
Thus \(i_f\) and \(\pi_A\) are inverse comparison-preserving maps on the
admissible candidate preorders and do induce an order isomorphism.  Only
under that stronger certificate does the readout graph propagate exactly
the unresolved \(A\)-alternatives without selecting one.
An independently generated relation such as \(f(A)=g(D)\), with
\(D\) ranging over another unresolved family \(\mathfrak D\), instead replaces
the graph/product by intersection with the fibre-product incidence
\[
 \{(A,D)\in\mathfrak A\times\mathfrak D:f(A)=g(D)\},
\]
while every other joint admissibility condition remains. The result may be
a proper subset or empty.  Such a reduction is licensed because
it comes from the independent frontier datum, not from the readout graph
alone.
\end{proof}

\begin{remark}[A readout graph need not preserve the admissible preorder]
\label{p2a:rem:v280-readout-graph-not-order-isomorphism}
The conservative-comparison hypothesis above cannot be replaced by the
bijection $A\mapsto(A,f(A))$.  For a precise conditional countermodel, let
$\mathfrak A=\{a,b\}$ be two candidate-internally admissible parent objects
and let a \emph{previously licensed} output constructor assign the singleton
comparison family $f(a)=\{0\}$ to $a$ and the unbounded family
$f(b)=\mathbb N$ to $b$.  Both graph elements exist as raw mathematical
pairs and projection is bijective, but if the \emph{complete} generated
output family is the consumed physical datum, GP1 excludes the $b$ branch.
No independent occurrence datum was needed for this exclusion: the same
licensed constructor exposed a new joint admissibility obligation.
Alternatively, a generated comparison mark can change which
whole-candidate comparison arrows are admitted.  Therefore the rigorous
no-selection statement is exactly the conservative-graph branch, not the
bare graph bijection.  This example is a countermodel to the invalid
inference, not an assertion that these candidate objects physically occur.
\end{remark}

\begin{theorem}[Generated exact-solution moduli and common-content inheritance]
\label{p2a:thm:generated-exact-solution-moduli}
Let \(P\) be a frozen typed presentation, let
\(D=(d_1,\ldots,d_n)\) be a jointly occurring/generated tuple, and let
\(\mathsf Q\) be a target sort declared before any proposed output is
evaluated.  Suppose \(\Delta_D\) is a finite-exact typed diagram already
generated from \(D\).  Define
\[
 \operatorname{Sol}^{\rm ex}(\Delta_D)
 :=\left\{
 q:\ q\text{ is a complete GP1-admissible realization of }\Delta_D,
 \; \mathfrak D^{\rm GP2}_{\rm joint}(D,q)=0
 \right\}
\]
Require a certified word \(w_{\Delta}\) in the frozen grammar whose predeclared
output in \(\mathsf Q\) is this complete exact-solution family:
\[
 \operatorname{Orig}_{\Sigma}(w_{\Delta};D,\mathsf Q),
 \qquad
 \operatorname{ReqOut}_{\Sigma}(w_{\Delta},D)
 =\operatorname{Sol}^{\rm ex}(\Delta_D).
\]
Thus the completion, transport, quotient, or continuation operation actually
used is licensed before the family is formed.  Let
\[
 \mathfrak M_{\Delta_D}
 :=\operatorname{Sol}^{\rm ex}(\Delta_D)/(\cong,\sim_{\rm cal}),
\]
where structured isomorphism and the optional calibration congruence
\(\sim_{\rm cal}\) are themselves licensed steps of \(w_{\Delta}\), frozen
before the quotient is formed.  Assume only that the resulting structured
presentation lies in the stated GP1 finite/exact domain; naturality,
conservativity, and type preservation are checked properties of the certified
word rather than substitutes for its license.

Then the following statements hold.
\begin{enumerate}[label=(\roman*),leftmargin=2em]
\item If \(\mathfrak M_{\Delta_D}\) is nonempty, it is generated role-level
      family content of \(D\) and
      Theorem~\ref{p2a:thm:rooted-generated-image-corealization} co-realizes it
      in \(\mathsf Q\).  If it is empty, the requested complete exact role is
      FRC-obstructed and GP3 cannot create a solution.
\item If the nonempty family has a unique GP3-least structured isomorphism
      class, GP3 retains it.  If it does not, GP3 retains the unresolved
      irreducible family and no member is selected.
\item For nonempty \(\mathfrak M_{\Delta_D}\), let
      \(S_{\rm GP3}\) be its nonempty complete GP3-retained class family
      after all currently licensed generated frontier obligations have been
      consumed, and let
      \(R:S_{\rm GP3}\to\mathsf Y\) be a licensed, congruence-respecting
      generated readout into a predeclared sort.  If
      \(R(S_{\rm GP3})=\{y_*\}\), then \(y_*\) is common generated content
      even when the source family remains unresolved.  This derives the
      common readout and does not select a source member.
\item When joint admissibility factors coordinatewise, no generated
      cross-role relation exists, and no newly enabled common constructor
      adds an obligation, the joint moduli are the relation-free product.
      If the owner additionally generates the common-value equation for
      typed maps \(r_i:\mathfrak M_i\to\mathfrak B\), the complete joint locus is
      \[
        \mathfrak Z_{\rm joint}
        =A_{\rm joint}\cap R_{\rm prior}\cap
          \{(q_1,q_2):r_1(q_1)=r_2(q_2)\},
      \]
      where \(A_{\rm joint}\) records every joint GP1/typing/constructor
      condition and \(R_{\rm prior}\) every other retained incidence.
      This is the ordinary fibre product precisely when these other
      conditions are redundant on it. Comparison maps alone do not impose
      equality. If the family is specified as a parent image, equality with
      the entire fibre product additionally requires the joint-lift criterion
      of Lemma~\ref{p2a:lem:repair-joint-lift-completeness}.
\end{enumerate}
None of these conclusions actualizes every member of an unresolved family, a
particular solution history, or an output in a sort not declared by the
constructor schema.
\end{theorem}

\begin{proof}
The certificate for \(w_{\Delta}\) places the exact-solution assignment and its
quotient in the constructor-relative closure; its candidate set is not an
ambient solution space formed by bare mathematical permission.  By
construction every member is GP1-admissible and lies in the complete joint
zero-GP2-defect locus.  Quotient by the licensed generated congruence gives,
when nonempty, the generated orbit/family of
Corollary~\ref{p2:cor:generated-orbit-family-inheritance}.  If the locus is
empty, Corollary~\ref{p2:cor:joint-unresolved-family-exhaustion} gives the
stated obstruction.  The rooted generated-image theorem gives physical
co-realization on the nonempty branch.  This proves~(i).

Part~(ii) is exactly the whole-candidate least-class rule of
Postulate~\ref{p2:post:GP3}; neither the existence of a quotient nor the
number of its parameters supplies an additional selector.  For~(iii), first
consume all licensed generated GP1/GP2 frontier obligations and GP3
comparisons and form the nonempty complete retained class family
\(S_{\rm GP3}\).  Its licensed readout has singleton image
\(R(S_{\rm GP3})=\{y_*\}\) by hypothesis, so \(y_*\) is the same generated
content for every surviving class and distinguishes none of them.  The raw
readout graph
\[
 \Gamma_R^{\rm raw}=\{(s,R(s)):s\in S_{\rm GP3}\}
\]
has a bijective set projection to \(S_{\rm GP3}\).  This bijection \emph{does
not} imply an isomorphism of complete admissible GP3 comparison preorders.
The stronger order-isomorphism conclusion requires a conservative extension
for GP1/GP2 candidate status together with preservation \emph{and reflection}
of complete GP3 comparisons, exactly as in
Corollary~\ref{p2:cor:joint-unresolved-family-exhaustion}.  That additional
certificate is not needed for the singleton common-value conclusion.  A
licensed output that changes admissibility or comparison relations has
already been consumed when forming the final survivors; the common value
itself supplies no outcome-driven member selector.

For~(iv), the relation-free composition clause gives the product on
the stated factored joint domain and forbids insertion of an equality.
Otherwise exact admissibility is the conjunction of all generated
conditions. A generated common-value incidence contributes its equalizer
locus without deleting another condition. The universal property of that
locus identifies the ordinary fibre product only on the stated redundant-
constraint branch. A parent image exhausts it exactly when all compatible
tuples have common lifts. The final boundary is
Corollary~\ref{p2a:cor:g0-witness-transport-multiroot-boundary} together with
the declared-output actualization theorem.
\end{proof}

\begin{theorem}[Finite critical stationary relation with retained physical kernel]
\label{p2a:thm:v3110-finite-critical-stationary-relation}
Let $S(w,z)$ be an already generated real action in a finite local coordinate
chart on its admitted $C^2$ domain, with retained coordinates $w$ and interior
coordinates $z\in V$.  Apply only gauge quotients already licensed on this
domain, and let $(w_0,z_0)$ belong to its interior stationary relation,
$d_zS(w_0,z_0)=0$.  Let $G_0$ be the positive-definite kinetic inner product
on $V$ inherited from the physical owner at that point.  Use its fixed Riesz
map in this chart to define $\nabla_z^0S$, and define the self-adjoint
interior Hessian operator $C$ by
\[
 \langle Cv_1,v_2\rangle_{G_0}
 =D_z^2S(w_0,z_0)[v_1,v_2],\qquad
 N:=\ker C,\qquad E:=N^{\perp_{G_0}}.
\]
This fixed splitting is a coordinate construction from the retained metric;
it does not replace any varying kinetic metric in the action.  Write
$z=z_0+n+e$ with $n\in N$, $e\in E$, and let $P_E$ be the
$G_0$-orthogonal projector onto $E$.

There are neighborhoods of $(w_0,0,0)$ and a unique $C^1$ map
$e=\psi(w,n)$, with $\psi(w_0,0)=0$, such that within these neighborhoods
\[
 P_E\nabla_z^0S(w,z_0+n+e)=0
 \quad\Longleftrightarrow\quad e=\psi(w,n).
\]
The reduced action
\[
 \Phi(w,n):=S(w,z_0+n+\psi(w,n))
\]
is $C^2$, and the full interior stationary relation has the exact local
presentation
\begin{equation}
\label{p2a:eq:v3110-critical-stationary-relation}
 \boxed{
 d_zS(w,z)=0
 \quad\Longleftrightarrow\quad
 e=\psi(w,n),\qquad d_n\Phi(w,n)=0.}
\end{equation}
The graph map $(w,n)\mapsto(w,z_0+n+\psi(w,n))$ and its inverse
$(w,z)\mapsto(w,P_N(z-z_0))$ identify the two local solution sets, even
when those sets are singular and no single stationary graph $z=z_*(w)$
exists.  In particular, all surviving critical branches remain in the
relation rather than being selected or discarded.
\end{theorem}

\begin{proof}
Self-adjointness of $C$ implies $C(E)\subseteq E$: for $e\in E$ and
$n\in N$, $\langle Ce,n\rangle_{G_0}=\langle e,Cn\rangle_{G_0}=0$.
The kernel of $C|_E$ is zero, so finite dimensionality makes $C|_E$ invertible.
The map
\[
 F(w,n,e):=P_E\nabla_z^0S(w,z_0+n+e)
\]
is $C^1$, satisfies $F(w_0,0,0)=0$, and has
$D_eF(w_0,0,0)=C|_E$.  The finite-dimensional implicit-function theorem
therefore gives precisely the displayed local graph and its uniqueness.

On that graph $d_zS$ vanishes on $E$.  Since the derivatives of $\psi$ take
values in $E$, the chain rule gives, for $\dot n\in N$ and any $\dot w$,
\[
 d_n\Phi[\dot n]
 =d_zS[\dot n+D_n\psi\,\dot n]=d_zS[\dot n],
 \qquad
 d_w\Phi[\dot w]
 =d_wS[\dot w]+d_zS[D_w\psi\,\dot w]=d_wS[\dot w].
\]
All derivatives of $S$ on the right are evaluated on the graph.  They are
$C^1$ functions of $(w,n)$ because $S$ is $C^2$ and $\psi$ is $C^1$;
hence $\Phi$ is $C^2$.  The direct sum $V=N\oplus E$ now shows that
$d_zS=0$ is equivalent to the complement equation together with
$d_n\Phi=0$.  The stated inverse follows by the fixed orthogonal projection.
Neither implication assumes that the critical relation is a manifold or
that one of its branches is preferred.
\end{proof}

\begin{corollary}[Same-owner reuse of the finite critical relation]
\label{p2a:cor:v3110-critical-relation-owner-reuse}
When the action, interior variation, kinetic metric and stationary-relation
target of Theorem~\ref{p2a:thm:v3110-finite-critical-stationary-relation}
are already licensed, its local re-presentation is an internal theorem
closure after the stated $C^2$ domain is certified, in the sense of
Definition~\ref{p2a:def:v368-theorem-closure-status}.  It is therefore
available as a theorem-generated dependency frontier under
Corollary~\ref{p2a:cor:v176-licensed-theorem-generated-frontier}, with no
additional physical field, smoothing profile, normalization or occurrence
token.  Any later licensed whole-family constructor may consume this same
stationary relation while retaining all of its critical branches and every
other generated constraint and owner incidence.

The kernel coordinate $n$ remains coupled inside the original owner.  The
theorem asserts no autonomous kernel carrier or restriction of other host
operations to $N$; such a restriction would require its own invariance and
typing proof.  GP3 may identify equivalent presentations but supplies no
choice of critical point, local minimum or history branch.  The local
relation result does not establish that an occurring history reaches the
reference point, follows any particular branch, continues globally, or has
uniform control under infinite refinement.
\end{corollary}

\begin{proof}
The theorem retains exactly the already owned stationary solutions in its
local domain through mutually inverse coordinate maps.  All additional
owner relations are retained by transport or intersection with their original
constraints, not removed by the partial elimination.  The operation is thus
a re-presentation of the licensed variational relation rather than a new
physical target or selector.  The cited theorem-closure rule applies after
its analytic hypotheses are met; any certified whole-family use preserves
the same solution alternatives.  No assertion about invariance under other
host operations, actual traversal or global continuation enters the local
implicit-function argument, which establishes the stated scope.
\end{proof}
\begin{proposition}[Necessary second and third constraint jets]
\label{p2a:prop:v257-constraint-obstructions}
For $C:U\subseteq\mathbb R^n\to\mathbb R^m$ of class $C^3$ and a $C^3$ solution curve with initial velocity $v$,
acceleration $a$ and third derivative $b$, put $J=DC_q$. Then
\[
 Jv=0,\qquad Ja+D^2C[v,v]=0,\qquad
 Jb+3D^2C[v,a]+D^3C[v,v,v]=0.
\]
In particular $[D^2C[v,v]]\in\operatorname{coker}J$ must vanish.
These are necessary conditions, not a sufficiency theorem for curve
existence. Even all smooth jets need not detect an isolated zero.
\end{proposition}
\begin{proof}
Differentiate $C(\gamma(s))=0$ successively. For the last warning use
$C(0)=0$, $C(x)=e^{-1/x^2}$ for $x\ne0$: every jet at zero vanishes,
but the zero set is the singleton.
\end{proof}

\begin{proposition}[Effective quartic of the complete critical reduction]
\label{p2a:prop:v257-effective-critical-quartic}
On the finite domain of
Theorem~\ref{p2a:thm:v3110-finite-critical-stationary-relation}, additionally
assume the complete action is $C^4$. Fix the retained coordinates and
center the stationary point at zero. Use the inherited fixed metric and
write $N=\ker C$, $E=N^\perp$, $C_E=C|_E$, $T=D^3S(0)$,
$Q=D^4S(0)$ and
\[
 \langle t(n,n),e\rangle=T[n,n,e].
\]
Then its exact complementary stationary graph obeys
\[
 \psi(n)=-\tfrac12C_E^{-1}t(n,n)+O(|n|^3),
\]
and its reduced action has expansion
\[
 \Phi(n)=S(0)+\tfrac16T[n,n,n]+\tfrac1{24}Q[n,n,n,n]
 -\tfrac18\langle t(n,n),C_E^{-1}t(n,n)\rangle+o(|n|^4).
\]
Thus
\[
 Q_{\rm eff}[n^4]=Q[n^4]-3\langle t(n,n),C_E^{-1}t(n,n)\rangle.
\]
If $T|_{N^3}=0$ and $Q_{\rm eff}[n^4]\ge c|n|^4$ for some $c>0$,
the stationary point is locally isolated. If also $C_E>0$, it is a strict
local minimum. No positive quadratic mass or invertible full Hessian is
implied.
\end{proposition}
\begin{proof}
The complement equation has leading terms $C_Ee+t(n,n)/2$ and gives
the graph expansion by the implicit-function theorem. In the action,
the order-four complement terms are
$\langle C_Ee_2,e_2\rangle/2+T[n,n,e_2]/2$ with
$e_2=-C_E^{-1}t/2$. They sum to $-\langle t,C_E^{-1}t\rangle/8$.
The reduced action is $C^4$ because complement stationarity removes
derivatives of $\psi$ from its first derivative. Therefore
$n\cdot d\Phi(n)=Q_{\rm eff}[n^4]/6+o(|n|^4)>0$ for small nonzero $n$.
No such point is stationary. Positivity of the complementary Hessian
also gives a positive quadratic displacement from its graph and proves
the minimum statement. All licensed gauge reductions precede this
construction; no physical kernel is discarded.
\end{proof}

\begin{remark}[Nonlinear isolation and a Hessian zero mode]
\label{p2a:rem:v257-critical-quartic-control}
The exact two-coordinate action
\[
 S(n,e)=\tfrac12e^2+gen^2+\tfrac\lambda4n^4
\]
has complementary stationary graph $e=-gn^2$ and reduced action
$\Phi(n)=(\lambda-2g^2)n^4/4$. If $\lambda\ne2g^2$, its stationary
point is isolated despite the Hessian kernel. If $\lambda>2g^2$, it is
a strict local minimum; if $\lambda<2g^2$, it has a negative quartic
direction. When $\lambda=2g^2$, this exact model has the stationary
curve $e=-gn^2$. In a general action, a vanishing effective quartic
requires higher orders and does not establish a stationary curve.
\end{remark}



\begin{theorem}[Rooted generated-parent closure cascade]
\label{p2a:thm:rooted-generated-parent-cascade}
Let \(D_\star\) be one jointly occurring/generated root tuple in a frozen
typed presentation.  Consider a finite dependency diagram whose sector
parents are obtained either as sibling descendants
\[
 P_a=\mathsf F_a(D_\star)
\]
or recursively as
\begin{equation}
\label{p2a:eq:rooted-generated-parent-cascade}
 P_{n+1}
 :=\operatorname{Cl}^{\min}_{\rm GP2}
 \bigl(P_n,\operatorname{EssIm}Q_n\bigr).
\end{equation}
For each sibling require a frozen certificate
\[
 \operatorname{Orig}_{\Sigma}(w_a;D_\star,\mathsf P_a),
 \qquad \operatorname{ReqOut}_{\Sigma}(w_a,D_\star)=P_a,
\]
and for each recursive step require
\[
 \operatorname{Orig}_{\Sigma}
   (v_n;(P_n,Q_n),\mathsf P_{n+1}),
 \qquad
 \operatorname{ReqOut}_{\Sigma}(v_n;(P_n,Q_n))=P_{n+1}.
\]
Thus every \(\mathsf F_a\), output sort of \(Q_n\), and closure operation is
licensed and declared before candidate evaluation, with retained outputs
available for reuse and a GP1-finite exact presentation.  Naturality,
conservativity, and type preservation are properties checked on these words,
not sources of constructor license.  The symbol \(Q_n\) is required to be
already generated as either

\begin{enumerate}[label=(\alph*),leftmargin=2em]
\item the unique GP3-retained exact output at stage \(n\); or
\item a licensed readout that is constant on the complete surviving stage-
      \(n\) moduli and hence is common generated content by
      Theorem~\ref{p2a:thm:generated-exact-solution-moduli}.
\end{enumerate}


For every complete exact-solution-moduli output asserted at a node,
also require the certified family word and licensed quotient steps of
Theorem~\ref{p2a:thm:generated-exact-solution-moduli}, with its complete
input tuple already generated at that node and its target declared before
evaluation. A family already retained as a licensed generated output
satisfies this requirement. Certificates for the parents alone establish
the parent cascade; they do not license a new solution-family constructor.

Then every displayed \(P_a\), \(P_n\), and every complete exact-solution
moduli output certified as above is generated from \(D_\star\).  On every nonobstructed step, the occurrence of
\(D_\star\) co-realizes the corresponding descendant role; no second
independent occurrence is required for a certified-word output.

Distinct sector parents need not be equal, and their laws, clocks, rates,
carriers, and defects retain their declared relations and joint admissibility.
A generated common-value incidence intersects the candidate locus with a
fibre product; any other joint conditions remain imposed.  An unresolved
member-dependent output may not be inserted into
\eqref{p2a:eq:rooted-generated-parent-cascade}; doing so would silently select
its source member.

The cascade has no general global-termination conclusion.  For a fixed finite
role lattice it terminates whenever every nonstationary step strictly enlarges
the retained generated subpresentation.  In general it stops at the first
stage where the exact moduli are empty, no unique/common output exists, the
proposed descendant is not projectable or correctly typed, or the required
constructor-origin certificate is absent.
\end{theorem}

\begin{proof}
The first parent is generated by its certificate $w_a$ and
Theorem~\ref{p2a:thm:rooted-generated-image-corealization}.  Suppose inductively
that \(P_n\) is generated from \(D_\star\).  In case~(a), \(Q_n\) is the
GP3-retained exact class; in case~(b), its value is common generated content
and does not select a member of the unresolved source family.  The certificate
\(v_n\) then places the displayed closure in
\(\operatorname{Gen}_{\Sigma}(P_n,Q_n)\); GP2 saturation retains
\(P_{n+1}\).  The rooted generated-image theorem gives co-realization at every
nonobstructed step. Finite induction proves the parent cascade. At each
asserted family output, composition with its additional certified family
word and Theorem~\ref{p2a:thm:generated-exact-solution-moduli} supplies the
moduli conclusion. Corollary~\ref{p3f:cor:v257-classical-family-provider}
instantiates this family word on the complete classical-action branch.

The exact joint-incidence and conditional fibre-product assertion is part~(iv) of
Theorem~\ref{p2a:thm:generated-exact-solution-moduli}.  If \(Q_n\) varies with
an unresolved class, inserting one value would violate whole-family GP3
exhaustion.  The listed stopping conditions are exactly the failures of the
solution-moduli, common-content, typing, projectability, and constructor-origin
gates.  Strict enlargement on a finite role lattice cannot occur
indefinitely; GP1 finite presentation alone does not make the global lattice
finite, so no stronger termination statement follows.
\end{proof}

\begin{corollary}[Whole-family parent reuse without member selection]
\label{p2a:cor:v173-whole-family-parent-reuse}
Let \(D_\star\) be one jointly occurring/generated root tuple and let
\(\mathfrak M_{\Delta_D}\) be a nonempty exact-solution family generated in a
predeclared family-valued sort by
Theorem~\ref{p2a:thm:generated-exact-solution-moduli}.  Suppose a later frozen
constructor word
\[
 v_{\rm fam}:
 (P_n,\mathfrak M_{\Delta_D})\longrightarrow P_{n+1}^{\rm fam}
\]
has a certified origin and required output, consumes the complete family node
in its native sort, and has a GP1-finite exact presentation.  Require that any
member-dependent output be retained fibrewise as a graph or family over
\(\mathfrak M_{\Delta_D}\), and that the word import no section, preferred
member, probability measure, amplitude, or new calibration.

Then \(P_{n+1}^{\rm fam}\) is a rooted generated descendant of \(D_\star\).
On every nonobstructed realization, occurrence of \(D_\star\) co-realizes that
whole-family output without selecting or actualizing a member of
\(\mathfrak M_{\Delta_D}\).  In particular, an unresolved family is not a
cascade stop merely because it has no unique GP3-least member; the stop occurs
only when the required whole-family constructor, output sort, typing, or exact
presentation is absent.

This conclusion does not identify alternative fibres as simultaneous physical
sectors.  Such an interpretation requires an independently generated internal
decomposition, common owner, and dynamics preserving the corresponding
projectors.
\end{corollary}

\begin{proof}
Theorem~\ref{p2a:thm:generated-exact-solution-moduli} places the complete
family node in the constructor-relative closure and retains it in its declared
sort.  The origin certificate for \(v_{\rm fam}\) permits that retained node to
fill a later input port without choosing an element.  Composition of the two
certified words is therefore rooted in \(D_\star\), and GP2 saturation retains
its required output.  The rooted generated-image theorem gives co-realization
on the nonobstructed branch.  Fibrewise graph retention preserves every source
alternative; evaluation at one fibre, a section of the family, or a measure on
it would be a different constructor and is excluded by hypothesis.
\end{proof}


\begin{corollary}[Generated constrained owner-family principle]
\label{p2a:cor:v186-generated-constrained-owner-family}
Let $D_\star$ be a jointly occurring/generated root tuple and let
$\Delta_D$ be a finite typed constraint diagram whose target family sort is
declared before solving.  Suppose a certified constructor word
$w_\Delta$ has required output
\[
 \operatorname{ReqOut}_{\Sigma}(w_\Delta,D_\star)
 =\operatorname{Sol}^{\rm ex}(\Delta_D)
\]
in the sense of Theorem~\ref{p2a:thm:generated-exact-solution-moduli}.  Then a
nonempty complete zero-defect solution family is itself generated
\emph{owner-family content}.  It may be consumed by every later certified
whole-family constructor without selecting a numerical member.

In particular, failure of uniqueness inside the family is not a physical
constructor gap.  A residual gap exists only when a downstream claim requires
one of: (i) a section/member of the family, (ii) a probability or occurrence
law on the family, (iii) a new comparison into a different target sort, or
(iv) a constructor word not already certified by the frozen grammar.  Common
hosting, an ambient isomorphism, or mathematical naturality alone still do
not supply any of these four items.
\end{corollary}
\begin{proof}
Apply Theorem~\ref{p2a:thm:generated-exact-solution-moduli} to retain the
complete nonempty solution family in its predeclared sort, then apply
Corollary~\ref{p2a:cor:v173-whole-family-parent-reuse}.  The four listed
residual cases are exactly the operations excluded by the hypotheses of that
corollary and therefore require their own certified words or G0 occurrence.
\end{proof}


\begin{theorem}[GP-derived generated joint-owner saturation]
\label{p2a:thm:v183-generated-joint-owner-saturation}
Let $D_\star$ be one occurring/generated rooted tuple in a frozen presentation
with constructor grammar $\Sigma$.  Suppose a finite licensed constructor word
\[
 w_{\mathcal O}:D_\star\longrightarrow\mathcal O
\]
has certified origin, declared terminal target sort $\mathcal O$, and consumes
a finite family of generated/occurring ports $X_1,\ldots,X_n$ together with
every declared incidence $R_1,\ldots,R_m$.  Assume the complete word is
GP1-admissible and GP2-exact on the branch under consideration.  Then the
occurrence of $D_\star$ co-realizes the complete joint owner
\[
 \boxed{\operatorname{CoReal}
 (\mathcal O;X_1,\ldots,X_n;R_1,\ldots,R_m\mid D_\star,w_{\mathcal O}).}
\]
No additional G0 occurrence or binding premise is required for any
nonredundant generated port or incidence of that owner.  Omitting such an
output, retyping it, or dropping a declared incidence is a GP2 defect.  If
several complete GP1/GP2-exact owners survive, GP3 acts only after this
joint-owner closure and either retains its unique least class or preserves the
unresolved complete owner family.

The theorem does not license $w_{\mathcal O}$ from common hosting,
co-realization, abstract isomorphism, a nonzero Hom-space, or mathematical
naturality.  Constructor provenance and the physical target sort remain prior
requirements.
\end{theorem}

\begin{proof}
Write the frozen word as a finite composition of licensed constructors.
Rooted generated-image co-realization gives the base step.  At each later
step all required input ports are already retained; Corollary~\ref{p2:cor:no-idle-licensed-constructor}
therefore forces the declared nonredundant output, while GP2 typing,
no-insertion, and relation faithfulness retain it in its declared sort with
its declared incidences.  Finite induction yields the complete owner.
GP1 supplies the obstruction alternative.  GP3 is subsequent and hence
cannot manufacture a missing port or constructor.  The final firewall is the
constructor-origin/typed-cross-relation firewall already imposed above.
\end{proof}

\begin{corollary}[Joint-owner relation closure and the fibre-product case]
\label{p2a:cor:v183-joint-owner-fibre-product}
Under Theorem~\ref{p2a:thm:v183-generated-joint-owner-saturation}, let
$X=\prod_iX_i$ be the declared candidate atlas. Let $A_{\rm joint}\subseteq X$
express the full joint GP1, typing and enabled-constructor admissibility
conditions, and let $R_\alpha\subseteq X$ be all retained incidences. The
complete admissible joint locus is
\[
 \boxed{J=A_{\rm joint}\cap\bigcap_\alpha R_\alpha.}
\]
Suppose, in addition to typed maps $r_i:X_i\to C$, the frozen owner generates
the common-value incidence $r_1(x_1)=\cdots=r_k(x_k)$. Then
\[
 J=A_{\rm joint}\cap R_{\rm prior}\cap
       \{(x_i):r_1(x_1)=\cdots=r_k(x_k)\},
\]
where $R_{\rm prior}$ retains every other incidence. The joint locus is the
ordinary fibre product precisely when those other conditions impose no
additional restriction on that fibre product. A common comparison codomain
alone does not supply the common-value incidence.

When joint admissibility factors coordinatewise, no cross-relation exists,
and no newly enabled joint constructor adds an obligation, the relation-free
Cartesian product is retained. A proposed empty complete locus is an
obstruction. An unresolved nonempty family may be reused only by a licensed
whole-family constructor, without a section or member selection.
\end{corollary}
\begin{proof}
Exact joint admissibility is the conjunction of the complete predicates and
incidences, hence their intersection. A generated equality incidence adds
exactly the displayed equalizer predicate; it cannot delete another retained
constraint. Its locus has the fibre-product universal property by the unique
map $z\mapsto(f_i(z))$ for every compatible cone. This universal property
identifies the licensed object and is not an independent constructor license.
The relation-free case is the GP2 relation-free-composition clause on its
stated domain. GP3 cannot manufacture a point of an empty locus or a preferred
member of an unresolved family.
\end{proof}

\begin{lemma}[Common-parent image and joint-lift completeness]
\label{p2a:lem:repair-joint-lift-completeness}
Let generated maps $a_i:P\to X_i$ and $r_i:X_i\to C$ satisfy
$r_i a_i=c:P\to C$. Then
\[
 \operatorname{im}(a_1,\ldots,a_k)\subseteq\prod_C X_i.
\]
Equality holds if and only if every compatible tuple $(x_i)$ has a common
lift $p\in P$ with $a_i(p)=x_i$ for all $i$. The induced map from $P$ to the
fibre product is bijective if and only if such lifts exist uniquely.
Structural-isomorphism status additionally requires the declared structure
to be preserved by that bijection and its inverse.
\end{lemma}
\begin{proof}
The common comparison gives inclusion. Equality of the image with the
fibre product is exactly surjectivity onto the compatible tuples, which is
the existence of the displayed lifts. Adding uniqueness gives injectivity.
Separate surjectivity of the legs is not joint-lift completeness.
\end{proof}


\begin{theorem}[Typed Parent Correspondence and Action Descent]
\label{p2a:thm:v187-pcad}
Let $R_P\subseteq X_1\times\cdots\times X_n$ be a generated parent relation,
and let $\pi_i:X_i\to D_i$ be licensed maps on the domains actually used by
the word. With $\pi=\prod_i\pi_i$, the licensed image constructor retains
\[
 \boxed{R_D=\pi(R_P).}
\]
This is existential image descent. A representative-independent predicate
with $R_P=\pi^{-1}(R_D)$ requires exactly the saturation condition
$\pi^{-1}\pi(R_P)=R_P$ on the retained domain.

For a binary relation, functionality defines a map from its input projection
support. Totality on the declared input requires that projection to be onto.
Functionality in both directions gives a bijection of the two projection
supports; a full-domain bijection requires both projections to be onto.
For an $n$-ary relation the input-port grouping and output port must first be
specified. An isomorphism in a structured category is retained only when
the forward and inverse maps preserve its declared structure. In particular,
a full two-sided functional subgroup or linear-subspace graph gives the
corresponding group or linear isomorphism; a closed two-sided functional
relation between compact Hausdorff spaces gives a homeomorphism of supports.

Equivariant descendant maps carry invariant parent relations to invariant
descendant relations. A scalar action descends through a quotient only when
it is constant on its fibres. Variations require the admitted chain-rule
domain. A nonlinear Hessian pullback retains its gradient--second-derivative
term unless that term is proved zero. Stationary Schur descent requires an
invertible eliminated block or a separately certified compatible quotient.
For a linear quotient $q$, an operator graph descends to a map exactly
when the operator preserves $\ker q$; for subspace restriction it must preserve
that subspace. Lemma~\ref{p2a:lem:repair-operator-graph-descent} distinguishes
these conditions. These conditions are not consequences of relation
functionality alone.
\end{theorem}
\begin{proof}
The image consists exactly of tuples with a common parent lift. Pulling it
back gives the saturation hull, proving the predicate-descent criterion.
On a projection support, existence of a related point is built into the
support and uniqueness follows from functionality; reverse functionality
then gives a bijection. Subgroup or subspace closure proves preservation of
the corresponding operations. A continuous bijection from a compact space
to a Hausdorff space is a homeomorphism, which applies to the compact graph
projections in the stated closed-relation branch.
Equivariance transports each parent-related tuple and its inverse action
within the same relation. Fibre constancy is precisely the condition for
defining a scalar on a quotient without choosing a representative.
The variation and Hessian statements are the ordinary first and second chain
rules. Differentiating the eliminated stationary equation gives the Schur
formula on its invertible-block domain. Every physical operation used here
remains subject to the frozen constructor-origin and chronology requirements.
\end{proof}

\begin{corollary}[Chronology firewall]
\label{p2a:cor:v187-pcad-chronology-firewall}
PCAD cannot be used retroactively.  The relation $R_P$ must be generated before
the descendant split, or independently generated later by a licensed
constructor.  Common hosting, equal dimension, an abstract isomorphism, a
nonzero Hom-space, or co-realization without a retained parent relation does
not manufacture $R_P$.
\end{corollary}

\begin{definition}[Frozen momentum-seed dependency frontier]
\label{p2:def:momentum-seed-frontier}
After the round momentum parent
\[
 M_p\cong SU(2)\cong S^3(\Rmom)
\]
and its derived \(\Spin(4)\) symmetry have been frozen, the initial Paper-II
momentum-seed dependency ledger records the frontier
\[
 \mathcal F_p^{\rm seed}
 :=\{E_0,\ S_{0,+},\ S_{0,-},\ E_1,
       \ \nabla,\ [\,\cdot,\cdot\,]_{\Gamma(TS^3)}\}.
\]
Here \(E_0\) is the scalar unit, \(E_1\) is the first nonconstant scalar
harmonic object, \(S_{0,\pm}\) are the two intrinsic lowest spinor
objects, and \(\nabla\) and the vector-field Lie bracket are the round-host
natural constructors.  The tangent-section codomain is typed by
Definition~\ref{p2:def:parent-geometric-object-occurrence} as the same-role
parent object sort \(\Gamma(TS^3)_{\rm sec,obj}\), while the induced action on
smooth functions is a separate derivation occurrence.  The constituted
composite constructor is therefore
\[
 \beta:E_1\times E_1\longrightarrow\Gamma(TS^3)_{\rm sec,obj},\qquad
 \beta(f,g)=[\nabla f,\nabla g].
\]
Its alternation permits the mathematical factorization
\[
 E_1\times E_1\longrightarrow\Lambda^2E_1
 \xrightarrow{\,\bar\beta\,}\Gamma(TS^3)_{\rm sec,obj},
\]
but \(\Lambda^2\) is classificatory here and is not added to the frozen
physical constructor grammar.  Paper~II-B must verify from the previously
established round geometry the stated round-host provenance and the finite
canonical image before using it.

The seed frontier contains no finite level, component ceiling, tensor-word
boundary, candidate-specific projector, compression map, or unrestricted
tensor/polynomial constructor.  In particular, ordinary pointwise
multiplication of scalar harmonics remains available in the smooth parent
mathematics but is not, merely by that availability, a finite momentum-carrier
constructor.  A later stage may constitute such a constructor through a new
dependency frontier, in which case the recursive GP calculus is reapplied to
the enlarged signature.  Completeness of \(\mathcal F_p^{\rm seed}\) is
therefore relative to the declared initial momentum-seed dependency ledger in
the sense of Definition~\ref{p2:def:recursive-gp-application}.
\end{definition}

\begin{remark}[Later discharge does not alter theorem ownership]
\label{p2:rem:recursive-gp-discharge}
Definition~\ref{p2:def:recursive-gp-application} does not require each paper
or each layer to discharge the antecedents of its own conditional theorems.
If stage \(n\) proves \(H\Rightarrow T\), and a later layer independently
generates content whose recursive GP closure establishes \(H\), then the
earlier implication supplies \(T\).  The earlier stage owns the conditional
theorem; the later stage owns only the discharge of its antecedent.  A
specialized conclusion \(T\) may not be used to generate the later datum that
is then invoked to establish \(H\).
\end{remark}

\begin{remark}[Static-algebra scope guard]
\label{p2:rem:static-algebra-scope-guard}
A finite represented static \(*\)-algebra does not determine which injective
spectral parametrization, parent differential order, locality class, principal
symbol, or dispersion law generated it.  The precise reparametrization no-go
and selector corollary are deferred to
Paper~II-F, Section~\ref{p2:app:selector-limitations}.  This is why the paired realization
below retains the filtered parent grammar and declared compression/quantization
map rather than attempting to reconstruct them from one static finite algebra.
\end{remark}

\begin{definition}[Paired parent/finite realization]
\label{p2:def:paired-realization}
For each constituted role \(R\) and finite level \(m\), Paper~II uses the
following concrete structured realization of
Definition~\ref{p2:def:structured-role-realization}, augmented by its parent
presentation data:
\[
 \mathfrak R_{R,m}
 :=\left(
     \mathfrak F_{R,m},
     \mathfrak P_R,
     \operatorname{Fil}^{\bullet}\operatorname{Op}(\mathfrak P_R),
     Q_{R,m}
   \right).
\]
Here \(\mathfrak F_{R,m}\) is the finite ontic realization,
\(\mathfrak P_R\) is the smooth or equivalently filtered parent grammar,
\(\operatorname{Fil}^{\bullet}\operatorname{Op}(\mathfrak P_R)\) records
locality, differential order, ellipticity, and principal symbols, and
\(Q_{R,m}\) is the declared compression or quantization map.  Parent-side
operator selection occurs before compression.  The parent is not an
infinite-capacity physical stage.  Whenever \(\mathfrak P_R\) contains an
exact-operation group, that group means the closed exact operation image;
when its generated finite state/test action is affine, it is unitarized in
that finite faithful representation by
Theorem~\ref{p2:thm:frc-automatic-compact-lie-closure}, not a
nonclosed dense parametrized subgroup of a larger compact group.
\end{definition}

\begin{definition}[Paper-II theorem-status convention]
\label{p2:def:paper-ii-status-convention}
The following statuses are available for reading any Paper-II result, and are
used wherever a result's dependency class is not evident from its statement:
\begin{enumerate}[label=(\roman*),leftmargin=2em]
\item \([\mathrm{G0}{-}\mathrm{OCC}]\): a physical occurrence, preparation,
test, process, trajectory, or cross-role bridge is represented only after a
contingent token, a geometrically licensed claim-relative invariant
distinction, and result-independent provenance are supplied as in
Definition~\ref{p2:def:g0-occurrence-types}; the tag alone asserts no
occurrence;
\item \([\mathrm{G0}{-}\mathrm{CAL}]\): an already constituted invariant is
expressed through a declared comparison with a realized standard or readout
map.  If the standard and map are already constituted, this is
instantiational and fixes numerical presentation only.  A new standard,
comparison law, or interface is constitutive and enters through its own
downstream dependency frontier; neither branch alters the carrier;
\item \([\mathrm{PARENT}]\): exact parent-grammar theorem, with no ontic
infinity claim;
\item \([\mathrm{HOST}]\): exact constructor-relative host-closure theorem;
its carrier content is selected from the generated host signature rather than
from a preassigned finite level;
\item \([\mathrm{FIN}{-}m]\): exact for arbitrary finite \(m\);
\item \([\mathrm{SPEC}{-}m_0]\): mathematical specialization at a named
finite level, with no physical selection;
\item \([\mathrm{COND}{-}H]\): physical-regime conclusion conditional on a
registered hypothesis \(H\);
\item \([\mathrm{BAND}{-}B]\): exact stabilization on a fixed band above a
finite threshold;
\item \([\mathrm{ASYM}{-}B,\varepsilon]\): controlled finite-level
approximation with an explicit remainder;
\item \([\mathrm{DEFERRED}]\): discharge, normalization, physical
identification, or dynamics outside the theorem being stated.
\end{enumerate}
A conditional theorem is a complete kinematic implication.  Only the truth of
its antecedent remains to be discharged.
\end{definition}

\begin{remark}[Paper-II parent, seed, and unfolding architecture]
\label{p2:rem:carrier-seed-unfolding}
Paper~II owns the parent geometry, the constructor-relative momentum seed, the
generic finite tower used where explicitly constituted, every finite
specialization, and every quantified smooth/local recovery implication:
\[
\boxed{\begin{aligned}
 \Sigma_{\rm primitive}+\mathrm{FRC}+\mathrm{GP2}+\mathrm{GP3}
 &\Longrightarrow U(1)\cong S^1,\quad SU(2)\cong S^3,\\
 \text{II-B round-host closure}
 &\Longrightarrow \mathcal H_p\cong\mathcal H_1
    \Longrightarrow (\mathrm{LVL}_p^1),\\
 \begin{gathered}\text{GP2 reciprocal-role closure}\\
   +\text{intrinsic }E_1\text{ spectrum}\end{gathered}
 &\Longrightarrow S^3_B,\quad
   \mathcal P_x(m)=\bigoplus_{k=0}^{m}E_k^x,\\
 \mathcal P_x(1)=E_0^x\oplus E_1^x
 &\Longrightarrow m_x[0]=1,\\
 \begin{gathered}\text{chronological persistence}\\
   \text{of the frozen momentum object}\end{gathered}
 &\Longrightarrow m_p=1\text{ remains fixed},\\
 \text{generated finite frontier }F
 &\Longrightarrow m_x^{\rm res}(F)
   =\min\{m:F\subseteq\mathcal P_x(m)\},\\
 (\mathrm{RES}_x{:}B,\varepsilon)
 &\Longrightarrow \text{smooth-parent kinematics on }B,\\
 (\mathrm{RES}_x)+(\mathrm{FLAT}_x)+(\mathrm{PAIR}_{xp})
 &\Longrightarrow \mathsf K_{\rm local}^{B,L,\varepsilon}.
\end{aligned}}
\]
Here the second line is a theorem-ownership statement, not an additional
antecedent: II-B imports the frozen momentum-seed frontier of
Definition~\ref{p2:def:momentum-seed-frontier}, verifies its round-host
provenance, and derives the momentum host closure using the finite
generated-image inheritance
Lemma~\ref{p2:lem:gp2-finite-generated-image-inheritance} together with the
same-role no-splitting Lemma~\ref{p2:lem:gp2-same-role-no-splitting}.  Only afterward may it attach
the word-length label \(m_p=1\).  The generic word tower remains available only
for roles whose constructor grammar explicitly contains it.  Paper~II-C
instead derives the scalar position filtration and its five-dimensional seed
from the intrinsic coordinate spectrum.  Its least-resolution theorem assigns
a unique finite level to every generated finite position frontier; nontrivial
growth occurs exactly when later generated content has a nonzero harmonic
component above the current level.  No such higher component, and hence no
nonconstant trajectory, is asserted here.  The reciprocal-role line uses
Theorem~\ref{p2:thm:gp2-generated-reciprocal-role-closure}: downstream work
must generate the reciprocal separating spectrum and prove the
\((\mathrm{BR}_{\rm ind})\) and \((\mathrm{IND}_{xp})\) antecedent statuses,
but those statuses are not additional inputs once proved.  If one or more
later antecedents require dynamics, Paper~II does not defer the kinematics: it
proves the corresponding implications conditionally, and dynamics may
discharge those antecedents non-circularly.
\end{remark}

\begin{definition}[Constructor-relative GP2 generated closure]
\label{p2:def:gp2-generated-closure}
Let \(\Sigma\) be the frozen constituted host signature and let \(X\) be the
content constituted at the current stage.  The \emph{GP2 generated closure}
\(\mathcal G_\Sigma(X)\) is the smallest typed subpresentation containing
\(X\) and closed under:
\begin{enumerate}[label=(\roman*),leftmargin=2em]
\item every operation and constructor schema actually constituted in
\(\Sigma\) whenever its typed inputs are available;
\item the canonical compositions and partial evaluations required by those
schemata;
\item their generated output images in the declared target sorts, including
finite generated-image inheritance from already constituted parent object
sorts and canonical-image objectification when the antecedents of GP2 apply;
\item generated subobjects and quotients only when the corresponding
subobject, congruence, or quotient construction is itself generated; and
\item universal relation-free composition of independently generated roles
when the signature supplies no cross-relation.
\end{enumerate}
Arbitrary direct sums, tensor products, symmetric powers, exterior powers,
isotypic decompositions, polynomial constructions, or theorem consequences
are not members of \(\mathcal G_\Sigma(X)\) merely because they are
mathematically available.  They enter physical generation only when the
corresponding constructor schema belongs to \(\Sigma\).  Representation theory may nevertheless be used after generation to classify
the type of a canonical image without thereby constituting the full
representation-theoretic closure.  Simultaneous realization of several
independently generated invariant carrier sectors in one carrier is assembly of
the exact type-support required by the host signature, not a new direct-sum
constructor that generates additional representation types.
\end{definition}

\begin{definition}[Budgeted generated closure and the level normalization]
\label{p2:def:budgeted-finite-closure}
For the current \(\Spin(4)=SU(2)_L\times SU(2)_R\) representation grammar put
\[
 F_L:=V_{\frac12,0},\qquad F_R:=V_{0,\frac12}.
\]
When the fundamental pair has been established as the atomic constructor seed
of the frozen parent, define intrinsic tensor-word length by
\[
 \ell(V_{j_L,j_R})
 :=\min\{r:\ V_{j_L,j_R}\prec(F_L\oplus F_R)^{\otimes r}\}.
\]
The representation-theoretic word-budget argument gives
\[
 \ell(V_{j_L,j_R})=2j_L+2j_R.
\]
The half-level convention is therefore the normalization
\[
 m:=\ell/2,
 \qquad \deg F_L=\deg F_R=\frac12,
\]
not an independent numerical selector.  For
\(m\in\frac12\mathbb Z_{\ge0}\), define the budgeted word object and its
multiplicity-free irreducible support by
\[
 \mathcal W_{\le m}
 :=\bigoplus_{0\le r\le 2m}(F_L\oplus F_R)^{\otimes r},
 \qquad
 \mathcal C_R(m):=\operatorname{IrrSupp}(\mathcal W_{\le m}).
\]
Thus \(\mathcal C_R(m)\) contains one unmarked copy of every irreducible
isomorphism class generated within intrinsic word budget \(2m\).

This definition is representation-theoretic bookkeeping.  It separates the
atomic constructor seed, the intrinsic word-length grading, and any later
physical decision to constitute a tensor-word grammar.  The unfiltered
mathematical tensor closure remains infinite, but revised GP2 does not identify
that entire closure with physical generation.  If a later role explicitly
constitutes the tensor-word constructor grammar, then a generated or registered
boundary \(\mathsf P_{\rm gbd}\) may supply a finite scope and Proposition~
\ref{p2:prop:gp1-gp2-bounded-generated-exactness} exactifies that scope.  For
the initial momentum seed, the word tower may instead be used only to classify
the constructor length of a carrier selected independently from the host
grammar.
\end{definition}

\begin{proposition}[Word budget equals the total-weight lower ideal]
\label{p2:prop:word-budget-lower-ideal}
For every \(m\in\frac12\mathbb Z_{\ge0}\),
\[
 \mathcal C_R(m)
 =\bigoplus_{j_L+j_R\le m}V_{j_L,j_R}
\]
as a multiplicity-free \(\Spin(4)\)-module.  Hence the intrinsic
word-length sublevel \(\ell\le2m\) and total-weight sublevel
\(j_L+j_R\le m\) are the same generated support once the atomic constructor
grammar of Definition~\ref{p2:def:budgeted-finite-closure} is available.
\end{proposition}

\begin{proof}
A word with \(r_L\) left and \(r_R\) right fundamentals has
\(r=r_L+r_R\), and every irreducible constituent satisfies
\(j_L\le r_L/2\) and \(j_R\le r_R/2\).  Thus every constituent of a word of
length at most \(2m\) has \(j_L+j_R\le m\).  Conversely,
\(V_{j_L}\) occurs as the highest-spin summand of
\((V_{1/2})^{\otimes 2j_L}\), and similarly on the right.  Therefore
\(V_{j_L,j_R}\) occurs in a word of length
\(2j_L+2j_R\le2m\).  Taking one unmarked copy of each isomorphism class proves
the equality.
\end{proof}

\begin{definition}[Typed structural-status, occurrence, bridge, and recovery registry]
\label{p2:def:level-tags}
The following registry separates theorem-status certifications, actual G0
occurrences, open bridge/calibration interfaces, and recovery conditions.  A
label in this registry is therefore not automatically a selector assumption.
Generated/co-realized carriers and constituents are theorem outputs and need no
second occurrence tag unless a genuinely new contingent token is introduced.

\subparagraph{Derived structural-status labels.}
\begin{enumerate}[label=(\roman*),leftmargin=2em]
\item \((\mathrm{BR}_{\rm ind})\): theorem-status certification that a
generated reciprocal separating test spectrum satisfies the generated-role
objectification antecedent of GP2: it carries its own faithful operation role
and autonomous variance.  This is not an independent constitutive hypothesis.
By Theorem~\ref{p2:thm:gp2-generated-reciprocal-role-closure}, once the status
is proved GP2 itself constitutes the reciprocal spectrum as an autonomous
physical role.
\item \((\mathrm{IND}_{xp})\): theorem-status certification that the frozen
relation slot contains no prior generated joint-support constraint linking the
constituted position and momentum roles.  A total comparison map or generated
role isomorphism is not such a constraint unless a graph/equality condition on
joint values is separately generated.  This is not an independent
constitutive hypothesis: once the status is proved, GP2 supplies the universal
relation-free composition by
Theorem~\ref{p2:thm:gp2-generated-reciprocal-role-closure}.  The status does
not supply a reciprocal normalization or local physical interpretation map.
\item \((\mathrm{LVL}_p^1)\): \(m_p=1\).  This tag records a theorem output or
downstream interface; under the revised momentum constitution it is not a
primitive selector premise.
\item \((\mathrm{LVL}_{x,0}^1)\):
\(m_x(\widetilde t_0)=1\).  Paper~II-C derives this as the least resolution
of the intrinsic \(E_0^x\oplus E_1^x\) seed; it is not a primitive selector.
\item \((\mathrm{SEED}_{11})\): the conjunction
\((\mathrm{LVL}_p^1)\wedge(\mathrm{LVL}_{x,0}^1)\).
\item \((\mathrm{FRZ}_p)\): shorthand for the theorem conclusion
\(m_p(\widetilde t)=1\) throughout the specified branch or epoch, obtained
from \((\mathrm{LVL}_p^1)\) together with chronological persistence of the
already frozen momentum object.  It is not an independent constitutive
premise.
\item The canonical first-symbol position--transvection pairing is a
Paper-II-C theorem and a generated physical-role comparison.  Its legacy
\((\mathrm{OBS}^{(1)}_{xp})\) label denotes only a particular performed
comparison token, not the structural pairing itself.
\item On any retained finite position-frontier history,
\((\mathrm{MON}_x)\) is only a theorem-status shorthand for
\[
  s\preceq t\Longrightarrow F_x(s)\subseteq F_x(t)
  \Longrightarrow m_x(s)\le m_x(t).
\]
It is not a separate premise or selector on that branch.
\end{enumerate}

\subparagraph{Generated/co-realized structural content.}
The reciprocal role, its relation-free product with the frozen parent, every
nonredundant generated constituent of an occurring composite, and the
constructor-relative momentum carrier are structural outputs rather than
selector entries.  Their role occurrence is inherited by
Theorems~\ref{p2:thm:generated-role-corealization} and
\ref{p2:thm:generated-constituent-corealization} whenever the stated typed
constructor hypotheses hold.  They therefore require no duplicate G0 witness
and no memberwise selector merely because a later calculation focuses on one
co-realized constituent.  A new contingent preparation, performance, history,
or calibration remains a new occurrence token when it is not generated from
the existing root.

\subparagraph{Actual G0 occurrence labels.}
\begin{enumerate}[label=(\roman*),leftmargin=2em,resume]
\item \((\mathrm{HIST}_x^{\rm fin})\): an outcome-independent physical
history of finite-dimensional generated position-test frontiers.  In the
Paper-II-C notation used here,
\[
 \mathcal A_x^{\rm alg}
 :=\operatorname{alg}_{\mathbb C}^{*}(E_0^x\oplus E_1^x)
 =\bigcup_{m\ge0}\mathcal P_x(m),
\]
and the history is
\[
  \widetilde t\longmapsto F_x(\widetilde t),
  \qquad
  F_{x,0}\subseteq F_x(\widetilde t)\subset\mathcal A_x^{\rm alg}.
\]
Its level is not independently chosen but is the derived least scalar
harmonic-resolution statistic
\[
  m_x(\widetilde t)
  :=m_x^{\rm res}(F_x(\widetilde t))
  =\min\{m:F_x(\widetilde t)\subseteq\mathcal P_x(m)\}.
\]
FRC permits \(\sup_{\widetilde t}m_x(\widetilde t)=\infty\) across a whole
history while forbidding an infinite realized stage.  A sparse frontier and
the complete band \(\mathcal P_x(m_x)\) with the same resolution are not
identified.  The legacy symbol \((\mathrm{TRJ}_x)\), where encountered in an
older dependency ledger, is retired as an independent integer trajectory and
may refer only to this derived statistic of an occurring
\((\mathrm{HIST}_x^{\rm fin})\).
\item A particular \((\mathrm{OBS}^{(1)}_{xp})\) performance, preparation,
test, outcome, empirical record, or calibration is a contingent G0 token even
when its role-level carrier, comparison, or effect space is already generated.
\end{enumerate}

\subparagraph{Open bridge/calibration interfaces.}
\begin{enumerate}[label=(\roman*),leftmargin=2em,resume]
\item The residual tag \((\mathrm{PAIR}_{xp})\) denotes a separately proved
stronger finite-band or continuum local/global comparison identifying the
declared position harmonics with normalized local momentum observables on its
stated band, topology, and error bound.  This remains a genuine downstream G0
bridge/calibration certificate; G0 types but does not prove its witness.
\item A claim depending on a particular ambient/raw coordinate action, omitted
Wigner--Eckart coefficients, an absolute protected radius, raw coordinate
normalization, or coordinate--coordinate commutator magnitude must state the
required completion \(\mathsf C\in\mathfrak M_{\rm coord}\) or derive an
independent discriminator.  The protected even \(\Cl(0,3)\) algebra on
\(\varrho_{\rm sp}>0\) is excluded from this completion-selection interface,
because current Paper~II-C proves it independent of the nonzero scale and
relative chiral phase.
\end{enumerate}

\subparagraph{Optional finite-role grammar and recovery conditions.}
\begin{enumerate}[label=(\roman*),leftmargin=2em,resume]
\item \((\mathrm{TWR}_R)\): the realized role \(R\) is represented by a
stagewise bounded member of an intrinsic generated finite-role family
\(\{\mathfrak F_{R,m}\}_{m<\infty}\).  This is an optional role-specific
interface, not the default meaning of GP2 generation.  It applies when the
host signature has actually constituted the corresponding finite-role grammar.
The grading of Definition~\ref{p2:def:budgeted-finite-closure} may remain
useful mathematically even when \((\mathrm{TWR}_R)\) is not a constitutional
premise.
\item \((\mathrm{SPEX}_R)\): for a role whose host signature has explicitly
constituted a finite generated-depth scope, GP2 exactness retains the complete
generated support inside that scope.  In the tensor-word implementation the
multiplicity-free mathematical support is \(\mathcal C_R(m)\), and
Proposition~\ref{p2:prop:word-budget-lower-ideal} identifies it with the
total-weight lower ideal.  This tag remains useful for later tower roles but is
not required when GP2 constructor-relative host closure already fixes the
carrier support directly.
\item \((\mathrm{FIT}_R^1)\): the complete generated role content of an
already constituted finite-level grammar is faithful at level one and not at
the preceding level.  This remains a useful diagnostic for generic tower
roles.  It is no longer the constitutional route intended for the initial
momentum carrier: there Paper~II-B is assigned the direct host-closure theorem,
and \(m_p=1\) is read off only after \(\mathcal H_p\cong\mathcal H_1\) has been
derived.
\item \((\mathrm{RES}_x{:}B,\varepsilon)\): for a declared spectral or
observational band \(B\) and \(\varepsilon>0\),
\(m_x\ge M(B,\varepsilon)\), where \(M(B,\varepsilon)\) is the threshold of
the controlled-asymptotic schema displayed above.
\item \((\mathrm{FLAT}_x{:}L,\varepsilon)\): a separate curvature/radius
estimate makes regions of size \(L\) approximately Euclidean to accuracy
\(\varepsilon\).  Large \(m_x\) alone does not imply this condition.  This
and the preceding recovery bound remain analytic/geometric theorem-domain
conditions; G0 does not derive their estimates.
\end{enumerate}
\end{definition}

\begin{definition}[Typed continuum-recovery contract]
\label{p2:def:typed-continuum-contract}
Let \(P_B\) project onto a finite parent band and let
\(J_{R,m,B}:\operatorname{Ran}P_B\to\mathfrak F_{R,m}\) be the declared
band inclusion when \(m\) contains \(B\).  Put
\(\mathcal O_{R,m}=Q_{R,m}(\mathcal O_{\rm par})\) whenever the finite
operator is the compression of the parent operator.  Every cross-level claim
is stated in the strongest of the following three forms that holds; the forms
are nested, so the strongest available is the one to state.
\begin{enumerate}[label=(\roman*),leftmargin=2em]
\item \emph{Finite-exact}: an identity holds at every finite \(m\).
\item \emph{Band-stabilized}:
\[
 J_{R,m,B}^*\mathcal O_{R,m}J_{R,m,B}
 =P_B\mathcal O_{\rm par}P_B
 \qquad(m\ge M(B)).
\]
\item \emph{Controlled asymptotic}:
\[
 \bigl\|J_{R,m,B}^*\mathcal O_{R,m}J_{R,m,B}
       -P_B\mathcal O_{\rm par}P_B\bigr\|
 \le r_B(m),
 \qquad r_B(m)\to0.
\]
\end{enumerate}
Thus the phrase ``\(m_x\) very large'' is inadmissible without an explicit
threshold or remainder.  The generic schema is
\[
 \forall B,\varepsilon>0\ \exists M(B,\varepsilon)\ \forall m_x\ge M:
 \operatorname{Err}_B(m_x)\le\varepsilon.
\]
\end{definition}

\begin{remark}[Typed discharge/occurrence ledger and non-circularity]
\label{p2:rem:level-discharge-ledger}
The registry distinguishes theorem-status results, generated/co-realized
content, contingent G0 occurrences, and genuinely open bridge, completion,
calibration, and recovery interfaces.  The source of an open downstream
antecedent is not fixed in advance.
The momentum statement \((\mathrm{LVL}_p^1)\) is not in that open ledger:
Paper~II-B owns its discharge by first proving the direct round-host theorem
\(\mathcal H_p\cong\mathcal H_1\) and then reading off \(m_p=1\).
Likewise \((\mathrm{BR}_{\rm ind})\) and \((\mathrm{IND}_{xp})\) are not open
physical hypotheses: they are theorem-status certifications of antecedents in
Theorem~\ref{p2:thm:gp2-generated-reciprocal-role-closure}.  A later
reciprocal-spectrum construction may prove those statuses, but no additional
constitutive postulate is then discharged.  Paper~II-C derives
\((\mathrm{LVL}_{x,0}^1)\) from the intrinsic reciprocal spectrum, while
Paper~II-A's chronological-persistence theorem supplies
\((\mathrm{FRZ}_p)\); neither belongs to the open ledger.  Paper~II-C
additionally closes the role-level first-symbol comparison.  It also derives
the least-resolution functor for every later generated finite position
frontier and proves monotonicity from retained inclusion.  Thus the genuinely
open downstream ledger contains the stronger finite-band/local--global G0
bridge \((\mathrm{PAIR}_{xp})\), the actual frontier-history occurrence
\((\mathrm{HIST}_x^{\rm fin})\), completion/calibration data only for claims
that genuinely depend on them, and the separate recovery conditions.  The
legacy \((\mathrm{TRJ}_x)\) is not an open selector, and
\((\mathrm{MON}_x)\) is not a separate premise on retained histories.  If
kinematics does not suffice, a structural or dynamical theorem may generate
higher frontier content, relate the actual frontiers to physical time, or
discharge the remaining bridge/recovery antecedents.
In every case the discharge proof may consume only parent theorems and the
generic finite family \(\mathsf K(m_x,m_p)\), never a specialization whose
proof already assumes the conclusion being discharged.  Formally,
\[
 \mathsf K_{\rm gen}\longrightarrow\mathsf{DYN}
 \longrightarrow H\longrightarrow\mathsf K_{\rm spec}(H),
\]
not \(\mathsf K_{\rm spec}(H)\to\mathsf{DYN}\to H\).

A large-position theorem must state its strength: existence of one branch,
branch-conditional growth, or forced growth for every admissible solution.
Only the last licenses an unconditional statement that dynamics forces
\(m_x\) large.  Finally, \((\mathrm{SEED}_{11})\) does not by itself imply
\(R_x(\widetilde t_0)=\ell_{\rm B}\).  The independently allocated scalar
\(H_{\rm succ}\cong U(1)\) is not identified with a circle inside the
nonabelian transfer parent.  On the scalar-succession branch,
Definition~\ref{p2:def:hopf-succession-provenance} instead constitutes a
typed provenance arrow from the already frozen Paper-I-A Hopf-flow occurrence;
Theorem~\ref{p2:thm:transported-succession-orientation} transports its
orientation without a carrier or subgroup identification.
\end{remark}





\begin{remark}[Covariant-readout existence versus physical occurrence]
\label{p2a:rem:v270-covariant-readout-ledger}
Paper~II-C, Theorems~\ref{p2c:thm:covariant-stabilizer-capacity}
and~\ref{p2c:thm:covariant-direction-complete-povm}, prove a
\emph{mathematical} Spin(4)-covariant effect-capacity hierarchy
$65/107/137$ for position, position plus an unoriented tangent axis,
and position plus a directed tangent. At radius $r>\sqrt3\Rmom$,
$8/8/6$ passive homogeneous-orbit labels are sharp for saturating
those spans with the existing intrinsic first moments.  In the
directed case full $M_{11}\oplus M_4$-native tomography is thus
available without an all-band map into the frozen momentum carrier.
This is not a new G0-typed physical observation: neither a detector
incidence, a prepared state, nor a CP instrument is forced by the
constructor merely because the corresponding positive POVM exists.
The local first-symbol role is already generated; transporting its
unit tangent bundle as a \emph{mathematical outcome geometry} does not
constitute a physical detector, nor does the distinct left-Hopf mark
choose a global tangent direction.  All GP3, frontier-occurrence,
charge and calibration ledger qualifications remain in force.
\end{remark}

\begin{definition}[Geometrically constituted primitive-role existence]
\label{p2:def:geometric-role-existence}
For a primitive physical role \(R\), Geometric Primacy constitutes the
role-geometric existence assertion whose actual occurrence is certified under
G0,
\[
 \boxed{\mathsf E_R:\quad
 \exists\,\mathfrak M\;\operatorname{Geo}_R(\mathfrak M).}
 \tag{\ensuremath{\mathsf E_R}}
\]
The predicate \(\operatorname{Geo}_R\) is a strong condition on a whole
candidate realization, not bare set-theoretic existence.  It requires:
\begin{enumerate}[label=(\roman*),leftmargin=2em]
\item a nonempty comparison geometry \((X_{\mathfrak M},d_{\mathfrak M})\);
\item an identity-rooted continuous local reapplication constructor containing
a full connected identity neighbourhood, with its finite sequential
reapplications spanning the candidate's generated connected operation role;
\item an action \(\alpha_{\mathfrak M}\) of those reapplications by comparison
isometries, with an effective occurrence
\[
 \exists a\ \exists x\in X_{\mathfrak M}:\qquad
 \alpha_{\mathfrak M}(a)x\ne x;
\]
and
\item the intrinsic role diagrams which compare the two well-typed sequential
reapplications \(a\circ b\) and \(b\circ a\), together with whatever
comparison/test grammar the geometry itself generates.
\end{enumerate}
A role-faithful candidate must realize those diagrams exactly.  Thus the
geometry, rather than a free label, determines whether the order-swap diagram
is identically equal or is distinguishable somewhere.  No operation, point,
axis, or witness pair is named or marked.  The role-indexed predicate is read
with the actual equalities and inequivalences of its intrinsic diagrams; it is
not a disjoint union of incompatible diagram outcomes.  Consequently two
candidates for the same \(\operatorname{Geo}_R\) cannot choose different
order-swap branches, although the labels \(\mathrm{ab}\) and \(\mathrm{nc}\)
are introduced only after exactification.

The assertion \(\mathsf E_R\) does not choose a completed ambient group or
freeze one candidate's multiplication table before GP3.  It contains no free
labels for nontriviality, connectedness, abelianity, or noncommutativity and no
chosen dimension, rank, quotient relation, representation, metric scale, or
global form.  Nevertheless, clauses (ii)--(iv) explicitly constitute
connected local generation, effectiveness, and the order-swap diagram.  A
faithful realization therefore has nontrivial/connected statuses and one
classified order-swap branch; GP1--GP2 realize and preserve this content rather
than creating it.  If the geometrically primary role genuinely contains
additional intrinsic directions, brackets, relations, or tests, they remain
inside \(\operatorname{Geo}_R\) and every role-faithful comparison must preserve
them; the primitive endpoint theorem below then applies only if its endpoint is
still a realization of that same role.
\end{definition}

\begin{lemma}[Connected exact-operation closure without a Lie hypothesis]
\label{p2a:lem:v281-connected-generated-group}
Let \(K\) be a compact topological group and let \(C\subset K\) be a
connected subset containing the identity.  If
\[
 K=\overline{\bigcup_{n\geq1}C^n},
\]
then \(K\) is connected.  In particular, for a primitive
\(\operatorname{Geo}_R\) candidate whose generated closed exact parent
is exhausted by the finite reapplications of its constituted connected
identity-neighbourhood and has no separately generated discrete
operation seed, the GP1 compact reversible closure is a connected
compact \emph{topological} group, even if no faithful finite affine
state/test representation has yet been certified.  Its algebraic
commutator branch (trivial or nontrivial) is well-defined without a
Lie hypothesis.  Effectiveness of the constituted reapplication
makes that group nontrivial.
\end{lemma}
\begin{proof}
Each finite product \(C^n\) is connected as the continuous image of
\(C^n\) under multiplication, and every \(C^n\) contains the
identity; the nested union \(\bigcup_{n\geq1}C^n\) is connected.
Its closure is connected, hence equals a connected \(K\) by the
stated generation hypothesis.  In the constituted primitive role,
clause (ii) of Definition~\ref{p2:def:geometric-role-existence}
provides the generating connected identity neighbourhood and GP2
no-insertion prevents adding a distinct discrete seed.  Compactness
and reversal are supplied by
Theorem~\ref{p2:thm:gp1-bounded-isometry-closure}; an effective
reapplication remains nonidentity by typed host faithfulness.  The
commutator equal/non-equal dichotomy exists for every group, not just
Lie groups.
\end{proof}

\begin{theorem}[Faithful realization and classification of constituted role geometry]
\label{p2:thm:geometric-existence-exactification}
Let \(\mathfrak M\) satisfy \(\operatorname{Geo}_R(\mathfrak M)\), GP1, and
candidate-internal GP2 exactness, and let \(H_{\mathfrak M}\) be its closed
exact-operation parent.  Then:
\begin{enumerate}[label=(\roman*),leftmargin=2em]
\item \(H_{\mathfrak M}\) is a nontrivial connected compact topological group;
\item if the role supplies a faithful finite state-affine operation port,
Theorem~\ref{p2:thm:frc-automatic-compact-lie-closure} additionally
identifies \(H_{\mathfrak M}\) as compact Lie; the same follows
from any independently certified Lie-regularity input; and
\item exactification of the intrinsic order-swap diagram has exactly one of
the exhaustive statuses
\[
 \epsilon_R=\mathrm{ab}:
 [H_{\mathfrak M},H_{\mathfrak M}]=\{e\},
 \qquad\text{or}\qquad
 \epsilon_R=\mathrm{nc}:
 [H_{\mathfrak M},H_{\mathfrak M}]\ne\{e\}.
\]
Every role-faithful candidate for the same
\(\operatorname{Geo}_R\) has the same value of \(\epsilon_R\), although its
completed operation group and multiplication table may differ.  On the second
branch, if the constituted physical test doctrine separates transformations,
GP2 faithfulness also yields an endogenous order-sensitive test distinction
\[
 \exists r,s\in H_{\mathfrak M},\ \exists b\in\mathcal B_{\mathfrak M}:
 \qquad [r,s]^*b\ne b.
\]
Metric-curry closure is one sufficient generated grammar for this last
operational witness; it is not needed for the algebraic branch status.
\end{enumerate}
Thus the provenance is
\[
 [\mathrm{CONST}]\ \operatorname{Geo}_R
 \longrightarrow[\mathrm{REAL}]\ H_{\mathfrak M}
 \longrightarrow[\mathrm{CLASS}]\
 (\text{nontrivial},\text{ connected},\epsilon_R).
\]
The connectedness and algebraic branch statuses follow already from
the strong constituted predicate and compact recurrence.  The Lie upgrade
uses the separate faithful affine state/test port (or another valid Lie
certificate); it is not needed to define the algebraic branch and is not
a new independent branch label.  The branch status is common role content entering the later
\([\mathrm{MIN}]\) GP3 comparison, while proof witnesses and each candidate's
completed self-presentation remain unfrozen.
\end{theorem}

\begin{proof}
GP1-TB and the exact comparison-isometry action give the compact
reversible parent by
Theorem~\ref{p2:thm:gp1-bounded-isometry-closure}.
The primitive role's connected local reapplication word and absence
of a separately licensed discrete seed satisfy the generation
hypotheses of
Lemma~\ref{p2a:lem:v281-connected-generated-group}; this proves
connectedness without assuming Lie structure.  GP2 faithful retention
of the effective reapplication proves nontriviality.
If the candidate additionally has the faithful affine operation port
of Theorem~\ref{p2:thm:frc-automatic-compact-lie-closure}, its closed
group is compact Lie.  An independent valid Lie certificate may
supply the same topological upgrade, but it does not create a
Hilbert-space lift or a new physical readout.

Finally, the exact multiplication either makes every instance of the intrinsic
order-swap diagram equal or it does not; these alternatives are mutually
exclusive and exhaustive.  Role faithfulness preserves equality and
inequality of that diagram, so the status cannot change between candidates of
the same \(\operatorname{Geo}_R\).  In the second case choose a commutator
\(c=[r,s]\ne e\) only inside the proof.  If the constituted test doctrine
separates transformations, it supplies some
\(b\in\mathcal B_{\mathfrak M}\) with \(c^*b\ne b\); canonical metric profiles
do so under metric-curry closure by
Theorem~\ref{p2:thm:canonical-metric-profile-separation}.  Neither \(r,s\), nor
\(b\) is retained as a marked datum.  Hence the branch is classified after
faithful exact realization of the constituted role geometry rather than
inserted as a primitive label, while the physical
B-witness retains its stated conditional scope.
\end{proof}

\begin{definition}[Cumulative role-allocation chronology]
\label{p2:def:cumulative-role-allocation}
Let \(P_{<R}\) be the inventory frozen before the role \(R\), and let
\(\mathsf E_R\) be the geometrically primary assertion of
Definition~\ref{p2:def:geometric-role-existence}.  Define the GP1/GP2-admissible
class by
\[
 \operatorname{Adm}_R
 :=\left\{\mathfrak M:
 \begin{array}{l}
 \operatorname{Geo}_R(\mathfrak M),\quad
 \mathrm{GP1}(\mathfrak M),\\[-1pt]
 \operatorname{Pres}_{\mathfrak M}(\mathfrak M)
 =\operatorname{Gen}_R(\mathfrak M),\quad
 \mathfrak M\text{ is faithful to }P_{<R}
 \end{array}\right\}.
 \tag{\ensuremath{\mathrm{Adm}_R}}
\]
For each \(\mathfrak M\in\operatorname{Adm}_R\), put
\[
 \Sigma_R[\mathfrak M]
 :=\operatorname{Pres}_{\mathfrak M}(\mathfrak M)
 =\operatorname{Gen}_R(\mathfrak M).
\]
This is the candidate's complete exact self-presentation; there is no single
completed \(\Sigma_R\) frozen across candidates before GP3.  Applying
Theorem~\ref{p2:thm:geometric-existence-exactification} candidatewise
to every admitted compact topological parent derives one common status
\(\epsilon_R\in\{\mathrm{ab},\mathrm{nc}\}\) from the intrinsic role
diagram.  A separate continuous faithful finite state-affine operation port,
where actually generated, upgrades that individual parent to compact Lie.  The two values are exhaustive theorem branches, not
candidate choices or independent retained bits.  The objects of
\(\operatorname{Adm}_R\) are paired exact-operation models
\[
 \mathfrak M=
 \bigl(\mathfrak F_{\mathfrak M};
       X_{\mathfrak M},d_{\mathfrak M},H_{\mathfrak M},
       \alpha_{\mathfrak M},Q_{\mathfrak M}\bigr)
\]
of the same geometrically primary role \(R\).  They agree on
\(P_{<R}\), \(\operatorname{Geo}_R\), and the derived branch, while each has
its own complete exact presentation \(\Sigma_R[\mathfrak M]\).  Here
\(\mathfrak F_{\mathfrak M}\) contains the finite ontic carrier together with
its represented finite-dimensional operation envelope, so the paired model is
a concrete structured realization in the sense of
Definition~\ref{p2:def:structured-role-realization}.  The operation parent
\(H_{\mathfrak M}\) is the \emph{closed} exact image obtained by taking the
closure of the exact composition monoid.
Theorem~\ref{p2:thm:geometric-existence-exactification} identifies the role
parent with the connected component generated by the local grammar.
If the candidate carries a constituted continuous faithful finite-state
affine operation action of its \emph{entire closed parent}, that action is
unitarized on the state/test span by
Theorem~\ref{p2:thm:frc-automatic-compact-lie-closure}.  Neither GP1
finite capacity nor GP2 host faithfulness alone proves such an action for
every literal member of \(\operatorname{Adm}_R\).  This conditional Lie
upgrade is not an extra role selector or an admissibility hypothesis.  The model
also satisfies FRC-a and FRC-c, together with every relevant FRC-b
scalar or rate bound.  The tuple
\((X_{\mathfrak M},d_{\mathfrak M},H_{\mathfrak M},\alpha_{\mathfrak M})\)
is the generated parent/presentation comparison and action, and
\(Q_{\mathfrak M}\) is the declared parent-to-finite realization map.  Its
constituted comparison image satisfies GP1-TB through the finite-operational
uniform bridge of Corollary~\ref{p2:cor:gp1-tb}.

Every model is faithful to every relation and distinction in \(P_{<R}\) and
\(\operatorname{Geo}_R\), is internally exact to
\(\Sigma_R[\mathfrak M]\), and obeys the no-extra-marked-datum clause
\ref{p2:sr:no-extra-marked-datum}.  When the candidate carries the constituted faithful state-affine
operation port, its closed parent is automatically compact Lie by
Theorem~\ref{p2:thm:frc-automatic-compact-lie-closure}; otherwise
compactness alone is established and the Lie status requires an independent
certificate.  When a local Lie type
is exposed, a compact source integration that already admits such a faithful
paired FRC model is retained with no defining equality beyond the congruence
generated by the geometrically primary local source.  If the source integration is not globally FRC-b-bounded, FRC retains
only closed compact globalizations that themselves admit a faithful paired
model.  GP2 then applies its prior-signature relation-exactness rule.  In the
compact-reductive decomposition used below, a pure semisimple central kernel
whose removal preserves the same projected full abelian lattice and leaves a
compact paired model is an additional ungenerated quotient relation and is
excluded; Lemma~\ref{p2:lem:no-gratuitous-semisimple-kernel} proves that
case.  Mixed lattice--center identifications that participate in the
compactifying lattice are not removed by that lemma, and genuinely different
surviving compactification classes are not selected from one another.  Only
after this semantic filter is applied is a unique surviving unmarked
isomorphism class forced; a nonsingleton relation-exact class remains
unresolved and supplies no preferred compactifying quotient.
Proposition~\ref{p2:prop:admissible-globalization-operator} characterizes
this rule for compact-Lie globalizations below rather than defining
admissibility by the result it is used to prove.

Admissibility is role-relative and cannot be certified by a candidate's own
later automorphisms or canonical quotients.  This does not identify the point
sets, multiplication tables, or full profile algebras of two candidates.  Each
is retained by candidate-internal exactness, but none is prior data to which
every other candidate must be faithful.

For each local type whose source integration is compact, GP2 retains that
source form and permits no additional discrete quotient relation unless the
role geometry generates it.  Candidate-specific operation elements,
multiplication tables, and metric profiles belong to
\(\Sigma_R[\mathfrak M]\); they are exact inside that candidate but are not
part of \(P_{<R}\).  The
full profile family generated candidatewise by currying
\(d_{\mathfrak M}\) detects the kernel of each candidate's parent geometric
action.  It is a physical B-test family only under metric-curry closure and
does not become a prior distinction between candidates.  In this unmarked
primitive-role domain, GP3 compares the whole already-filtered paired-model
class by role-faithful comparisons induced through
\(\operatorname{Skel}_R\) as specified in
Definitions~\ref{p2:def:primitive-role-skeleton} and
\ref{p2:def:global-gp3-preorder}; carrier, represented
operations, relations, and tests remain jointly subject to role faithfulness
and candidate-internal exactness throughout that comparison.  The
primitive-domain preorder does
not claim a unique point carrier, finite presentation, or metric scale.  Only
after GP3 closure---a unique least structured class when one exists, or else
the common content of the unresolved irreducible family---does the selected
candidate's full multiplication, topology, representation theory, and
canonical tensor structure enter the chronological inventory:
\[
 \mathsf E_R\ \longrightarrow\
 \operatorname{Adm}_R\ \longrightarrow\
 \mathcal M_R:=\operatorname{Close}_{\rm GP3}(\operatorname{Adm}_R)\ \longrightarrow\
 P_R\ \longrightarrow\ \operatorname{Th}(H_R),
\]
with
\[
 P_R=P_{<R}\cup
 \operatorname{Pres}_{\mathfrak M_R}(\mathfrak M_R)
 =P_{<R}\cup\operatorname{Gen}_R(\mathfrak M_R)
\]
when a unique least \(\mathfrak M_R\) exists.  For an unresolved family,
\(P_R\) contains only \(P_{<R}\) and the generated items proved common to all
irreducible members.
In the primitive allocation proved below, \(H_R\) denotes the common endpoint
skeleton of every GP3-irreducible structured realization; the unmarked
faithful carrier presentation may remain unresolved.  The GP3 comparison is
set-sized and well posed without assuming every ambient parent is Lie.
The compact metric comparison has a countable dense set and is therefore
coded up to isometry by a real distance matrix on \(\mathbb N\); every
isometry is determined by its values on that dense set.  Its closed
operation subgroup, action and finite role-presentation data are therefore
coded by subsets and functions of sets, so the resulting candidate
isomorphism classes and preorder form a set rather than a proper class.
Replacing \(\operatorname{Geo}_R\) by the complete self-description of a
chosen \(H_R\) before the GP3 comparison reverses the chronology and is
inadmissible.  Conversely, after selection, discarding any element of
\(\operatorname{Gen}_R(\mathfrak M_R)\) would violate GP2 no-omission.
\end{definition}

\begin{definition}[Unmarked applied primitive-role domain]
\label{p2:def:unmarked-primitive-role-domain}
The master theorem applies at the first allocation of a primitive role, before
a global operation group has been selected.  Its chronological input is
\(P_{<R}\) together with the geometrically primary assertion \(\mathsf E_R\)
of Definition~\ref{p2:def:geometric-role-existence}.  The word
\emph{unmarked} introduces no further hypothesis: it means exactly that GP2
no-insertion is evaluated relative to this prior content rather than relative
to the self-description of a proposed candidate.

Candidatewise GP1--GP2 closure realizes the constituted effective, connected
local-generation data as a nontrivial connected exact parent and classifies the
intrinsic order-swap diagram by one common status
\[
 \epsilon_R=\mathrm{ab}:
 [H_{\mathfrak M},H_{\mathfrak M}]=\{e\},
 \qquad\text{or}\qquad
 \epsilon_R=\mathrm{nc}:
 [H_{\mathfrak M},H_{\mathfrak M}]\ne\{e\}.
\]
Neither status is separately postulated and no witness pair is designated.  A
physical test witness on the second branch additionally requires a constituted
transformation-separating test doctrine, as stated in
Theorem~\ref{p2:thm:geometric-existence-exactification}.

A named operation element, completed multiplication table, candidate point
set, full candidate profile algebra, dimension, rank, quotient relation,
representation, or scale may occur inside
\(\Sigma_R[\mathfrak M]\) and is then retained there by GP2 no-omission.  It
cannot be inserted into \(P_{<R}\), used as its own selector, or imposed on
every other candidate merely because it belongs to \(\mathfrak M\).  This is
GP2 no-insertion, not a separate local-role assumption.  If
\(\operatorname{Geo}_R\) itself genuinely requires richer directions,
brackets, relations, or tests, every role-faithful candidate must preserve
them; a proposed \(U(1)\) or \(SU(2)\) endpoint that fails that condition is
outside \(\operatorname{Adm}_R\), and the primitive endpoint theorem abstains.
\end{definition}

\begin{definition}[Primitive-role skeleton of a structured realization]
\label{p2:def:primitive-role-skeleton}
For \(\mathfrak M\in\operatorname{Adm}_R\) in the unmarked applied
primitive-role domain, let \(\operatorname{Skel}_R(\mathfrak M)\) be the
bookkeeping image retaining the connected operation-role parent and local germ,
the relation-exact source/globalization data, the intrinsic diagrams required
by \(\operatorname{Geo}_R\), and the GP2-derived exact branch status.  It
forgets only presentation coordinates irrelevant to comparison---for example
a carrier basis, an unmarked equivalent faithful carrier presentation, a
metric scale, named proof-witness elements, and profile coordinatizations.

This forgetful image is not a physical truncation of \(\mathfrak M\): the full
candidate presentation \(\Sigma_R[\mathfrak M]\) remains internally exact.
Every skeleton comparison used by GP3 must have complete structured endpoints
whose carriers, represented operations, relations, and tests still satisfy
GP1, candidate-internal GP2 exactness, and role faithfulness.  After
Definition~\ref{p2:def:global-gp3-preorder}, write
\(S_1\preceq_R^{\rm sk}S_2\) for the induced skeleton order: there exist
structured lifts \(\mathfrak M_i\) with
\(\operatorname{Skel}_R(\mathfrak M_i)\cong S_i\) and
\(\mathfrak M_1\preceq_R\mathfrak M_2\).  Thus the primitive preorder may
factor through skeleton data without minimizing carrier and operations
independently, deleting source-candidate operations, or promoting forgotten
presentation parameters into selection criteria.
\end{definition}

\begin{definition}[Primitive-role specialization of the global GP3 preorder]
\label{p2:def:global-gp3-preorder}
For \(\mathfrak M_1,\mathfrak M_2\in\operatorname{Adm}_R\) in the unmarked
domain, let \(\preceq_R\) be the primitive-role specialization of the universal
role-faithful candidate comparison of GP3.  Write
\(\mathfrak M_1\preceq_R\mathfrak M_2\) when, before either candidate's
self-presentation is frozen, a finite chain of the following elementary
comparisons takes \(\operatorname{Skel}_R(\mathfrak M_2)\) to
\(\operatorname{Skel}_R(\mathfrak M_1)\), and every step has complete
GP1-admissible, candidate-internally GP2-exact-and-faithful structured
endpoints:
\begin{enumerate}[label=(\roman*),leftmargin=2em]
\item passage to an effective continuous connected homomorphic image;
\item construction of the complete exact candidate generated by a closed
connected local subrole;
\item deletion of a direct factor unused by \(\operatorname{Geo}_R\);
\end{enumerate}
provided every target still realizes \(\operatorname{Geo}_R\), preserves
\(P_{<R}\), and has the same GP2-derived exact branch status.  The target is
exact to its own \(\operatorname{Gen}_R\); it is not required to retain every
operation or test generated only by the source candidate.  If a source
direction, bracket, relation, or test is required by \(\operatorname{Geo}_R\),
however, the target must preserve it and the comparison is otherwise
inadmissible.  Before the endpoint of any image, subrole, or factor-deletion
step is admitted to the preorder, its retained local germ is passed through
\[
        E_{\rm adm}(\text{retained local germ}),
\]
the FRC--GP2 admissible globalization of Proposition~
\ref{p2:prop:admissible-globalization-operator}.  A compact admissible source
is retained without an added kernel relation.  A noncompact source is first
restricted by the closed compact-parent requirement and then by GP2 relation
exactness, including removal of every gratuitous pure semisimple central
kernel as proved in
Lemma~\ref{p2:lem:no-gratuitous-semisimple-kernel}.  If the remaining compact
relation-exact globalizations form one unmarked isomorphism class, that class
is forced, and if several classes survive, the endpoint is their unresolved
family rather than a preferred quotient.  On a set-valued branch, \(\preceq_R\) is read as
holding when the displayed chain exists for some member, but no member is
selected.  In the two applications below the operator is single-valued:
Lemma~\ref{p2:lem:abelian-compact-globalization} gives \(U(1)\) on the abelian
line, and the compact source gives \(SU(2)\) on the noncommutative branch.

The preorder is the reflexive transitive closure of these complete-candidate
comparisons,
modulo natural isomorphism of the corresponding relation-exact structured
role models.  GP3 retains a unique least structured isomorphism class when one
exists.  If there is no unique least class, distinct irreducible classes remain
unresolved whether they are incomparable or preorder-equivalent but
non-isomorphic.  The preorder contains no dimension, rank, root datum, named
group, carrier-dimension functional, or presentation scale.  The three
elementary comparisons displayed here are the primitive-role specialization of
the universal role-faithful GP3 preorder; richer role geometries may prove
other concrete comparison lemmas, but they remain instances of the same GP3
rather than new minimization principles.  A unique least role skeleton may
nevertheless be common to an unresolved family of faithful structured lifts;
in that case only the common skeleton is frozen, not an unmarked carrier
representation or scale.
\end{definition}

\begin{lemma}[Primitive-role skeleton bookkeeping]
\label{p2:lem:skeleton-factorization}
For the geometrically primary role condition \(\operatorname{Geo}_R\):
\begin{enumerate}[label=(\alph*),leftmargin=2em]
\item a skeleton comparison with no complete GP1-admissible,
candidate-internally GP2-exact-and-role-faithful structured endpoints is not a
GP3 comparison;
\item distinct faithful structured lifts of the same skeleton are not selected
from one another unless the role geometry or frozen inventory supplies an
additional structural comparison distinguishing them; absent such a comparison they remain an
unresolved presentation family over the common skeleton.
\end{enumerate}
\end{lemma}

\begin{proof}
Clause~(a) is part of Definition~\ref{p2:def:global-gp3-preorder}.  Clause~(b)
follows because the primitive preorder contains no rule comparing lifts that
differ only in unmarked equivalent presentation data; adding such a comparison
would itself require a new generated role relation.
\end{proof}

\begin{remark}[Endpoint comparison rather than source-operation transport]
\label{p2:rem:primitive-endpoint-reachability-scope}
The induced order \(\preceq_R^{\rm sk}\) does not assert that every operation
of a source candidate is transported into an endpoint candidate.  It asserts
that both are complete exact realizations of the same geometrically primary
role and that a licensed role-faithful comparison places the endpoint below
the source.  The required quantifier is candidatewise: from every admissible
candidate there is such a comparison to the least endpoint.  That statement is
proved directly in
Theorem~\ref{p2:thm:primitive-endpoint-lift-reachability} after the
compact-globalization lemmas.  Distinct faithful presentations over the common
endpoint skeleton can therefore remain unresolved even though every
GP3-irreducible presentation has that skeleton.
\end{remark}

\begin{remark}[Illustration of the exactification step]
\label{p2:rem:exactification-illustration}
Outside the definition, and using names the preorder itself does not contain:
an image such as \(SU(2)\to SO(3)\) is used only to expose its local
\(\mathfrak{su}(2)\) germ and exactifies back to \(SU(2)\); the added kernel
relation is not a strict GP3 comparison.
\end{remark}

\begin{remark}[Three controls on the candidate/frozen distinction]
\label{p2:rem:three-candidate-chronology-controls}
The following cases test three different clauses of the chronology.
\begin{enumerate}[label=(\roman*),leftmargin=2em]
\item A flat \(T^2\) candidate retains its full two-torus operation law by
candidate-internal GP2 no-omission.  A complete \(U(1)\) candidate nevertheless
lies below it for the primitive scalar role.  If a second toral direction is
intrinsic to \(\operatorname{Geo}_R\), that comparison is forbidden.
\item A compact \(SU(3)\) candidate likewise retains every generated direction
and bracket.  A complete root-\(SU(2)\) candidate lies below it for the
primitive order-sensitive role, but not when an \(SU(3)\)-specific direction,
bracket, relation, or test belongs to \(\operatorname{Geo}_R\).
\item \(SO(3)\) is not removed as surplus structure by GP3.  It is excluded at
the relation-exact GP2 stage because it imposes the additional global equality
\(-1=1\) on an already compact admissible source integration.  Thus the
relation-exact type-\(A_1\) candidate is \(SU(2)\).
\end{enumerate}
The first two clauses compare different complete exact candidates; the third
excludes an ungenerated quotient relation.  None deletes an operation from a
selected or frozen geometry.
\end{remark}

\begin{definition}[Paper-II primitive-role scope declaration]
\label{p2:def:primitive-role-unmarked-applicability}
The primitive roles allocated in Paper~II-A are the roles of
Definition~\ref{p2:def:unmarked-primitive-role-domain}.  Each is introduced by
the strong geometrically primary assertion \(\mathsf E_R\), not by a completed
global candidate.  Here \emph{primitive} means that the first-allocation role
content is generated by the single effective identity-rooted local
reapplication constructor and its intrinsic comparison and order-swap diagrams
of Definition~\ref{p2:def:geometric-role-existence}; there is no independently
constituted second role direction, bracket schema, quotient relation, or test
doctrine beyond consequences of that role geometry.  This is a domain
definition, not an additional selector.  If the actual geometry contains such
richer content, GP2 no-omission retains it and the role leaves this primitive
endpoint domain.

The word ``single'' here refers to one undecomposed role sort, not to a
one-parameter family, a one-dimensional tangent space, or Cartan rank one.

The allocation theorem has an explicit feasibility branch.  If
\(\operatorname{Adm}_R=\varnothing\), the role is FRC-obstructed and no parent
skeleton is allocated.  The endpoint conclusions concern the complementary
non-obstructed case \(\operatorname{Adm}_R\ne\varnothing\).  This nonemptiness
condition is feasibility of exact finite/bounded realization, not a structural
selector among surviving candidates.

After GP1 and GP2 have been applied, the name \emph{scalar
branch} denotes \(\epsilon_R=\mathrm{ab}\), while
\emph{directed-transfer branch} denotes the derived order-sensitive status
\(\epsilon_R=\mathrm{nc}\).  Nontriviality, connectedness, and branch type are
therefore theorem statuses, not primitive marks.  No operation pair is named.

Before GP3, no candidate's global parent, complete multiplication table,
profile algebra, dimension, rank, quotient relation, representation, or scale
belongs to \(P_{<R}\).  Such data may and must occur internally in an exact
candidate, but cannot serve as prior selection marks.  This is the content of
GP2 no-insertion, not an independent ``unmarked local-role'' assumption.  It
is also not a theorem that GP1--GP3 erase richer geometry: if
\(\operatorname{Geo}_R\) requires additional intrinsic structure, every
candidate comparison must preserve it and the endpoint theorem applies only
after a role-faithful endpoint candidate is proved.  The symbols
\((X_{\mathfrak M},H_{\mathfrak M})\) range over candidates and are not
earlier physical constants.
\end{definition}

\section{Principle-separated compact parent selection}
\label{paper2a:sec:organized-selection}

\subsection{Principle-separated proof chain}
\label{p2:subsec:principle-separated-chain}

This subsection contains the unique load-bearing logical order.  For each
candidate model \(\mathfrak M\in\operatorname{Adm}_R\), let
\((X_{\mathfrak M},d_{\mathfrak M})\) be its comparison carrier, let
\(S_0(X_{\mathfrak M})\) be the identity-containing forward monoid generated
by its exact comparison-preserving self-maps, and let
\(\mathcal B(X_{\mathfrak M})\) be its generated tests and readouts.  In the
candidatewise theorems below we suppress the subscript \(\mathfrak M\).
The initial grammar contains identity and composition, but no inverse symbol.
At this point an element of \(S_0(X)\) is only an isometric embedding of
\(X\) into itself; bijectivity, reversal, topology, smoothness, dimension,
connectedness, and a sphere are not assumed.

In the FRC typing, \(X\) is the generated comparison
\emph{presentation image} of a finite realized role.  It may have continuum
cardinality only in that parent/presentation sense.  The ontic carrier is the
finite \(\mathcal H_{R,s}(\mathcal D)\) and its finite operator algebra.
Accordingly, compact completions of \(X\) and the groups acting on them are
parent operation geometries unless a finite realization is displayed
separately.

Until Theorem~\ref{p2:thm:master-compact-role-allocation}, \(X\) denotes one
generic candidate in \(\operatorname{Adm}_R\).  The completion,
operation closure, and metric-profile theorems below are candidatewise.  They
do not place the complete operation set or profile algebra of this placeholder
into \(P_{<R}\) or the prior role assertion.  Only the GP2-derived exact branch
status of the intrinsic role diagram, never a witness pair or candidate
multiplication table, is held fixed across candidates;
if a unique least structured class exists its full constituted data are
frozen after GP3; otherwise only the prior/generated data common to the
unresolved irreducible family are frozen.

\paragraph{FRC-b comparison and rate specializations}

\begin{definition}[GP1-TB: the comparison specialization of FRC-b]
\label{p2:def:intrinsic-operational-boundedness}
For the generated metric geometry \((X,d)\), write
\[
 D_X:=\operatorname{diam}_dX
     =\sup_{x,y\in X}d(x,y).
\]
GP1-TB is exactly the condition
\[
 N_X(\varepsilon)<\infty\qquad(\varepsilon>0),
\]
which implies \(D_X<\infty\) but not conversely.  For geometric operations use
the uniform comparison
\[
 D_\infty(a,b):=\sup_{x\in X}d(ax,bx).
\]
If an elapsed scalar cover \(T\subseteq\mathbb R\) and the full family
\(\mathcal O_{\rm prop}\) of timed propagation operations have also been
constituted, each \(\Phi\in\mathcal O_{\rm prop}\) has the generated invariant
rate
\[
 v(\Phi)
 :=
 \sup_{\substack{x\in X\\s,t\in T,\ s\ne t}}
 \frac{d(\Phi_t x,\Phi_s x)}{|t-s|}
 \in[0,\infty].
\]
Here \(\sup\varnothing:=0\) for a degenerate elapsed cover.
The operation comparison includes this rate evaluation.  Thus GP1
(clause (b)) requires
every constituted propagation operation to have finite \(v(\Phi)\) and
applies total boundedness to the scalar rate image
\[
 \mathcal V_{\rm prop}
 :=\{v(\Phi):\Phi\in\mathcal O_{\rm prop}\}
 \subset[0,\infty)
\]
with its ordinary invariant distance.  In an FRC-b-admissible timed role the
extended value \(+\infty\) is therefore excluded rather than treated as a
point of that metric geometry.  More generally, for any constituted
scalar readout \(q\), boundedness means that \(q(X)\) is a bounded subset of
its invariant value geometry.  These conditions use the already generated
comparison.  Replacing an
unbounded metric by \(\min\{d,1\}\), bounding coordinates through a function
such as \(\tanh\), or observing only that every individual value is a real
number does not establish GP1-TB.  Total boundedness also supplies no positive
minimum distance: a compact interval is the elementary counterexample.
\end{definition}

\begin{theorem}[GP1-TB gives a canonical compact completion]
\label{p2:thm:gp1-bounded-compactness}
If \((X,d)\) satisfies GP1-TB, its metric
completion \((\widehat X,\widehat d)\) is compact.  The completion is unique up
to a unique isometry fixing \(X\), so admitting it introduces no branch,
quotient, lattice, or coordinate choice.
\end{theorem}

\begin{proof}
A metric space is totally bounded precisely when its completion is totally
bounded.  It is also complete, being a metric completion.  A complete totally
bounded metric space is compact.  Uniqueness of
metric completion is canonical and functorial for isometric embeddings.
\end{proof}

\begin{corollary}[Resolution completeness follows from GP1-TB]
\label{p2:cor:gp1-resolution-completeness-derived}
Under GP1-TB, \(X\) and \(\widehat X\) have finite
\(\varepsilon\)-nets for every \(\varepsilon>0\).  Resolution completeness is
therefore built into the total-boundedness statement, while compactness of the
canonical completion is derived.  Neither statement makes the comparison presentation finite
or discrete; FRC-a separately makes the ontic carrier finite.
\end{corollary}

\begin{proof}
The first assertion is the definition of total boundedness.  A finite
\(\varepsilon/2\)-net in \(X\) is an \(\varepsilon\)-net in its completion.
\end{proof}

\begin{remark}[Completion boundary]
\label{p2:rem:gp1-completion-boundary}
The original totally bounded comparison presentation \(X\) need not be
complete or compact.  Only its canonical completion \(\widehat X\) is compact.
Every later compact-parent argument is therefore stated on \(\widehat X\), or
on a separately derived compact operation parent \(H\).
\end{remark}

\begin{theorem}[FRC-b operational propagation ceiling]
\label{p2:thm:gp1-operational-speed-ceiling}
Let \(X\), \(T\), and the nonempty full family
\(\mathcal O_{\rm prop}\) be a constituted timed
propagation role as in
Definition~\ref{p2:def:intrinsic-operational-boundedness}.  Then
\[
 c_{\rm op}:=\sup\mathcal V_{\rm prop}<\infty,
\]
and every \(\Phi\in\mathcal O_{\rm prop}\) obeys
\begin{equation}
\label{p2:eq:gp1-finite-propagation-cone}
 d(\Phi_t x,\Phi_s x)\le c_{\rm op}|t-s|
 \qquad(x\in X,\ s,t\in T).
\end{equation}
The same estimate is stable under pointwise limits of histories with the
common bound \(c_{\rm op}\), and under compatible concatenation of histories.
Here a compatible concatenation is a history \(\Psi\) obtained from a finite
or locally finite partition, with junctions in \(T\), of an elapsed interval
by histories on the same carrier \(X\), with the same metric \(d\), elapsed
cover \(T\), and common bound \(c_{\rm op}\), whose pointwise values agree at
every junction; the partition is finite on each compact elapsed subinterval.
Call a timed operation \emph{stationary} when \(d(\Phi_tx,\Phi_sx)=0\) for
every \(x\in X\) and all \(s,t\in T\).  If the role contains a nonstationary
operation, then \(c_{\rm op}>0\).

The ceiling is universal only relative to the constituted full timed
propagation role and need not be attained.  FRC-b requires its finiteness but
not its unit-dependent numerical value: rescaling the elapsed-time unit
rescales the number assigned to \(c_{\rm op}\).  This theorem alone proves no
field or light-channel identity.  Paper~II-A does, however, prove the general
inheritance result
Theorem~\ref{p2a:thm:undeformed-host-principal-inheritance}: after a later
channel's support and undeformed principal ancestry are established, its
front speed is \(c_{\rm op}\).  A measurement identifying an observed signal
with the later vacuum photon channel calibrates that common rate; it does not
create another structural speed.
\end{theorem}

\begin{proof}
By FRC-b the subset
\(\mathcal V_{\rm prop}\subset\mathbb R_{\ge0}\), with its ordinary metric,
is totally bounded.  Every totally bounded subset of \(\mathbb R\) is bounded,
so \(c_{\rm op}<\infty\).  The defining supremum of \(v(\Phi)\) gives
\[
 d(\Phi_t x,\Phi_s x)
 \le v(\Phi)|t-s|
 \le c_{\rm op}|t-s|,
\]
which proves \eqref{p2:eq:gp1-finite-propagation-cone}.  If histories
\(\Phi^{(n)}\) with the same bound converge pointwise, joint continuity of \(d\)
passes the inequality to the limit.  For a compatible concatenation and
\(s<t\), let
\[
 s=u_0<u_1<\cdots<u_N=t
\]
be \(s,t\) together with the finitely many intervening junctions on the
compact elapsed interval \([s,t]\).
Pointwise junction matching and the triangle inequality give
\[
 d(\Psi_t x,\Psi_s x)
 \le \sum_{i=0}^{N-1}d(\Psi_{u_{i+1}}x,\Psi_{u_i}x)
 \le c_{\rm op}\sum_{i=0}^{N-1}(u_{i+1}-u_i)
 =c_{\rm op}(t-s).
\]
Thus finite and locally finite compatible concatenations inherit the same
bound.  A nonstationary operation has
\(d(\Phi_t x,\Phi_s x)>0\) for some \(x,s,t\), hence \(v(\Phi)>0\) and
\(c_{\rm op}>0\).

A change
\(t\mapsto at\), \(a>0\), sends
\(c_{\rm op}\mapsto c_{\rm op}/a\), while preserving
finiteness and the invariant cone
\(d(\Phi_tx,\Phi_sx)\le c_{\rm op}|t-s|\).  Therefore FRC-b requires a finite ceiling for
the constituted full role
but no unit-dependent numerical value.
\end{proof}

\begin{proposition}[Elapsed concatenation and synchronized iterates]
\label{p2a:prop:clocked-composition-rate-scope}
Let $X$ be the generated metric carrier, and let
$\Phi_t$, $t\in\mathbb R$, be a flow with
$0<v(\Phi)<\infty$, using the full elapsed cover in
Definition~\ref{p2:def:intrinsic-operational-boundedness}.  For each positive
integer $n$, the synchronized iterate
$\Psi^{(n)}_t:=(\Phi_t)^n=\Phi_{nt}$ has
\[
 v(\Psi^{(n)})=n\,v(\Phi).
\]
Consequently an FRC-b-admissible constituted timed role cannot contain this
entire family with unchanged clock accounting and the same rate comparison.
This restriction does not apply to the compatible elapsed concatenations of
Theorem~\ref{p2:thm:gp1-operational-speed-ceiling}, which add the elapsed
durations of their pieces.
\end{proposition}

\begin{proof}
Put $u=nt$ and $w=ns$ in the defining supremum.  Since the elapsed cover is
$\mathbb R$, the substitution ranges over all $u\ne w$ and gives
\[
 \sup_{x,s\ne t}\frac{d(\Phi_{nt}x,\Phi_{ns}x)}{|t-s|}
 =n\sup_{x,u\ne w}\frac{d(\Phi_ux,\Phi_wx)}{|u-w|}.
\]
The right side is unbounded as $n$ varies.  FRC-b requires the full
constituted rate image, rather than each member separately, to be bounded.
By contrast, compatible concatenation uses the sum of the piece durations,
exactly as in the preceding proof.  Thus constructor closure must retain its
actual elapsed typing: it cannot silently license synchronized acceleration
by applying an untimed composition rule to clocked histories.
\end{proof}

\begin{corollary}[Compact rate closure and completion of timed histories]
\label{p2:cor:gp1-timed-history-completion}
In the domain of
Theorem~\ref{p2:thm:gp1-operational-speed-ceiling},
\[
 \overline{\mathcal V_{\rm prop}}\subset\mathbb R_{\ge0}
 \quad\text{is compact},\qquad
 c_{\rm op}=\max\overline{\mathcal V_{\rm prop}}.
\]
For every \(\Phi\in\mathcal O_{\rm prop}\) and \(x\in X\), the history
\[
        T\longrightarrow X,\qquad t\longmapsto\Phi_t x
\]
extends uniquely to a \(c_{\rm op}\)-Lipschitz map from the metric completion
\(\widehat T\) of \(T\) into \(\widehat X\).  Moreover,
\[
 c_{\rm op}=0
 \quad\Longleftrightarrow\quad
 \text{every history in the constituted role is stationary on \(T\)}.
\]
The maximum in the first display belongs to the closure of the rate image;
an operation attaining it need not exist.
\end{corollary}

\begin{proof}
Total boundedness of \(\mathcal V_{\rm prop}\subset\mathbb R\) makes its
closure compact.  The continuous identity function attains its maximum on
that closure, and the maximum equals
\(\sup\mathcal V_{\rm prop}=c_{\rm op}\).
Theorem~\ref{p2:thm:gp1-operational-speed-ceiling} makes each displayed
history \(c_{\rm op}\)-Lipschitz.  A Lipschitz map sends Cauchy sequences to
Cauchy sequences, so the universal property of metric completion gives the
unique extension \(\widehat T\to\widehat X\), with the same bound.
Finally \(c_{\rm op}=0\) makes every pairwise displacement vanish.  The
converse follows because a stationary history has rate zero, including under
the convention \(\sup\varnothing=0\).
\end{proof}

\begin{remark}[Why endpoint compactness is not the speed theorem]
\label{p2:rem:endpoint-compactness-not-speed}
The load-bearing FRC-b object here is the constituted \emph{timed} operation
comparison, including its invariant scalar rate image.  Compactness of the
endpoint operation group alone would not suffice.  For example, on the unit
circle the histories
\[
\Phi_t^{(n)}(z)=e^{int}z
\]
all take values in the compact group \(U(1)\), while
\(v(\Phi^{(n)})=n\) is unbounded.  FRC-b excludes this as one unbounded full
timed-operation family precisely because
\(\{v(\Phi^{(n)})\}\) is not totally bounded.  Thus the theorem neither
confuses finite spatial diameter with finite speed nor inserts
\(c_{\rm op}\): it
applies the same total-boundedness principle to the generated rate readout of
operations.
\end{remark}

\begin{theorem}[Forward exact isometries derive a compact reversible closure]
\label{p2:thm:gp1-bounded-isometry-closure}
Arzel\`a--Ascoli makes the isometry group of a compact metric space compact in
the uniform topology, with composition and inversion continuous; this
compactness is used twice below.  Every \(s\in S_0(X)\) extends uniquely to an
isometric self-map
\(\widehat s\) of \(\widehat X\), and every such extension is surjective.
Let \(S\) be the forward submonoid of
\(\operatorname{Isom}(\widehat X,\widehat d)\) formed by these extensions.
Then
\[
 G:=\overline S^{\,D_\infty}
 \leq\operatorname{Isom}(\widehat X,\widehat d)
\]
is a compact topological group.  In particular \(g^{-1}\in G\) for every
\(g\in G\): reversibility is derived in the compact closure and was not
assumed in the forward syntax.  The uniform metric is bi-invariant on \(G\).
\end{theorem}

\begin{proof}
An isometric embedding extends uniquely to the completion, and the extension
is again an isometry, since \(d\) is jointly continuous and the two sides
agree on a dense subspace.  If that extension
\(f:\widehat X\to\widehat X\) were not onto, compactness would make
\(f(\widehat X)\) closed.  Choose \(y\notin f(\widehat X)\) and put
\(\delta=d(y,f(\widehat X))>0\).  For \(n>m\),
\[
 d(f^my,f^ny)=d(y,f^{n-m}y)\ge\delta,
\]
so \(\{f^ny:n\ge0\}\) would be an infinite \(\delta\)-separated subset of a
compact metric space, a contradiction.  Thus every extension is a bijective
isometry, and \(S\) is a submonoid of the compact group
\(\operatorname{Isom}(\widehat X)\).

Distinct maps differ at some point, so \(D_\infty\) is nondegenerate.  Left
invariance,
\(D_\infty(ca,cb)=D_\infty(a,b)\), needs only that \(c\) is an isometry;
right invariance,
\(D_\infty(ac,bc)=D_\infty(a,b)\), additionally needs \(c\) surjective, so
that the supremum reindexes over all of \(\widehat X\).  Hence
\[
 D_\infty(ac,bc)=D_\infty(a,b)
 =D_\infty(ca,cb).
\]
This bi-invariance uses only the extension and surjectivity established above,
so it is available for the next paragraph.

Its closure \(G\) is a compact semigroup; cancellativity is not needed, and
the Numakura route via cancellation is not the one taken here.  For completeness, fix
\(g\in G\).  Compactness of the closure of \(\{g^n:n\ge0\}\) gives arbitrarily
large \(n>m\) with \(g^n\) and \(g^m\) arbitrarily close.  Bi-invariance in the
ambient isometry group then makes \(g^{n-m}\) arbitrarily close to \(e\).
Hence a subnet \(g^{k_i}\to e\), where the \(k_i\ge2\) may be taken
arbitrarily large, and \(g^{k_i-1}\to g^{-1}\) with
\(g^{k_i-1}\in G\).  Closedness gives \(g^{-1}\in G\), so \(G\) is a
group.
No reversibility,
homogeneity, manifold structure, or Lie theorem was assumed.
\end{proof}

\begin{theorem}[Optional \((\mathrm{FD})+(\mathrm{LR})\) regularity package]
\label{p2:thm:fd-lr-regularity-summary}
Let \(G\) be the compact parent-operation closure obtained from GP1-TB.
If, independently, \((\mathrm{FD})\) and \((\mathrm{LR})\) hold, then
\(G\) is a finite-dimensional compact Lie group.  Its uniform comparison
has a finite Ad-invariant first-order norm \(F\), is locally bi-Lipschitz to
the associated bi-invariant Finsler distance on \(G^0\), and the John
ellipsoid of the \(F\)-unit ball gives a canonical Ad-invariant quadratic
comparison.  This conditional package is not used to derive the primitive
endpoint allocation or the FRC finite-operation Lie closure.
\end{theorem}

\begin{proof}
The definitions, packing criterion, Lie-regularity proof, Finsler-rate
construction, and John-ellipsoid argument are collected without change in
Paper~II-F, Section~\ref{p2:app:fd-lr-technical-chain}; see in particular
Theorems~\ref{p2:thm:finite-mode-lie-regularity},
\ref{p2:thm:gp1-propagation-derived}, and
\ref{p2:thm:operational-john-metric}.
\end{proof}

The finite-representation proof metric, standard compact-geometric
consequences, and the compact-resolvent implication for bounded elliptic
spectral ranges are supplementary consequences, not steps in the primitive
allocation.  Paper~II-F, Section~\ref{p2:app:compact-lie-consequences} states their exact
hypotheses and proves respectively the curvature sign, completeness/Haar
package, and finite-dimensionality of each bounded elliptic spectral subspace;
see Theorems~\ref{p2:thm:finite-mode-trace-metric},
\ref{p2:thm:gp1-compact-geometric-consequences}, and
\ref{p2:thm:gp1-bounded-spectrum-finite-carrier}.

\begin{corollary}[FRC does not select a numerical bound]
\label{p2:cor:gp1-boundedness-no-scale}
FRC is invariant under a common positive rescaling of every constituted
dimensionful comparison.  Its FRC-b clause therefore asserts boundedness but
does not select the value of a diameter, curvature radius, momentum ceiling,
energy ceiling, operation norm or rate ceiling, or any other unit-dependent
positive number.  Separately registered \(\mathsf P_{\rm gbd}(q,N)\)
can constrain a realized generated presentation to a specified finite intrinsic
sublevel after the constructor grammar and its length have been constituted.
This registered, indexed boundary is not supplied by GP1 and does not by
itself derive a minimal or physically selected depth \(N\).  Thus GP1
supplies finite/bounded admissibility of carrier and operation structure,
not a numerical or discrete minimization rule.
\end{corollary}

\begin{proof}
If \(d\) is an admissible dimensionful comparison, then \(c d\) is again
bounded and totally bounded for every \(c>0\), with
\(N_{(X,cd)}(\varepsilon)=N_{(X,d)}(\varepsilon/c)\).  Scalar, operation, and
rate bounds rescale in the corresponding units, so FRC distinguishes finite
or bounded realization from an unbounded one but no positive coordinate
value.  Definition~\ref{p2:def:generated-boundary-hypothesis} instead
registers the \emph{indexed} predicate \(\mathsf P_{\rm gbd}(q,N)\):
realized generated support for slot \(q\) lies in its specified sublevel
\(\Gamma_q^{\le N}\).  When such a predicate is actually registered for
some finite \(N\), it entails the weaker existential statement
\[
 \exists N_R(s,\mathcal D)<\infty:
 \operatorname{scope}(R,s,\mathcal D)\subseteq\Gamma_R^{\le N_R};
\]
in the matching typed role.  The converse is not an inference to any
\emph{specified} \(\mathsf P_{\rm gbd}(q,N_0)\), nor does GP1 supply the
registration.  Nested sublevels make a valid upper bound nonunique in
general; the indexed hypothesis itself contains no leastness or comparison
rule between different depths.  GP2 may exactify a separately constituted
bound and GP3 may compare bounded structured realizations only through a
proved role-faithful structural comparison; neither step is contained in
GP1 itself.
\end{proof}

\begin{remark}[Optional comparison regularity is genuinely additional]
\label{p2:rem:gp1-regularity-countermodels-pointer}
GP1-TB, nonisolation of the identity, and canonical separating metric profiles
do not by themselves imply either optional parent-comparison regularity
condition (FD) or (LR).  Explicit compact countermodels separating the two
conditions are given in Paper~II-F, Section~\ref{p2:app:selector-limitations}.  These
countermodels concern parent/presentation comparisons and do not contradict the
automatic compact Lie closure of faithful exact operations on a finite carrier.
\end{remark}



\subparagraph{Scope guard for GP2 alone.}
GP2 fixes prior-geometric origin, exact generated content, host/sort
faithfulness, and constructor-relative closure.  It does not by itself insert
an external frame, orientation, topology, numerical normalization, or a
restricted doctrine not already generated by the host signature.  The
standalone scope diagnostics and local counterexamples, none of which are used
as premises later in the constitutional chain, are collected in
Paper~II-F, Section~\ref{p2:app:gp2-scope-diagnostics}.

\paragraph{Canonical effectiveness and irreducible/least GP3 comparison}

\begin{remark}[Host faithfulness precedes GP3 comparison]
\label{p2:rem:gp3-faithful-reduction-alone}
For a structured realization with supplied action
\[
 \rho:\mathcal A(X)\longrightarrow\operatorname{Aut}(\mathcal B(X)),
\]
GP2 host faithfulness requires the represented operation action to have the
kernel allowed by the frozen inventory and role geometry, the retained tests
to separate every state or operation distinction assigned to the test doctrine,
and the complete
carrier--operation--constructor structure to realize all prior laws jointly.
GP3 then asks only which of these already exact-and-faithful models are
irreducible and whether one is least.  It never minimizes carrier and
operations separately and never inserts a numerical complexity functional.

Definition~\ref{p2:def:global-gp3-preorder} gives the primitive-role
specialization of the universal role-faithful whole-candidate comparison.
Richer domains may require new comparison \emph{lemmas} proving that an image,
local-subrole candidate, quotient, or factor deletion realizes their richer
geometry, but do not require a new minimization principle.  If no least structured class is
proved, the irreducible family remains unresolved.  Once retained, the
resulting structured presentation is frozen.
\end{remark}

\begin{definition}[Metric-partial-evaluation closure of the B grammar]
\label{p2:def:metric-partial-evaluation-closure}
For a constituted binary comparison \(d:X\times X\to\mathbb R_{\ge0}\), say
that the generated B-test grammar is \emph{metric-curry closed} when, for
every constituted reference value \(z\in X\), it contains the partial
evaluation
\[
        b_z=d(z,\cdot).
\]
On the canonical completion the same clause uses
\(\widehat b_z=\widehat d(z,\cdot)\).  This closure introduces no preferred
reference point or new comparison law.  Under G0 it is a test-occurrence
certificate for an apparatus or comparison channel that realizes the full
profile doctrine; it is an operational grammar
clause: GP2 guarantees that an admitted test must have prior-geometric
origin; it does not conversely guarantee that every canonical geometric
function is physically available.  The mathematical profile theorem below
does not require this clause.  Its conclusion becomes a statement about the
physical B-test doctrine only when that G0 certificate is supplied.  G0 does
not require every possible apparatus to realize every mathematical profile.
\end{definition}

\begin{theorem}[The comparison metric canonically generates separating mathematical profiles]
\label{p2:thm:canonical-metric-profile-separation}
Let \((X,d)\) satisfy GP1-TB and let
\((\widehat X,\widehat d)\) be its canonical compact parent completion.  The binary
comparison extends uniquely to \(\widehat d\), and currying that already
constituted comparison canonically generates, without choosing a preferred
point, the full distance-profile family
\[
        b_z(x):=\widehat d(z,x),\qquad z,x\in\widehat X.
\]
Then
\[
 \lVert b_x-b_y\rVert_\infty=\widehat d(x,y).
\]
Hence the profile family is a functorial mathematical rewriting of the
comparison metric, requiring no further geometric choice, and separates
states isometrically.  Its physical test availability is conditional on
Definition~\ref{p2:def:metric-partial-evaluation-closure}.  It also separates
distinct isometries of the completion: if
\(R\ne S\), choose any \(x\) with \(Rx\ne Sx\) and put
\(z=Rx\); then
\[
        (R^*b_z)(x)=0,
        \qquad
        (S^*b_z)(x)=\widehat d(Rx,Sx)>0.
\]
Thus, on each geometric operation candidate used in this paper, the full
mathematical metric-profile family is faithful.  Under metric-curry closure,
the same family is a faithful B-test family.  More generally, for an abstract operation
presentation
\(P\to\operatorname{Isom}(\widehat X,\widehat d)\), its canonical effective parent
is
\[
 P_{\rm eff}=P/\ker(P\curvearrowright\widehat X),
\]
and the displayed metric profiles detect exactly that kernel: equal images in
\(\operatorname{Isom}(\widehat X)\) give equal profile pullbacks, and distinct
images are separated by the isometry-separation clause above.  This
effectiveness theorem precedes GP3.  It ensures that the subsequent global comparison is among faithful candidates;
it does not itself establish a least class.
\end{theorem}

\begin{proof}
Uniqueness of metric completion gives the unique extension \(\widehat d\).
The triangle inequality gives
\(|\widehat d(x,z)-\widehat d(y,z)|\le\widehat d(x,y)\) for all \(z\).
Taking \(z=x\) gives equality,
which proves the first formula.  The second calculation is immediate from the
definition of pullback.  The entire family is functorial under isometries, so
its mathematical construction inserts no distinguished point.  Whether all
of these functions are admitted as physical tests is exactly the separate
closure clause of
Definition~\ref{p2:def:metric-partial-evaluation-closure}.  Restriction to
the dense image of \(X\) recovers the original profiles \(d(z,\cdot)\).
\end{proof}

\begin{lemma}[Peter--Weyl finite-image comparison witnesses for unmarked roles]
\label{p2:lem:peter-weyl-finite-image-reduction}
Let \(H\) be the compact connected exact-operation parent obtained from one
candidate realization of \(\mathsf E_R\).
\begin{enumerate}[label=(\roman*), leftmargin=2em]
\item If \(\epsilon_R=\mathrm{ab}\), \(H\) has an effective
continuous image isomorphic to \(U(1)\).
\item If \(\epsilon_R=\mathrm{nc}\), \(H\) has a finite-dimensional connected
noncommutative compact Lie image.
\end{enumerate}
The image in either clause is only a comparison witness.  A complete candidate
with that local Lie type must pass its own GP1--GP2 exactification before it
enters the GP3 preorder.
\end{lemma}

\begin{proof}
Peter--Weyl separation~\cite{Peter1927} supplies finite-dimensional unitary
representations that distinguish points of every compact group.  In the
abelian case, choose a nontrivial one.  Its commuting unitary matrices are
simultaneously diagonalizable, and at least one diagonal character is
nontrivial.  Connectedness makes the image of that character a nontrivial
compact connected subgroup of \(U(1)\), hence all of \(U(1)\).

In the noncommutative case choose, only inside the candidate and only to
witness the GP2-derived branch status, \(r,s\in H\) with
\(c=[r,s]\ne e\).  Peter--Weyl supplies a finite-dimensional unitary
representation \(\rho\) with \(\rho(c)\ne I\).  Its connected image
\(\rho(H)\) is therefore a noncommutative compact Lie group.  Neither the
chosen character, commutator pair, nor representation is retained as prior
role data.
\end{proof}

\begin{lemma}[Least local comparison type of the noncommutative role]
\label{p2:lem:a1-least-local-reduction}
Every connected noncommutative compact Lie candidate has a closed
connected local subrole whose Lie algebra is
\(\mathfrak{su}(2)\).  Conversely, the local type
\(\mathfrak{su}(2)\) has no strict image, connected-local-subrole, or
factor-deletion comparison that preserves the geometrically primary
order-sensitive role.  It is therefore the least local role type in the
declared structural preorder once the target is completed and exactified
candidate-internally.
\end{lemma}

\begin{proof}
The Lie algebra \(\mathfrak l\) of a connected noncommutative compact Lie
group is reductive with nonzero semisimple derived algebra.  A root of that
semisimple algebra gives a compact real root subalgebra
\(\mathfrak s_\alpha\cong\mathfrak{su}(2)\), and the compact root-subgroup
theorem integrates it to a homomorphism
\[
        \varphi_\alpha:SU(2)\longrightarrow L
\]
with connected image and differential image
\(\mathfrak s_\alpha\)~\cite{Hall2015lie}; the image is closed because
\(SU(2)\) is compact and a continuous image of a compact set in a Hausdorff
group is closed.  This exposes an allowed connected-local-subrole comparison
witness and preserves the GP2-derived noncommutative branch status.  The
target candidate is completed separately; no chosen root or embedding is
frozen.

Every proper Lie subalgebra of \(\mathfrak{su}(2)\) is abelian, while
simplicity gives no nonzero proper Lie-algebra quotient and no nontrivial
direct-sum factor.  Thus each of the three elementary comparisons in
Definition~\ref{p2:def:global-gp3-preorder} either leaves the local type
unchanged or loses the derived noncommutative branch status.
\end{proof}

\begin{lemma}[Faithful compact paired models and relation-exact endpoints]
\label{p2:lem:relation-exact-endpoint-models}
Let \(C\) be any compact connected Lie group.  Give \(C\) any positive
bi-invariant Riemannian metric \(g\), let
\(d_g\) be its geodesic distance, and choose a faithful finite-dimensional
unitary representation \(\rho:C\to U(V)\).  On
\(\mathcal H_C:=\mathbb C\oplus V\), use \(1\oplus\rho\).  Conjugation on
\(\operatorname{End}(\mathcal H_C)\) is a faithful exact reversible action,
while
\[
        \mathfrak P_C:=(C,d_g,C,L)
\]
is its smooth parent self-presentation under left translation.  Hence
\[
 \mathfrak F_C
 :=\bigl(\mathcal H_C,\operatorname{End}(\mathcal H_C),
          \operatorname{Ad}_{1\oplus\rho}\bigr),
 \qquad
 (\mathfrak F_C,\mathfrak P_C)
\]
is a paired FRC-admissible role model: the ontic carrier is finite, the parent
comparison is compact, and the full mathematical parent distance-profile
family separates parent points and operations.  No marked element, axis,
basis, or particular proof representation is promoted into the role
signature; the positive metric parameters remain free presentation data.

In particular, when \(C\) is the compact source-simply-connected
integration exposed by a finite-dimensional comparison witness---including \(SU(2)\),
and with the abelian line handled by its FRC-forced compact globalization
\(U(1)\)---exactification of a quotient endpoint is represented by this
faithful source-form paired model.  It does not mean that a restored central
kernel acts trivially on an old quotient carrier, and it does not make the
smooth group an infinite ontic carrier.
\end{lemma}

\begin{proof}
Every compact Lie group has a faithful finite-dimensional unitary
representation.  If conjugation by \(1\oplus\rho(g)\) is trivial on the full
matrix algebra, then \(1\oplus\rho(g)\) is scalar.  Its trivial summand forces
that scalar to be \(1\), and faithfulness of \(\rho\) gives \(g=e\).  Thus
the finite reversible action is faithful, including on the center.  A compact
Lie group admits a positive bi-invariant metric by averaging an inner product
over the adjoint action.  Compactness of \((C,d_g)\) gives the bounded parent
comparison, and
Theorem~\ref{p2:thm:canonical-metric-profile-separation} gives profile
separation for the regular parent action.  Varying the positive metric changes
no structural role datum.

For \(SU(2)\), the source integration is already compact, so this paired
model is the relation-exact endpoint.  For the abelian local line, the source
integration \(\mathbb R\) is not FRC-b bounded as a global parent; every
one-dimensional lattice globalization is abstractly \(U(1)\), with its metric
scale still free.  Replacing a quotient endpoint by the paired admissible model
constructs another complete exact member of \(\operatorname{Adm}_R\); it never
reuses a quotient action with a newly invisible kernel.
\end{proof}

\begin{lemma}[Compact globalizations of an unmarked abelian local type]
\label{p2:lem:abelian-compact-globalization}
Let \(\mathfrak a\) be a real abelian Lie algebra of dimension \(k>0\).  Up to
unmarked Lie-group isomorphism, there is exactly one compact connected Lie
group with Lie algebra isomorphic to \(\mathfrak a\), namely
\[
 T^k\cong\mathbb R^k/\mathbb Z^k.
\]
This uniqueness selects no lattice basis, period, metric scale, or marked
identification of \(\operatorname{Lie}(T^k)\) with \(\mathfrak a\).
\end{lemma}

\begin{proof}
The simply connected integration of \(\mathfrak a\) is the additive group
\((\mathfrak a,+)\).  Every connected globalization is
\(\mathfrak a/\Lambda\) for a discrete subgroup \(\Lambda\), and the quotient
is compact exactly when \(\Lambda\) is a full lattice of rank \(k\).  Given
two full lattices, a linear isomorphism sending a lattice basis of one to a
lattice basis of the other descends to an isomorphism of the quotients.  Every
linear isomorphism of \(\mathfrak a\) is a Lie-algebra automorphism because the
bracket is zero.  The descended isomorphism is not canonical, which is why no
marked lattice or normalization follows.
\end{proof}

\begin{lemma}[GP2 removes gratuitous pure semisimple kernel relations]
\label{p2:lem:no-gratuitous-semisimple-kernel}
Under the presentation-theoretic relation-exact semantics of GP2 fixed in
Postulate~\ref{p2:post:GP2}, let \(H\) be a compact connected
exact-operation parent in the unmarked primitive GP1/GP2-admissible class,
with Lie algebra
\[
        \mathfrak h=\mathfrak z\oplus\mathfrak s,
\]
where \(\mathfrak z\) is abelian and \(\mathfrak s\) is compact semisimple.
Write the simply connected integration as
\[
        \widetilde H=\mathbb R^r\times S,
\]
where \(S\) is compact, simply connected, and semisimple, and write
\[
        H\cong \widetilde H/\Gamma
\]
for a discrete central subgroup \(\Gamma\).  Then
\[
        \Gamma\cap\bigl(\{0\}\times Z(S)\bigr)=\{e\}.
\]
Consequently the canonical homomorphism \(S\to H\) is injective.  Thus every
semisimple source factor is present as an actual closed subgroup of the
GP2-exact parent; FRC compactification may impose the lattice relations needed
to compactify the abelian covering directions, but, in this GP2 domain, it
does not license an additional pure semisimple central quotient relation.
\end{lemma}

\begin{proof}
Because \(\Gamma\) is central,
\(\Gamma\subseteq\mathbb R^r\times Z(S)\).  Let
\[
 p:\mathbb R^r\times S\longrightarrow\mathbb R^r
\]
be projection and put
\[
 K:=\Gamma\cap\bigl(\{0\}\times Z(S)\bigr),
 \qquad \Lambda:=p(\Gamma).
\]
The center \(Z(S)\) is finite.  Hence \(\Lambda\) is discrete: if a sequence
of distinct points of \(\Lambda\) converged in \(\mathbb R^r\), choose lifts
in \(\Gamma\); after passing to a subsequence their finite \(Z(S)\)-components
agree, so the differences of the lifts converge to the identity, contradicting
discreteness of \(\Gamma\).

The quotient map \(H\to\mathbb R^r/\Lambda\), induced by projection to the
abelian factor, is continuous and onto.  Compactness of \(H\) therefore makes
\(\mathbb R^r/\Lambda\) compact.  Thus \(\Lambda\) is a full lattice and, as
an abstract group, \(\Lambda\cong\mathbb Z^r\).  There is an exact sequence
of abelian groups
\[
 1\longrightarrow K\longrightarrow\Gamma
   \xrightarrow{\ p\ }\Lambda\longrightarrow1.
\]
Because \(\Lambda\) is free abelian, this sequence splits.  Choose a subgroup
\(\Gamma_0\le\Gamma\) on which \(p\) is an isomorphism onto \(\Lambda\).
Then
\[
 \Gamma=K\Gamma_0,
 \qquad K\cap\Gamma_0=\{e\}.
\]
Set
\[
 H_0:=(\mathbb R^r\times S)/\Gamma_0.
\]
To see directly that \(H_0\) is compact, choose a compact fundamental
parallelepiped \(F\subset\mathbb R^r\) for the full lattice \(\Lambda\).
For every \((x,s)\in\mathbb R^r\times S\), the unique element of
\(\Gamma_0\) over a suitable lattice translate moves the first coordinate
into \(F\).  Hence the quotient image of the compact set \(F\times S\) is all
of \(H_0\).

Every compact Lie group admits a faithful finite-dimensional unitary
representation, so the construction of
Lemma~\ref{p2:lem:relation-exact-endpoint-models} gives \(H_0\) a faithful
paired FRC realization.  The natural quotient homomorphism
\[
        H_0\longrightarrow H
\]
has kernel canonically isomorphic to \(K\).  Thus a nontrivial \(K\) imposes
additional central equalities on a compact paired realization that already
satisfies FRC.  The geometrically primary local source and previously frozen
inventory generate no such pure semisimple quotient relation.  By the stated
GP2 no-insertion semantics,
\(K\ne\{e\}\) is therefore inadmissible.  Hence \(K=\{e\}\).  The image in
\(H\) of \(\{0\}\times S\) then has trivial kernel, so \(S\hookrightarrow H\).
\end{proof}

\begin{proposition}[FRC--GP2 admissible-globalization operator]
\label{p2:prop:admissible-globalization-operator}
For an unmarked connected compact-reductive local Lie type
\(\mathfrak g\) exposed by an FRC exact-operation parent, let
\(\mathcal K(\mathfrak g)\) be the set of unmarked isomorphism classes of
\emph{closed, connected} globalizations of \(\mathfrak g\) that admit a
faithful paired FRC model.  Write the simply connected integration in the
standard form
\[
       \widetilde G=\mathbb R^r\times S,
\]
with \(S\) compact, simply connected, and semisimple.  For a compact
class represented as \((\mathbb R^r\times S)/\Gamma\), define
\(\mathcal K_{\rm ex}(\mathfrak g)\subseteq\mathcal K(\mathfrak g)\) by the
GP2 relation-exact condition
\[
       \Gamma\cap\bigl(\{0\}\times Z(S)\bigr)=\{e\}.
\]
Lemma~\ref{p2:lem:no-gratuitous-semisimple-kernel} proves that every
GP2-exact compact globalization lies in this subset.  The cumulative FRC--GP2
rule is therefore the following partial operator:
\[
 E_{\rm adm}(\mathfrak g)=
 \begin{cases}
 \text{FRC-obstructed},
   & \mathcal K_{\rm ex}(\mathfrak g)=\varnothing,\\[0.3em]
 [\widetilde G],
   & [\widetilde G]\in\mathcal K_{\rm ex}(\mathfrak g),\\[0.3em]
 \text{the unique class in }\mathcal K_{\rm ex}(\mathfrak g),
   & [\widetilde G]\notin\mathcal K_{\rm ex}(\mathfrak g),\quad
     |\mathcal K_{\rm ex}(\mathfrak g)|=1,\\[0.3em]
 \mathcal K_{\rm ex}(\mathfrak g)\text{ unselected},
   & |\mathcal K_{\rm ex}(\mathfrak g)|>1.
 \end{cases}
\]
The empty branch is the existing FRC obstruction verdict.  The source branch
is GP2 no-insertion: no unnecessary kernel relation is added to an already
admissible source.  On a noncompact covering source, FRC may force a lattice
in the abelian covering directions, but GP2 removes any additional pure
semisimple central kernel before the singleton/unresolved comparison.  The
singleton branch is FRC-forced compactification, not a freely chosen quotient.
The last branch records a genuine unresolved relation-exact family rather than
choosing extra data.

In particular,
\[
 E_{\rm adm}(\mathfrak u(1))=[U(1)],
 \qquad
 E_{\rm adm}(\mathfrak{su}(2))=[SU(2)].
\]
\end{proposition}

\begin{proof}
The input \(\mathfrak g\) of this operator is already a
finite-dimensional compact-reductive \emph{Lie algebra} by hypothesis.
Its compact globalizations are compact Lie groups and those admitted by
\(\mathcal K(\mathfrak g)\) have faithful paired models by definition.
This says nothing about Lie regularity of the possibly non-Lie ambient
parent from which a finite Lie-image germ was exposed.
If no admissible compact globalization exists, this local-germ comparison
branch is FRC-obstructed by definition.
If \(\widetilde G\) itself is compact, Lemma~
\ref{p2:lem:relation-exact-endpoint-models} supplies a faithful paired model,
and every proper connected globalization \(\widetilde G/K\) adds a nontrivial
kernel relation.  GP2 therefore retains \(K=\{e\}\).  If the source is not an
admissible compact parent, FRC first restricts the comparison to compact
members of \(\mathcal K(\mathfrak g)\), and GP2 then removes the gratuitous
pure semisimple kernels by
Lemma~\ref{p2:lem:no-gratuitous-semisimple-kernel}.  This leaves
\(\mathcal K_{\rm ex}(\mathfrak g)\); uniqueness permits no choice, whereas
more than one surviving isomorphism class supplies no selector.

For \(\mathfrak u(1)\), the semisimple factor is trivial, so
\(\mathcal K_{\rm ex}(\mathfrak u(1))=\mathcal K(\mathfrak u(1))\), and
Lemma~\ref{p2:lem:abelian-compact-globalization} gives the singleton
\(\{[U(1)]\}\).  For \(\mathfrak{su}(2)\), the source is the compact group
\(SU(2)\), so the source branch applies.  A dense winding copy of \(\mathbb R\) in a torus is not a
counterexample: it is not closed, and its closure has the Lie algebra of the
ambient subtorus rather than the one-dimensional local type unless that
subtorus is a circle.
\end{proof}

\begin{theorem}[Primitive role endpoint candidate comparison]
\label{p2:thm:primitive-endpoint-lift-reachability}
Let \(\mathfrak M\in\operatorname{Adm}_R\) be any complete
GP1-admissible, candidate-internally GP2-exact-and-role-faithful realization of
an unmarked applied primitive role.  Then the role-local construction below
produces a second complete exact candidate \(\mathfrak M_*\) and a
role-faithful comparison
\[
       \mathfrak M_*\preceq_R\mathfrak M
\]
with
\[
 \operatorname{Skel}_R(\mathfrak M_*)\cong
 \begin{cases}
 U(1),&\epsilon_R=\mathrm{ab},\\[2pt]
 SU(2),&\epsilon_R=\mathrm{nc}.
 \end{cases}
\]
provided the displayed endpoint realizes every intrinsic diagram of
\(\operatorname{Geo}_R\).  This proviso is automatic in the Paper-II primitive
domain; if Geometric Primacy constitutes richer toral directions, brackets,
relations, or tests not realized by the endpoint, the comparison is
inadmissible and the theorem returns ``not applicable'' rather than deleting
them.  If the skeleton of \(\mathfrak M\) is not already the displayed
endpoint, the comparison is strict.  No carrier basis, root, circle embedding,
or metric normalization chosen in the proof is frozen into \(P_{<R}\).
\end{theorem}

\begin{proof}
Let \(H=H_{\mathfrak M}\) be the nontrivial compact
connected \emph{topological} parent supplied by
Theorem~\ref{p2:thm:geometric-existence-exactification}.  No faithful
finite-dimensional action of the \emph{entire} \(H\) is used.  In the primitive
preorder of Definition~\ref{p2:def:global-gp3-preorder}, each effective
homomorphic-image or closed connected local-subrole step enters GP3 only
after its retained germ passes through \(E_{\rm adm}\) and receives a
complete GP1/GP2-exact, role-faithful paired endpoint.  This exactification
is \emph{not} asserted to be a homomorphic image of the original parent.

If \(\epsilon_R=\mathrm{ab}\), the Peter--Weyl reduction
Lemma~\ref{p2:lem:peter-weyl-finite-image-reduction} yields a nontrivial
continuous character \(\chi:H\twoheadrightarrow U(1)\): its compact
connected nontrivial image is the entire circle.  This is an allowed
branch-preserving image comparison.  Since
\(E_{\rm adm}(\mathfrak u(1))=[U(1)]\),
Lemma~\ref{p2:lem:relation-exact-endpoint-models} gives a new complete
faithful finite paired endpoint \(\mathfrak M_{U(1)}\).  Its regular
comparison geometry realizes the primitive commuting diagrams, without
transporting the old carrier or freezing the character.  Thus
\(\mathfrak M_{U(1)}\preceq_R\mathfrak M\).

If \(\epsilon_R=\mathrm{nc}\), choose only as a proof witness
\(r,s\in H\) with \(c=[r,s]\ne e\).  Peter--Weyl supplies a finite
unitary representation \(\rho\) with \(\rho(c)\ne I\).  The compact,
connected, noncommutative image \(K=\rho(H)\) is a Lie group, with
\[
 \mathfrak k=\operatorname{Lie}(K)=\mathfrak z\oplus\mathfrak s,
 \qquad\mathfrak s\ne0.
\]
The \emph{image germ}, not an unqualified raw quotient of the final source
form, is the witness for comparison (i).  It has at least one admissible
compact exact globalization: writing the simply connected integration
as \(\mathbb R^r\times S\), with \(S\) compact simply connected
semisimple, choose a full lattice \(\Lambda\subset\mathbb R^r\).  Then
\[
 C_0=(\mathbb R^r\times S)/(\Lambda\times\{e\})
       \cong T^r\times S
\]
is compact, connected, has germ \(\mathfrak k\), and satisfies
\(\Gamma\cap(\{0\}\times Z(S))=\{e\}\).  Its faithful finite paired
model is provided by
Lemma~\ref{p2:lem:relation-exact-endpoint-models}; therefore
\(\mathcal K_{\rm ex}(\mathfrak k)\ne\varnothing\).
If \(E_{\rm adm}(\mathfrak k)\) returns several unmarked compact classes,
none is physically selected.  Every allowed class
\(C=(\mathbb R^r\times S)/\Gamma\) has
\(\Gamma\cap(\{0\}\times Z(S))=\{e\}\), so the compact factor \(S\)
embeds as a closed subgroup of \(C\) by
Lemma~\ref{p2:lem:no-gratuitous-semisimple-kernel}.
Choose a simple factor of \(S\) and a simple root \(\alpha\) only in the
proof.  Its root map \(SU(2)\to S\) is injective: its possible kernel is
central, and the simple coroot is primitive in the cocharacter lattice of
the simply connected semisimple factor, so \(-1\) is not killed.  Thus
\emph{every} allowed exact globalization contains a closed embedded
\(SU(2)\) root subrole.

The resulting complete exact intermediate paired candidate \(\mathfrak M_C\)
is below \(\mathfrak M\) by the image-germ and exactification rule (i).
Applying the closed connected local-subrole rule (ii), followed by
\(E_{\rm adm}(\mathfrak{su}(2))=[SU(2)]\), supplies a complete exact
paired endpoint \(\mathfrak M_{SU(2)}\preceq_R\mathfrak M_C\).
Transitivity therefore gives
\(\mathfrak M_{SU(2)}\preceq_R\mathfrak M\).
Neither \(H\to C\) nor a full-parent finite linear port has been
asserted.  Every intermediate and final candidate is internally exact;
the unmarked primitive role's noncommuting diagram, previous inventory
and effective connected reapplication are preserved.  A richer intrinsic
bracket, extra toral direction or separately generated test would require a
fresh role-faithfulness check and may invalidate the comparison.

It remains to justify \emph{strictness}, not infer it from distinct group
names.  Starting at the \(U(1)\) endpoint, every nontrivial connected
image or connected local subrole retaining the commuting primitive role has
the same \(\mathfrak u(1)\) germ; exactification returns \(U(1)\), and
there is no unused connected direct-product factor.  Starting at \(SU(2)\),
any noncommutative connected image has the same \(\mathfrak{su}(2)\) germ
(including a raw \(SO(3)\) image, exactified back to \(SU(2)\)); a proper
closed connected subgroup is abelian and fails the order-sensitive branch,
and no nontrivial connected product factor can be deleted.  Consequently a
finite branch-preserving chain starting at either exact endpoint cannot
reach a different skeleton, including a non-Lie source.  Hence a source
with a different skeleton has no reverse comparison and the displayed
endpoint comparison is strict.  No choice of character, root, lattice,
carrier or metric is frozen into \(P_{<R}\).
\end{proof}

\begin{corollary}[Unconditional endpoint skeleton of every primitive GP3 irreducible]
\label{p2:cor:primitive-irreducible-endpoint-skeleton}
Every GP3-irreducible complete structured realization in the primitive
endpoint-comparison domain has skeleton isomorphic to \(U(1)\) on the abelian
branch and to \(SU(2)\) on the noncommutative branch.  Distinct unmarked
faithful structured lifts over that common skeleton may remain unresolved.
\end{corollary}

\begin{proof}
If an irreducible structured realization had a different skeleton,
Theorem~\ref{p2:thm:primitive-endpoint-lift-reachability} would supply a strict
role-faithful comparison to a complete exact endpoint candidate, contradicting
GP3 irreducibility.  The last sentence is
Lemma~\ref{p2:lem:skeleton-factorization}\,(b).
\end{proof}

\begin{theorem}[Master unmarked compact-role allocation chain]
\label{p2:thm:master-compact-role-allocation}
\label{p2:thm:gp-least-role-chain}
\label{p2:thm:gp-minimal-role-chain}
Let \(\mathsf E_R\) be the unmarked geometrically primary primitive-role
assertion of Definition~\ref{p2:def:geometric-role-existence}, and form
\(\operatorname{Adm}_R\) by applying GP1 and GP2 candidatewise.  By
definition, if \(\operatorname{Adm}_R=\varnothing\), the role is
FRC-obstructed and no endpoint parent is allocated.  Assume henceforth that
\(\operatorname{Adm}_R\ne\varnothing\).  By
Theorem~\ref{p2:thm:geometric-existence-exactification}, every surviving
candidate has a nontrivial connected compact \emph{topological} parent;
that parent is Lie when the full-parent finite affine-operation port or an
independent Lie certificate has actually been supplied.  All candidates
realize the same classified exact status
\(\epsilon_R\in\{\mathrm{ab},\mathrm{nc}\}\) of the constituted geometry.
No designated witness pair or candidate self-presentation is inserted into
\(P_{<R}\).  Every candidate in the resulting exact branch is a paired FRC model
of Definition~\ref{p2:def:cumulative-role-allocation}.  Its generated parent
comparison supplies the canonical mathematical metric profiles, which detect
the geometric action kernel; under metric-curry closure they also make the
candidate's B-test action faithful.  The profiles are not identified across
candidates and are not promoted into the prior inventory.  Apply GP3
in the order of Definition~\ref{p2:def:cumulative-role-allocation}, using the
preorder of Definition~\ref{p2:def:global-gp3-preorder}.  The endpoint clauses
below apply exactly when the complete endpoint candidate realizes the same
\(\operatorname{Geo}_R\); richer geometrically primary content is preserved and
causes the endpoint comparison to abstain rather than be erased.
\begin{enumerate}[label=(\roman*), leftmargin=2em]
\item Candidatewise, every generated parent comparison has a canonical
compact completion.  Whenever a continuous faithful finite-dimensional
operation port of the \emph{entire closed parent} has been licensed, that
closed parent is compact Lie and its exact reversible action is unitarizable.
Neither conclusion inserts such a port into a candidate lacking one.
\item If \(\epsilon_R=\mathrm{ab}\), its globally unique least nontrivial connected
primitive-role skeleton has parent
\[
 H_R\cong U(1)\cong S^1.
\]
\item If \(\epsilon_R=\mathrm{nc}\), the unique least
primitive-role skeleton has faithful local Lie type
\[
 \operatorname{Lie}(H_R)\cong\mathfrak{su}(2).
\]
The GP2-admissible representative of this local type has the relation-exact
globalization
\[
 H_R\cong SU(2)\cong S^3.
\]
In both the abelian and noncommutative branches the displayed parent
operation-role class is the unique least primitive-role skeleton in the sense
of Definition~\ref{p2:def:primitive-role-skeleton}.

\item Every GP3-irreducible complete structured realization in the endpoint
comparison domain has the corresponding displayed skeleton.  This follows
from the candidatewise whole-candidate comparison of
Theorem~\ref{p2:thm:primitive-endpoint-lift-reachability}, not from a claim that
the endpoint retains every self-operation of every rejected candidate.
Distinct unmarked
faithful structured lifts over the common skeleton may nevertheless remain
unresolved.  In particular, the theorem does not select a carrier basis,
faithful carrier representation, or metric scale unless a later signature
marks it.

\item On the noncommutative branch the negative Killing form
\[
 g_K(X,Y):=-\alpha_K B_{\mathfrak{su}(2)}(X,Y),\qquad \alpha_K>0,
\]
is the unique automorphism-invariant positive quadratic form up to scale.
Its translated metric is the round metric on \(S^3(\Rmom)\).  No GP fixes
\(\Rmom\), and no claim is made that the primitive parent comparison equals
its geodesic distance.
\end{enumerate}
\end{theorem}

\begin{proof}
For part (i), compact completion of each constituted parent comparison is
Theorem~\ref{p2:thm:gp1-bounded-compactness}; the compact-Lie and
Haar-unitarity upgrades follow from
Theorem~\ref{p2:thm:frc-automatic-compact-lie-closure} only when its
faithful finite-operation antecedent holds.  The canonical metric
profiles detect the parent geometric action kernel by
Theorem~\ref{p2:thm:canonical-metric-profile-separation}; their physical
B-test use additionally invokes
Definition~\ref{p2:def:metric-partial-evaluation-closure}.

For the remaining parts, GP3 is evaluated after candidate-internal GP2
exactification and relation filtering.  Role faithfulness means preservation
of \(P_{<R}\), every intrinsic diagram of \(\operatorname{Geo}_R\), and the
constituted-and-realized nontriviality, connectedness, and classified abelian
or noncommutative branch; each comparison endpoint remains exact to its own
generated geometry.
The endpoint existence and, crucially, the universal
``below every candidate'' quantifier are supplied directly by
Theorem~\ref{p2:thm:primitive-endpoint-lift-reachability}.

On the abelian branch, Lemma~\ref{p2:lem:relation-exact-endpoint-models}
supplies an admissible paired \(U(1)\) model, and
Theorem~\ref{p2:thm:primitive-endpoint-lift-reachability} gives, from
each admissible primitive abelian candidate, a complete role-faithful endpoint
candidate whose skeleton is \(U(1)\).  Hence the \(U(1)\)
skeleton class lies below every candidate skeleton in the induced order.  The
group \(U(1)\) has no proper nontrivial connected homomorphic image or
connected subrole and no nontrivial connected direct-product factor to delete.
Therefore this skeleton is irreducible and is the unique least abelian
primitive-role skeleton; no single common structured lift is asserted.  Lemma~\ref{p2:lem:peter-weyl-finite-image-reduction} and
Proposition~\ref{p2:prop:admissible-globalization-operator} give an independent
local/global replay of the same endpoint classification, but they are not used
to replace this candidatewise quantifier.

On the noncommutative branch,
Theorem~\ref{p2:thm:primitive-endpoint-lift-reachability} gives, from each
admissible primitive noncommutative candidate, a complete role-faithful
endpoint candidate whose skeleton is \(SU(2)\).  Hence the \(SU(2)\)
skeleton class lies below every candidate
skeleton.  Existence of admissible paired endpoint lifts is
Lemma~\ref{p2:lem:relation-exact-endpoint-models}.  It remains only to check
that no strict elementary comparison below the endpoint preserves the
geometrically primary order-sensitive role.  A connected homomorphic image of \(SU(2)\) that still satisfies
the derived nontrivial noncommutative status is either \(SU(2)\) or \(SO(3)\); the
trivial image is inadmissible because it loses nontriviality and
noncommutativity.  GP2 exactification sends the \(SO(3)\) quotient back to the
source-global \(SU(2)\) class because the additional equality \(-1=1\) is
ungenerated.  Every proper closed connected subgroup of \(SU(2)\) is trivial
or one-dimensional abelian and therefore fails the derived noncommutative
branch status and intrinsic order-swap diagram.
Finally \(\mathfrak{su}(2)\) is simple, so connected \(SU(2)\) has no
nontrivial connected direct-product factor to delete.  Thus the endpoint skeleton is irreducible and, being represented below every
candidate skeleton, is the unique least noncommutative primitive-role
skeleton.  Again, this does not assert a single common structured lift.

Corollary~\ref{p2:cor:primitive-irreducible-endpoint-skeleton} now follows
with the same candidatewise quantifier: every GP3-irreducible complete
structured realization has the corresponding \(U(1)\) or \(SU(2)\) skeleton.
Distinct unmarked faithful presentations over that common skeleton need not be
comparable and may remain an unresolved presentation family.

Finally, on the compact simple algebra \(\mathfrak{su}(2)\), every
automorphism-invariant symmetric bilinear form is a scalar multiple of the
Killing form.  Indeed, the real adjoint module is absolutely irreducible: its
complexification is the irreducible three-dimensional
\(\mathfrak{sl}_2(\mathbb C)\)-module.  Equivalently, an endomorphism commuting
with every element of
\(\operatorname{Aut}(\mathfrak{su}(2))\cong SO(3)\) is scalar.  Real Schur
then forces any invariant form to be proportional to a fixed one.  Its
negative is positive definite and its translated metric on
\(SU(2)\) is round.  Positive rescaling leaves all three GPs invariant, so
the radius remains free.
\end{proof}

\begin{corollary}[Rank and center diagnostics]
\label{p2:cor:rank-center-diagnostics}
For the \(U(1)\) and \(SU(2)\) \emph{selected endpoints} of
Theorem~\ref{p2:thm:master-compact-role-allocation}, finite-dimensional Lie
regularity follows from the explicit compact Lie endpoint models, while
operational Cartan rank one and a one-parameter ray of positive invariant
Riemannian metrics are consequences of the least-skeleton role allocation.  On the noncommutative branch
this ray is the negative-Killing ray \(g_K\).  On the abelian branch the
Killing form vanishes, and the metric ray is instead the unique positive
invariant quadratic ray on \(\mathfrak u(1)\).  On the noncommutative branch a
faithful representation contains center-odd, equivalently half-integral,
sectors exactly when it is sensitive to the source center.  This diagnoses
the already selected \(SU(2)\) form; it does not select that form.
\end{corollary}

\begin{proof}
Rank one is immediate for \(\mathfrak u(1)\) on the abelian branch and
is read off \(\mathfrak{su}(2)\) on the noncommutative branch, whose local
selection is Lemma~\ref{p2:lem:a1-least-local-reduction}; the metric ray is the Killing
uniqueness established in the proof of
Theorem~\ref{p2:thm:master-compact-role-allocation}, while on the abelian
branch the Killing form vanishes and every positive form on a one-dimensional
space is a positive multiple of every other.  The central-character statement
is proved with
Theorem~\ref{p2:thm:relation-minimal-globalization}.
\end{proof}

\begin{corollary}[Conditional comparison of primitive and Killing metrics]
\label{p2:cor:primitive-killing-metric-comparison}
Suppose a retained faithful realization over the derived least \(SU(2)\)
skeleton carries a primitive comparison satisfying (FD)+(LR), so that
\(D_\infty\), \(F\), and \(g_F\) refer to that same operation group.  Fix any representative \(g_K\) of the
canonical negative-Killing ray, with geodesic distance \(d_K\).  Then there
are constants \(\beta>0\) and, for every faithful finite-dimensional unitary
proof representation \(\rho\), \(a_\rho>0\), such that
\[
 F(X)=\beta\sqrt{g_K(X,X)},\qquad
 g_F=\beta^2g_K,\qquad
 g_\rho=a_\rho g_K.
\]
Consequently,
\[
 D_\infty\le\beta d_K\le C_R D_\infty
 \qquad\text{near the identity},
\]
and for some \(C_K<\infty\),
\[
 D_\infty\le\beta d_K\le C_KD_\infty
 \qquad\text{on all of \(SU(2)\)}.
\]
This proves conditional bi-Lipschitz equivalence, not equality of the
primitive and Killing distances, and fixes no metric scale.
\end{corollary}

\begin{proof}
The norm \(F\) is Ad-invariant by
Corollary~\ref{p2:cor:gp1-first-order-norm-derived}.  The adjoint action of
\(SU(2)\) is \(SO(3)\), which is transitive on every \(g_K\)-unit sphere.
Thus \(F\) is constant, say \(\beta\), on that sphere; positive homogeneity
gives the first formula.  Its unit ball is already an ellipsoid, so the John
construction returns that same ball and \(g_F=\beta^2g_K\).  The trace metric
\(g_\rho\) is a positive Ad-invariant symmetric form on the simple algebra
\(\mathfrak{su}(2)\), hence is \(a_\rho g_K\) with \(a_\rho>0\).

The Finsler distance of \(F=\beta\sqrt{g_K}\) is \(\beta d_K\).
Theorem~\ref{p2:thm:gp1-propagation-derived} gives
\(D_\infty\le D_F\le C_RD_\infty\) near the identity.  The first inequality
holds globally because the intrinsic length metric dominates its defining
metric.  Away from a fixed identity neighborhood, the continuous ratio
\(D_F(e,g)/D_\infty(e,g)\) is bounded on a compact set; bi-invariance
transports the resulting bound to every pair of points.  Enlarging the local
constant if necessary gives \(C_K\).
\end{proof}

\begin{definition}[One-dimensional operational factorization]
\label{p2:def:operational-rank-one-factorization}
Let \(R\) be a generated operation role and set
\[
 H_R:=\rho_A(\langle R\rangle)
 \leq\operatorname{Aut}(\mathcal B(X)).
\]
It is \emph{operationally one-dimensionally factorable} when
there exist a connected one-dimensional Lie operation parent \(C\), a
homomorphism \(\phi:H_R\to C\), and an action
\(\sigma:C\to\operatorname{Aut}(\mathcal B(X))\) such that
\[
 \iota_{H_R}=\sigma\circ\phi,
\]
where \(\iota_{H_R}\) is the displayed inclusion, and this factorization
preserves every generated transformation--test
distinction.  It has a \emph{one-dimensional nonfactorizing residue} when no
such factorization exists.  This notion concerns factorization of an action.
It is not Cartan rank.  The primary compact-role classification is instead
the global least-skeleton theorem
\ref{p2:thm:master-compact-role-allocation}; (OR1) is a consequence on the
derived least parent-role skeletons.
\end{definition}

\begin{theorem}[Nonabelian commutator equivalence and conditional endogenous B-witness]
\label{p2:thm:gp2-commutator-dichotomy}
\label{p2:thm:gp2-generated-rank-one-obstruction}
Let \(H_R\) be the generated transformation parent of a role \(R\), and let
\[
 \rho_R:H_R\longrightarrow\operatorname{Aut}(\mathcal B_R)
\]
be its action on the generated test doctrine.  The following are equivalent
provided the constituted doctrine is transformation-separating, equivalently
\(\rho_R\) is faithful:
\begin{enumerate}[label=(\roman*), leftmargin=2em]
\item \(H_R\) is nonabelian;
\item there exist \(r,s\in H_R\) with \([r,s]\ne e\);
\item there is an endogenous order-sensitive witness
\[
 \exists r,s\in H_R,\ \exists b\in\mathcal B_R:
 \qquad [r,s]^*b\ne b.
 \]
\end{enumerate}
For a primitive role arising from \(\mathsf E_R\), clauses (i)--(ii) give the
two exhaustive algebraic statuses derived in
Theorem~\ref{p2:thm:geometric-existence-exactification}; neither is a
primitive role-type assumption.
Clause (iii) is the additional operational form under the displayed separation
hypothesis.  When \(H_R\) acts by the geometric isometries retained in this paper and the
B-test grammar is metric-curry closed, the witness in (iii) may be chosen
from the canonical distance-profile family of
Theorem~\ref{p2:thm:canonical-metric-profile-separation}.  Without
transformation separation, no universal physical \(b\in\mathcal B_R\) is
claimed.  When separation is constituted but metric-curry closure is absent,
the theorem supplies some \(b\) but does not identify it with a distance
profile.  When these conditions hold, every factorization through an abelian
one-dimensional parent loses a generated distinction and therefore fails GP2
faithfulness; no GP3 comparison may reintroduce that loss.
Thus the noncommutative branch is unconditional inside an admissible
order-sensitive exactification, while its physical B-witness is conditional on
the constituted separation doctrine.
\end{theorem}

\begin{proof}
The equivalence of (i) and (ii) is the definition of nonabelianity.  If (iii)
holds, then \([r,s]\ne e\), because the identity fixes every test, so (iii)
implies (ii).  Conversely, suppose \(c=[r,s]\ne e\).  Both \(c\) and \(e\)
are generated transformations.  By the explicit faithfulness hypothesis on
\(\rho_R\), the supplied B-test doctrine separates them: there is
\(b\in\mathcal B_R\) with
\[
        c^*b\ne e^*b=b.
\]
If the geometric B grammar is metric-curry closed,
Theorem~\ref{p2:thm:canonical-metric-profile-separation} permits \(b\) to be
chosen from the canonical distance-profile family.  GP3 preserves this
already constituted distinction under comparison.  This proves (ii)
implies (iii) in the stated domain.  Finally, every one-dimensional connected Lie
group is abelian, so a homomorphism into it kills \([r,s]\); the displayed
test distinction would be lost.
\end{proof}

\begin{corollary}[GP3 does not generate its own evidence]
\label{p2:cor:gp3-no-generation-alone}
GP3 cannot create a missing operation, test, topology, bound, geometric
origin, or rank-one obstruction.  If transformation separation is required by
the constituted role but the supplied pairing is not separating, the candidate
is not GP2-faithful and never enters the GP3 comparison.  If separation is not
part of the role grammar, the algebraic noncommutative branch remains available
but no physical B-test witness is inferred.  Within a transformation-separating
GP2-admissible pairing the separator exists by
Theorem~\ref{p2:thm:gp2-commutator-dichotomy}; this is a faithfulness
consequence, not test insertion.
\end{corollary}

\begin{proof}
Postulate~\ref{p2:post:GP2} requires every role distinction assigned to the
test doctrine to be represented faithfully before GP3 begins.  A missing
required separator therefore removes a candidate from the exact/faithful
class; it is not created by GP3 comparison.
Theorem~\ref{p2:thm:gp2-commutator-dichotomy} applies this definition to a
nontrivial commutator and yields exactly the stated alternative.
\end{proof}

\paragraph{Additional consequences of GP1 and GP2}

\begin{theorem}[Relation-exact and FRC-forced globalization]
\label{p2:thm:relation-minimal-globalization}
Under the presentation-theoretic semantics of GP2 fixed in
Postulate~\ref{p2:post:GP2}, let \(\mathfrak g\) be a candidate connected
local Lie type in the GP2-admissibility filter, and let \(\widetilde G\) be
its source-simply-connected integration.  The globalizations compared here
are integrations of this one common local path presentation, not unrelated
abstract groups.  Every connected globalization is
\[
 G\cong\widetilde G/K
\]
for a discrete central kernel \(K\).  Thus each nontrivial \(k\in K\)
imposes the global defining relation \(k=e\) on the common source
presentation.  If that equality is absent from the geometrically primary
local source and the previously frozen inventory, and
\(\widetilde G\) is compact and admits the faithful paired FRC model of
Lemma~\ref{p2:lem:relation-exact-endpoint-models}, GP2 relation exactness
forces \(K=\{e\}\).  The parent component of that source model is
\[
 \mathfrak P_{\widetilde G}
 :=(\widetilde G,d_g,\widetilde G,L),
 \qquad L_h(k)=hk,
\]
for any positive bi-invariant metric \(g\); its finite ontic component is the
faithful matrix action supplied by the same lemma.  Exactification therefore
replaces a quotient endpoint by a paired source-form model and never lets a
restored kernel act trivially on the old quotient carrier.

If \(\widetilde G\) is not an admissible compact parent, the closed-parent FRC
filter is applied before GP2.  Lemma~\ref{p2:lem:no-gratuitous-semisimple-kernel}
then removes every pure semisimple central kernel that is not required for
compactness.  When all compact connected relation-exact globalizations that
survive this two-stage filter are isomorphic as unmarked groups, their common
class is forced and no lattice or quotient representative is selected.  For
an abelian \(k\)-dimensional local type the semisimple filter is vacuous and
the common class is \(T^k\) by
Lemma~\ref{p2:lem:abelian-compact-globalization}; in particular the scalar
line globalizes to \(U(1)\).  If several nonisomorphic relation-exact compact
classes survive, the operator of Proposition~
\ref{p2:prop:admissible-globalization-operator} returns the unresolved family.

In particular, for the GP3-least noncommutative parent-role skeleton,
\(\mathfrak g\cong\mathfrak{su}(2)\) and \(\widetilde G=SU(2)\) is compact.
Hence
\[
        G\cong SU(2),\qquad\text{not}\qquad
        SO(3)=SU(2)/\{\pm1\}.
\]
The latter adds the ungenerated relation \(-1=1\).  In every faithful unitary
realization of the selected source form, center-oddness is equivalent to the
occurrence of an honest half-integral \(SU(2)\) summand.  It is therefore a
representation-theoretic diagnostic of the derived global form, not an
independent datum selecting it.
\end{theorem}

\begin{proof}
The classification of connected integrations of a finite-dimensional Lie
algebra gives \(G\cong\widetilde G/K\) with \(K\) discrete and central.
Quotienting is exactly the imposition of the relations \(k=e\), \(k\in K\).
GP2 admits such a relation only if prior geometry generated it or FRC requires
a compact paired realization unavailable on the source.  When
\(\widetilde G\) is already compact, Lemma~\ref{p2:lem:relation-exact-endpoint-models}
supplies that realization, so neither condition holds for nontrivial \(K\)
and relation exactness leaves the source integration unchanged.

When the source is not compact, closed compact-parent admissibility is a prior
FRC restriction rather than a discretionary GP2 permission.  GP2 then removes
pure semisimple kernel relations by
Lemma~\ref{p2:lem:no-gratuitous-semisimple-kernel}, after which Proposition~
\ref{p2:prop:admissible-globalization-operator} selects a class only when the
remaining relation-exact unmarked class is a singleton.  The abelian assertion
is exactly Lemma~\ref{p2:lem:abelian-compact-globalization}; its proof also
shows why the resulting group is unique without selecting a particular
lattice.

For \(\mathfrak{su}(2)\), the source integration and its center are
\(SU(2)\) and \(\{\pm1\}\).  Thus the only proper connected global quotient
is \(SO(3)\), obtained by imposing \(-1=e\).  The central character formula
says that on a spin-\(j\) irreducible, \(-1\) acts by \((-1)^{2j}\), proving
the final diagnostic statement.
\end{proof}

\begin{remark}[The \(S^1/S^3/SO(3)\) allocation as generated-defect fidelity]
\label{p2:rem:rank-one-defect-fidelity-example}
The primitive allocation already provides the canonical worked example of
Definitions~\ref{p2:def:intrinsic-generated-defect} and
\ref{p2:def:gp2-realization-defect}.  Let \(\delta_{\rm swap}\) denote the
typed intrinsic defect signature of the geometrically primary order-swap
diagram.  On the abelian branch \(\delta_{\rm swap}=0\); on the
order-sensitive branch \(\delta_{\rm swap}\ne0\).  A candidate that abelianizes
the latter branch therefore has
\[
 \mathfrak D_{\rm rel}\ne0,
\]
because it fails to preserve the generated nonzero distinction.  Conversely,
the source-global \(SU(2)\) realization preserves that order-swap signature.
Passing instead to
\[
 SO(3)=SU(2)/\{\pm1\}
\]
imposes the additional equality \(-1=1\), which is not generated by the local
source role, and hence produces an insertion/relation realization defect.
Thus GP2 first filters to candidates that represent every intrinsic defect
signature and every generated relation exactly.  GP3 then compares the
remaining zero-realization-defect complete candidates and supplies the least
primitive-role skeleton.  Vanishing GP2 realization defect by itself is not a
substitute for GP3 leastness.
\end{remark}

\begin{theorem}[Nontrivial connected exact role derives positive Lie dimension]
\label{p2:thm:gp1-gp2-positive-dimensionality}
Let \(G\) be the compact Lie exact-operation closure obtained from FRC at a
fixed role, region, and stage.  If it belongs to a candidate satisfying the
primitive assertion \(\mathsf E_R\), then
\[
        \dim G^0\ge1.
\]
Consequently every GP1/GP2-admissible realization with an independently
certified compact-Lie parent (for example a licensed full-parent faithful
finite operation port) has a nonzero Lie algebra.  This holds in particular
for the derived primitive endpoint \(U(1)\) or \(SU(2)\), whose complete
paired models are already constructed.  An arbitrary non-Lie compact
connected ambient candidate instead has a positive-dimensional finite Lie
\emph{image} by Peter--Weyl; the assertion does not attach a Lie algebra
to that ambient parent.  No separate (FD) or (LR) hypothesis is needed
on the certified Lie branch.
\end{theorem}

\begin{proof}
Theorem~\ref{p2:thm:geometric-existence-exactification} makes \(G\)
nontrivial and connected.  Hence \(G=G^0\).  A zero-dimensional connected Lie
group is trivial, contradicting the effective occurrence required by
\(\operatorname{Geo}_R\).  Therefore \(\dim G^0\ge1\).
For an ambient compact connected parent not known to be Lie, the
nontrivial finite image of
Lemma~\ref{p2:lem:peter-weyl-finite-image-reduction} is a connected
nontrivial compact Lie group, hence has positive Lie dimension; this is
an image statement, not a Lie theorem about the ambient parent.
\end{proof}

\begin{theorem}[Compact connected operational parent under FRC+GP2]
\label{p2:thm:gp1-gp2-compact-natural-alternatives}
Apply FRC to a faithful exact reversible candidate satisfying the
geometrically primary assertion \(\mathsf E_R\), with a generated
faithful affine finite-state/test action or an independent compact-Lie
certificate as specified in
Theorem~\ref{p2:thm:geometric-existence-exactification}.
Theorem~\ref{p2:thm:frc-automatic-compact-lie-closure}
gives a compact Lie operation group; exact realization of the constituted
local reapplication grammar restricts the primitive parent to its generated
identity component, and
Theorem~\ref{p2:thm:gp1-gp2-positive-dimensionality} gives positive dimension.
The resulting parent is therefore a nontrivial compact connected Lie group.
These principles impose no Cartan-rank bound and no simply-connectedness
condition.  Before irreducible/least GP3 comparison, higher tori and
higher-rank compact Lie groups remain possible.  The optional (FD)+(LR)
conditions concern the
regularity of a primitive parent comparison and the derivation of a
first-order displacement norm; they are not needed for Lie closure.
\end{theorem}

\begin{proof}
Compact Lie closure is the direct finite-carrier theorem.  Nontriviality and
connectedness follow from application of GP2 to \(\mathsf E_R\) by
Theorem~\ref{p2:thm:geometric-existence-exactification}, and the preceding
theorem gives positive dimension.  Neither FRC nor GP2 restricts maximal commuting
subgroups or simple rank.  Faithful finite unitary representations of
\(T^2\) and \(SU(3)\), paired with their compact regular parent presentations,
give explicit FRC+GP2 candidates against any stronger pre-GP3 conclusion.
They are exact to their own operation geometries but are not least in the
primitive endpoint-comparison domain.  If their extra directions are required
by \(\operatorname{Geo}_R\), they instead define a richer domain and cannot be
removed.
\end{proof}

\begin{remark}[Exact FRC+GP2 capacity-and-origin criterion]
\label{p2:rem:gp1-gp2-bounded-origin-criterion}
A proposed physical datum passes FRC+GP2 exactly when it belongs to the
constructor-relative generated host closure, admits a finite/bounded structured
realization in the appropriate slots, and is represented with exact
sort-faithfulness.  GP2 forbids insertion, forbids omission, and retains every
nonredundant canonical image generated by a constituted constructor.  A
registered boundary \(\mathsf P_{\rm gbd}\) is relevant only in a role whose
frozen grammar has actually constituted the corresponding generated-depth
scope; Proposition~\ref{p2:prop:gp1-gp2-bounded-generated-exactness} then
exactifies that richer scope.  A bounded parent comparison or operation family
without prior-generated origin is inserted structure; an origin statement
without a GP1-admissible exact/faithful realization gives at most a parent
construction or an FRC obstruction.
\end{remark}

\begin{remark}[Family-selection scope]
\label{p2:rem:gp1-gp2-family-selection-pointer}
Finite/bounded admissibility plus common GP2 origin does not select a member of
an otherwise exact/faithful family.  A member is fixed only when prior geometry
generates the distinction or a role-faithful GP3 comparison makes one class
least.  The formal corollary is deferred to
Paper~II-F, Section~\ref{p2:app:selector-limitations}.
\end{remark}

\paragraph{Why GP3 is load-bearing in primary role allocation}

For each candidate geometric action, GP1-TB gives the compact completion of
its comparison presentation, the finite-carrier theorem gives compact Lie
closure of its faithful reversible operations, and the canonical mathematical
metric profiles detect the action kernel, giving the
geometric effective quotient.  Their status as physical B-tests remains
conditional on metric-curry closure.  Those facts do not compare
\emph{different} effective candidates.  GP3 supplies exactly that missing
global comparison, applied after candidate-internal GP2 exactification and
before candidate-specific self-data are frozen.  This distinction is what
places \(T^2\) and \(SU(3)\) above the primitive endpoints without deleting
their internally exact operations or inserting (OR1).  If their additional
structure belongs to \(\operatorname{Geo}_R\), the endpoint comparison is
inadmissible.  GP3 must not be replaced by an irreducibility argument internal
to an already chosen group.

\paragraph{Additional consequences of GP2 and GP3}

\begin{proposition}[GP2 exact/faithful closure followed by GP3 leastness]
\label{p2:prop:gp2-gp3-origin-faithful-moduli}
GP2 excludes every ungenerated prior mark and requires each candidate's entire
generated content to be exact and faithful in its generated sorts.  GP3 orders
the resulting complete candidates by role-faithful structural comparisons
before candidate-specific data are frozen.  If a unique least
isomorphism class exists, GP3 retains it; if no unique least class exists,
distinct irreducible classes remain unresolved whether incomparable or
preorder-equivalent but non-isomorphic.  GP3 cannot authorize a datum that
failed GP2, and GP2 cannot use naturality to declare a nonfaithful realization
closed.  Once the least class is retained, its complete exact generated
presentation is frozen; theorem consequences remain available without
automatically becoming new presentation items.
\end{proposition}

\begin{proof}
The two filters concern different operations: GP2 completes and exactifies
each generated candidate, while GP3 orders the admissible candidates by
role-faithful comparison.  Reversing the order would allow an irreducible but
inserted object to survive; freezing candidate self-data first
would make every candidate trivially indispensable.
\end{proof}

\begin{corollary}[Conditional GP2+GP3 reciprocal test obligations]
\label{p2:cor:gp2-gp3-reciprocal-obligations}
Whenever the constituted role assigns transformation separation to its
generated test doctrine, every transformation distinction imposes the
corresponding separation obligation before GP3.  Every GP3 comparison must
preserve that supplied separation.  Neither
principle may insert a topology, reciprocal manifold, numerical normalization,
or new test to discharge that obligation.
\end{corollary}

\begin{proof}
Under the stated role assignment, GP2 host faithfulness requires two
distinguished transformations to induce different pullbacks on at least one
test in the supplied doctrine.  GP2 also permits only geometrically generated
tests and relations.  Thus each transformation distinction in that typed
pairing has a separation obligation, while adding any of the listed structures
solely to manufacture a separator would violate GP2\,(i).  Without the stated
assignment, the algebraic distinction remains but this corollary asserts no
physical test witness.
\end{proof}

\paragraph{Rank-one classification as a consistency corollary}

\begin{theorem}[Dependency-complete least-skeleton classification]
\label{p2:thm:dependency-complete-parent-allocation}
Apply the cumulative chronology of
Definition~\ref{p2:def:cumulative-role-allocation} to an unmarked primitive
role assertion \(\mathsf E_R\) in the endpoint-comparison domain of the master
theorem.  Generate
each candidate's full mathematical metric profiles canonically from its
comparison; if they are used as physical B-tests, assume metric-curry
closure.
Then
\[
 H_R\cong
 \begin{cases}
 U(1)\cong S^1,
 &\epsilon_R=\mathrm{ab},\\[2pt]
 SU(2)\cong S^3,
 &\epsilon_R=\mathrm{nc}.
 \end{cases}
\]
Equivalently, the resulting first branch has
\([H_R,H_R]=\{e\}\), while the second has
\([H_R,H_R]\ne\{e\}\).
Here \(\epsilon_R\) is classified after GP2 faithfully realizes the constituted
order-swap diagram, as in
Theorem~\ref{p2:thm:geometric-existence-exactification}; it is not an
independent primitive label.  Both branches have operational Cartan rank one:
\[
 [X,Y]=0\quad\Longrightarrow\quad X\wedge Y=0
 \qquad(X,Y\in\mathfrak h_R).
 \tag{OR1}
\]
Thus (OR1) is a theorem about the derived least parent-role skeletons, not an independent
selection hypothesis.  On the noncommutative branch \(SO(3)\) remains a
useful local-form comparison, but GP2 rejects its additional relation
\(-1=1\).  A center-odd faithful sector diagnoses the resulting source form and is not
supplied as another global-form bit.  The displayed groups are the unique least
primitive-role skeletons, and
Corollary~\ref{p2:cor:primitive-irreducible-endpoint-skeleton} shows that every
irreducible complete structured realization in that domain has the
corresponding skeleton.  Distinct unmarked lifts over a common skeleton remain
unresolved without further generated structure.
\end{theorem}

\begin{proof}
The branch allocation is
Theorem~\ref{p2:thm:master-compact-role-allocation}; the global-form assertion
is Theorem~\ref{p2:thm:relation-minimal-globalization}.  The Lie algebras
\(\mathfrak u(1)\) and \(\mathfrak{su}(2)\) each have Cartan rank one, so the
displayed commutator condition follows.  The central-character assertion is
Corollary~\ref{p2:cor:rank-center-diagnostics}.
\end{proof}

\begin{theorem}[Role-local torus obstruction]
\label{p2:thm:role-local-torus-obstruction}
Let \(H_R\) be one nontrivial compact connected primitive operation parent.
Then the following are equivalent:
\[
\begin{split}
&[X,Y]=0\Longrightarrow X\wedge Y=0
   \quad(X,Y\in\mathfrak h_R),\\
&\operatorname{rank}H_R=1,\\
&T^2\not\hookrightarrow H_R .
\end{split}
\]
This is a condition on one undecomposed role \(R\), not on a product of
distinct roles.  If \(R_1,\ldots,R_n\) are separately constituted roles, then
\[
 \operatorname{rank}\!\left(\prod_i H_{R_i}\right)
 =\sum_i\operatorname{rank}H_{R_i}.
\]
In particular, two scalar roles each isomorphic to \(U(1)\) legitimately
produce
\[
 U(1)_{R_1}\times U(1)_{R_2}\cong T^2,
\]
although neither primitive factor contains an internal \(T^2\).
\end{theorem}

\begin{proof}
In a compact Lie algebra every abelian subalgebra is contained in a maximal
abelian subalgebra, and all such subalgebras are conjugate of dimension equal
to the rank.  (The corresponding statement for commuting \emph{group}
elements is false: the Klein four-group in \(SO(3)\) is commuting but lies in
no maximal torus.)  Thus an independent commuting tangent pair exists exactly
when a maximal torus has dimension at least two.  Such a torus contains a
closed subgroup \(T^2\), and conversely an embedded \(T^2\) differentiates to
an independent commuting pair.  This proves the three-way equivalence.
Ranks add under direct products because maximal tori do.  The final statement
is the special case of two rank-one abelian factors.
\end{proof}

\subparagraph{Further compact-Lie consequences.}
Once the compact connected parent has been derived, its standard toral--simple
decomposition, maximal-torus structure, recurrence properties, invariant
measure, homogeneous/symmetric metric geometry, representation-theoretic
orthogonality, and elliptic spectral consequences are available without adding
new physical hypotheses.  They are consistency and import results rather than
premises of the parent-allocation theorem, so their full statements and proofs
are collected in Paper~II-F, Section~\ref{p2:app:compact-lie-consequences}.

\clearpage
\section{Global-cover spectral discriminators and their constitutional order}
\label{p2a:sec:v352-global-cover-tests}
The following round-geometry statements are independent mathematical
comparisons.  They are not substituted for the earlier GP2 constructor
or for the constitution of an actually occurring primitive role.

\begin{theorem}[Complete connected $A_1$ globalizations and first scalar test]
\label{p2a:thm:v352-a1-global-scalar}
For Lie algebra $\mathfrak{su}(2)$, the only connected Lie-group
globalizations, up to isomorphism, are $SU(2)$ and
$SO(3)=SU(2)/\{\pm I\}$.  With the locally isometric round metrics of
radius $R>0$, their first positive scalar Laplace eigenvalues,
complex eigenspace dimensions, and scale-free ratios are
\[
\begin{array}{c|c|c|c}
 &\lambda_1&\dim_{\mathbb C}\ker(-\Delta-\lambda_1)&
 \lambda_1/\operatorname{Scal}(g)\\\hline
 SU(2)\simeq S^3&3/R^2&4&1/2\\
 SO(3)\simeq\mathbb{RP}^3&8/R^2&9&4/3.
\end{array}
\]
Thus a \emph{previously licensed} first scalar harmonic comparison
with the left-hand signature excludes the projective quotient.
\end{theorem}
\begin{proof}
Every connected integration of a Lie algebra is a quotient of its
simply connected group by a discrete central subgroup.  The center of
$SU(2)$ is $\{\pm I\}$.  Scalar harmonics of degree $k$ on the round
$S^3$ have eigenvalue $k(k+2)/R^2$ and multiplicity $(k+1)^2$.
Under the antipodal deck map they acquire sign $(-1)^k$; therefore
only even $k$ descend to $\mathbb{RP}^3$.  The first nonconstant
degrees are $k=1$ and $k=2$ respectively.  Both metrics have
$\operatorname{Scal}=6/R^2$, proving the table.
\end{proof}

\begin{remark}[Chronological and semantic firewall for global selection]
\label{p2a:rem:v352-a1-chronology}
A test assembled \emph{after} choosing the $SU(2)$ harmonic or
spinorial atlas cannot retroactively justify that original choice.
The source predicates, test and occurrence must be licensed before the
cover selection for the displayed comparison to become a GP2 selector.
Likewise the bare assertion of a nontrivial complete geometry alone
has compact connected abelian models (e.g. a circle), so it cannot
by itself establish the occurrence of a noncommuting primitive role.
This observation neither changes G0 nor GP1--GP3, nor negates the
stronger constructor-relative theorems established elsewhere in this
paper; it isolates their load-bearing domain hypotheses.
\end{remark}

\section{Endpoint scope and exact theorem ownership}
\label{paper2a:sec:organized-ownership}

\subsection{Scope and theorem ownership}
\label{p2:sec:ownership-boundary}

The preceding proof chain is the sole constitutional derivation in this paper.
This section states the load-bearing scope limits that must remain visible in
the main text and records theorem ownership across Paper~II.  Detailed selector
no-go arguments and legacy finite-level arithmetic are deferred to
Paper~II-F, Section~\ref{p2:app:selector-limitations}.

\begin{theorem}[Boundedness does not select a realization label]
\label{p2:thm:no-cutoff-from-finiteness}
Let \(\{\mathfrak R_\lambda\}_{\lambda\in\Lambda}\) be inequivalent
realizations of the same constituted parent and operations, each satisfying
the relevant FRC admissibility and boundedness conditions.  If no earlier natural construction
or generated operation singles out one member, GP1 selects no value of
\(\lambda\).  In particular it selects no cutoff, truncation integer,
representation label, least or greatest family member, or preferred
finite-dimensional subcarrier.
\end{theorem}

\begin{proof}[Proof deferred]
See Paper~II-F, Section~\ref{p2:app:selector-limitations}.
\end{proof}

\begin{corollary}[No primitive-stage numerical carrier scale]
\label{p2:cor:no-numerical-scale-from-gp1-gp2}
The nonabelian least-skeleton classification fixes every
automorphism-invariant quadratic
comparison on \(\mathfrak{su}(2)\) only up to a positive common scale.  In the
absence of an earlier dimensionful datum, the primitive operation-role
allocation of GP1--GP3 selects no numerical value of
\[
 \Rmom,\qquad \Punit=\hbar/\Rmom,\qquad P_{\rm ceil},
 \qquad\text{or}\qquad P_{\rm ceil}/\Punit .
\]
This does not preclude a constructor-relative GP2 host closure from fixing a
finite carrier support directly, nor a later role-faithful GP3 comparison or
non-circular dynamical theorem from fixing a dimensionless level in a richer
role.  If such a carrier is selected, its sharp
actual maximum \(\Pmax=m_p\Punit\) is a derived readout; common rescaling still
leaves the absolute scale free.
\end{corollary}

\begin{proof}
On the compact simple algebra \(\mathfrak{su}(2)\), invariant symmetric
bilinear forms are multiples of the Killing form.  Positivity leaves one
positive scaling orbit.  All structural hypotheses above are preserved by a
common positive rescaling, so none chooses a point of that orbit.
Theorem~\ref{p2:thm:no-cutoff-from-finiteness} excludes converting the
primitive-stage unfixed auxiliary ratio into a selected truncation label.  A
later dimensionless carrier label may instead follow from a generated host
closure or a role-faithful structural comparison, but not from this scale
ratio.
\end{proof}

\begin{remark}[Legacy tower-capacity arithmetic]
\label{p2:rem:legacy-module-capacity-pointer}
For roles that separately invoke the legacy \(\mathsf P_{\rm sli}\) or an
explicit admissibility ceiling, the exact dimension and floor-function
capacity formulas for the declared Paper-I-A family remain valid.  They are not
part of the constitutional momentum-seed proof and are recorded in
Paper~II-F, Section~\ref{p2:app:selector-limitations}.  In particular, the derived label
\((\mathrm{LVL}_p^1)\) is distinct from an independently asserted admissibility
ceiling \(m_{\rm adm}=1\).
\end{remark}

\begin{lemma}[No canonical scalar branch inside the transfer group]
\label{p2:lem:no-su2-u1-character}
For \(H=SU(2)\) or \(SO(3)\), every continuous homomorphism
\(\chi:H\to U(1)\) is trivial.  Moreover, no nontrivial proper connected
subgroup of \(H\) is characteristic.  Hence a nontrivial scalar succession
law is neither a character of the nonabelian transfer group nor a canonically
selected circle inside it.
\end{lemma}

\begin{proof}
Both groups are connected and perfect, so every homomorphism to an abelian
group factors through the trivial abelianization.  A connected characteristic
subgroup would differentiate to a characteristic ideal of the simple algebra
\(\mathfrak{su}(2)\), hence is trivial or all of \(H\).  Although maximal
tori exist, they are mutually conjugate; choosing one requires an additional
axis or torus selector.
\end{proof}

\begin{corollary}[Circle homomorphisms into the nonabelian rank-one parent]
\label{p2:cor:v193-u1-su2-cocharacters}
Let \(A_{\rm v}\cong SU(2)\) be the unmarked nonabelian rank-one parent.
Every continuous homomorphism
\[
 \lambda:U(1)\longrightarrow A_{\rm v}
\]
is conjugate to
\[
 \lambda_n(z)=\operatorname{diag}(z^n,z^{-n}),\qquad n\in\mathbb Z.
\]
Conjugation by the Weyl normalizer identifies \(n\) and \(-n\), so the
conjugacy classes are indexed by \(|n|\).  For \(n\ne0\), the image is a
maximal torus and
\[
 \ker\lambda_n\cong C_{|n|}.
\]
Hence \(\lambda_n\) is faithful if and only if \(|n|=1\).  In particular,
an abstract scalar circle does not select a circle subgroup of the unmarked
nonabelian parent: a faithful inclusion requires a maximal-torus occurrence
together with a primitive choice.  The already retained Hopf inclusion of
Lemma~\ref{p2:lem:chosen-hopf-vector-parent-inclusion} is such a faithful
primitive inclusion on the separately constituted Hopf occurrence.
The circle classification is Paper~1B,
Proposition~\ref{p1b:result:circle-cocharacters}~\cite{Kugelmann2026Paper1B},
applied to the rank-one parent already established here; the signed degree
records the retained primitive orientation before Weyl conjugacy.
\end{corollary}

\begin{remark}[Classification versus Hopf/succession provenance]
\label{p2:rem:v193-u1-su2-provenance-firewall}
Corollary~\ref{p2:cor:v193-u1-su2-cocharacters} classifies the mathematical
circle homomorphisms available once the group types are known.  It does not
create a new Hopf/succession bridge.  The physical provenance remains the
separately constituted \(\mathsf{HS}_t\) arrow \(\kappa_t\) of
Definition~\ref{p2:def:hopf-succession-provenance}, together with the already
chosen Hopf occurrence and its typed inclusion \(\iota_H\).
\end{remark}

\begin{theorem}[Operational allocation of the least compact parent-role skeletons]
\label{p2:thm:finite-realization-momentum-selection}
For every non-FRC-obstructed unmarked primitive role assertion in the
endpoint-comparison domain of
Theorem~\ref{p2:thm:master-compact-role-allocation}, after the cumulative
GP1--GP3 chronology has been applied and with its canonical effective
metric-profile action,
\[
 \epsilon_R=\mathrm{ab}
 \quad\Longrightarrow\quad H_R\cong U(1),
\]
whereas
\[
 \epsilon_R=\mathrm{nc}
 \quad\Longrightarrow\quad
 H_R\cong SU(2)\cong S^3_A(\Rmom).
\]
The resulting parent operation groups satisfy \([H_R,H_R]=\{e\}\) and
\([H_R,H_R]\ne\{e\}\), respectively.
The branch status \(\epsilon_R\) is classified after GP2 faithfully realizes
the constituted order-swap diagram and is not an independent primitive label.
An endogenous physical commutator witness on the second
branch additionally requires a transformation-separating B-test doctrine.
OR1 is derived for either least skeleton.  On the second
line \(SO(3)\) is excluded because it adds \(-1=1\), while a center-odd
half-integral sector is a later diagnostic of the selected source form.
These statements concern the common least parent-role skeletons.  Every
GP3-irreducible complete structured realization in the endpoint-comparison
domain has the corresponding skeleton
by Theorem~\ref{p2:thm:primitive-endpoint-lift-reachability}; a unique complete
unmarked structured lift is not asserted.
\end{theorem}

\begin{proof}
This is Theorem~\ref{p2:thm:dependency-complete-parent-allocation} with the
transfer notation of
Definition~\ref{p2:def:momentum-branch-specialization}.
\end{proof}

\begin{definition}[Paper-II derived momentum-parent allocation]
\label{p2:def:momentum-branch-specialization}
Let \(A_{\rm v}\) denote the geometrically constituted primitive
directed-transfer role.  GP1--GP2 exactification of its intrinsic order-swap
diagram derives
\(\epsilon_{A_{\rm v}}=\mathrm{nc}\); no independent branch bit or candidate
group is supplied.  Theorems~\ref{p2:thm:master-compact-role-allocation} and
\ref{p2:thm:relation-minimal-globalization} derive the common least kinematic
parent-role skeleton and its operation geometry
\[
        M_p:=H_{A_{\rm v}},\qquad
        L_h(k)=h\cdot k,\qquad
        M_p\cong SU(2)\cong S^3_A(\Rmom).
\]
The symbol \(M_p\) denotes this regular self-parent of operation
values, supplied by
Theorem~\ref{p2:thm:regular-homogeneity-closure}; when a later role explicitly constitutes the paired tower, its stagewise
physical realizations may be represented by finite objects
\(\mathfrak F_{p,m}\); the initial momentum host is instead selected directly
and only afterward classified as the \(m=1\) tower member.  It neither proves
\(X_{\mathfrak M}\cong H_{A_{\rm v}}\) for a generic candidate action nor
identifies the primitive comparison metric with the geodesic distance of
the derived negative-Killing metric.
Here \(h\cdot k\) is abstract group multiplication.  Under the radius-
\(\Rmom\) quaternion realization of
Paper~I-A~\cite{Kugelmann2026a}, write \(q=\Rmom u\) and \(r=\Rmom v\) with
\(u,v\in S^3(1)\).  Paper~I-A's scaled-unit-quaternion proposition supplies
the transported group product and regular action
\[
        q\odot r:=\Rmom^{-1}qr,
        \qquad
        L_q(r)=q\odot r,
\]
whose identity is \(\Rmom\in\HH\) and whose inverse is
\(q^{-1}_{\odot}=\bar q\).  Thus \(S^3_A(\Rmom)\) means the
radius-\(\Rmom\) round sphere equipped with this transported operation law;
unrescaled quaternion multiplication of two norm-\(\Rmom\) elements is not
being used.
The radius \(\Rmom>0\) is not fixed.  A faithful representation may test the
source center; an honest half-integral sector is center-odd and thereby
confirms the selected \(SU(2)\) form.  It is not a premise of this
allocation.  The abelian class \(H_R\cong U(1)\cong S^1_A\) is separately
derived for a scalar role and is exported without time or energy
interpretation.  No downstream reciprocal-role construction enters the
Paper-II-A proof.
\end{definition}

\begin{definition}[Hopf-provenanced scalar-succession occurrence]
\label{p2:def:hopf-succession-provenance}
Let \(U(1)_H\), with positive primitive Reeb lattice generator
\(\tau_H^+\), be the oriented Paper-I-A Hopf-flow occurrence, and let
\(H_{\rm succ}\cong U(1)\) denote the separately allocated abelian endpoint
parent for the scalar-succession role.  The role-specific Geometric Primacy
predicate for scalar succession is the strong predicate
\[
 \operatorname{Geo}^{H+}_{\rm succ}(\mathfrak M_t,\kappa_t)
\]
which, in addition to the primitive scalar-role diagrams, constitutes an
actual host-natural faithful continuous group-homomorphism provenance arrow
\[
 \boxed{
 \kappa_t:(U(1)_H)_{\rm flow}
          \longrightarrow (H_{\rm succ})_{\rm succ}.}
 \tag{\ensuremath{\mathsf{HS}_t}}
\]
The target is the autonomous scalar-succession operation role.  It is not a
subgroup of the nonabelian transfer parent and is not an object occurrence on
the momentum carrier.  No target orientation is supplied independently: the
target receives its orientation from this actual provenance arrow.
\end{definition}

\begin{theorem}[Transported scalar-succession orientation]
\label{p2:thm:transported-succession-orientation}
Assume the Hopf-provenanced scalar-succession occurrence of
Definition~\ref{p2:def:hopf-succession-provenance}.  Then \(\kappa_t\) is an
isomorphism of the two compact rank-one operation parents and induces an
isomorphism of period lattices.  The scalar primitive is therefore
\[
 \boxed{e_t:=(\kappa_t)_*\tau_H^+.}
\]
It is primitive and its positive direction is uniquely inherited from the
Paper-I-A Hopf orientation.  GP2 retains the provenance arrow and forbids a
post-constitutional replacement \(\kappa_t\mapsto I\circ\kappa_t\) by scalar
circle inversion.  No physical duration unit or Hamiltonian is fixed.
\end{theorem}

\begin{proof}
A faithful continuous homomorphism between compact connected circles is an
isomorphism.  Its induced map on period lattices sends primitive generators
to primitive generators.  Since the target had no independently marked
orientation, the image of \(\tau_H^+\) defines its positive primitive.
Composing the retained provenance arrow with inversion changes the structured
role occurrence and is not an automorphism of the same frozen datum.
\end{proof}

\begin{lemma}[Typed inclusion of the chosen Hopf circle in the vector parent]
\label{p2:lem:chosen-hopf-vector-parent-inclusion}
Let \(A_{\rm v}\cong SU(2)\cong S^3_A(\Rmom)\) be the vector operation parent
of Definition~\ref{p2:def:momentum-branch-specialization}, represented by the
same scaled-unit-quaternion model used for the oriented Paper-I-A Hopf-flow
occurrence \(U(1)_H\).  Then the already chosen Hopf circle has a faithful
smooth subgroup inclusion
\[
 \iota_H:U(1)_H\hookrightarrow A_{\rm v},
 \qquad
 \exp_H(sX_H)\longmapsto \exp_{A_{\rm v}}(sX_H).
\]
Because \(U(1)_H\) is abelian, the same underlying map is also a group
homomorphism
\[
 \boxed{\iota_H^{\rm op}:U(1)_H\hookrightarrow A_{\rm v}^{\rm op}.}
\]
Both inclusions preserve the positive primitive
\(\tau_H^+=2\pi X_H\).  They introduce no additional maximal-torus axis,
handedness choice, character of the unmarked parent, or physical occurrence:
the axis and orientation are exactly those already contained in the frozen
Paper-I-A Hopf-flow occurrence used by \(\mathsf{HS}_t\).
\end{lemma}

\begin{proof}
The vector-parent specialization transports scaled quaternion multiplication
from the Paper-I-A sphere to \(A_{\rm v}\).  In that same model the chosen
Hopf flow is the one-parameter subgroup generated by \(X_H\), so its inclusion
is smooth and injective.  If \(u,v\in U(1)_H\), then, writing
\(\star_{\rm op}\) for multiplication in \(A_{\rm v}^{\rm op}\),
\[
 \iota_H^{\rm op}(uv)=uv=vu
 =\iota_H^{\rm op}(u)\star_{\rm op}\iota_H^{\rm op}(v).
\]
Thus the opposite inclusion is a homomorphism.  Preservation of the primitive
is immediate from the displayed exponential, and no new selector is used.
\end{proof}

\begin{remark}[Minimality and typing of the provenance occurrence]
\label{p2:rem:hopf-succession-provenance-scope}
Without \(\mathsf{HS}_t\), simultaneous inversion
\(e_t\mapsto-e_t\), \(e_E\mapsto-e_E\) preserves the abstract scalar circle
and its reciprocal pairing.  Hence the provenance occurrence, not a second
sign postulate, is minimal.  This does not contradict
Lemma~\ref{p2:lem:no-su2-u1-character}: the arrow starts at the already chosen
Paper-I-A Hopf-flow occurrence and ends at an autonomous scalar role; it is not
a character of, or a canonical circle extracted from, the unmarked
nonabelian transfer parent.  A full countermodel is recorded in
Paper~II-F, Section~\ref{p2:app:selector-limitations}.
\end{remark}


\subsubsection{Theorem ownership across Paper II}
\label{p2:subsec:theorem-ownership}
Paper~II-A owns FRC, the exact/host-faithful GP2 architecture, the
irreducible/least GP3 architecture, primitive parent-skeleton allocation,
the G0 actualization/readout/calibration semantics,
the role-specific Hopf-to-succession provenance occurrence, transported
scalar orientation, and typed source-Hopf inclusion in the vector parent and
its opposite,
relation-exact source globalization, the canonical Killing geometry, the
constructor-relative generated-closure rule, typed finite generated-image
inheritance, the GP2 reciprocal-role objectification/relation-free-composition
schema of Theorem~\ref{p2:thm:gp2-generated-reciprocal-role-closure}, the
paired-realization contract, and the typed status/occurrence/bridge registry.
Paper~II-A also owns the typed rooted structural-realization equivalence, the
no-duplicate structural-actuality corollary, and typed generated-identity
deduplication.  They are consequences of the existing G0/GP1--GP3
architecture and add no constitutional principle or constructor.
Paper~II-B owns the finite equivariant word tower as mathematical
classification and, under the revised target-sort inheritance rule, the
unconditional momentum host theorem
\(\mathcal G^{\rm car}_{\Sigma_p}\cong\mathcal H_1
\Rightarrow m_p=1\), with the full-endpoint carrier comparison separately
conditional on Theorem~\ref{p2b:thm:momentum-carrier-rigidity}.
II-C owns generic finite
coordinate/operator kinematics and its level-one specialization, the finite
position-frontier history type and least-resolution statistic, retained-history
monotonicity, and the protected even \(\Cl(0,3)\) invariant on the
nondegenerate protected stratum.  II-D owns parent spectral grammar, exact
finite spectral statements, and quantified recovery interfaces; the merged
II-E component owns its typed geometric and protected-block exports.  Later
nonconstant position-frontier growth, recovery, raw coordinate completion,
absolute normalization, and dynamical branch conditions remain conditional
unless separately generated or discharged.  A later reciprocal-spectrum paper owns the concrete
proof of the \((\mathrm{BR}_{\rm ind})\) and \((\mathrm{IND}_{xp})\)
antecedent statuses for its generated spectrum; once proved, GP2 supplies the
generated reciprocal role and its relation-free product rather than treating
those labels as independent physical assumptions.  At the primitive-role level the exported parent groups
are least-skeleton results in the stated primitive domain, and candidatewise
whole-candidate endpoint comparison makes those skeletons common to every GP3-irreducible complete
structured realization.  Distinct unmarked lifts over a common skeleton remain
unresolved.  No later result is a premise of the parent-allocation theorem.


\subsubsection{Stagewise scalar characters and carrier-role scope}
\label{p2a:sec:v180-stagewise-scalar-character-support}

\begin{corollary}[Stagewise finite scalar character support]
\label{p2a:cor:stagewise-scalar-character-support}
Let a realized scalar role carry the finite faithful unitary circle
representation supplied by GP1--GP2.  Then for a finite set
$F_s\subset\mathbb Z$,
\[
 \rho_s(e^{i\theta})=\sum_{n\in F_s}e^{in\theta}P_{n,s},\qquad
 Q_s=\sum_{n\in F_s}nP_{n,s},
\]
with mutually orthogonal nonzero $P_{n,s}$ summing to the identity.  Hence
$N_s:=\max_{n\in F_s}|n|<\infty$ and
\[
 \prod_{n\in F_s}(Q_s-nI)=0.
\]
Faithfulness is equivalent to
$\gcd\{|n|:n\in F_s,\ n\ne0\}=1$.
\end{corollary}
\begin{proof}
Finite-dimensional unitary representations of $U(1)$ split into integer
characters.  The polynomial follows from the spectral decomposition, and the
kernel is trivial exactly when the nonzero character weights have greatest
common divisor one.
\end{proof}

\begin{remark}[Scalar skeleton, fixed momentum carrier, and comparison band]
\label{p2a:rem:scalar-skeleton-carrier-comparison}
The primitive group skeleton does not identify all of its faithful
representations.  In particular the already derived compact momentum carrier
has $m_p=1$, whereas the symbol $m_x$ labels harmonic resolution of the
Born-derived reciprocal-position window
$\mathcal P_x(m_x)=\bigoplus_{k=0}^{m_x}E_k^x$ on the smooth reciprocal
position role.  It does not label a distinct microscopic position carrier.
The later position-side comparison level $m_{\rm cmp}=2$ is selected by an
exact generated-algebra/quartic-evaluation comparison; it is neither a
replacement for $m_p=1$ nor a new GP3 rule that minimizes arbitrary scalar
comparison bands.  Higher occurring reciprocal-position bands remain retained
in their own role.
\end{remark}

\subsection{Scoped application certificates}
\label{p2a:sec:v2210-lit-application-certificates}
The literature-assisted results of Papers~I-B and II-B--II-E narrow several
application calculations without changing G0, GP1--GP3 or the primitive
role domain.  They provide explicit defect and equivalence tests rather
than selecting a new physical member.

\begin{corollary}[Claim-relative zero defect is the transferred obligation]
\label{p2a:cor:v2210-lit-visible-defect}
For a licensed application at a frozen presentation, let $D$ be the actual
finite comparison defect and let $\{\ell_j\}$ be its predeclared consumed
source/effect tests.  If the claim asserts only these readouts, exact
transfer is equivalent to $\ell_j(D)=0$ for every consumed test.
If the tests separate the actual difference space, this is equivalent to
$D=0$ on that space.  Without that separation, zero observed tests cannot
be promoted to full operator, source-jet or carrier equivalence.
\end{corollary}
\begin{proof}
Subtracting the two supplied readouts gives $\ell_j(D)$.  Separation is
exactly injectivity of their evaluation map on the occurring difference
space.  Paper~II-C, Theorem~\ref{p2c:thm:v2210-lit-effect-frame}, gives a
positive finite Gram/rank test for that injectivity.  This uses the
existing claim-relative G0 semantics and inserts no full-tomography axiom.
\end{proof}
For a product word, Paper~II-B,
Theorem~\ref{p2b:thm:v2210-lit-word-defect}, computes $D$ from its omitted
images.  For a stronger reciprocal comparison, Paper~II-D,
Theorem~\ref{p2d:thm:v2210-lit-positive-recovery}, separates reconstruction
from multiplication error.  These are certificates on the actual licensed
maps; they do not generate those maps merely because an auxiliary atlas
contains them.

\begin{remark}[Structural outputs and surviving physical applications]
\label{p2a:rem:v2210-lit-retained-frontier}
The native exceptional packing/subgroup implications, first-symbol
reciprocal bridge, coordinate fingerprints, protected Clifford block,
rooted trimer/family maps and unstable $M_8$ classifications stay closed
on their current domains.  An arbitrary energy-minimization rule does not
select the Hurwitz support, an auxiliary Fourier/quantum-metric limit does
not change $m_p=1$, and a positive Born-connection theorem does not replace
the physical Lorentzian solder.  The new finite boundary response is
generated by the complete supplied Hermitian block; its physical boundary
ownership remains the action/domain calculation.  Determinant scale-germ
and finite-part uses retain the distinct readout tests in
Theorem~\ref{p2d:thm:v2210-lit-zeta-scale}.  Actual apparatus/history tokens
and unresolved source members remain G0 occurrences; none is a missing
mathematical consequence manufactured by these transfers.
\end{remark}


\section{Typed finite operations and generated descendant closure}
\label{paper2a:sec:organized-applications}

\subsection{Typed graded-filtered parent-operation descent}
\label{p2a:sec:v189-graded-filtered-parent-operations}
The parent-saturation and PCAD rules retain not only an endpoint object but
also the typed constructor path by which that endpoint is generated.  For a
licensed word $w$ rooted at a constituted parent $P$, write the path signature
\[
 \Sigma_P(w)=
 \bigl(P;\tau_{\rm src}\to\tau_{\rm tgt};\rho;
 b_B,b_{\rm slot},b_Z;g,h;d\bigr),
\]
where entries are recorded only on sorts on which they are defined.  Here
$\rho$ is constructor-word order, $b_B$ Born parity, $b_{\rm slot}$ input-slot
exchange parity, $b_Z$ the spin/central parity, $g$ and $h$ the Clifford and
harmonic filtrations, and $d$ the derived-commutator depth.  Undefined entries
are not replaced by artificial numerical values.

\begin{theorem}[Path-relative graded-filtered PCAD]
\label{p2a:thm:v189-graded-filtered-pcad}
Let $w_1,w_2$ be composable licensed parent words.  On every common typed
domain on which the displayed metadata are defined,
\[
 \rho(w_2\circ w_1)=\rho(w_1)+\rho(w_2),
 \qquad
 b(w_2\circ w_1)=b(w_1)+b(w_2)\pmod2
\]
for each multiplicative parity $b$.  Metric contraction may lower displayed
Clifford or harmonic degree by an even integer, so those degrees are
filtrations with an exact parity residue rather than unrestricted additive
gradings.  Parent Correspondence and Action Descent transports the complete
path signature to every licensed descendant relation.

If $\mathsf D^0(H)=H$ and
$\mathsf D^{d+1}(H)=[\mathsf D^d(H),\mathsf D^d(H)]$, then the canonical
nested-commutator provenance tree at derived depth $d$ has $2^d$ primitive
leaves.  This is a statement about that constructor path; it is not a lower
bound on the length of every unrelated word representing the same endpoint.
\end{theorem}

\begin{proof}
Composition concatenates constructor words and therefore adds path order.
Multiplicative $\mathbb Z_2$ characters multiply and hence add modulo two.
The Clifford relation $uv+vu=2g(u,v)$ and the spherical relation
$|x|^2=R^2$ contract degree by two, proving preservation of parity without
preservation of raw degree.  PCAD retains the originating parent word by
construction.  Finally a commutator of two depth-$d$ trees has twice their
number of leaves, giving $2^{d+1}$ inductively.
\end{proof}

\begin{remark}[Endpoint-degree firewall]
\label{p2a:rem:v189-endpoint-degree-firewall}
The same endpoint may admit several constructor histories.  For example
$-1\in Q_8$ occurs both as the commutator $[i,j]$ and as the square $i^2$.
Thus degree belongs to the triple
$(\text{root},\text{constructor path},\text{endpoint})$, not to the endpoint
alone.  No new GP is introduced by recording this provenance.
\end{remark}


\subsubsection{Conditional constitution-atlas status}
\label{p2a:subsec:v236-constitution-atlas}

\begin{definition}[Atlas status labels]
\label{p2a:def:v236-atlas-status}
The symbols
\[
 \mathrm{CND},\qquad \mathrm{GEN},\qquad \mathrm{DL},\qquad
 \mathrm{ROOT},\qquad \mathrm{OBS}
\]
are bookkeeping labels for the chronology already enforced by G0 and
GP1--GP3.  They mean, respectively: a conditional realization branch; a
role generated/co-realized from an occurring ancestor; a branch or member
whose occurrence is discharged only by a later independently constituted
result; an actually asserted physical member with no upstream discharge on
the current claim; and an empty or incompatible branch.  These labels add no
new GP and never turn a later discriminator into an earlier GP3 selection.
\end{definition}

\begin{definition}[Typed finite-output atlas]
\label{p2a:def:v236-finite-output-types}
Finite realizations are typed before comparison.  We distinguish
\[
 \mathrm{GeoSupp}(\Gamma),\qquad
 \mathrm{GrpSupp}(\Gamma),\qquad
 \mathrm{Tors}(\Gamma),\qquad
 \mathrm{Abs}(\Gamma),\qquad
 \mathrm{Cmp}_m.
\]
The first two are, respectively, a finite geometric support and an embedded
finite subgroup in the constituted compact parent; \(\mathrm{Tors}(\Gamma)\)
is an unbased regular torsor, \(\mathrm{Abs}(\Gamma)\) an abstract group type,
and \(\mathrm{Cmp}_m\) an exact finite presentation of the retained
comparison band.  No equality of cardinalities or abstract group types
identifies these sorts.  In particular an exact presentation need not be a
physical point support, and an abstract modular realization need not be an
embedded momentum support without a typed owner/action bridge.
\end{definition}

\begin{remark}[Prospective later discharge]
\label{p2a:rem:v236-prospective-discharge}
If GP3 retains an unresolved atlas family, a later independently constituted
physical theorem may prove that one previously classified branch occurs.  Its
status is then \(\mathrm{DL}\) for the later claim.  The earlier theorem
remains an unresolved conditional classification and is not rewritten as
upstream uniqueness.  Content constant on every retained branch may descend
without selecting a member.
\end{remark}



\subsection{Geometric discharge of comparison premises}
\label{p2a:sec:v233-premise-closure}

\begin{theorem}[Conditional full-quadratic support discharge]
\label{p2a:thm:v233-quadratic-support-discharge}
Suppose an already occurring role is generated by an admitted finite
constructor word on occurring roots, and its frozen requirements are:
\begin{enumerate}[label=(\roman*),leftmargin=2em]
\item exact evaluation of every inner product on the full scalar
comparison band \(\mathcal P_2=E_0\oplus E_1\oplus E_2\);
\item a finite group-closed point support in the constituted \(SU(2)\)
parent;
\item comparison of its unselected candidate supports by subgroup
restriction, faithful to every other retained requirement of this
same role.
\end{enumerate}
Then its group candidates are exactly \(2T,2O,2I\), and the unique
least unmarked type is \([2T]\). The generated support output
co-realizes without a second occurrence certificate.
Uniform weights, antipodality, a separate axis-free predicate and
exceptional classification are not additional premises on this branch.
No labelled vertex, frame or trimer is selected.
\end{theorem}
\begin{proof}
Paper~I-C, Theorem~\ref{p1b:thm:v233-subgroup-sampler-classification},
proves the complete weighted candidate range and its subgroup leastness.
In particular all candidates contain \(-1\), and no cyclic or binary
dihedral candidate passes the quadratic tests. Requirement (iii)
makes that mathematical order an admissible comparison for the same
physical role. The admitted word supplies its ancestry and output
typing, so Theorem~\ref{p2:thm:generated-role-corealization} applies.
The output is an orbit; selecting an embedded representative adds a mark.
\end{proof}

This specializes the orbit-level route of
Corollary~\ref{p2a:cor:two-conditional-2t-discharge-routes}.
It does not derive its three role requirements from generic point
resolution. Even an antipodal, identity-anchored exact quadratic
sampler need not contain a Hurwitz subset, by
Proposition~\ref{p1b:prop:v233-hopf72}. If fourth powers of quadratic
fields belong to this same role, \(2I\), rather than \(2T\), is required.
Restriction preserves an exact band image but can discard frozen point
tests or a full permutation action
(Proposition~\ref{p1b:prop:v233-sampling-restriction}).


\begin{theorem}[Finite operation/comparison exceptional frontier]
\label{p2a:thm:v236-qop-exceptional-frontier}
Fix \(m\ge2\).  Let \(\mathsf{QOp}_m\) be an already occurring role whose
frozen physical content is: a finite subset \(Z\subset SU(2)\) with
\(1\in Z\) and \(ZZ\subseteq Z\); exact realization of the retained complete
scalar band \(\mathcal P_m\); and only the group-word structure together with
the exact sampled \(\mathcal P_m\) image and its transported operators.  No
individual node tests, full point algebra, marked incidence data or full
left-regular permutation action are frozen.

Then \(Z\) is automatically a finite subgroup.  For exact subgroup candidates,
Paper~I-C's sampling-restriction map is unitary on the retained band image and
preserves every frozen group word and transported band operator, so subgroup
restriction is faithful to this same role.  Consequently
\[
 \boxed{
 \begin{array}{c|c}
 m & \operatorname{Irr}\mathfrak M_{\mathsf{QOp}_m}\\ \hline
 2 & [2T]\\
 3 & \{[2O],[2I]\}\\
 4,5 & [2I]\\
 m\ge6 & \varnothing
 \end{array}}
\]
where the \(m=3\) pair is retained because \(2O\) and \(2I\) are
incomparable.  These are conditional realization statements; they do not
assert that \(\mathsf{QOp}_m\) occurs.
\end{theorem}
\begin{proof}
For \(z\in Z\), finiteness of the sequence \(1,z,z^2,\ldots\) gives
\(z^r=z^s\) for some \(r<s\), hence \(z^{s-r}=1\) in the ambient group and
\(z^{-1}=z^{s-r-1}\in Z\).  Thus \(Z\le SU(2)\).  Paper~I-C,
Theorem~\ref{p1b:thm:v233-subgroup-sampler-classification}, gives the sharp
complete-band ranges for \(2T,2O,2I\), and
Proposition~\ref{p1b:prop:v233-sampling-restriction} proves exact unitary
restriction on the sampled band image.  Since the frozen role contains no
extra node-specific data, subgroup inclusion preserves all remaining group
words.  GP3 therefore selects the unique least class only in the rows where
one exists; the empty row is the exact obstruction.
\end{proof}

\begin{corollary}[Native quadratic construction for \(\mathsf{QOp}_2\)]
\label{p2a:cor:v236-native-qop2}
Suppose the occurring finite operation role retains \(E_1\) and licenses the
pointwise metric--Gram constructor
\[
 (f,h)\longmapsto R^2\langle\nabla f,\nabla h\rangle
\]
as a function-valued physical comparison.  Then
\[
 q_aq_b=\delta_{ab}-R^2\langle\nabla q_a,\nabla q_b\rangle,
 \qquad
 \operatorname{Sym}^2E_1=E_0\oplus E_2,
\]
so its retained scalar comparison image is the full
\(\mathcal P_2=E_0\oplus E_1\oplus E_2\).  If its finite support is
identity-rooted and multiplication-closed, the preceding theorem yields the
unmarked \([2T]\) core without separately assuming inverse closure,
antipodality, exceptionality or a \(2T\) occurrence.
\end{corollary}

\begin{lemma}[Recentring an occurring optimal \(24\)-cell]
\label{p2a:lem:v236-optimal24-recentering}
Let an occurring same-owner regular \(24\)-cell support in
\(S^3\simeq SU(2)\) be written \(Z=aHb^{-1}\), where \(H\cong2T\).  For any
\(z=ahb^{-1}\in Z\),
\[
 z^{-1}Z=bH b^{-1}.
\]
Hence the unmarked embedded subgroup orbit is independent of which support
point is used to recenter.  The geometric \(24\)-cell support therefore
generates \(\mathrm{GrpSupp}(2T)\) at orbit level once the ambient
multiplication/inversion is the already constituted owner operation; no
physical distinguished vertex is selected.
\end{lemma}
\begin{proof}
Since \(h\in H\), \(h^{-1}H=H\), and
\(z^{-1}Z=bh^{-1}a^{-1}aHb^{-1}=bHb^{-1}\).  Quaternionic inversion also
preserves \(H\), so orientation reversal does not change the unmarked output.
\end{proof}
\begin{theorem}[Observable-sensitive exceptional subgroup frontier, with physical occurrence firewall]
\label{p2a:thm:v268-observable-exceptional-frontier}
On the already constituted round momentum parent, let a \emph{specified,
already occurring} finite operation/comparison role require a finite
nonempty multiplication-closed support $\Gamma\subset SU(2)$ and exact
uniform quadrature of a declared family of comparison observables.
If the frozen role permits the subgroup-restriction comparison of
Theorem~\ref{p2a:thm:v236-qop-exceptional-frontier}, the following
mathematical candidate classifications hold:
\begin{enumerate}[label=(\roman*),leftmargin=2em]
\item \textbf{Complete quadratic pairing or quadratic endpoint readout.}
Exact $L^2$ inner products on $\mathcal P_2$, exact Dirichlet pairings
on $\mathcal P_2$, or the complete quadratic-symbol endpoint readout
$\Phi_{a,\zeta}|_{\mathcal P_2}$ each admit precisely
$[2T],[2O],[2I]$.  In this role the unique subgroup-least
unmarked class is $\boxed{[2T]}$, of order $24$.
\item \textbf{All cubic products or all compressed quadratic multipliers.}
Exact integrals of all $fgh$ with $f,g,h\in\mathcal P_2$, or
exact realization of every $P_2M_fP_2$ with $f\in\mathcal P_2$,
admit exactly $\boxed{\{[2O],[2I]\}}$.
These are \emph{two incomparable subgroup-minimal classes}; in
particular GP3 subgroup leastness does not select $[2O]$.
The group $2O$ is the unique \emph{smallest-cardinality} candidate
($48<120$) if that is the separately licensed comparison order.
If the same role also requires a nontrivial homomorphism
$\Gamma\to\{\pm1\}$, then the only candidate is $[2O]$, provided
this character is part of the frozen observable and is preserved by
admissible comparison maps.
\item \textbf{All quartic scalar field actions on $\mathcal P_2$.}
Exact integration of $\phi^4$ for \emph{every real}
$\phi\in\mathcal P_2$ (and hence of its mode derivatives) admits
precisely $\boxed{[2I]}$, of order $120$.
\end{enumerate}
These statements classify \emph{possible finite-group implementations}
of already specified observables.  Neither the canonical quadratic
operator submodule of
Theorem~\ref{p2c:thm:v268-native-jordan-quadratic-copy}, nor the
smooth pointwise metric--Gram identity, nor a continuous quartic
integral, nor the intrinsic first-moment POVM by itself asserts
occurrence of a physical finite subgroup or an action consumer of
all the listed polynomial tests.  In particular no unconditional
$2T$, $2O$, or $2I$ physical-support occurrence is concluded.
\end{theorem}
\begin{proof}
A nonempty finite subset of a group closed under multiplication
contains the identity and inverses by finite repetition of powers;
it is therefore a subgroup.  Theorem~\ref{p1b:thm:v233-subgroup-sampler-classification}
classifies all finite-subgroup designs, with exact strengths $5,7,11$
for $2T,2O,2I$.  Its proof excludes all cyclic and binary dihedral
subgroups from $\mathcal P_2$ isometry.  A direct exclusion for the
Dirichlet-only condition is as follows: each such subgroup preserves
a decomposition $\mathbb R^4=A\oplus B$ into two orthogonal planes,
and its points lie in $S(A)\cup S(B)$ (cyclic groups use only one
circle).  The harmonic
$h(q)=\|P_Aq\|^2-\|P_Bq\|^2\in E_2$
obeys $\nabla_{S^3}h=0$ at every such point, while
\[
 \int_{S^3}|\nabla_{S^3}h|^2\,d\sigma
 =8\int_{S^3}h^2\,d\sigma=\frac83>0.
\]
Thus none of the nonexceptional subgroups meets this requirement,
regardless of point weights.  Sufficiency for all three exceptional
groups follows from degree-four exactness.  For the full endpoint
quadratic readout, Theorem~\ref{p2c:thm:v268-endpoint-quadratic-iff-design4}
proves the converse to the degree-four sufficiency statement.
The $2T$ inclusions into $2O$ and $2I$ give the first least class.

For (ii), products of three members of $\mathcal P_2$ span
$\mathcal P_6$.  The matrix entries of
$P_2M_fP_2$ are exactly $\langle g,fh\rangle$ with
$f,g,h\in\mathcal P_2$, so \emph{all} such compressed multiplication
tests are exact if and only if the rule is a six-design.
The witness
\[
 \tau_{2T}\big(q_1^2(q_1^2-q_2^2)^2\big)
 -\sigma\big(q_1^2(q_1^2-q_2^2)^2\big)
 =\frac1{12}-\frac1{16}=\frac1{48}
\]
excludes $2T$.  The remaining classes $2O$ and $2I$ are
incomparable since $48\nmid120$; neither is subgroup-least.
Their abelianizations give
$\operatorname{Hom}(2O,\{\pm1\})\cong C_2$
and $\operatorname{Hom}(2I,\{\pm1\})=0$.

For (iii), real polarization reconstructs all four-linear products
of $\mathcal P_2$ fields from their fourth powers, and those products
span $\mathcal P_8$.  Thus exactness of \emph{all} quartic scalar
field actions is equivalent to degree-eight cubature.  The degree-eight
moment for $u=q_1^2-q_2^2$ gives the sharp checks
\[
 \tau_{2T}(u^4)=\frac16,\qquad
 \tau_{2O}(u^4)=\frac5{48},\qquad
 \int_{S^3}u^4\,d\sigma=\tau_{2I}(u^4)=\frac3{40}.
\]
The finite-subgroup classification leaves only $2I$.
The final occurrence statement follows from the G0 action-consumer and
preaction no-self-witness rules and the representative no-go in
Corollary~\ref{p2a:cor:exceptional-support-status-after-no-go}.
\end{proof}

\begin{remark}[No conflation of group occurrence, endpoint realization, or orbit means]
\label{p2a:rem:v268-problem3-ledger}
The $14$-dimensional $\mathcal P_2$ bandwidth and the canonical native
operator copy are derived from the fixed smooth carrier and its
intrinsic first-symbol tuple; these are unconditional geometric/operator
facts.  The $[2T]$ statement is an exact conditional consequence of an
occurring role consuming the \emph{complete} quadratic test family
and allowing role-faithful subgroup restriction.  Eight-fibre
$Q_8$ isometry and first moments remain exact weaker operations.
The descendant weights $(1/8,1/7,1/9)$ require the separately
constituted octet, septet, and pair transitive permutation actions of
Corollary~\ref{p2:cor:paper1-octet-descendant-weight-closure}; they
are not free-standing conclusions from the abstract class $[2T]$.
\end{remark}






\begin{theorem}[Finite positive convolution repeatability derives a subgroup]
\label{p2a:thm:v273-positive-repeatable-convolution-group}
On $G=SU(2)$, let
$\mu=\sum_{s\in S}w_s\delta_s$ have finite nonempty support,
$w_s>0$, and $\sum_sw_s=1$. Define the left averaging operation
$A_\mu f(x)=\int f(g^{-1}x)\,d\mu(g)$.
Then the following are equivalent:
\begin{enumerate}[label=(\roman*),leftmargin=2em]
\item $A_\mu^2=A_\mu$ on all $C(G)$;
\item $\mu*\mu=\mu$ as Borel probability measures;
\item $S\le G$ is a finite subgroup and
      $\mu=|S|^{-1}\sum_{s\in S}\delta_s$.
\end{enumerate}
For an \emph{already occurring, licensed} finite-output convolution
role, full-domain repeatability therefore discharges the separate
multiplication-closure and weight-uniformity conditions.
It neither creates the finite-output role from the smooth parent nor
identifies a POVM with a convolution action.
\end{theorem}
\begin{proof}
Convolution composition and evaluation at the identity against all
continuous functions yield (i)$\Leftrightarrow$(ii).
Positivity gives $\operatorname{supp}(\mu*\mu)=SS$;
thus (ii) implies $SS=S$, and a nonempty finite multiplicatively
closed subset of a group contains the identity and each inverse by
finite repetition of powers. Let $a=\max_{s\in S}w_s=w_{s_0}$.
Then idempotence gives
\[
 a=\sum_{g\in S}w_gw_{g^{-1}s_0}
       \le a\sum_{g\in S}w_g=a.
\]
All $w_g$ are strictly positive, so equality forces
$w_{g^{-1}s_0}=a$ for every $g$, hence uniformity on $S$.
Uniform subgroup averages are idempotent, proving the converse.
\end{proof}

\begin{proposition}[Native finite-band repeatability cannot certify group closure]
\label{p2a:prop:v273-finite-band-repeatability-countermodel}
Let $H=2T$ and $g\in SU(2)\setminus H$.
The translated positive measure $\mu_{gH}$ supported on $gH$
is not supported on a subgroup, yet
\[
 A_{\mu_{gH}}|_{\mathcal P_5}
 =A_{\rm Haar}|_{\mathcal P_5},\qquad
 A_{\mu_{gH}}^2|_{\mathcal P_5}=A_{\mu_{gH}}|_{\mathcal P_5}.
\]
Thus no inference from quadratic sampling, even with five-band
idempotence, to an occurring $\operatorname{GrpSupp}(2T)$ is valid.
\end{proposition}
\begin{proof}
Each translated $2T$ coset remains a spherical five-design;
its convolution Fourier coefficient vanishes on all nonconstant
$E_k$ for $1\le k\le5$ and fixes constants. The support
$gH$ omits the identity, so cannot be a subgroup.
\end{proof}

\begin{remark}[Two-axis finite-exceptional comparison and physical occurrence]
\label{p2a:rem:v273-superposition-consumer-firewall}
The exact local finite frames and coset mixtures of
Theorems~\ref{p1b:thm:v273-cosetwise-resolution-groupoid} and
\ref{p1b:thm:v273-triple-icosahedral-nineteen-design}
are mathematical presentations of the constituted smooth scalar
comparison owner, not new physical nodes. Cubature exactness and
Reynolds invariant-sector projection are distinct observable roles:
a positive nineteen-design need not reproduce a rank-thirteen
$2I$ projector on $E_{12}$. Co-realization preserves the existing
physical carrier and its licensed native observables without
requiring finite points, but GP2 cannot turn a comparison frame,
positive mixture, or an unmarked projector family into a physically
occurring finite multiplication action. A frozen finite-output
constructor and its actual observable consumer are required for
that stronger target sort. In particular $2O$ and $2I$ remain
incomparable for the degree-six exactness role; cardinality
minimization is not silently substituted for subgroup leastness.
\end{remark}


\begin{theorem}[First leaked harmonic exactly detects exceptional-coset repeatability]
\label{p2a:thm:v274-first-leaked-harmonic-coset}
Let $\Gamma\in\{2T,2O,2I\}$ be embedded in $SU(2)$, let
$\mu_{g\Gamma}$ be the uniform positive probability on a left coset,
and put
\[
 (k_\Gamma,|\Gamma|)=(6,24),(8,48),(12,120)
 \quad\hbox{respectively}.
\]
For $\rho_k=\operatorname{Sym}^k(\mathbb C^2)$ define
$T_{g,k}=\widehat\mu_{g\Gamma}(k)=\rho_k(g)P_\Gamma^{(k)}$.
At $k=k_\Gamma$ the fixed subspace of $\rho_k(\Gamma)$ is a
single line $\mathbb Cv$ and
\[
 \boxed{T_{g,k}^2-T_{g,k}
    =\bigl(\langle v,\rho_k(g)v\rangle-1\bigr)T_{g,k},}
 \qquad\|v\|=1.
\]
Therefore
\[
 \boxed{T_{g,k}^2=T_{g,k}\iff g\in\Gamma.}
\]
On each $1\le j<k_\Gamma$ the Fourier coefficient of the coset
measure vanishes regardless of $g$, so $k_\Gamma$ is the first
single-harmonic degree that distinguishes every nontrivial coset
translation by its repeatability defect. This test presupposes an
already specified \emph{uniform coset} and does not construct a
finite-output physical role.
\end{theorem}
\begin{proof}
The character sum $|\Gamma|^{-1}\sum_{\gamma\in\Gamma}
U_k(\operatorname{Re}\gamma)$ gives the first nonconstant fixed
ranks: one at $k=6,8,12$ for $2T,2O,2I$ respectively and zero at
all smaller positive degrees. These are the first degrees beyond
their exact design strengths $5,7,11$, respectively.
Writing $P_\Gamma^{(k)}=|v\rangle\langle v|$ yields the displayed
rank-one product identity. Its vanishing is equivalent to
$\langle v,\rho_k(g)v\rangle=1$, which by unitarity is equivalent
to $\rho_k(g)v=v$. The stabilizer of $v$ contains $\Gamma$ and is
a closed subgroup. It cannot be $SU(2)$ because $k>0$, and no
proper positive-dimensional closed subgroup of $SU(2)$ contains any
of these nonabelian binary polyhedral groups. Thus it is finite.
For $2O$ and $2I$ the finite-subgroup classification has no larger
finite supergroup; for $2T$ the only larger finite possibilities are
$2O$ and $2I$, neither of which has an invariant in $\rho_6$.
Hence the stabilizer is exactly $\Gamma$, proving the equivalence.
\end{proof}

\begin{theorem}[Two polynomial tests suffice for any finite positive convolution]
\label{p2a:thm:v274-two-scalar-repeatability-certificate}
Let $\mu=\sum_{s\in S}w_s\delta_s$ be a positive probability
measure with $S\subset SU(2)$ finite and nonempty, $w_s>0$, and
$N=|S|$. In quaternionic $S^3$ coordinates define
\[
 d_s(q)=1-\langle q,s\rangle,\quad
 p_S(q)=\prod_{s\in S}d_s(q),\quad
 \ell_s(q)=\prod_{t\in S\setminus\{s\}}
             \frac{d_t(q)}{d_t(s)},\quad
 h_\mu(q)=\sum_{s\in S}w_s\ell_s(q).
\]
Here $p_S\ge0$ has degree $N$, $h_\mu\ge0$ has degree at most
$N-1$, $p_S^{-1}(0)=S$, and $h_\mu(s)=w_s$ on $S$.
The following two \emph{scalar} conditions are necessary and
sufficient for $\mu$ to be the uniform probability on a finite
subgroup:
\[
 \boxed{\int p_S\,d(\mu*\mu)=0,\qquad
 \int h_\mu\,d(\mu*\mu)=\int h_\mu\,d\mu.}
\]
For the left-averaging convention
$A_\mu f(x)=\int f(g^{-1}x)\,d\mu(g)$, these two displayed
moment conditions are equivalently identity-point repeatability of the
\emph{inverse-pulled} tests $p_S\circ\mathrm{inv}$ and
$h_\mu\circ\mathrm{inv}$, since
\[
 A_\mu^2(f\circ\mathrm{inv})(e)=\int f\,d(\mu*\mu),\qquad
 A_\mu(f\circ\mathrm{inv})(e)=\int f\,d\mu.
\]
The same identity-point claim for the unmodified functions $p_S$ and
$h_\mu$ is not valid in general.  Inversion preserves their polynomial
degrees on quaternionic $S^3$; no idempotence assumption on all of
$C(SU(2))$ is needed.
Once the first test holds, let $w_{\min}=\min_Sw_s$ and
$u_s=1/N$. The nonnegative second-test defect satisfies
\[
 \boxed{\delta_h:=\int h_\mu\,d\mu-
       \int h_\mu\,d(\mu*\mu)
 \ \ge Nw_{\min}\sum_{s\in S}(w_s-u_s)^2.}
\]
Thus the same certificate controls nonuniform weights quantitatively.
For the first test, if $SS\setminus S\ne\varnothing$ and
$c_S=\min_{q\in SS\setminus S}p_S(q)>0$, then
\[
 \Pr_{(s,t)\sim\mu\otimes\mu}[st\notin S]
 \le\frac{\int p_S\,d(\mu*\mu)}{c_S}.
\]
These are finite mathematical/action tests only when their polynomial
observables are licensed on the same occurring averaging owner.
\end{theorem}
\begin{proof}
For each $s\in S$, $d_s(q)\ge0$ on the unit sphere and vanishes
only at $q=s$, whence the support and interpolation statements.
The first test reads
$0=\sum_{s,t\in S}w_sw_t p_S(st)$, a sum of nonnegative terms
with all weights positive. It is therefore equivalent to $SS\subseteq S$.
A nonempty finite multiplicatively closed subset of a group is a
subgroup, so $S$ is then a finite subgroup and both $\mu$ and
$\mu*\mu$ are supported on it.
Writing the weight vector as $w\in\ell^2(S)$, interpolation gives
\[
 \int h_\mu\,d\mu=\|w\|_2^2,\qquad
 \int h_\mu\,d(\mu*\mu)=\langle w,w*w\rangle.
\]
Convolution by any probability measure is a convex combination of
unitary translations on $\ell^2(S)$, so $\|w*w\|_2\le\|w\|_2$.
Cauchy--Schwarz shows that equality of the two displayed integrals
forces $w*w=w$. The positive idempotent-measure theorem
\ref{p2a:thm:v273-positive-repeatable-convolution-group} then forces
uniform Haar weights; its converse proves necessity of both tests.

For stability put $x=w-u$ and $\alpha=Nw_{\min}$. If $\alpha<1$,
write $w=\alpha u+(1-\alpha)v$, where $v$ is another probability
vector. Convolution by $u$ kills zero-mean vectors, and convolution
by $v$ is an $\ell^2$ contraction, hence
$\|w*x\|_2\le(1-\alpha)\|x\|_2$.
Since $w*w=u+w*x$, it follows that
$\delta_h=\langle x,x-w*x\rangle\ge\alpha\|x\|_2^2$.
For $\alpha=1$ the weights are already uniform. The first-test
leakage bound follows by retaining just the products outside $S$,
where each term $p_S(st)$ is at least $c_S$.
\end{proof}

\begin{corollary}[Exact support-observable selector on an occurring quadratic role]
\label{p2a:cor:v268-observable-sixth-moment-2t}
On the \emph{already occurring} finite, identity-rooted,
multiplication-closed $\mathsf{QOp}_2$ role of
Theorem~\ref{p2a:thm:v236-qop-exceptional-frontier}, let
$\mathcal A_6^{\rm supp}$ be the canonical unweighted support
invariant of Paper~I-C,
Theorem~\ref{p1b:thm:v268-sixth-moment-discriminator}.  Then
\[
 \mathcal A_6^{\rm supp}>0
 \quad\Longleftrightarrow\quad
 [Z]=[2T],\qquad
 \mathcal A_6^{\rm supp}=0
 \quad\Longleftrightarrow\quad
 [Z]\in\{[2O],[2I]\}.
\]
Thus an actually consumed positive sixth-order support readout would
select the $2T$ \emph{type} without using a leastness decision, a
chosen quaternion axis, or an embedded member mark.  No such readout,
finite-role occurrence, or embedded support is generated by this
corollary.  The antecedent must be established on the same physical
action/measurement owner; its weighted effects cannot silently be
replaced by the uniform-support measure.
\end{corollary}



\begin{remark}[Native constructor and presentation boundary]
\label{p2a:rem:v233-native-constructor-boundary}
The momentum seed of Definition~\ref{p2:def:momentum-seed-frontier}
retains its scalar, spinor, gradient and section-bracket output sorts.
It does not include unrestricted scalar-product or tensor generation.
Mathematically,
\[
 \operatorname{Sym}^2 E_1=E_0\oplus E_2,\qquad
 R^2\langle\nabla q_a,\nabla q_b\rangle=\delta_{ab}-q_aq_b .
\]
The second formula generates quadratic scalar functions as physical
tests only if a position-test role licenses these pointwise metric
comparisons and retains their function-valued outputs.
For any output sort absent from the roots and from the targets of
all admitted constructors, induction on finite terms proves that no
term has that sort. This is a statement about the current grammar;
a later licensed constructor can change it.
The group-word restriction theorem
\ref{p2a:thm:v188-group-word-restriction} operates on an already
generated subgroup; it supplies neither an initial finite group nor
a normalizer constructor.

By contrast, an operator already constituted on the retained carrier
needs no new physical point support merely to be written in sampled
coordinates: Corollary~\ref{p1:cor:v233-eight-fibre-sampling} gives
an exact geometric presentation. A claim using support incidence,
point labels or a performed instrument still consumes that role.
The alternative frame-normalizer and three-torsion formulas in
Proposition~\ref{p1b:prop:v233-alternative-2t-constructors} identify
their outputs only on their declared constructor domains.
\end{remark}

\begin{corollary}[Canonical loop coefficients without a preparation premise]
\label{p2a:cor:v233-loop-coefficients}
For the already licensed canonical duality closure on a finite Hilbert
space \(H\),
\[
 \operatorname{ev}\circ(1\otimes T)\circ\operatorname{coev}
       =\operatorname{Tr}_H T .
\]
If \(T=\bigoplus_\lambda t_\lambda I_{H_\lambda}\), its coefficient
is \(\sum_\lambda d_\lambda t_\lambda\), where
\(d_\lambda=\dim H_\lambda\); normalizing the same closure by its
value at the identity gives weights \(d_\lambda/\dim H\).
These coefficients require no prepared-state choice or assumption
of invariance under an enlarged unitary group.
\end{corollary}
\begin{proof}
The canonical coevaluation is the basis-independent identity tensor
\(\sum_i e^i\otimes e_i\); evaluation after \(1\otimes T\)
gives \(\sum_i e^i(Te_i)\). This is
Lemma~\ref{p2:lem:canonical-finite-duality-trace}, specialized to
the block decomposition. A retained insertion \(B\) instead gives
\(\operatorname{Tr}(BT)\) and is a different word.
The result does not choose a general invariant preparation.
\end{proof}




\subsection{Generated co-consumption of complete constructor outputs}
\label{p2a:sec:v3101-generated-coconsumption}

The existing generated-constituent, generated-role, no-idle and joint-owner
theorems already imply a stronger reusable statement for a constructor whose
complete output has been frozen before evaluation.  The following theorem
names that consequence; it is not a fourth generating principle.

\begin{theorem}[Generated co-consumption of a complete constructor output]
\label{p2a:thm:v3101-generated-coconsumption}
Let $P$ be an occurring rooted input and let a result-independent licensed
constructor have frozen typed output
\[
 \mathfrak C(P)=Q=(q_1,\ldots,q_N),
\]
where the $q_i$ are the nonredundant required constituents of the declared
output.  Then the occurrence of $P$ co-realizes the complete output $Q$ and
every $q_i$ in its native generated sort.  Any later licensed constructor
whose declared input is the complete $Q$ retains that complete typed input jointly:
no $q_i$ may be independently deleted, rescaled, rebound, or subjected to a
second occurrence condition while claiming the same generated output.

The conclusion preserves all frozen coefficients, incidences,
normalizations and target sorts carried by $\mathfrak C$.  It does not create
a tuple constructor that was not licensed, combine independently generated
objects merely because they share a host, select a member of an unresolved
family, or promote mathematical availability to physical consumption.
\end{theorem}

\begin{proof}
Apply Theorem~\ref{p2:thm:generated-constituent-corealization} to the frozen
finite presentation of $Q$, together with generated-role co-realization and
the no-idle licensed-constructor corollary.  Deleting, rescaling or rebinding a
nonredundant $q_i$ changes the generated presentation and violates GP2 exact
faithfulness.  A later constructor declared on the complete $Q$ therefore
inherits that complete generated input.  This is a typed-availability
claim, not a proof that every downstream readout depends essentially on
each constituent value; Proposition~\ref{p2a:prop:v281-complete-input-factorization}
separates those assertions.  The final paragraph is the existing
typed-cross-relation firewall: GP2 preserves licensed composition but does not
insert a new composition or selected member.
\end{proof}

\begin{corollary}[Common readouts of a co-consumed unresolved family]
\label{p2a:cor:v3101-coconsumed-family-readout}
If $Q$ is a generated unresolved family and a later licensed readout is
constant on every surviving member, that value is common generated content.
If the readout is member-dependent, the complete image family is retained;
co-consumption supplies no hidden section or selector.
\end{corollary}
\begin{remark}[Constructor origin, missing certificate and consumer]
\label{p2a:rem:constructor-origin-and-consumer}
A generated constituent is located at its earliest retained output in the
actual certified word, even if a later paper first names its use.
A claim records separately that origin, its first unproved input,
incidence/domain/comparison arrow, and its downstream consumer.
This is a constructor-relative partial order, not a distance inferred from
paper order, a hierarchy rung or the completed \(M_8\) host.
Theorems~\ref{p2a:thm:rooted-generated-image-corealization}
and~\ref{p2a:thm:v183-generated-joint-owner-saturation} co-realize a complete
licensed parent and its required constituents; no second occurrence or
normalization is added at the consumer.
A certified whole-family constructor may instead consume the retained
nonempty complete family, by
Corollary~\ref{p2a:cor:v186-generated-constrained-owner-family}, and retain its
joint image.  A common value needs no member selection; a varying image
remains a family.  A required new physical incidence, member probability,
target-sort map or analytic continuation is recorded at the arrow that
actually consumes it.  Hosting alone does not license that arrow.
\end{remark}

\begin{proposition}[Complete generated input does not imply essential readout dependence]
\label{p2a:prop:v281-complete-input-factorization}
Let \(Q\subseteq Q_1\times\cdots\times Q_N\) be the surviving
complete output family of a licensed frozen constructor, let
\(F:Q\to Y\) be a later licensed readout, and for
\(S\subseteq\{1,\ldots,N\}\) write \(\pi_S:Q\to\pi_S(Q)\)
for coordinate restriction.  There exists a unique map
\(\widetilde F:\pi_S(Q)\to Y\) such that
\(F=\widetilde F\circ\pi_S\) if and only if
\[
 \forall q,q'\in Q:\quad
 \pi_S(q)=\pi_S(q')\ \Longrightarrow\ F(q)=F(q').
\]
In particular, the downstream readout is \emph{value-dependent} on a
coordinate \(i\) relative to the others precisely when there exist
\(q,q'\in Q\) agreeing away from \(i\) but with
\(F(q)\neq F(q')\).  If those comparison pairs do not exist, the
port may be indispensable to the \emph{typed generated presentation}
without being informationally indispensable to \(F\).

Consequently GP2 co-realization of the complete input and its obligatory
transport under a licensed \(Q\)-typed constructor do not imply that
\emph{every} constituent is mathematically necessary to calculate each
later readout.  Nor does factorization through \(\pi_S\) itself license a
new physical deletion/compression constructor.  Proof-dependency arrows
must distinguish a required \emph{typed input} from an \emph{essential
value premise}.
\end{proposition}
\begin{proof}
If \(F=\widetilde F\circ\pi_S\), equality of projections forces
equality of outputs.  Conversely, if \(F\) is constant on fibers of
\(\pi_S\), define \(\widetilde F(z)=F(q)\) for any
\(q\in Q\) satisfying \(\pi_S(q)=z\).  The condition makes this
well-defined and uniqueness follows from surjectivity of
\(\pi_S:Q\to\pi_S(Q)\).  The single-coordinate criterion is the
negation of the same fiberwise condition for
\(S=\{1,\ldots,N\}\setminus\{i\}\).

For sharpness, take \(Q=\{0,1\}^2\) with both coordinate sorts fixed as
required output ports, and \(F(q_1,q_2)=q_1\).  All complete tuples
remain co-realized, while the second coordinate does not influence the
readout.  If \(Q\) has cross-coordinate constraints, several distinct
minimal sufficient coordinate subsets may exist; GP2 does not select one
as a licensed replacement presentation.  These are mathematical
factorizations, not physical constructor occurrences.
\end{proof}



\subsection{Bounded required readouts of a complete realization}
\label{p2a:sec:bounded-generated-readouts}

\begin{theorem}[Bounded whole-family readouts without a second occurrence]
\label{p2a:thm:bounded-required-readout-family}
Let \(P_{\mathcal R}\) be one complete occurring rooted datum on a declared
stage or region \(\mathcal R\). Let a result-independent licensed constructor,
with its hypotheses certified on that datum, have required typed output
containing the real readout families \(f_1,\ldots,f_k\) over the whole claimed
\(\mathcal R\). Retain their frames, units, source and parameter marks.
Then those required families are co-realized from \(P_{\mathcal R}\), and
\[
 \boxed{\text{for each }j\text{ there is }M_{\mathcal R,j}<\infty
 \text{ such that }|f_j(s)|\leq M_{\mathcal R,j}
 \quad(s\in\mathcal R).}
\]
The descendant families need no independent occurrence token. A restriction
of the parent cannot delete a required family while retaining the same
complete claimed output.
\end{theorem}
\begin{proof}
The no-idle licensed-constructor corollary, generated-constituent and
generated-role co-realization, and
Theorem~\ref{p2a:thm:v183-generated-joint-owner-saturation} retain the complete
required output in its declared sort. GP2 forbids an omission or retyping of
one of those constituents. GP1, Postulate~\ref{p2:post:GP1}, makes the image
of each constituted real scalar/parameter/rate family totally bounded;
in \(\mathbb R\) it is bounded. Apply this to each \(f_j\).
\end{proof}

\begin{corollary}[Positive compression bound from an owned inverse readout]
\label{p2a:cor:bounded-inverse-readout-compression}
On the theorem's domain, suppose a positive coordinate/readout \(y(s)\)
and a required nonnegative family \(f(s)\) obey the proved same-owner relation
\[
 f(s)\geq C\,y(s)^{-p},\qquad C>0,\quad p>0,
\]
with the units and positive coefficient retained. For any positive finite
GP1 bound \(M_{\mathcal R}\) on that constituted \(f\) family,
\[
 \boxed{y(s)\geq(C/M_{\mathcal R})^{1/p}>0.}
\]
\end{corollary}
\begin{proof}
Combine \(C y^{-p}\leq f\leq M_{\mathcal R}\) and invert the positive
inequality. The relation between \(f\) and \(y\), rather than a common host
or numerical coincidence, licenses this inference.
\end{proof}

The constants are family- and stage-dependent; the theorem asserts neither
attainment of a supremum nor one number common to all realizations. It also
does not construct a future region. Its metric/action consumers in
Paper~VI-A derive the density, spectrum and curvature relations before using
this corollary. A nonempty matched continuation still consumes its existing
solution word and analytic domain; the empty-family alternative of
Theorem~\ref{p2a:thm:generated-exact-solution-moduli} remains intact.





\section{Application inputs and owner-resolved action interfaces}
\label{p2a:sec:clean-application-inputs}
\subsection{Generated application inputs and explicit residual obligations}
\label{p2a:sec:application-obligation-record}
A later consumer of a Paper-II export uses a licensed map from its complete
action, state, process or history parent. The generated constituent and
co-consumption results already proved here retain its output in the proper
sort. If the map and its hypotheses are certified on that parent, its input
is a descendant, not a second occurrence or structural selector. If the
output varies across an admitted family, the licensed image family is
retained; a scalar readout can close by being common on that family.

For an application, the explicit record states (i) the target and domain;
(ii) its required input; (iii) the existing parent and licensed generating
map, with a source theorem; (iv) the part already generated and its retained
marks; (v) the first remaining predicate, with its kind and certificate;
and (vi) the consumer activated by that predicate and any later dependent
steps. A predicate is classified as a construction or identification proof,
physical input, domain or limit proof, evaluation, or metrology. Optional
branches first declare their consumed target. These are bookkeeping types,
not additional constitutive principles.

For example, an occurring sharp spectral frontier generates the finite
position history by Paper~III-F,
Corollary~\ref{p3f:cor:sharp-frontier-generates-position-history}; its
physical frontier word, scale and any claimed threshold reach remain
visible. The native fresh cyclic root state and descendants instead use
Paper~VI-A, Theorem~\ref{p8:thm:v2114-fresh-root-neutral-qroot}, and
Paper~VI-B, Corollary~\ref{p9:cor:v187-k16-state-inheritance}, without a
second K16 preparation. Neither example inserts an ungenerated physical
state or supplies the domain of every later interacting or limiting claim.

On a specified consumer domain, ``first remaining'' means the first
uncertified consumed arrow after the already licensed input chain.
Alternative sufficient routes are alternatives; shared certificates are
reused; a new experimental input is not a missing universal law. A broad
undeclared target has an activation question rather than a located proof
gap.

An application coverage contract is not automatically a
native construction target.  In particular, an \texttt{ACTIVATE\_CLAIM},
\texttt{PHYSICAL\_INPUT}, or \texttt{METROLOGY} entry is an application
interface at that consumer, not evidence for a missing independent physical
mechanism.  If an unconditional prediction is intended, follow the recorded
generator upstream until the first uncertified generated arrow is reached;
if a conditional application is intended, supply the named instance without
reopening the native constructor.

Accordingly, a claim of application closure must identify the constituted
geometric role, the licensed complete action/state/process parent, the already
generated prefix, the first unproved arrow, and a finite pass/fail certificate
for that arrow.  A synthetic model, generic receiver, or downstream mechanism
name can diagnose a route but does not advance the frontier unless it evaluates
that actual arrow.  This is an execution rule for the existing Geometric
Primacy hierarchy, not a new constitutive principle.

The geometric target includes the retained source incidence, action and
state response, not only the momentum carrier.  In the electroweak example,
the occurring rooted odd solder generates its gauge stabilizer by
Paper~V-A, Theorem~\ref{p5:thm:v377-occurring-solder-em-stabilizer};
the positive kinetic Gram retains the same kernel independently of scalar
stationarity by Paper~III-F,
Corollary~\ref{p3f:cor:v111-gauge-gram-no-stationarity}.  A scalar equilibrium,
its Hessian or an actual symmetry-breaking history is required only for a
claim consuming that separate output, as scoped by Paper~V-C,
Corollary~\ref{p7v2:cor:ewsb-status}.  Thus the generated stabilizer is not
held open by an unused vacuum-construction target.

Section~\ref{p2a:app:later-application-interfaces} records the Paper-II
handoffs and later providers. The current series application map records the
same fields for each scoped contract, with exact source anchors and the
closure/application distinction.





\subsection{Generated charged-pole inputs and ordered application obligations}
\label{p2a:sec:v2168-charged-pole-inputs}
The charged-pole application now supplies a concrete example of
generated inputs and first unproved consumption. The source occurrence,
current tags and terminal coefficients are existing V-A/V-B outputs.
III-F supplies same-action kinetic normalization and proves actual
component-band retention of the internal current. IV supplies full finite
source dynamics, memory, a charge-sector pole test and a state-family
visibility test. V-C evaluates the native internal vertices and coherent
weak return. Each theorem retains its declared domain.

\begin{corollary}[Reuse of generated charged response inputs]
\label{p2a:cor:v2168-charged-pole-input-reuse}
After the same-owner source/action/constraint chart and any stationary/rest
dictionary consumed by the target have been certified, the complete
$H,\rho,q_r$ word generates $A,B,C_{\rm mem}$, the weak receiver and all
marked receiver kernels. They need no independent source, interaction,
covariance or mass-selection input. If a pole-location readout is common
on an already generated faithful stationary family, that location can be
retained without selecting one member. Residues retain their state weights. A root calculated on the exact actual
stationary support already includes visibility there; a faithful-carrier
root or a family admitting smaller supports retains its support test.
\end{corollary}
\begin{proof}
Apply the generated co-consumption and normalization-inheritance theorems
to the complete action/source word. IV,
Proposition~\ref{p4:prop:v2168-stationary-source-quotient}, gives the actual
quotient and Schur blocks; V-C,
Theorem~\ref{p5c:thm:v2168-complete-weak-memory}, constructs every marked
return from those blocks. IV,
Theorem~\ref{p4:thm:v2168-faithful-stationary-poles}, proves common
finite pole locations on the faithful stationary family. The existing
common-family readout rule then retains this particular common output.
It does not make the full response or residue common when they vary.
\end{proof}

\begin{center}
\begin{tabular}{p{.24\textwidth}p{.30\textwidth}p{.34\textwidth}}
\textbf{Consumed input}&\textbf{Existing generator}&\textbf{First unproved application}\\\hline
Terminal coefficients and occurring source
&V-A/V-B direct compiler and hierarchy synthesis
&Free route only: actual source unit-support/return equality.\\
Physical $H,q_r$, kinetic frame and constraint chart
&III-A complete action; III-F slot transport; IV quantization
&Certify the named chart and its stationary/rest dictionary if not already licensed.\\
Full marked memory
&IV source quotient; V-C native vertex receiver and full complement resolvent
&Evaluate the complete source-visible root at the fixed terminal energy.\\
Physical state visibility
&IV positive residue effect and generated stationary support family
&A root on the exact actual support includes visibility; a faithful-carrier root or smaller-support family needs its own support test. The named particle interpretation remains.\\
Physical particle and limits
&III-F/VII-C source/residue and physical observable transport
&Prove the named particle dictionary and each limit actually consumed.
\end{tabular}
\end{center}

The full dynamical route is sufficient for a pole-only claim and need not
close the stronger complete free-response certificate. Conversely a
successful geometric comparison lift does not identify the actual source
exporter. Owner-normalization inheritance is applied after owner ancestry;
it is not a proof of that ancestry. Once the finite complete root on the exact actual support is
established, visibility is included and existing pole transport supplies its finite residue with the
full derivative norm. Physical mass identification and limiting survival
remain visible at their consuming arrows. A05 remains open there, with
A01 closed and A12's separate actual K16 handoff unchanged.





\subsection{Exact generated inputs and the first unproved charged applications}
\label{p2a:sec:v2169-charged-construction-reuse}

The complete finite action now generates a terminating tagged source
polynomial; the actual stationary support generates its physical quotient.
This replaces an unspecified memory input by a finite construction and
removes density weights from the pole-location input. It does not prove
that the fixed terminal coefficient is a root.

\begin{corollary}[Reuse on the actual support, without a duplicate visibility task]
\label{p2a:cor:v2169-actual-support-reuse}
On the already certified same-owner finite chart, let $P$ be the exact
support of the actual stationary state or an already licensed stationary
constructor. The complete word $(H,q_r,P)$ generates $p_{r,P}$ and a
positive residue effect. The finite terminal-pole statement is exactly
\[
 \boxed{p_{r,P}(\varepsilon_r)=0.}
\]
This equality includes source visibility on that same exact support for
every stationary density having it. The weights within $P$ are not an
additional input for this pole location. A containing support, a family
allowing smaller supports, or a faithful-full-carrier calculation alone
still requires its own actual-support visibility certificate.
\end{corollary}
\begin{proof}
IV, Theorem~\ref{p4:v2169:b7:thm:closure}, proves exact finite closure and
Theorem~\ref{p4:v2169:b10:thm:support} proves the common null space and roots
on the exact physical support. The divided polynomial supplies its residue
effect. Apply the preterminal relation and common-family readout rules to
these already generated words. No independent source, covariance or state
member is chosen. The effect's weights and particle interpretation are not
inferred to be common from the common roots.
\end{proof}

\begin{center}
\begin{tabular}{p{.24\textwidth}p{.31\textwidth}p{.33\textwidth}}
\textbf{Required input/output}&\textbf{Existing generator}&\textbf{First unproved consumption}\\\hline
$\varepsilon_r,R_{\rm act},U,q_r$
&V-A/V-B source occurrence, terminal compiler and current; III-F/IV source transport
&Actual free return only if that stronger free-response route is claimed.\\
Complete finite $H$ and chart
&III-A complete action; III-F kinetic and slot transport; IV occupation/constraint quantization
&Named physical chart and consumed rest dictionary where not already licensed.\\
Actual stationary support $P$
&Occurring state/history or IV's licensed complete Gibbs vacuum; actual Slater/unit constructor on its own domain
&For a state not already stationary, evaluate the complete coefficient/support stationarity test.\\
Conserved fermion degree
&Balanced complete normal grammar proves $[N_F,H]=0$
&Only actual unbalanced words need cancellation or a larger cyclic space; no second conservation premise on the balanced branch.\\
Complete tagged polynomial
&IV's full commutators or actual adjacent-degree blocks; V-C's weak density and occupation words
&Assemble the actual complete blocks and calculate the source-visible roots of $p_{r,P}$.  If the terminal coefficient itself is claimed as the pole, then test $p_{r,P}(\varepsilon_r)=0$; the constructor's existence is closed.\\
Physical particle and limits
&III-F/VII-C same-source residue and observable transport
&Identify the named particle and prove each refinement, resonance or charged infrared limit used.
\end{tabular}
\end{center}


The evaluation order is therefore forward: $(H,q_r,P)$ first determines the
actual polynomial and its tagged root set.  Equation
$p_{r,P}(\varepsilon_r)=0$ is the certificate for the stronger
terminal-coordinate--equals--physical-pole claim; it is not an input used to
select a root.  If the generated tagged root is shifted, that shift is retained
as a downstream dynamical output and does not back-generate or fit a new
$\varepsilon_r$.

The action support test is IV,
Theorem~\ref{p4:v2169:b11:thm:normalstate}. Balanced typing and adjacent
sectors are IV, \eqref{p4:eq:v2169-b12-n0} and
Theorem~\ref{p4:v2169:b12:thm:adjacentnumber}.
On an actual empty-fermion support the density direction is reduced exactly,
but the bosonic occupation word remains. A filled negative-energy ground
uses its actual occupied degree, or all degrees of the licensed degenerate
ground support; it is not replaced by the empty unit. The native coordinate
port acts nontrivially on its free unit. The filled Slater--free-coordinate
state has weak leakage $3\pi^2\kappa/2>0$ on the stated chiral constituent;
complete-action stationarity needs the assembled coefficient cancellation or
an actual dressed support. An already licensed complete Gibbs ground is
stationary and bypasses this Slater test.

The general polynomial route, the whole-response free/fidelity route and
III-F's optional reference-mode protection route are alternatives. They
must not become cumulative obligations. Source-normalization inheritance
retains $q_r=\mathcal QN^{1/2}e_r$; neither projected CAR nor an unchanged
basis tag is assumed after a mixing Gram normalization. Primitive $m_p=1$
and all fixed terminal coefficients are retained. A01 stays closed;
A05 keeps its physical identification/limit status; A12 and A21 retain
their separate first unproved applications.




\subsection{Clocked action geometry and generated evolution}
\label{p2a:sec:v2190-clocked-action-evolution}

The following is an abstract representation theorem for an already licensed
action/evolution word.  It neither assumes a downstream action exists nor
adds a fourth constitutional principle.  Its physical use consumes the
frozen constructor origin of Definition~\ref{p2a:def:frozen-constructor-origin-certificate}
and the declared-target discipline of
Corollary~\ref{p2a:cor:declared-output-actualization-boundary}.

\begin{definition}[Regular clocked phase-action word]
\label{p2a:def:v2190-clocked-phase-word}
On a finite local chart of an already generated phase role \(\mathcal P\),
let \(t\) be its retained oriented clock coordinate.  A complete phase-action
word retains a real one-form \(\vartheta_t\), a real scalar \(\mathcal H_t\),
all source/constraint/normalization marks, and its actual allowed path
variations.  Write
\[
 \alpha=\vartheta_t-\mathcal H_t\,dt,\qquad
 \Omega_t=-d_{\mathcal P}\vartheta_t,
 \qquad S[\sigma]=\int_\sigma\alpha.
\]
The regular branch has nondegenerate \(\Omega_t\) after only already licensed
constraint and gauge treatment.  The coefficients are \(C^2\) locally in
phase coordinates and time, or have the corresponding regularity needed by
the claim.  The action, its variational zero relation, and any physical
evolution target are licensed before their evaluation.  The chart is a
presentation of the retained role; it is not a new autonomous carrier.
\end{definition}

\begin{theorem}[Clock-normalized characteristic and variational equivalence]
\label{p2a:thm:v2190-clocked-characteristic}
On Definition~\ref{p2a:def:v2190-clocked-phase-word}, put
\[
 \eta_t=\partial_t\vartheta_t+d_{\mathcal P}\mathcal H_t,
 \qquad \beta=-d\alpha=\Omega_t+\eta_t\wedge dt.
\]
There is a unique vector field \(X=\partial_t+Y_t\) with
\[
 \boxed{\iota_X\beta=0,\quad dt(X)=1,
 \qquad \iota_{Y_t}\Omega_t=\eta_t.}
\]
Its clocked integral curves are precisely the stationary paths of \(S\)
for compactly supported vertical variations, or for the retained endpoint
conditions that kill the actual endpoint term.  If \(\vartheta=p_i\,dq^i\)
in a time-independent canonical chart, this is exactly
\[
 \dot q^i=\partial_{p_i}\mathcal H,
 \qquad \dot p_i=-\partial_{q^i}\mathcal H.
\]
On every common existence domain its local evolution satisfies
\[
 F_{t,s}^{*}\Omega_t=\Omega_s,
 \qquad
 \frac{d}{dt}\mathcal H_t(\sigma(t))
 =\partial_t\mathcal H_t-(\partial_t\vartheta_t)(Y_t).
\]
For each initial phase point there is a unique local solution within the
regular chart.  A solution continues while it remains in a compact subset
of that chart over a compact clock interval.
\end{theorem}
\begin{proof}
The time-space split of \(-d\alpha\) gives the displayed \(\beta\).
Nondegeneracy uniquely solves \(\iota_Y\Omega=\eta\).
Since \(\eta(Y)=\Omega(Y,Y)=0\), the remaining temporal component of
\(\iota_X\beta\) vanishes.  Conversely the spatial component of that
contraction is the same nondegenerate equation, proving uniqueness.
For a variation \(V\) of a path,
\[
 \delta S=\int d\alpha(V,\dot\sigma)
               +[\alpha(V)]_{\rm endpoints}.
\]
After the actual endpoint term vanishes, arbitrary allowed vertical tests
force the displayed characteristic equation.  Full vertical testing is
used only on the regular phase role stated here; a restricted variation
sort retains its corresponding tested relation instead.

Since \(d_{\mathcal P}\Omega_t=0\) and
\(\partial_t\Omega_t=-d_{\mathcal P}\partial_t\vartheta_t\), Cartan's formula
gives \(\partial_t\Omega_t+\mathcal L_{Y_t}\Omega_t=0\), proving transport
of the form.  The energy identity follows from
\(d_{\mathcal P}\mathcal H(Y)=-\partial_t\vartheta(Y)\).
The coefficients of \(Y\) are locally Lipschitz, because the inverse of
the nondegenerate coefficient matrix is smooth on this chart.  On a closed
ball and a sufficiently short interval, the integral map
\(z\mapsto z_s+\int_s^tY_u(z(u))\,du\) maps the ball into itself and is a
contraction.  Iteration proves existence and uniqueness; covering compact
subsets by such neighborhoods proves the continuation statement.
\end{proof}

\begin{theorem}[Unmarked clock-comparison rigidity]
\label{p2a:thm:v2190-clock-isotropy-rigidity}
Let a licensed intrinsic clock-operation role already have a smooth
clock-normalized generator \(X_{\rm op}\), with
\(dt(X_{\rm op})=1\), over the regular comparison geometry above.  Compute
the unmarked isotropy group \(G_m\) at each point from the complete frozen
pre-generator signature, including every retained mark; no evaluated
generator or desired history may be used as a selector.  Suppose the
operation is equivariant under those actual owner automorphisms.  Set
\(V_m=\ker dt_m\).  If
\[
 \boxed{V_m^{G_m}=0,}
\]
then \(X_{\rm op}(m)=X(m)\), the characteristic direction, without a
separate postulate of stationary action.  A licensed isotropy element
acting as \(-I\) on \(V_m\) is sufficient.  More generally the invariant
clock-normalized tangent directions are exactly
\(X(m)+V_m^{G_m}\); this is a tangent classification, not an assertion
that all those directions are generated physical laws.
\end{theorem}
\begin{proof}
Every isotropy element preserves \(\beta\) and \(dt\).  It therefore
preserves their unique clock-normalized kernel vector \(X(m)\) pointwise.
Equivariance makes \(X_{\rm op}(m)\) fixed as well.  Their difference lies
in \(V_m^{G_m}\), proving the claim and the affine classification.
A transverse \(-I\) has no nonzero fixed vector.  On a native same-owner
generated branch the equivariance antecedent is the existing
Theorem~\ref{p2a:thm:v368-same-owner-unmarked-equivariance}.
\end{proof}

This is a sufficient geometric derivation of the governing direction,
not a claim that every marked physical action owner has that isotropy.
The bare regular \((\beta,dt)\) comparison geometry has such local
isotropy: straightening its characteristic flow and using a Darboux chart
on a time slice makes \(\beta\) the constant transverse symplectic form,
so transverse inversion preserves both tensors.  Passing from a complete
marked parent to that bare comparison role requires a licensed
factorization of the operation.  A coordinate basis change alone is not
an owner automorphism preserving its source, field and boundary marks.
If the complete signature separately retains a nonzero transverse energy
differential, its symplectic dual is already an invariant transverse
vector, so the zero fixed-space criterion cannot hold on that owner.
Paper~II-F, Section~\ref{p2a:prop:v2190-retained-mark-fixed-vector} proves this
boundary; applying the bare criterion there requires an independently
licensed operation factorization.

\begin{theorem}[Same-root clock-evolution closure and identification test]
\label{p2a:thm:v2190-clock-evolution-closure}
Let a certified word have the complete output of
Definition~\ref{p2a:def:v2190-clocked-phase-word}, with the stationary-path
relation and its local-solution presentation licensed on the stated domain.
The characteristic field and the local Cauchy graph of
Theorem~\ref{p2a:thm:v2190-clocked-characteristic} are representations of
that same relation.  Rooted generated-image co-realization retains these
representations without an additional governing-law occurrence.  A retained
parent family gives a family of graphs; common readouts descend without
selecting a member.  Initial points and actual histories retain their G0
status.  The local mathematical existence result supplies nonempty local
Cauchy fibres, not an arbitrary boundary-value or global extension.

If an independently retained clock-operation field \(X_{\rm op}\) is
already licensed on the same phase role and \(dt(X_{\rm op})=1\), then
\[
 \boxed{X_{\rm op}=X
 \quad\Longleftrightarrow\quad
 \mathfrak r_{\rm op}:=\iota_{X_{\rm op}}(-d\alpha)=0.}
\]
Testing this covector on a retained finite tangent basis is an exact finite
certificate.  If no typed map places \(X_{\rm op}\) in that role, its
comparison with \(X\) is not yet formed; provenance alone cannot supply it.
\end{theorem}
\begin{proof}
The variation identity proves equality of the already licensed stationary
relation and the characteristic relation.  The local initial-value result
re-presents its actual Cauchy branch.  Composition and rooted-image
co-realization apply to these retained outputs, exactly as in
Theorem~\ref{p2a:thm:generated-exact-solution-moduli}; they do not create a
new target sort by naming a formula.  The residual criterion is the
uniqueness of the clock-normalized characteristic direction.  A covector
on a finite tangent space vanishes exactly when its values on a basis do.
The typed comparison requirement is
Theorem~\ref{p2a:thm:typed-cross-relation-firewall}.
\end{proof}

The theorem closes an already licensed action/evolution relation.  It does
not infer stationary action from GP2 alone or identify every geometric
clock operation with that relation.  The native diagnostic in
Paper~II-F, Section~\ref{p2a:app:v2190-exactness-stationarity} makes that scope concrete.
Time dependence of \(\vartheta_t\) is retained; dropping
\(\partial_t\vartheta_t\) changes the governing relation.



\subsection{Assembly with an explicitly retained origin certificate}
\label{p2a:sec:v228-assembly-origin}
\begin{proposition}[An assembly map retains its input-generation antecedent]
\label{p2a:prop:v228-assembly-origin}
Suppose a licensed map $\mathfrak A$ consumes a complete generated packet
$Q=(q_r)_r$, and the only available export of an entry $q_r$ is conditional
on an independently generated input arrow $a_r$. Then generated-image
co-realization applies to $\mathfrak A(Q)$ on the domain where those input
origins are certified. Defining $q_r=\Phi(a_r)$ and printing
$\mathfrak A(Q)$ do not, by themselves, discharge the origin antecedent
for $a_r$. The required certificate is an existing licensed constructor
on the occurring parent, not an additional G0 token.
\end{proposition}
\begin{proof}
Theorems~\ref{p2a:thm:v187-pcad} and
\ref{p2a:thm:v3101-generated-coconsumption} transport generated inputs
through already licensed maps. Their hypotheses include the complete
input and its retained relations. Composition with $\Phi$ or
$\mathfrak A$ leaves any undischarged input hypothesis in place.
This is the constructor-origin partial order of
Remark~\ref{p2a:rem:constructor-origin-and-consumer}.
The statement does not exclude a separate origin proof elsewhere;
it specifies the certificate needed before this conditional chain is used.
\end{proof}
\begin{corollary}[Common readout without an input-member selector]
\label{p2a:cor:v228-constant-source-fibre}
Once the actual licensed output is proved to lie in a retained complete
family, a readout constant on that family is generated without selecting
a member. A family of merely available comparison operators cannot replace
that output-domain proof. Simultaneous changes of a shared carrier frame
transport every consumed map; applying a relative operation to one entry
is decided by its constructor and is not forced by a coordinate choice.
\end{corollary}
\begin{proof}
Apply Corollary~\ref{p2a:cor:v3101-coconsumed-family-readout}
and GP2 preservation of the complete typed diagram.
\end{proof}


\subsection{Later providers of the Paper-II application interfaces}
\label{p2a:app:later-application-interfaces}
This ledger explains later realizations without using them as premises of
the Paper-II seed proof. Each entry retains its provider's branch and domain;
optional targets become obligations only when a consumer claims them.

\subparagraph{K01 --- Exceptional support entry.}
On MobCl the $2I/2T/2O/C_6$ system and unmarked active family are generated. Certify the actual licensed finite frame-moving role if this support claim consumes it; a marked orientation is only needed by an orientation-sensitive output. Provider statements: \ref{p3c:def:v237-mobile-clifford-role}, \ref{p3c:thm:v237-generated-clifford-exceptional-resolution}, \ref{p3c:cor:v237-mobile-active-family-ancestry}.

\subparagraph{K02 --- Physical position-frontier history.}
$\mathsf{SF}_{\Lambda_F}$ generates $\mathrm{HIST}_x^{\rm fin}$, shell thresholds and retained $m_x$; no second m-history. Establish the actual sharp-frontier word and its scale, and threshold reach for any nonconstant growth claim. $\Lambda_F$=action cutoff is an optional separate matching condition. Provider statements: \ref{p3:def:sharp-scalar-resolution-frontier}, \ref{p3:thm:exact-spectral-resolution-staircase}, \ref{p3f:cor:sharp-frontier-generates-position-history}, \ref{p3f:thm:spectral-frontier-event-process}.

\subparagraph{K03 --- Stronger operator pairing.}
Local and harmonic comparisons are generated. A universal faithful all-band target in fixed $\operatorname{End}(\mathcal H_1)$ is excluded. Declare an admissible stronger comparison before proving its actual map. The excluded universal target is not a pending dynamical theorem. Provider statements: \ref{p2c:cor:harmonic-pairxp-generated}, \ref{p2c:thm:operator-germ-xp-pairing}, \ref{p3a:thm:canonical-generator-germ-localization}.

\subparagraph{K04 --- Reach of recovery bands.}
Covered finite/guarded recovery is exact; SF plus nondecreasing unbounded expansion generates exhaustion. On the actual physical frontier/history verify the needed threshold/guard. Prove uniform/interacting/continuum estimates only for a stronger claim; do not vary $m_p=1$. Provider statements: \ref{p2d:thm:scalar-fixed-band-recovery}, \ref{p2d:cor:guard-banded-nonlinear-recovery}, \ref{p3:cor:spectral-resolution-thresholds}, \ref{p8:prop:large-mx-smooth-limit}.

\subparagraph{K05 --- Absolute normalization and units.}
The intrinsic momentum curvature and neutral hierarchy are retained. The
$U=(3/\pi)\mathcal K_p$ and momentum-rooted $E_0$ rewrites additionally
consume the evaluated $(\mathrm{CELL}_{K\to U})$ comparison of
Eq.~\eqref{p3e:eq:v2211-momentum-eh-port}, together with the actual
energy-owner/kinetic dictionary where a physical energy is claimed. Physical
Newton normalization also consumes the complete scalar/joint matching.
After those actual ports are certified, numerical instrument/unit comparison
is metrology; calibration cannot supply a missing cross-owner identity. Provider statements: \ref{p3e:cor:typed-planck-normalization}, \ref{p6:thm:momentum-rooted-shell-normalization}.

\subparagraph{K06 --- States, preparations and instruments.}
Native root state and descendants close without second K16 preparation; state/effect rules are generated. For a particular new preparation or measurement supply/generate its actual apparatus process and required domain. Effects alone do not determine post-measurement instruments. Provider statements: \ref{p4:thm:v2114-neutral-fock-unit}, \ref{p8:thm:v2114-fresh-root-neutral-qroot}, \ref{p9:cor:v187-k16-state-inheritance}, \ref{p4:prop:effect-instrument-nonselection}.

\subparagraph{K07 --- Field/event envelope extension.}
Connected licensed bridges/tests yield their stated envelopes; native finite Born/event calculus is closed. Certify the actual bridge and connectivity on the target domain. Distinguish graded $M_{15}$ field envelope from parity-preserving $M_{11}\oplus M_4$ observables and from a full accessible event algebra. Provider statements: \ref{p4:thm:test-graph-env-discharge}, \ref{p4:thm:odd-bridge-field-envelope}.

\subparagraph{K08 --- Physical action, pole and stationary vacuum.}
Native action and source ownership and electroweak solder do not require a second scalar vacuum; finite pole transport is proved. For a physical pole certify its actual rest/pole dictionary and limit. For a stationary-equilibrium claim certify the effective action/state and stable admissible stationary family. Provider statements: \ref{p3f:thm:v174-generated-joint-gauge-phase-space-action}, \ref{p5:cor:v377-no-independent-ew-vacuum-token}, \ref{p3f:thm:v1140-finite-pole-source-transport}, \ref{p7v2:cor:ewsb-status}.

\subparagraph{K09 --- Critical/global continuation.}
Native local root, lineage and finite arrival close on their domains; full critical relation is retained. Prove the additional nonempty global/critical/strong-field continuation and reach actually consumed. A common readout can close without a unique complete history member. Provider statements: \ref{p2a:thm:v3110-finite-critical-stationary-relation}, \ref{p8:thm:complete-cap-root-cauchy-reclosure}, \ref{p8:thm:v2111-cyclic-k16-finite-arrival}.

\subparagraph{K10 --- Raw completion or complete structured uniqueness.}
Native $\mathsf C_D$ and protected even algebra close; a common scalar readout needs no singleton complete lift. Declare the optional raw target, then derive a licensed discriminator or provide its completion and prove compatibility. Do not infer an isomorphism from equal dimensions or demand complete uniqueness for a common output. Provider statements: \ref{p2c:cor:intrinsic-joint-comparison-role-export}, \ref{p2a:thm:compatible-family-observable-closure}.

\subparagraph{Remaining moment, charge and trace identifications.}
Paper~II-E, Corollary~\ref{p2e:cor:moment-prequantization-joint-locus},
gives $\mu_0=\varepsilon_{\rm red}\hbar n_{\rm pq}$; it selects no
occurring moment level. Equality to a separately owned Hopf/charge integer
requires the actual line/representation bridge. Later local cotangent,
operational Lorentzian and null-ray constructions and finite matter/anomaly
coefficients do not prove every global phase-space or quantum anomaly domain.
The $15/24$ supported trace identity remains in its own Born--McKay owner;
the excluded natural-$2T$ carrier identification is not a missing theorem.



\subsection{Generated source inheritance in the conormal response application}
\label{p2a:sec:v229-conormal-inheritance}
The exact conormal source map and covariance are supplied by
Proposition~\ref{p1b:prop:v229-affine-conormal}. On its stated whole
homothetic domain, Theorem~\ref{p3e:thm:v229-conormal-cap} evaluates
the complete flux response and its required absolute weight.
The existing generated-member and owner-normalization theorems preserve
the normalization of an identified native insertion; they do not
identify this insertion with a different action source without the
actual map. Proposition~\ref{p3e:prop:v229-source-counterfamily}
provides an explicit family sharing the listed geometric anchors when
that incidence is unspecified. Once the affine conormal prescription
is fixed, those deformations leave its class.
This is a concrete application of the existing target-member rule.
It adds no G0, source-value or particle-occurrence postulate and does
not reopen the already constructed rooted primitive source.

\begin{remark}[Constitutional scope of marked gravitational source transport]
\label{p2a:rem:v3028-marked-cap-incidence}
On the conditional four-port gravitational comparison of
Definition~\ref{p3e:def:v223-cap-data}, the fixed-reference current,
the determinant-normalized current and the positive-isometric current
are different source-constructor candidates, even when the last two
produce the same scalar vacuum term. GP2 preserves the exact current,
normalization, comparison frame and source measure of an \emph{identified}
native action; GP3's structured least-realization comparison does not
select a current solely for making vacuum energy constant or because
an isometry exists. The source-incidence certificate of
Definition~\ref{p2a:def:v3100-source-incidence-certificate}
must locate the native Dirac current and its gravitational action owner
before the proposed transport is promoted to an occurring response.
Corollary~\ref{p2a:cor:v239-observable-minimal-action}
allows an already licensed scalar contraction to be certified without
reconstructing unobserved matrix entries, but it does not license the
physical source or choose its transport.  No additional constitutional
principle, momentum-radius variation or independent vacuum field is
introduced by this scope distinction.
\end{remark}





\section{Physical occurrence and observable-readout limitations}
\label{p2a:sec:clean-observable-limits}
\subsection{Principle-theoretic motivation and empirical status}
\label{p2:app:axiom-motivation}

Geometric Primacy and Minimal Constitution begin with G0, geometric
actualization of distinguishability, and organize the constitutional triptych as
Finite/Bounded, Exact/Faithful, and Irreducible/Least.  G0 supplies the common
semantic type of physical occurrence, preparation, test, process, trajectory,
bridge, and calibration: each must exhibit a contingent token of a
geometrically licensed, claim-relative invariant distinction, with every
comparison and readout map fixed independently of its result.  A new role
enters only through a declared downstream frontier and relative GP1--GP3
closure.  It does not prove which admissible token occurs.  It does, however,
co-realize every typed canonical descendant of an occurring root; in
particular the transfer-operation parent and its generated momentum carrier
do not require an independent G0-MOM witness.  GP1 is
Finite-and-Bounded Realization Closure.  It applies to the whole structured
role: discrete carrier and represented-operation capacity is
finite, generated comparison/scalar/operation/rate images are bounded or
totally bounded in their constituted comparisons, and the total retained
composition and occupation envelope has finite realization.  These are typed
capacity statements, not one common metric notion of boundedness.  When an
intrinsic locally finite atomic construction grammar has already been
generated, a bounded generated sublevel is instead supplied only by the
separately registered hypothesis \(\mathsf P_{\rm gbd}\); it is not part of
GP1.  Continuum parent
operation groups and phase variables may therefore occur as parent or
presentation objects while their represented physical operation envelope and
all constituted real-valued operation/rate images obey the appropriate FRC
clause.

The comparison specialization GP1-TB is total boundedness of a constituted
comparison geometry with its finite-operational uniform bridge.  GP1-TB alone
does not imply a shortest distance or finite parent dimension.  Independently,
GP1 finite-realization capacity plus GP2 continuous faithful finite linear
representation unitarizes the compact exact-operation closure and gives
compact Lie structure before GP3.  \(\mathsf P_{\rm gbd}\) becomes available only after an
intrinsic atomic grammar and construction length have already been generated;
it supplies a bounded scope, not a selected depth.  Thus GP1 selects neither a
numerical scale nor a level.  The same clauses are reapplied to every later
structured role.  Consequently, after the elapsed scalar cover and a timed
propagation role exist, their generated rate image is totally bounded and has
the finite ceiling of Theorem~\ref{p2:thm:gp1-operational-speed-ceiling}.
Conditions (FD) and (LR) are optional ultraviolet regularity hypotheses on the
primitive untimed comparison.
The countermodels in
Proposition~\ref{p2:prop:gp1-clause-sharpness} show why both must remain
visible whenever claims about the freely parametrized first-order group-path
rate are made; those countermodels violate FRC-b once an unbounded timed
propagation-rate image is constituted.  Minimum-length and curved-momentum proposals provide motivation,
not empirical proof, for this ultraviolet picture
\cite{Garay1995,Hossenfelder2013,AmelinoCamelia2002,MagueijoSmolin2002,
Snyder1947,KowalskiGlikman2002DSR,RelativeLocality2011}.

GP2 is exact and host-faithful geometric primacy: later mathematical
presentations cannot retroactively constitute points, directions, frames,
scales, quotients, covers, representations, operation laws, or cross-role
relations as prior selection marks, while every item and canonical
nonredundant image generated inside a candidate must be retained in its proper
sort.  A selected candidate's complete exact presentation is frozen only after
GP3.  Physical generation follows the typed
constructor grammar rather than unrestricted tensor-category closure.
Canonical-image objectification, finite generated-image inheritance,
generated-role objectification, and relation-free composition make the same
chronology apply to reciprocal roles and their joint finite kinematic
realization.  The legacy \((\mathrm{BR}_{\rm ind})\) and
\((\mathrm{IND}_{xp})\) labels therefore certify that the antecedents of this
GP2 closure have been proved; they are not extra principles or independent
physical handoffs.  The canonical first-symbol comparison is generated and
co-realized by Paper~II-C.  A dimensionful reciprocal normalization or
stronger finite-band/local--global momentum interpretation remains the
separate G0-typed \((\mathrm{PAIR}_{xp})\) bridge/calibration interface.
GP3 compares the whole class of complete candidate-internally exact and
faithful realizations before candidate-specific structure is frozen.  Its
role-faithful preorder preserves the frozen inventory and every intrinsic
diagram of \(\operatorname{Geo}_R\), but does not require one candidate to
contain all self-operations of another.  It retains an irreducible least class
when one exists and otherwise leaves the irreducible family unresolved;
profiles generated from each comparison metric detect the
candidate geometric action kernel mathematically, while physical B-test
faithfulness requires the separately constituted metric-curry closure.  Neither
replaces the GP3 whole-candidate comparison.

Geometric Primacy constitutes the transfer role through the strong condition
\(\operatorname{Geo}_{A_{\rm v}}\); application of GP1 and GP2 derives a
nontrivial connected exact candidate parent and exactifies the intrinsic
order-swap diagram to its noncommutative status.  In the Paper-II primitive
endpoint domain, GP3 then derives OR1 and the \(SU(2)\) global form.  The
low-energy success of
three-component rotational momentum is therefore a check of the result rather
than the premise selecting it.  The physical value of the compact momentum
scale and the optional primitive-metric regularity conditions remain subject
to phenomenological test
\cite{Mattingly2005,KosteleckyRussell2011}.



\subsection{Finite parent saturation and correspondence descent}
\label{p2a:sec:v187-parent-saturation-pcad}

The universal co-realization bridge is constructor incidence and provenance, not
in general a direct map between already separated descendants.  The following
finite prepass makes that principle operational without enlarging the
constitutional set $G0+GP1+GP2+GP3$.

\begin{definition}[Finite intrinsic parent saturation]
\label{p2a:def:v187-parent-saturation}
For a constituted compound parent $P$ with selected intrinsic type, let
$\Sigma_P$ be the finite registry of algebraic, quotient, differential,
relation and same-action constructors explicitly licensed for that type.  The
parent saturation is
\[
 \operatorname{Sat}_{\rm par}(P;\Sigma_P),
\]
obtained by applying only those registered constructors before $P$ is split
into descendant physical roles.  This is not closure under all mathematically
available maps or theorems.  An object not named by $\Sigma_P$ is not promoted
merely because it is canonical in some external category.
\end{definition}

The typed image relation, support/totality conditions, and action-transport
domains are stated in the main text in
Theorem~\ref{p2a:thm:v187-pcad}. This appendix retains the parent-registry
definition and its provenance interpretation.



\begin{remark}[Typed relation-hypergraph normal form]
\label{p2a:rem:v187-relation-hypergraph}
The natural provenance object is therefore a typed relation hypergraph.  Its
primitive relation kinds include \path{PARENT_RELATION},
\path{DESCENT}, \path{INDUCED_ACTION}, \path{QUOTIENT},
\path{DEFECT}, \path{SAME_ACTION_RESPONSE}, \path{OCCURRENCE},
\path{ANALYTIC_CERTIFICATE}, and \path{CALIBRATION}.  A binary arrow is
recorded only after functionality has been proved.  This refines the existing
constructor firewall; it does not weaken it.
\end{remark}


\begin{theorem}[Intrinsic group-word restriction in parent saturation]
\label{p2a:thm:v188-group-word-restriction}
Let a saturated parent contain a constituted group owner $G$, a generated
subgroup support $H\le G$, and an already retained intrinsic group word
\[
 w:G^n\to G
\]
built only from multiplication, inverse and the identity.  If $H$ is closed
under that word, then
\[
 \boxed{w_H:=w|_{H^n}:H^n\to H}
\]
is a licensed parent-saturation restriction.  No new cross-descendant map is
selected.  In particular inverse, powers and commutators restrict to every
generated subgroup support.

If $Z\triangleleft H$ is a generated central subgroup and $w_H$ is independent
of the choice of representatives modulo $Z$, then the corresponding factor
map on $H/Z$ is also a licensed descendant relation. Every requested
finite prefix of the derived series is a finite parent-saturation output
when its subgroup-generation steps and output retention are licensed. For
finite $H$, the sequence stabilizes after at most
$\lfloor\log_2|H|\rfloor$ strict decreases. For an infinite support no global
stabilization follows from finite presentation alone.
\end{theorem}

\begin{proof}
A word built from multiplication, inverse and the identity evaluates in $H$
whenever its inputs lie in $H$, because $H$ is a subgroup.  Restriction
therefore changes only the declared domain/codomain, not the constructor or
its chronology.  If changing representatives by elements of $Z$ leaves the
word value unchanged in the declared quotient codomain, the quotient
universal property gives the factor relation. The commutator image and its
licensed subgroup closure give $[H,H]$; finite induction gives each requested
prefix. On finite $H$, closure under multiplication and inverse terminates
because it stays within a fixed finite set. Each proper subgroup has index
at least two, proving the strict-decrease bound. Equality at one step makes
all subsequent steps equal; the stable subgroup may be nontrivial and
perfect. None of this creates another physical support or occurrence token.
\end{proof}

\begin{remark}[Finite-support occurrence firewall]
\label{p2a:rem:v188-finite-support-occurrence-firewall}
The preceding theorem licenses restriction of an already generated operation
to an already generated support.  It does not infer that the same finite
support occurs on another physical factor.  In particular the Cartan
correspondence $g\mapsto(g,g^{-1})$ may descend to a finite relation between
labelled position and momentum group roles without constituting a finite
position carrier.  PCAD retains the relation, not an unproved occurrence.
\end{remark}


\subsection{Response-role, process, and action-target synchronization}
\label{p2a:sec:v397-response-action-target-sync}

\subparagraph{Response-role/action-tangent firewall.}
A generated balanced-response block, compact operator module, cotangent/source
role, determinant character, or mathematically canonical boundary operator is
not thereby an independently varied physical field or an action owner.  A
promotion into a configuration tangent, physical process target, or action
sort requires a generated constructor with that declared target.  Equal
dimension, equivalent representations, a common hierarchy label, common
scalar support, or common parent ancestry are insufficient.

\subparagraph{Generated process projectability.}
Let $F_t:X\to X$ be a generated one-parameter process family and let
$\pi:X\to Y$ be a generated readout. Put $Y_0=\operatorname{im}\pi$.
If $\pi\circ F_t$ is constant on every $\pi$-fibre, there is a unique
descended map $\bar F_t:Y_0\to Y_0$ with
\[
 \pi\circ F_t=\bar F_t\circ\pi.
\]
Indeed $\bar F_t(\pi x)=\pi(F_tx)$ is well-defined and lies in $Y_0$.
Whenever $F_{t+s}=F_tF_s$, surjectivity onto $Y_0$ gives
$\bar F_{t+s}=\bar F_t\bar F_s$ there. A unique descent on all of $Y$
follows when $\pi$ is onto $Y$; the fibre condition alone does not
determine values outside $Y_0$. Continuity or smoothness requires the
corresponding quotient structure. Proposition~\ref{p2a:prop:v257-observable-process-target}
constructs an onto continuous finite state-space target for the licensed
invariant test algebra. A separately typed physical target is inherited
only when that target role and its output retention are already licensed
by the generating chain.

\subparagraph{Boundary-functional reclassification.}
If a physical action owner is already constituted and stationary Schur reduction
produces a boundary descendant, an exact formula for that same descendant is
a mathematical reclassification of an existing action-owned object.  By
contrast, the uniqueness, positivity, naturality, or gluing property of a
mathematical boundary functional does not by itself manufacture a new physical
action constructor.  This distinction is used downstream to keep the Paper-I-B
Gaussian boundary calculus separate from the local metric modality while still
allowing it to characterize any already action-owned collar reduction.

\subsection{Nonzero ports do not actualize nonzero role coordinates}
\label{p2a:sec:v399-port-coordinate-boundary}

\begin{corollary}[Nonzero-port/zero-coordinate firewall]
\label{p2a:cor:v399-nonzero-port-zero-coordinate}
Let $P\neq0$ be a generated normalized port in a declared role and let an
admissible member of a later physical role have the form
\[
 q=sP
\]
for a scalar or typed role coordinate $s$.  If $s=0$ preserves every frozen
relation and licensed constructor closure, then
\[
 \boxed{P\neq0\ \not\Longrightarrow\ q\neq0.}
\]
The nonzero member $q\neq0$ remains an occurrence/member statement unless a
licensed constructor has $q\neq0$ itself as a required generated output.
\end{corollary}

\begin{proof}
Apply Corollary~\ref{p2a:cor:zero-member-occurrence-test} to the target role of
$q$.  The nonvanishing of the normalized port $P$ does not change the fact
that the zero target member is obtained at $s=0$.  If that zero member
preserves the frozen presentation, GP2 has no omission defect to remove and
GP3 has no redundant marked datum to delete.  The typed cross-relation
firewall, Theorem~\ref{p2a:thm:typed-cross-relation-firewall}, also forbids
promoting the port relation into a distinct physical occurrence relation
without a licensed constructor.
\end{proof}

\begin{remark}[Constituent co-realization is one-way]
\label{p2a:rem:v399-corealization-one-way}
If a nonzero frozen composite actually occurs, GP2 retains every generated
nonredundant constituent of its declared presentation.  The converse is not a
constitutional rule: the mathematical availability or generation of all ports
does not by itself actualize a nonzero composite member.  This distinction is
controlling for the fermion-source and terminal-neutral applications.
\end{remark}


\subsection{Same-owner defect inheritance and separated-scale typing}
\label{p2a:sec:v3100-defect-scale-typing}

\begin{definition}[Owner-equivalent finite transport and defect signature]
\label{p2a:def:v3100-owner-equivalent-transport}
Let $T:E_R\to E_L$ be a constituted finite operator in a declared owner slot.
A descendant is called \emph{undeformed owner-equivalent transport} when the
generated grammar supplies invertible support maps
$L:E_L\overset{\sim}{\to}E'_L$ and
$R:E'_R\overset{\sim}{\to}E_R$, with no generated projection, quotient,
compression, support-dimension change, independent additive source, singular
threshold, or other defect-changing constructor.  In that case
\[
  T'=LTR.
\]
Define
\[
 \delta(T):=(\dim\ker T,\dim\operatorname{coker}T).
\]
\end{definition}

\begin{theorem}[Undeformed same-owner defect inheritance]
\label{p2a:thm:v3100-same-owner-defect-inheritance}
For the transport of Definition~\ref{p2a:def:v3100-owner-equivalent-transport},
\[
 \ker T'=R^{-1}(\ker T),\qquad
 \operatorname{im}T'=L(\operatorname{im}T),
\]
and therefore
\[
 \boxed{\delta(T')=\delta(T)},\qquad
 \boxed{\operatorname{rank}T'=\operatorname{rank}T}.
\]
For square operators, $\det T=0$ if and only if $\det T'=0$.
\end{theorem}
\begin{proof}
Invertibility of $L$ and $R$ gives
$LTRx=0\Leftrightarrow TRx=0\Leftrightarrow Rx\in\ker T$, proving the
kernel formula.  Likewise $\operatorname{im}(LTR)=L(\operatorname{im}T)$.
Finite dimensionality gives the rank and cokernel statements.  The theorem is
the zero-order analogue of the existing undeformed-principal-owner
inheritance theorem: host-faithful transport cannot collapse a generated
finite distinction unless a separately licensed defect-changing constructor
is present.
\end{proof}

\begin{corollary}[Defect-change firewall]
\label{p2a:cor:v3100-defect-change-firewall}
If the defect signature changes along a claimed physical chain, that change
must be localized to an explicitly generated defect-changing event, such as a
projection, quotient, compression, support change, noninvertible transport,
independent additive deformation, Schur/Feshbach reduction into a new owner,
or a singular endpoint.  Ordinary owner-equivalent transport cannot be cited
as the cause of the rank change.
\end{corollary}

\begin{corollary}[Uniform same-owner gap inheritance]
\label{p2a:cor:v3100-uniform-gap-inheritance}
For $T'=LTR$,
\[
 \boxed{\sigma_{\min}(T')\ge
 \sigma_{\min}(L)\,\sigma_{\min}(T)\,\sigma_{\min}(R)}.
\]
Hence a transport family with uniformly positive lower singular bounds on its
left and right factors inherits any positive finite source gap uniformly.  A
terminal infrared claim requires the quantitative condition
$\liminf \sigma_{\min}(T)>0$; full rank at every finite scale alone is not a
terminal-gap theorem.
\end{corollary}

\begin{lemma}[Separated-scale kernel elimination]
\label{p2a:lem:v3100-separated-scale-kernel-elimination}
Suppose UV data determine $X$ only modulo $K=\ker D_{\rm UV}$,
\[
 D_{\rm UV}X=u,\qquad X=X_0+k,\quad k\in K,
\]
and an independently generated IR/end-point condition gives
$B_{\rm IR}X=b$.  If $B_{\rm IR}|_K$ is injective, then the joint UV/IR
boundary-value problem has at most one solution.  Thus distinct-scale objects
need not be identified; the IR condition may remove precisely the UV kernel.
\end{lemma}
\begin{proof}
Two joint solutions differ by $k\in K$ and obey $B_{\rm IR}k=0$.
Injectivity of $B_{\rm IR}|_K$ forces $k=0$.
\end{proof}

\begin{remark}[Scale-matching firewall]
\label{p2a:rem:v3100-scale-matching-firewall}
A coefficient/source defined at one scale is not to be equated directly with
a response, running quantity, or pole defined at another scale unless a
generated matching/transport map supplies that relation.  This is the general
form of the separated-scale Born--Legendre repair and applies equally to
fermion source versus pole data and to K16 heat-boundary versus hard/plasma
couplings.
\end{remark}

\begin{definition}[Event-closure factorization]
\label{p2a:def:v3100-event-factorization}
For a physical interface/event $E$, keep the obligations
\[
 \{\mathrm{REACH}_E,\mathrm{DOMAIN}_E,\mathrm{STATE}_E,
   \mathrm{MATCH}_E,\mathrm{READOUT}_E,\mathrm{INFERENCE}_E\}
\]
type-distinct.  Closure of one entry does not close the others.
Deterministic propagation from an occurring state is a descendant, a predicate
on that descendant is not a second occurrence, and observational likelihood is
an inference layer rather than a physical-state constructor.
\end{definition}

\begin{definition}[Source-incidence certificate]
\label{p2a:def:v3100-source-incidence-certificate}
Promotion of a generated nonzero source port to an occurring source value must
identify, in one certificate, the occurring parent, licensed constructor,
target $\mathsf{SrcVal}$ sort, action owner, required-output status, and
nonzero proof.  Mathematical availability of the port or membership in a
larger owner family is insufficient.  This definition specializes the existing
cross-relation and zero-member firewalls; it introduces no new GP.
\end{definition}




\section{Rooted relations and same-owner descendant transport}
\label{p2a:sec:clean-rooted-relations}
\subsection{Preterminal relation preservation}
\label{p2a:sec:v3102-preterminal-relation-preservation}
Generated co-consumption applies to the complete declared constructor output.
A noninjective readout of individual constituents cannot be substituted for
that output before the relations needed by the claim have been formed.

\begin{theorem}[Preterminal relation preservation]
\label{p2a:thm:v3102-preterminal-relation-preservation}
Let one occurring parent generate channels $T_a$ and embeddings $E_a$ into a
common finite receiver $(Q,K_Q)$, and suppose the licensed action constructor
uses
\[
 B=\sum_aE_aT_a.
\]
If eliminating $Q$ is part of the declared action word, the generated return is
\[
 \boxed{
 B^\dagger K_Q^{-1}B
 =\sum_{a,b}T_a^\dagger
    \bigl(E_a^\dagger K_Q^{-1}E_b\bigr)T_b.}
\]
All cross relations present in this word are co-consumed with their inherited
coefficients.  Let $q$ be a later trace, norm, determinant, singular-value
map, support quotient, Moore--Penrose scalar readout, positive codiagonal or
other noninjective descendant.  Replacing the complete parent word by the
separate data $q(E_aT_a)$ before the displayed return is valid only if a
licensed theorem proves that $q$ is relation-faithful for every cross return
used by the physical claim.
\end{theorem}

\begin{proof}
Expand the complete quadratic return.  Generated co-consumption
(Theorem~\ref{p2a:thm:v3101-generated-coconsumption}) retains all generated
terms and mixed terms of the declared action.  A noninjective descendant can
be applied afterward on its stated target, but GP2 does not supply an inverse
section for information removed by that descendant.  Moving $q$ before the
complete contraction therefore changes the constructor unless
relation-faithfulness is separately proved.
\end{proof}

\begin{corollary}[Construct before terminalizing]
\label{p2a:cor:v3102-construct-before-terminalize}
For a physical claim whose owner incidence is carried by cross-channel returns,
the licensed chronology is
\[
 \boxed{\begin{aligned}
  &\text{generated parent channels}
    \to\text{complete action contraction}\\
  &\to\text{owner/family return}\\
  &\to\text{terminal scalar or support readout}.
 \end{aligned}}
\]
Independent scalarization followed by later ``reassembly'' is not a GP2
constructor and cannot create the discarded incidence.
\end{corollary}

\begin{remark}[No new constitutional principle]
\label{p2a:rem:v3102-no-new-gp}
No additional GP is needed: this is the chronological consequence of
generated co-consumption together with the no-hidden-section and
typed-target conditions.
\end{remark}

\subsection{Sibling readout and relation consumption}
\label{p2a:sec:v3103-sibling-output-consumption}
Preterminal relation preservation does not forbid a scalar descendant from
appearing in a later physical constructor.  It forbids discarding the sibling
relation data and then trying to reconstruct it from the scalar.

\begin{corollary}[Sibling scalar/relation output may be consumed jointly]
\label{p2a:cor:v3103-sibling-output-consumption}
Let a frozen constructor on an occurring root have the complete required
output
\[
 \mathfrak q=((m_r,\rho_r))_{r\in I},
\]
where each $m_r$ is a scalar or support readout and each $\rho_r$ is a retained
owner relation generated from the same rooted role.  If a later licensed
constructor has $\mathfrak q$ as its typed input and is frozen independently of
the evaluated values, then generated co-consumption retains every pair
$(m_r,\rho_r)$ and all relations used by the later word.  The scalar component
may therefore appear as a coefficient in that later output without becoming a
selector for $\rho_r$.

Conversely, replacing $\mathfrak q$ by $(m_r)_r$ and choosing or reconstructing
$\rho_r$ afterward is not the same constructor unless an independent
relation-faithfulness theorem supplies that reconstruction.
\end{corollary}

\begin{proof}
Apply Theorem~\ref{p2a:thm:v3101-generated-coconsumption} to the complete
output tuple and Theorem~\ref{p2a:thm:v3102-preterminal-relation-preservation}
to the later word.  No new target sort or cross-role relation is created: the
later constructor consumes only previously generated sibling outputs in their
native sorts.
\end{proof}

\begin{remark}[No new mass or source principle]
\label{p2a:rem:v3103-no-new-principle}
This corollary adds no GP, no numerical coefficient and no source-value target.
A downstream paper must still declare and license the physical target sort of
its own constructor.
\end{remark}


\subsection{Rooted relative-comparison reduction}
\label{p2a:sec:v3104-rooted-relative-comparison}

\begin{theorem}[Rooted identity-channel reduction]
\label{p2a:thm:v3104-rooted-identity-channel-reduction}
Let $X$ be a finite generated comparison orbit and let
\[
 \mathcal L_X:=\ell^2(X_L\times X_R)
\]
be an already licensed source--target relation carrier with its canonical
orthonormal label basis.  Suppose one occurring rooted member $x_0\in X$ is
already present before the relative comparison is evaluated.  If the licensed
relative-comparison defect is required to
\begin{enumerate}
\item vanish on the already-realized identity relation $(x_0,x_0)$, and
\item preserve every other generated relation label without rescaling or
mixing,
\end{enumerate}
then GP2 exactness and GP3 leastness force the unique operator
\[
 \boxed{Q_{\rm rel}=I-|x_0,x_0\rangle\langle x_0,x_0|.}
\]
For the invariant state
\[
 \Omega_X=\frac1{|X|}\sum_{x,y\in X}|x,y\rangle
\]
one has
\[
 \boxed{\langle\Omega_X,Q_{\rm rel}\Omega_X\rangle
 =1-|X|^{-2}.}
\]
\end{theorem}
\begin{proof}
The first requirement puts $|x_0,x_0\rangle$ in the kernel.  The second fixes
the operator to the identity on its orthogonal complement.  These two
orthogonal subspaces exhaust $\mathcal L_X$, so the displayed projector is
unique.  The invariant-state expectation follows from
$|\langle x_0,x_0|\Omega_X\rangle|^2=|X|^{-2}$.
\end{proof}

\begin{corollary}[No duplicate occurrence for a predeclared rooted relation carrier]
\label{p2a:cor:v3104-rooted-relation-corealization}
If an occurring finite support $I$ and rooted subgroup $H$ are already
co-realized by a licensed physical comparison whose predeclared target is the
relation carrier on $X=I/H$, then the relation labels and the reduction of
Theorem~\ref{p2a:thm:v3104-rooted-identity-channel-reduction} are generated
outputs of that comparison and require no second occurrence token.  This
corollary does not license an arbitrary mathematical quotient or relation
carrier when the target sort was not already declared.
\end{corollary}
\begin{corollary}[Co-consumption of the nested icosian and rooted comparison outputs]
\label{p2a:cor:v243-one-root-register-descent}
Assume the \emph{already occurring}, licensed relative-comparison
constructor of
Theorem~\ref{p2a:thm:v3104-rooted-identity-channel-reduction}
has $X=2I/2T$ and physically retains the $25$ ordered relation
labels together with their $24$-sample payloads.  Then its rooted
deleted-label word acts as $Q_{\rm rel}\otimes I_{24}$ on the
labelled payload carrier.  Its finite scalar comparisons coincide with
$\tau_{\rm root}=\tau_{2I}-\tau_{2T}/25$ by
Theorem~\ref{p1b:thm:v243-rooted-register-functional}.
If the licensed constructor also retains both the unrooted comparison
and the later rooted reduction as sibling outputs, the full
unrooted degree-$12$ source and the rooted degree-$6$ defect are
co-consumed without a second geometric selector; their harmonic
onsets are the nested projections of
Theorem~\ref{p1b:thm:v243-nested-reynolds-onsets}.
No independently unlicensed matrix coefficient, group-orbit span,
field embedding, source map or action coupling is produced by this
corollary.  In particular the result is conditional on occurrence of
the predeclared physical comparison and the stated retention/reuse
policy; it does not infer such occurrence from the ambient geometry.
\end{corollary}
\begin{proof}
The projection of Theorem~\ref{p2a:thm:v3104-rooted-identity-channel-reduction}
acts on the orthonormal label basis and preserves each other label.
Tensoring with the \emph{already retained} $24$-point cell payload
implements exactly the block deletion.  The fivefold group-element
incidence proves the corresponding scalar functional, independently
of how overlapping labelled cells are ordered.  Apply
Theorems~\ref{p2a:thm:v3101-generated-coconsumption} and
\ref{p2a:thm:v3102-preterminal-relation-preservation} only to the
hypothesized retained output tuple.  Nothing in this argument
constitutes a new target sort or inserts a new action word.
\end{proof}




\subsection{Terminal mass and spectral non-generation}
\label{p2a:sec:v3105-mass-ontology-chronology}

\begin{theorem}[No spectral back-generation across a frozen terminal readout]
\label{p2a:thm:v3105-no-spectral-backgeneration}
Let an occurring licensed constructor generate a terminal scalar or role-tagged
coefficient $m_r$ together with any retained sibling incidence data, and let a
later generated constructor form an operator $A$ from those already generated
outputs.  If the terminal map producing $m_r$ is frozen before $A$ is formed,
then GP2 exactness, PCAD chronology and the no-insertion rule forbid replacing
$m_r$ by an eigenvalue, singular value, Takagi value or other spectral
functional of $A$ unless a separately licensed relation theorem explicitly
identifies the two outputs.  A later spectral equality may be a derived
comparison dictionary; it is not an upstream mass-generation rule merely
because both quantities live on the same finite carrier.
\end{theorem}
\begin{proof}
The terminal coefficient and the later operator are descendants at different
positions in the frozen constructor order.  PCAD preserves generated parent
relations but forbids a descendant value from reconstructing or redefining an
upstream output after a noninjective readout.  GP2/no-insertion likewise
forbids adding an equality relation not produced by the constructor grammar.
Therefore only an independently generated relation can identify a later
spectral functional with the earlier terminal coefficient.
\end{proof}

\begin{corollary}[Interaction-coordinate dictionaries do not own native masses]
\label{p2a:cor:v3105-sm-dictionary-firewall}
A downstream Standard--Model coordinate such as a Yukawa matrix, a
CKM/PMNS representative, a Takagi factorization or a mass-matrix
representative may encode an already generated finite source/current packet.
It does not thereby become the owner or generator of a native terminal mass.
Conversely, a framework-native mixing claim must be phrased in terms of the
generated current/incidence and propagation maps actually consumed by the
action; a relative-eigenframe formula is admissible only as a derived
coordinate dictionary when its hypotheses are separately established.
\end{corollary}




\subsection{Physical-descendant exclusivity}
\label{p2a:sec:v3106-physical-descendant-exclusivity}

\begin{corollary}[Physical descendant exclusivity]
\label{p2a:cor:v3106-physical-descendant-exclusivity}
Let $S_\star$ be one occurring physical action and suppose a licensed physical
boundary/interior construction admits stationary interior Schur reduction,
\[
 S_\star^\partial[q_\partial;y]
 =S_\star[q_\partial,\bar q_I(q_\partial;y);y].
\]
Assume a downstream theorem has independently established that
$S_\star^\partial$ is the physical descendant consumed in a declared target
role.  If the frozen generated grammar contains no second owner, carrier or
role which reconstitutes the eliminated interior as an autonomous downstream
object, then GP2 same-role no-splitting and generated co-consumption forbid
using both $S_\star^\partial$ and an independently evolved copy of $q_I$ as
coequal physical descendants of that same occurrence.  The parent action and
its eliminated interior remain available as derivational data, comparison
histories and source ancestry, but not as a duplicated downstream physical
degree of freedom.
\end{corollary}

\begin{proof}
The stationary descendant is already a generated child of the occurring
parent action in the declared target role.  Reintroducing the eliminated
interior as a second coequal physical child requires an additional role or
owner distinction.  By hypothesis no such distinction is generated.  The
same-role no-splitting lemma therefore keeps the physical descendant in the
single constituted target role, while generated co-consumption prevents
discarding the actual reduced relation in favour of an independently
normalized duplicate.  The parent may still be used for derivation or
counterfactual comparison because this does not create a second physical
occurrence.
\end{proof}

\begin{remark}[Cosmological use]
\label{p2a:rem:v3106-k16-use}
When the K16 ordinary-particle branch is identified with a same-action boundary
descendant, this corollary forbids adding an independently normalized duplicate
of the homogeneous scalaron.  It does not suppress the original bulk scalaron:
that same-action mode retains its K16 Cauchy data and continues on the physical
post-K16 history, while stationary or dynamic Schur objects are reductions of
that mode.
\end{remark}


\subsection{Boundary-response owner firewall}
\label{p2a:sec:v3107-boundary-response-owner-firewall}

\begin{corollary}[Boundary-response owner firewall]
\label{p2a:cor:v3107-boundary-response-owner-firewall}
Let an occurring same-action construction produce a licensed physical boundary
descendant $S_\star^\partial$ on a declared boundary tangent $J_\partial$.
Let a downstream theorem independently constitute a physical particle carrier
$J_P$ and let an owner-preserving inclusion/restriction place $J_P$ inside the
retained response domain.  Schur reduction, pole extraction and
positive-frequency polarization may then identify the response carried by
$J_P$ with a restriction of the same-action boundary response.  They do not,
by themselves, promote the complete boundary stress, energy packet or exchange
current to ordinary-particle ownership.  In particular,
\[
 T_{\partial,\mathrm{dyn}}^{\mu\nu}=T_{\rm ord}^{\mu\nu}
 \quad\text{or}\quad
 Q_\partial=Q_{\rm ord}
\]
requires a generated common-owner relation or an equivalent typed theorem.
\end{corollary}

\begin{proof}
The boundary response and the particle carrier have distinct target sorts until
the licensed owner-preserving map is supplied.  The declared-output
actualization boundary forbids changing target sort merely because a canonical
restriction, Schur complement, pole space or positive-frequency subspace can be
formed.  Once the particle carrier is independently constituted, the response
operator on that carrier is licensed; equality of the complete boundary stress
with particle stress is an additional common-value relation and is not implied
by mathematical availability alone.
\end{proof}

\begin{remark}[Direction of physical-descendant exclusivity]
\label{p2a:rem:v3107-descendant-exclusivity-direction}
Corollary~\ref{p2a:cor:v3106-physical-descendant-exclusivity} is one-way.  Once
a downstream theorem has established a stationary descendant as the physical
object/process in a target role, the eliminated parent interior cannot be
reintroduced as a coequal duplicate on that branch.  The existence of a
stationary descendant does not conversely select its physical owner or assign
ordinary-particle stress to it.
\end{remark}

\subsection{Fixed-owner successor-cell inheritance}
\label{p2a:sec:v3108-fixed-owner-successor-cell}

\begin{corollary}[Fixed-owner successor-cell inheritance]
\label{p2a:cor:v3108-fixed-owner-successor-cell}
Let an occurring licensed same-action boundary constructor produce a response
family on one fixed normalized boundary datum $q$,
\[
 S_\partial^{(2)}[q;\rho]
 =\eta(\rho)S_\partial^{(2)}[q;0],
\]
with no constructor in the cell changing the target sort of $q$.  Suppose an
independent endpoint/common-parent theorem identifies $q$ with a physical
owner $P$ by an owner-preserving inclusion into the declared boundary response
domain.  Then every nonzero member of the same successor-cell response has the
same physical owner $P$.  In particular, neither an interior value of $\rho$
nor the scalar response fraction $\eta$ licenses rebinding the datum to a
second owner.

If the response normalization is inherited from the same generated owner,
$\eta$ is the complete generated cell-dependent scalar response factor.  A
second multiplicative realization/efficiency coefficient requires a separately
generated deformation or owner-changing relation.
\end{corollary}

\begin{proof}
The constructor schema fixes the datum $q$ and varies only its scalar response
coefficient.  Same-role no-splitting therefore forbids a target-sort change
inside the cell without a generated owner-changing constructor.  Once the
endpoint/common-parent theorem supplies the typed owner of $q$, GP2
no-omission carries that already generated owner relation through every member
of the same constructor family.  Generated-owner normalization inheritance
then prevents a duplicate scalar normalization.  The argument does not infer
an owner from the existence of the boundary response alone; that endpoint
owner relation is an explicit antecedent, consistent with
Corollary~\ref{p2a:cor:v3107-boundary-response-owner-firewall}.
\end{proof}

\subsection{Generated-member handoff to owner coordinates}
\label{p2a:sec:v3109-generated-member-owner-handoff}

The source-incidence certificate of Definition~\ref{p2a:def:v3100-source-incidence-certificate}
and generated co-consumption already imply a reusable closure pattern whenever
the physical target coordinate and its insertion map are declared before the
specific generated member is evaluated.  The following theorem names that
pattern without adding a new generating principle.

\begin{theorem}[Generated member / owner-coordinate co-realization]
\label{p2a:thm:v3109-generated-member-owner-corealization}
Let $P$ be an occurring rooted input and let a result-independent licensed
constructor have frozen nonredundant output
\[
 \mathfrak C_P(P)=Q=(q_1,\ldots,q_N).
\]
Let $\mathcal O$ be an independently occurring physical owner whose frozen
constructor schema already declares a target coordinate $\mathsf T$ and a
result-independent typed insertion
\[
 \iota_{\mathcal O}:Q\longrightarrow\mathsf T.
\]
Assume that on the joint constructor domain the image
$\iota_{\mathcal O}(Q)$ is a required nonredundant output of the same physical
owner $\mathcal O$.  Then
\[
 \boxed{
 \operatorname{Occ}(P)\wedge\operatorname{Occ}(\mathcal O)
 \Longrightarrow
 \operatorname{CoReal}_{\mathsf T}
 \bigl(\iota_{\mathcal O}(Q)\mid\mathcal O\bigr).}
\]
No second occurrence token, independent deletion, rescaling or owner rebinding
of a constituent of $Q$ is compatible with the same generated owner
presentation.

This theorem does not promote mere mathematical availability to physical
consumption.  The target $\mathsf T$, the insertion $\iota_{\mathcal O}$ and
the required-output status are antecedent parts of the licensed constructor
schema; if any of them is absent, the typed-cross-relation and zero-member
firewalls remain in force.
\end{theorem}

\begin{proof}
The occurrence of $P$ co-realizes the complete frozen output $Q$ by
Theorem~\ref{p2a:thm:v3101-generated-coconsumption}.  The owner $\mathcal O$
and the typed insertion are independently licensed, so no new target sort or
cross-role map is manufactured at evaluation time.  The declared-output
actualization boundary and no-idle licensed-constructor closure retain the
required image $\iota_{\mathcal O}(Q)$.  Deleting or rescaling one of its
nonredundant generated constituents changes the frozen owner presentation and
violates GP2 exact faithfulness.  The final paragraph is exactly the
source-incidence firewall: without the predeclared target/insertion the proof
does not start.
\end{proof}

\begin{corollary}[Finite pair-cell specialization]
\label{p2a:cor:v3109-finite-pair-cell-specialization}
Let an occurring finite quadratic dynamics owner have a predeclared fermionic
pairing coordinate $\Delta$ and let an occurring rooted constructor generate a
specific nonzero pair cell
\[
 Q_\Delta=(\Delta_\star,\Delta_\star^\dagger)
\]
in that coordinate.  If the frozen quadratic owner consumes the pair through
\[
 H_{\rm pair}(\Delta_\star)
 =\frac12\bigl(
   \mathfrak a^\dagger\Delta_\star\mathfrak a^\dagger
   +\mathfrak a\Delta_\star^\dagger\mathfrak a
  \bigr),
\]
then the complete Hermitian pair cell and $H_{\rm pair}(\Delta_\star)$
co-realize with that owner.  The zero pair and a rescaled
$\lambda\Delta_\star$ are different owner presentations unless the frozen
grammar separately generates the corresponding deletion or normalization
deformation.
\end{corollary}

\begin{remark}[Relation to the v3.100 source certificate]
\label{p2a:rem:v3109-source-certificate-relation}
The theorem is the owner-coordinate analogue of the rooted finite
$\mathsf{SrcVal}$ handoff used in the earlier fermion-source closures.  It
closes no unresolved owner family by itself: it fires only after the physical
owner coordinate and its result-independent insertion have already been
licensed.
\end{remark}

\subsection{K48 terminal-relation co-consumption}
\label{p2a:sec:v3111-k48-terminal-consumption}

\begin{corollary}[K48 terminal relation in a certified homogeneous action consumer]
\label{p2a:cor:v3111-k48-terminal-consumption}
Assume a downstream claim has already constituted the homogeneous $H$ owner of
one K40--K48 boundary-response cell and that the upstream terminal constructor
has generated the complete $C_3$-invariant $T_5\leftrightarrow S_7$ quadratic
relation before a noninjective scalar/Hilbert readout. Assume also the
actual licensed $H$-boundary action word has that complete terminal packet
as a declared input and consumes it as a required output in the fixed
action role, with its generated insertion, supports, coefficients and
incidences. This is the same-action incidence of
Definition~\ref{p2a:def:v229-action-incidence-certificate}, certified before
the terminal scalar is evaluated, or an already generated certificate
with this same required-output status. Then generated
co-consumption, preterminal-relation preservation, and fixed-owner
successor-cell inheritance require the occurring $H$-boundary action to consume
the complete terminal relation with its generated normalization.  It may not
replace that relation by a scalar readout, delete one nonredundant term, attach
a second $H$ owner, or insert an independent realization efficiency.
\end{corollary}

\begin{proof}
The licensed actual action word supplies the complete input incidence.
Theorem~\ref{p2a:thm:v3101-generated-coconsumption} retains its complete
required output, while Theorem~\ref{p2a:thm:v3102-preterminal-relation-preservation}
preserves every nonredundant preterminal relation before a noninjective
readout. The fixed owner retains the inherited normalization. These
rules preserve the certified action incidence; they do not manufacture
an action insertion from the terminal provider alone.
\end{proof}


\begin{remark}[Occurrence does not determine a new response jet]
\label{p2a:rem:v3111-k48-response-jet-firewall}
The preceding corollary preserves the certified action incidence and the already
generated data. It does not manufacture the insertion or a new target-coordinate law.  In
particular it does not by itself determine the metric/Hilbert derivative of the
reduced $H$-boundary action, a finite $H$-kernel shift, or a numerical abundance
coefficient.  Such a coefficient is fixed only by an independently licensed
same-action insertion/variation of the complete boundary action.  This is the
same distinction used by generated-member/owner-coordinate co-realization:
co-consumption forces an already generated member into a predeclared coordinate
but does not create a missing coordinate map.
\end{remark}



\subsection{Coordinate-generated operations and downstream tests}
\label{p2a:sec:v226-operation-propagation}
\begin{corollary}[The intrinsic operator envelope has no second carrier premise]
\label{p2a:cor:v226-native-operation-closure}
After the fixed-radius intrinsic coordinate class is generated, its unital
associative operation envelope is $\mathcal A_D=M_{11}(\mathbb C)\oplus
M_4(\mathbb C)$. Its rotations, separate geometric duality and protected
chirality projectors, and all native Hermitian tests are coordinate
polynomials. The represented word filtration stabilizes at degree six.
This is closure in the operation sort on the frozen momentum carrier;
it is not a new carrier constructor or a proof that every operation is
performed as a Hamiltonian control.
\end{corollary}
\begin{proof}
Apply Theorems~\ref{p1:thm:v226-coordinate-only},
\ref{p1:thm:v226-rotations} and~\ref{p1:thm:v226-chirality} to the already
generated intrinsic class. GP2 retains the licensed word output in its
operation/test sort. The actual coordinate-control Lie algebra has dimension
65 by Theorem~\ref{p1:thm:v226-control-lie}; associative completion alone
does not promote every Hermitian polynomial to a controllable process.
\end{proof}
The exact native coefficient and jet tests are
Theorem~\ref{p1b:thm:v226-dual-coefficients} and
Corollary~\ref{p1b:cor:v226-source-jets}. The universal optimal-radius
chirality obstruction is Theorem~\ref{p2c:thm:v226-chirality-radius}.
These results replace redundant operator-availability questions on their
specified native block. They do not reopen the K01--K10 role constitutions,
request another inherited root state or electroweak vacuum, choose the actual
physical source/action coefficients, or evaluate the individual principal
$SU(2)$ descent fibres. The latter retain the one carry bit and actual
pullback of Theorem~\ref{p1b:thm:v226-principal-two-types}.


\subsection{Readout closure can precede complete-parent selection}
\label{p2a:sec:v227-visible-closure}
\begin{corollary}[A finite certificate for a common response class]
\label{p2a:cor:v227-visible-closure}
For independently licensed finite Hermitian parents on one retained/source
target, with common kinetic normalization and complementary dimensions at
most $m$, equality of their complete responses at $m+1$ distinct upper-half-plane
energies proves equality of the whole response. A pole/readout depending
only on that response can therefore close on the retained family even when
an uncoupled dark complement is not uniquely selected.
\end{corollary}
\begin{proof}
Apply Theorem~\ref{p1b:thm:v227-finite-response}, then the minimal visible
realization theorem cited there. Every functional of the identical response
has the same value on its admitted common domain.
\end{proof}
The compared matrices and source ports are actual family members, not
independent adjustable fits. All consumed source jets use
Corollary~\ref{p1b:cor:v227-response-jets}. A claim that also consumes
absolute action, gravity, a different state or an omitted source still
requires those data. This supplies a finite sufficient execution of the
existing common-readout rule, not a physical member-selection axiom.




\section{Action-incidence and exact observable identification}
\label{p2a:sec:clean-action-certificates}
\subsection{Preterminal action-incidence certificates and observable-facing ownership}
\label{p2a:sec:v229-action-incidence}

\begin{definition}[Preterminal action-incidence certificate]
\label{p2a:def:v229-action-incidence-certificate}
Fix an occurring geometric provider $P$ and an observable-facing physical
owner $\mathcal O$.  An \emph{action-incidence certificate} is a typed and
source-traceable tuple
\[
 \mathsf{ActCert}_{P\to\mathcal O}
 =(S_\star,\iota_{P\to\mathcal O},\mathscr V_{\mathcal O}S_\star,
    \mathcal R_{\mathcal O},\operatorname{Read}_{\mathcal O})
\]
for which (i) the parent and owner occur, (ii) the preterminal insertion
$\iota$ is generated and licensed independently of terminal output values,
(iii) the \emph{actual} action word contains the indicated image with all
oriented supports, coefficients, mixed interior terms, kinetic, connection,
and boundary contacts required by its sort, (iv) variation and reduction are
performed in the order of that common action, and (v) the stated readout
follows on a declared domain.  The tuple is a \emph{certificate schema}, not
an additional generating axiom or a guarantee that such an insertion exists.
\end{definition}

\begin{proposition}[Availability, action occurrence, and readout are distinct]
\label{p2a:prop:v229-action-consumer-firewall}
A nonzero provider, nonempty typed morphism space, admissible candidate action,
or terminally correct scalar output does not in itself establish the
certificate of Definition~\ref{p2a:def:v229-action-incidence-certificate}.
Conversely, when its five entries are supplied from one already occurring
parent and the hypotheses of typed parent correspondence and action descent,
preterminal relation preservation, and source-incidence co-realization hold,
the claimed action readout is a licensed descendant and needs no independent
occurrence token.  This statement does not imply the existence of a
specific source insertion, chiral-frame assignment, or interacting pole.
\end{proposition}
\begin{proof}
The asserted positive direction is the composition of
Theorem~\ref{p2a:thm:v187-pcad},
Theorem~\ref{p2a:thm:v3102-preterminal-relation-preservation}, and the
source-incidence predicate of
Definition~\ref{p2a:def:v3100-source-incidence-certificate}.
For the converse, choose the zero insertion into an otherwise nonzero
provider space: its source readout is zero.  Or choose distinct licensed
preterminal words with different reduced cross terms.  Neither provider
existence nor its terminal scalar fixes the action's omitted image.
\end{proof}

\subparagraph{Three reusable proof targets.}
For charged fermions, the actual same-action source test is
\[
 \mathscr M_p=Z_{L,p}^{-1/2}
 \frac{\delta^2 S_\star}{\delta\bar\psi_{L,p}\,\delta\psi_{R,p}}
 Z_{R,p}^{-1/2},
\]
with its role, chirality, comparison and right-dressing marks retained.
For a regular local stationary gravitational reduction $W(R,z)$, its
curvature Hessian is the \emph{complete} Schur response
\[
 \chi_{\rm red}=H_{RR}-H_{Rz}H_{zz}^{-1}H_{zR};
\]
frequency-dependent or causal parents require their actual operator-valued
inverse and boundary domains rather than this local scalar abbreviation.
For the $K16$ normal owner, the required input is the restriction
$S_{16,{\rm normal}}=\mathcal R_{16}(S_{F,{\rm complete}})$, followed by
the inherited state and charge-current readout.  In all three cases the
observable is tested against the geometry; it is not prescribed by
substituting a preferred conventional mechanism.



\begin{corollary}[Observable-minimal action certification without a member selector]
\label{p2a:cor:v239-observable-minimal-action}
Let an already licensed preterminal constructor on an occurring parent
produce its \emph{actual} output $X$ in a typed operator space, and let
$F(X)$ be a generated relation invariant consumed by a licensed observable
readout $r=\widetilde r\circ F$. If the complete output is proved to lie
in a nonempty generated family $\mathfrak X$ on which $F$ is constant,
then $r(X)$ is uniquely determined without selecting a member of
$\mathfrak X$. The same holds if only the value of $F(X)$ is established
from the actual action contraction and no complete operator reconstruction
is made. Neither statement licenses the action constructor or proves that
its actual output belongs to $\mathfrak X$ merely from the geometric
availability of that family.
\end{corollary}
\begin{proof}
The observable factors through the proved invariant. The first case
is Corollary~\ref{p2a:cor:v228-constant-source-fibre} on the complete
preterminal output, with relation chronology from
Theorem~\ref{p2a:thm:v3102-preterminal-relation-preservation}.
The second is direct substitution into $\widetilde r$. The required
action origin remains the independent hypothesis of
Definition~\ref{p2a:def:v229-action-incidence-certificate}; constancy
of a readout cannot supply that missing constructor.
\end{proof}

\begin{proposition}[Finite sufficient certificates for an observable]
\label{p2a:prop:v257-observable-certificates}
Let $C_j(D,z)=0$ describe an already licensed nonempty real family, with
all its branch restrictions retained, and let $O$ be a licensed polynomial
readout. Write $f=O-o(D)$ and take integers $N\ge0$ and $k\ge1$.
Either exact identity
\[
 d^Nf=\sum_jh_jC_j\quad(d\ne0),\qquad\text{or}\qquad
 f^{2k}+\sum_\ell s_\ell^2=\sum_jh_jC_j\quad(k\ge1)
\]
proves $O=o(D)$ on the stated domain, without selecting a source member.
This supplies a sufficient constancy certificate for
Corollary~\ref{p2a:cor:v239-observable-minimal-action}; it does not
license an action insertion, readout or family constructor.
\end{proposition}
\begin{proof}
Evaluate the identity on a surviving member. In the first case divide
by the nonzero $d^N$. In the second, every real summand is nonnegative,
so $f=0$. Nonemptiness prevents vacuous prediction. There is no claim
that a certificate of bounded degree exists for every instance.
\end{proof}


\subsection{Exact finite converse certificates and observable identifiability}
\label{p2a:sec:v267-exact-converse-observables}
The certificates below strengthen the observable-minimal principle of
Corollary~\ref{p2a:cor:v239-observable-minimal-action}.  They are finite
mathematical equivalences on \emph{already constituted} operators and test
families.  They neither create an action constructor nor identify a geometric
comparison source with an actual mass or interaction action block.

\begin{theorem}[One positive-time heat--Bregman converse]
\label{p2a:thm:v267-heat-bregman-converse}
Let $B,D$ be finite Hermitian matrices and $\tau>0$.  Define
\begin{equation}\label{p2a:eq:v267-heat-bregman}
 {\cal H}_\tau(B,D)=\operatorname{Tr}(e^{-\tau B})
 -\operatorname{Tr}(e^{-\tau D})
 +\tau\operatorname{Tr}(e^{-\tau D}(B-D)).
\end{equation}
Then
\[
 \boxed{{\cal H}_\tau(B,D)\geq0,\qquad
 {\cal H}_\tau(B,D)=0\ \Longleftrightarrow\ B=D.}
\]
If $0\leq B,D\leq LI$, the exact quantitative comparison obeys
\[
 \frac{\tau^2e^{-\tau L}}2\|B-D\|_{\mathrm{HS}}^2
 \leq{\cal H}_\tau(B,D)
 \leq\frac{\tau^2}2\|B-D\|_{\mathrm{HS}}^2.
\]
For finitely many sectors, any positive weighted sum of these gaps
vanishes if and only if every included sector has $B_s=D_s$.
\end{theorem}
\begin{proof}
Put $X=B-D$, $H_u=D+uX$ and $f(u)=\operatorname{Tr}e^{-\tau H_u}$.
Duhamel differentiation gives
\[
 f''(u)=\tau^2\int_0^1\!\operatorname{Tr}
 \bigl(Xe^{-\tau(1-v)H_u}Xe^{-\tau vH_u}\bigr)\,dv.
\]
In a spectral basis of $H_u$ each summand is
$|X_{ij}|^2\int_0^1e^{-\tau[(1-v)h_i+vh_j]}dv$,
strictly positive if $X\neq0$.  Moreover
${\cal H}_\tau=f(1)-f(0)-f'(0)=\int_0^1(1-u)f''(u)du$.
Under the stated order bound, all exponential weights lie between
$e^{-\tau L}$ and $1$, proving the displayed stability inequalities.
A positive weighted sum cannot cancel a nonzero positive gap.
\end{proof}

\begin{corollary}[Canonical quadratic limit of the finite heat--Bregman defect]
\label{p2a:cor:v3026-hs-heat-quadratic-limit}
Let $B,D\in\operatorname{Herm}(\mathcal H)$ belong to one already licensed
finite complex carrier and put $X=B-D$.  For the \emph{same} canonical
unmarked trace closing word as in
Theorem~\ref{p2:thm:unmarked-natural-closure} and the heat--Bregman
functional of Theorem~\ref{p2a:thm:v267-heat-bregman-converse},
\[
 \boxed{\lim_{\tau\downarrow0}\frac{2}{\tau^2}
 {\cal H}_\tau(B,D)=\operatorname{Tr}(X^2)
    =\|B-D\|_{\mathrm{HS}}^2.}
\]
This is the infinitesimal quadratic comparison geometry; a finite positive
heat-time Bregman minimizer need not be the Hilbert--Schmidt minimizer.
For a general non-Hermitian $X$, the separately licensed adjoint/product
comparison is $\operatorname{Tr}(X^\dagger X)$; it is not obtained by
applying the Hermitian heat expression to arbitrary matrices.  Neither
identity selects an action member or supplies a dynamical dissipation law.
\end{corollary}
\begin{proof}
Finite-dimensional exponential expansion gives
$\operatorname{Tr}e^{-\tau B}=n-\tau\operatorname{Tr}B+
 \tau^2\operatorname{Tr}(B^2)/2+O(\tau^3)$, and likewise for $D$.
The linear correction is
$\tau\operatorname{Tr}[e^{-\tau D}(B-D)]
=\tau\operatorname{Tr}X-\tau^2\operatorname{Tr}(DX)+O(\tau^3)$.
The constant and linear pieces cancel.  The quadratic coefficient is
$\frac12\operatorname{Tr}(B^2-D^2)-\operatorname{Tr}(DX)
=\frac12\operatorname{Tr}X^2$, by cyclicity of the trace; dividing by
$\tau^2/2$ proves the formula.  Hermiticity gives
$\operatorname{Tr}X^2=\operatorname{Tr}X^\dagger X$.
\end{proof}

\begin{corollary}[Nearest compatible finite operator, with fixed comparison role]
\label{p2a:cor:v3026-hs-affine-projection}
Let $\mathcal C$ be a nonempty closed affine subspace of the real
Hilbert space $\operatorname{Herm}(\mathcal H)$ on one \emph{already
specified} finite carrier, and let $D$ be a same-sort target. Then
$A_*=\arg\min_{A\in\mathcal C}\|A-D\|_{\mathrm{HS}}^2$ exists,
is unique, and satisfies
\[
 \boxed{\|A-D\|_{\mathrm{HS}}^2
 =\|A-A_*\|_{\mathrm{HS}}^2+\|A_*-D\|_{\mathrm{HS}}^2
 \quad(A\in\mathcal C).}
\]
The family, carrier, target and comparison role must already be licensed;
this mathematical nearest representative is \emph{not} a replacement for
GP3's structural preorder, and need not be an occurring physical operator.
For a merely closed convex family the nearest point remains unique, but
only the projection variational inequality is automatic, not this
Pythagorean equality.
\end{corollary}
\begin{proof}
Writing $\mathcal C=A_0+V$ with $V$ a linear subspace, set
$A_*=A_0+P_V(D-A_0)$, using the real inner product
$\langle X,Y\rangle=\operatorname{Tr}(XY)$ on Hermitian matrices.
Then $D-A_*$ is orthogonal to every $A-A_*\in V$, proving the
identity and uniqueness. The distinction from physical least-realization
is a type/scope statement, not an extra variational postulate.
\end{proof}

\begin{theorem}[One Euclidean-frequency tagged converse]
\label{p2a:thm:v267-one-frequency-tagged}
Let $B\geq0$, $\|g\|=1$, $d=\langle g,Bg\rangle$ and $\omega>0$.
Writing $R_\omega=\langle g,(\omega^2+B)^{-1}g\rangle$,
\begin{equation}\label{p2a:eq:v267-tagged-jensen}
 \boxed{(\omega^2+d)R_\omega-1
 =\frac{\langle (B-d)g,(\omega^2+B)^{-1}(B-d)g\rangle}
 {\omega^2+d}\geq0.}
\end{equation}
Equality at \emph{one} $\omega>0$ is equivalent to $Bg=dg$.
If the same action independently evaluates $d=e^2$ for the unit left
source $g_e$, equality certifies its complete \emph{finite free}
tagged mass-only response $z/(z^2-e^2)$; no assertion of an interacting
physical pole follows.
\end{theorem}
\begin{proof}
Apply the spectral probability measure of $g$ to the identity
\[
 \frac{d-x}{\omega^2+x}=
 \frac{d-x}{\omega^2+d}
 +\frac{(d-x)^2}{(\omega^2+d)(\omega^2+x)}.
\]
The first expectation vanishes.  The remaining quadratic form is
strictly positive unless $(B-d)g=0$.  Squaring the finite off-diagonal
chiral Hamiltonian then gives the stated tagged resolvent.
\end{proof}

\begin{theorem}[One unmarked heat value after rooted first-moment certification]
\label{p2a:thm:v267-rooted-one-heat}
Let $g_1,\ldots,g_n$ be a \emph{complete} orthonormal left frame,
$B\geq0$, and $D=\sum_r\lambda_rg_rg_r^\dagger$ with $\lambda_r>0$.
Assume the \emph{actual source} has independently established
$\langle g_r,Bg_r\rangle=\lambda_r$ for every $r$. Then at any
$\tau>0$,
\[
 \boxed{E_\tau:=\operatorname{Tr}(e^{-\tau B})
 -\sum_re^{-\tau\lambda_r}\geq0,\qquad
 E_\tau=0\ \Longleftrightarrow\ B=D.}
\]
The rooted first-moment premise follows in particular from an
\emph{actual} anchored unit-row packet
$M=\sum_r\sqrt{\lambda_r}\,g_rr_r^\dagger$ with $\|r_r\|=1$
and $B=MM^\dagger$, even without orthogonality of the right rays.
Under $0\leq B,D\leq LI$ the quantitative inequalities of
Theorem~\ref{p2a:thm:v267-heat-bregman-converse} apply to $E_\tau$.
\end{theorem}
\begin{proof}
Strict scalar Jensen for the spectral measure of each $g_r$ gives
$\langle g_r,e^{-\tau B}g_r\rangle\geq
 e^{-\tau\langle g_r,Bg_r\rangle}$, with equality precisely if
$Bg_r=\lambda_rg_r$.  Sum the nonnegative gaps over the complete
frame.  The anchored unit-row expression has
$M^\dagger g_r=\sqrt{\lambda_r}r_r$, yielding the required diagonals.
Finally $\operatorname{Tr}(e^{-\tau D}(B-D))=0$ in the rooted basis,
so $E_\tau={\cal H}_\tau(B,D)$ exactly.
\end{proof}

\begin{proposition}[Hadamard determinant alternative]
\label{p2a:prop:v267-normalized-determinant}
On the anchored $n$-row source domain of
Theorem~\ref{p2a:thm:v267-rooted-one-heat}, with each row unit-normalized,
\[
 \boxed{\det(MM^\dagger)\leq\prod_r\lambda_r,\qquad
 \det(MM^\dagger)=\prod_r\lambda_r
 \ \Longleftrightarrow\ MM^\dagger=D.}
\]
\end{proposition}
\begin{proof}
In the complete rooted left frame $B=\Lambda^{1/2}G\Lambda^{1/2}$,
where $G_{rs}=\langle r_r,r_s\rangle$ is a positive Gram matrix
with unit diagonal. Hadamard's inequality gives $\det G\leq1$, with
equality exactly for $G=I$.  Positive $\lambda_r$ allow cancellation
of $\det\Lambda=\prod_r\lambda_r$.
\end{proof}

\begin{proposition}[Exact relation-fibre converse on a real algebraic domain]
\label{p2a:prop:v267-real-radical-fibre}
Let the complete nonempty real algebraic input domain be
$\mathcal F=\{x:C_1(x)=\cdots=C_m(x)=0\}$, and let $q(x)$ and $O(x)$
be polynomial retained-data and observable maps.  Form the real ideal
\[
 J=\langle C(x),C(y),q(x)-q(y)\rangle\subset\mathbb R[x,y].
\]
Then $O$ is constant on each $q$-fibre of $\mathcal F$ (equivalently,
$O=\widetilde O\circ q$ for a set map on $q(\mathcal F)$) if and only if
\[
 \boxed{O_i(x)-O_i(y)\in\sqrt[\mathbb R]{J}
 \quad\text{for every component }i.}
\]
This is a mathematical reverse-identifiability test, not an action
constructor.  Strict inequalities, analytic action domains and restricted
physical branches must be represented separately; they cannot silently be
replaced by this equality-only algebraic domain.
\end{proposition}
\begin{proof}
Two inputs lie in one retained-data fibre precisely when they solve
the simultaneous real polynomial equations defining $J$.
The factorization assertion says that each $O_i(x)-O_i(y)$ vanishes
on that real zero set.  The real Nullstellensatz identifies its real
vanishing ideal with $\sqrt[\mathbb R]{J}$.  Define $\widetilde O$
by the common value on each nonempty fibre.
\end{proof}

\begin{proposition}[Why scalar spectra alone do not identify rooted action rows]
\label{p2a:prop:v267-unmarked-heat-no-converse}
Let $D=\operatorname{diag}(a,m,b)$ with $a\ne b$ and let $L_\varphi$
unitarily rotate only the first and third left eigenvectors.  Put
$B_\varphi=L_\varphi D L_\varphi^\dagger$.  Then
$\operatorname{Tr}f(B_\varphi)=\operatorname{Tr}f(D)$ for every finite
spectral function $f$, whereas
\[
 \operatorname{Tr}(DB_\varphi)
 =\operatorname{Tr}(D^2)-(a-b)^2\sin^2\varphi
 <\operatorname{Tr}(D^2)
\]
for a nontrivial angle.  In particular no collection of \emph{unmarked}
heat traces can supply the rooted first-moment premise of
Theorem~\ref{p2a:thm:v267-rooted-one-heat} on a general action domain.
\end{proposition}
\begin{proof}
Functional calculus and trace are invariant under conjugation.
Expanding the $2\times2$ rotated block gives the displayed strictly
positive rooted mixed-moment defect.
\end{proof}

\paragraph{Action-consumer boundary.}
Theorem~\ref{p2a:thm:v267-rooted-one-heat} is an \emph{if-and-only-if}
statement only \emph{after} the actual rooted first moments are proved.
A single unmarked heat value by itself is not a general iff test;
Proposition~\ref{p2a:prop:v267-unmarked-heat-no-converse}
excludes that converse.  The normalized charged mass-source packet,
its preterminal origin, and the identification of its complete heat
factor must be supplied by the same action owner according to
Definition~\ref{p2a:def:v229-action-incidence-certificate}.



\section{Finite-band subgroup and exceptional-type certification}
\label{p2a:sec:clean-finite-certificates}
\begin{theorem}[Support-cardinality finite-band repeatability certificate]
\label{p2a:thm:v275-repeatability-PN-suffices}
Let $\mu=\sum_{s\in S}w_s\delta_s$ be a positive probability
on $N$ distinct points of $SU(2)$, all $w_s>0$.
For the left convolution averaging operator $A_\mu$, the following
conditions are equivalent:
\begin{enumerate}[label=(\roman*),leftmargin=2em]
\item $A_\mu^2=A_\mu$ on $\mathcal P_N$;
\item $\widehat\mu(k)^2=\widehat\mu(k)$ for $0\le k\le N$,
where $\widehat\mu(k)=\sum_sw_s\operatorname{Sym}^k(s)$;
\item $\mu$ is uniform Haar probability on the finite subgroup $S$.
\end{enumerate}
In particular the frame-independent, nonnegative Fourier scalar
\[
 \boxed{\mathfrak R_N(\mu)=
 \sum_{k=1}^{N}(k+1)
 \|\widehat\mu(k)^2-\widehat\mu(k)\|_{\rm HS}^2}
\]
vanishes if and only if $\mu$ is such a subgroup average.
The support-dependent tests of
Theorem~\ref{p2a:thm:v274-two-scalar-repeatability-certificate}
are a cheaper alternative when $S$ and its weights are known.
No occurring finite-output physical role is inferred from either test.
\end{theorem}
\begin{proof}
Peter--Weyl identifies the action of $A_\mu$ on each degree-$k$
scalar harmonic with multiplication by $\widehat\mu(k)$ (up to
unitary equivalence, inverse and adjoint conventions). Thus (i)
and (ii) are equivalent. Since $\mathcal P_N$ is closed under
$q\mapsto q^{-1}$, condition (i), evaluated at the identity
against $q\mapsto p_S(q^{-1})$, where
$p_S(q)=\prod_{s\in S}(1-\langle q,s\rangle)$, yields
\[
 0=\int p_S\,d(\mu*\mu)-\int p_S\,d\mu
  =\sum_{s,t}w_sw_t p_S(st).
\]
Every term is nonnegative, so $SS\subset S$. A finite nonempty
multiplicatively closed subset of a group is a subgroup.
Next, on this subgroup, the Lagrange polynomials
$\ell_s(q)=\prod_{t\ne s}(1-\langle q,t\rangle)/(1-\langle s,t\rangle)$
have degree $N-1$ and satisfy $\ell_s(t)=\delta_{s,t}$.
Applying (i) at the identity to their inverse pullbacks gives
$(\mu*\mu)(\{s\})=\mu(\{s\})$ for every $s$.
The positivity-and-maximum argument of
Theorem~\ref{p2a:thm:v273-positive-repeatable-convolution-group}
forces the weights to be uniform. Conversely, a uniform subgroup
average is idempotent on every band. The scalar is a sum of
nonnegative squared Hilbert--Schmidt norms; its vanishing is
precisely (ii), independent of a choice of representation bases.
\end{proof}

\begin{theorem}[Exact frame-independent eighth-moment exceptional signature]
\label{p2a:thm:v275-eighth-moment-3class}
For a finite probability $\nu$ on $S^3\subset\mathbb R^4$, put
\[
 M_8(\nu)=\int q^{\otimes8}\,d\nu(q),\qquad
 \mathcal A_8(\nu)=\|M_8(\nu)-M_8(\sigma)\|^2.
\]
For uniform exceptional subgroup measures the exact values are
\[
 \boxed{\mathcal A_8(2T)=\frac1{32},\qquad
        \mathcal A_8(2O)=\frac1{256},\qquad
        \mathcal A_8(2I)=0.}
\]
Thus, \emph{after} a finite-subgroup certificate and a restriction
to these three admissible types, the single axis-independent
scalar identifies their exceptional type. The same scalar is
unchanged by any quaternionic left translation:
$\mathcal A_8(\mu_{g\Gamma})=\mathcal A_8(\mu_\Gamma)$.
It therefore does not certify subgroup closure or an embedded
coset position. In particular
$\mathcal A_8(\xi_{H,n})=0$ for all $n\ge1$ although every
$\xi_{H,n}$ has non-subgroup support.
\end{theorem}
\begin{proof}
Orthogonality of rotations of $\mathbb R^4$ and Haar invariance
make $\mathcal A_8$ left-translation invariant. For a uniform
subgroup $\Gamma$, bi-invariance and quaternion multiplication give
\[
 \mathcal A_8(\mu_\Gamma)
   =\frac1{|\Gamma|}\sum_{g\in\Gamma}(\operatorname{Re}g)^8
      -\int_{S^3}q_0^8\,d\sigma(q),\qquad
 \int_{S^3}q_0^8\,d\sigma=\frac7{128}.
\]
The distributions of quaternionic real parts, with multiplicities,
are
\[
\begin{array}{c|l}
2T& (\pm1:2),\ (\pm1/2:16),\ (0:6)\\
2O& (\pm1:2),\ (\pm1/2:16),\ (\pm1/\sqrt2:12),\ (0:18)\\
2I& (\pm1:2),\ (\pm1/2:40),\
 (\pm\varphi/2:24),\ (\pm1/(2\varphi):24),\ (0:30)
\end{array}
\]
where each $\pm$ count is the total count of both signs and
$\varphi=(1+\sqrt5)/2$.
Their eighth-moment means are, respectively,
$11/128$, $15/256$, and $7/128$. Subtraction gives the claimed
numbers. Each $\xi_{H,n}$ is an exact $\mathcal P_{11}$ cubature
by Theorem~\ref{p1b:thm:v274-branch-symmetric-twin-convolution},
so its eighth tensor matches Haar exactly despite non-subgroup
support.
\end{proof}

\begin{corollary}[Two diagnostics classify exceptional cosets without confusing design with subgroup]
\label{p2a:cor:v275-coset-type-and-origin-test}
Suppose a positive probability is already known to be a uniform
left coset $g\Gamma$ of an unknown
$\Gamma\in\{2T,2O,2I\}$. Then $\mathcal A_8$ identifies the
type $\Gamma$. Setting $k_\Gamma=6,8,12$ accordingly, the
nonnegative first-leaked-degree scalar
\[
 \mathfrak r_{k_\Gamma}
   =\|\widehat\mu(k_\Gamma)^2
      -\widehat\mu(k_\Gamma)\|_{\rm HS}^2
\]
vanishes precisely when $g\in\Gamma$, that is, when the coset
itself is a subgroup. No coset hypothesis may be inferred from
these two scalar values alone.
\end{corollary}
\begin{proof}
Use Theorem~\ref{p2a:thm:v275-eighth-moment-3class} followed by
Theorem~\ref{p2a:thm:v274-first-leaked-harmonic-coset},
whose rank-one formula makes the second scalar equal to
$|\langle v,\rho_{k_\Gamma}(g)v\rangle-1|^2$.
\end{proof}
\begin{proposition}[Two-point sharpness of the support-cardinality harmonic band]
\label{p2a:prop:v278-two-point-band-sharpness}
For any $g\in SU(2)\setminus\{\pm1\}$, let
\[
 S_g=\{g,-g\},\qquad
 \mu_g=\tfrac12(\delta_g+\delta_{-g}).
\]
Then $\mu_g$ is a strictly positive probability on exactly two
points, and its left averaging operator $A_{\mu_g}$ is idempotent
on $\mathcal P_1=E_0\oplus E_1$, although $S_g$ is not a subgroup.
Thus degree $N-1$ does not suffice universally in the support-$N$
criterion of Theorem~\ref{p2a:thm:v275-repeatability-PN-suffices},
even for $N=2$.

More precisely, the degree-two nonnegative support polynomial
\[
 p_{S_g}(q)=(1-\langle q,g\rangle)(1-\langle q,-g\rangle)
            =1-\langle q,g\rangle^2
\]
has the exact repeatability defect
\[
 \boxed{\int p_{S_g}\,d(\mu_g*\mu_g)
            -\int p_{S_g}\,d\mu_g
       =1-(\operatorname{Re}g)^2>0.}
\]
This is a mathematical lower-band obstruction, not evidence that
any such two-point physical convolution role occurs.
\end{proposition}
\begin{proof}
Quaternionic antipodal points are distinct.  Every function in
$\mathcal P_1$ is a constant plus a real-linear coordinate function
on $S^3$.  For the left action, the two terms at $g^{-1}x$ and
$(-g)^{-1}x=-g^{-1}x$ cancel on every linear coordinate.  Hence
$A_{\mu_g}$ is identity on constants and zero on $E_1$; in particular
$A_{\mu_g}^2=A_{\mu_g}$ on $\mathcal P_1$.  But neither $g$ nor
$-g$ is the identity, so $S_g$ is not a subgroup.

Since $\mu_g*\mu_g=\tfrac12(\delta_{g^2}+\delta_{-g^2})$,
$p_{S_g}$ vanishes on $\operatorname{supp}\mu_g$ and has the value
$1-(\operatorname{Re}g)^2$ at both $g^2$ and $-g^2$:
indeed $\langle g^2,g\rangle=\operatorname{Re}(g^2\bar g)
=\operatorname{Re}g$.
For a unit quaternion the equality $|\operatorname{Re}g|=1$
holds only at $g=\pm1$, proving strict positivity of the defect.
\end{proof}




\section{Constitutional descent of the finite Petersen--Hopf packet}
\label{p2a:sec:v318-petersen-constitution}

\begin{theorem}[Generated five-core incidence under the occurring icosahedral role]
\label{p2a:thm:v318-petersen-same-owner-descent}
Suppose an already constituted, same-owner group-support role
$\mathsf{Supp}_{2I}$ has an occurring embedded representative in the
constituted round $SU(2)$ parent.  Its five conjugate binary-tetrahedral
subgroups $H_i\cong2T$, ten pairwise cyclic intersections
$C_{ij}=H_i\cap H_j\cong C_6$, and the fifteen unordered pairs
$\{C_{ij},C_{kl}\}$ with disjoint index sets are mathematically generated
and retained as one unmarked finite-incidence packet.  The last fifteen
pairs constitute the Petersen graph, with the group action factoring
through $2I/\{\pm1\}\cong A_5$.  As transitive $A_5$-sets, the core,
axis and edge orbits have sizes and stabilizer orders
\[
  (5,12),\qquad(10,6),\qquad(15,4).
\]
The graph and its exact linear-algebra descendants inherit the ancestry
of the occurring $\mathsf{Supp}_{2I}$ role: no second independent
finite-support witness is inserted.  No one $H_i$, cyclic-axis sign,
Petersen vertex, edge-channel basis, physical effect, or particular
preparation is thereby selected as an observed token.
\end{theorem}
\begin{proof}
The five conjugate $2T$ subgroups, their $C_6$ intersections, and the
intrinsic angular Petersen incidence are the finite conclusions of
Paper~I-B.  The stabilizers in $A_5$ are, respectively, $A_4$, the
order-six setwise stabilizer of a two-label subset, and the order-four
stabilizer of a partition into two pairs with one unused label.
All operations defining the graph and its real modules are finite,
relation-preserving descendants of the given subgroup support.  Apply
the existing generated-role co-realization theorem only within this
already-occurring antecedent.  The G0 mark/effect distinction is
unaffected by the mathematical descent.
\end{proof}

\begin{corollary}[Sharp exceptional-support test, with a frozen relation requirement]
\label{p2a:cor:v318-five-core-exceptional-test}
Suppose an already constituted finite-support comparison role has, as
an independently frozen \emph{same-owner} relation requirement, the
presence of five distinct tetrahedral subgroups with the full
pairwise $C_6$ intersection incidence, and its admissible embedded
finite-support candidates have already been reduced to
$[2T],[2O],[2I]$.  Then only $[2I]$ can realize that requirement.
This is a selector \emph{within that prequalified exceptional atlas},
not a universal selection of a discrete subgroup from the connected
$SU(2)$ parent.  It is circular to install the five-core requirement
solely because the $2I$ branch was considered.
\end{corollary}
\begin{proof}
The quotient groups are $A_4,S_4,A_5$.  The first contains one
$A_4$ subgroup (itself), the second contains its unique index-two
normal $A_4$ subgroup, and the third contains five conjugate $A_4$
subgroups with the required pairwise threefold-axis incidence.
Taking binary preimages gives precisely the asserted $2T$ families.
The independent role-premise and finite-candidate hypotheses are
indispensable for the stated GP3 comparison.
\end{proof}

\begin{proposition}[Nonconverse and exceptional-selection boundary]
\label{p2a:prop:v318-petersen-nonconverse}
The additional Petersen packet does not follow from generic point
resolution, from a bare $[2T]$ core-output class, or from the abstract
number fifteen.  It cannot be used to select the $[2I]$ branch before
its occurrence/comparison prerequisites have been established.
A labelled five-core packet can serve as an exact comparison witness
\emph{after} that branch has been generated; the labels remain
presentation marks.  In particular, the previously constituted
$SU(2)$ momentum parent, the $\mathcal H_1$ momentum seed, and the
conditional $2T/2O/2I$ support atlas do not change.
\end{proposition}


\section{Constitutional status of the projective rank-two finite packet}
\label{p2a:sec:v340-ranktwo-selection}
\begin{proposition}[Generated mathematical closure is not an occurring readout]
\label{p2a:prop:v340-ranktwo-nonactualization}
Conditioned on an \emph{already constituted} $2I$ support in the
same geometric parent, the sixty-point cyclic-coset incidence design,
its $722$ great polygons, its central rank-$25$ and rank-$16$
character projectors, the five-core $70\to80$ Reye completion and
its Petersen--dodecahedral second-neighbour operator are generated
mathematical descendants. No independent occurrence postulate is
needed to \emph{construct} them. Conversely, their existence in the
reference atlas does not establish G0 occurrence of that support,
does not make either projector an observed quantum effect, and
cannot supply a new GP3 universal uniqueness or momentum
multiplicity-one assertion.
\end{proposition}
\begin{proof}
Every incidence set and operator is defined by finite
multiplication, subgroup cosets, orthogonal projection and graph
relations on the supplied $2I$ input. Therefore all are invariant
finite descendants under that input. Constitution of the input and
selection of a same-owner effect, coupling, state and observable
remain separate obligations under the established G0/GP1--GP3 rules.
\end{proof}
\begin{remark}[Observable-first falsification criterion]
\label{p2a:rem:v340-ranktwo-observable-first}
A physical proposal using these projectors must state (i) which
source or detector observes a \emph{specific} mathematical block,
(ii) its same-owner embedding and normalization, (iii) a predicted
outcome distribution or pole and a quantitative discriminator against
an alternate support, and (iv) why the allowed preparation and effect
access that block. Equality of dimensions, a nonzero graph spectrum,
or a nearest-neighbour analogy is not a substitute. In particular
$T=S^2-3I$ is even in $S$ and alone cannot distinguish the two
rank-three signed-parity channels.
\end{remark}


\begin{remark}[Finite double-cover parity is not a new physical symmetry]
\label{p2a:rem:v341-desargues-occurrence-firewall}
The additional Desargues $GP(10,3)$ cross-incidence carrier and
coset-derived rank-fifteen projector of
Theorems~\ref{p1:thm:v341-desargues-rp2-parity} and
\ref{p1:thm:v341-coset-gram-hodge} are \emph{conditional finite
math} once a marked $2I$ packet is supplied.  The odd $S_5$
action exchanges the two $\mathbb{RP}^2$ face sets and changes
the signed Hodge operator $S_H$ to $-S_H$; this does not select
any physical parity effect.  An observable claim must supply
G0 occurrence, a GP2-frozen same-owner effect and preparation,
normalization, and a discriminator between the dodecahedral
holonomy and Desargues parity constructions.  A graph cover by
itself does not discharge GP3 or predict a mass or gauge mode.
\end{remark}


\section{Typed complete-octahedral particle-response realization}
\label{p2a:sec:v3222-octahedral-particle-role}
\begin{definition}[Working particle-response carrier, without a new primitive axiom]
\label{p2a:def:v3222-particle-octahedral-role}
For an already constituted local exceptional core $H\cong2T$ in the
round momentum parent, define the complete \emph{mathematical} local
response carrier
\[
 \mathsf C_{\rm oct}(H):=N_{\operatorname{Sp}(1)}(H)=2O(H),
 \qquad \mathsf C_{\rm oct}(H)=H\sqcup(2O(H)\setminus H).
\]
The \emph{working physical particle-carrier interpretation} is that
the finite elementary-particle response uses this complete $2O$,
not just its $2T$ subgroup and not the $2I$ mobile comparison packet.
The physical role has a separate occurrence/constructor condition:
the actual source and effect must be generated on a common carrier,
intertwine the momentum and reciprocal-position representations,
retain center/spin and all charges, and consume the required octahedral
reorientations.  This definition does not assert that the occurrence
condition has been discharged.
\end{definition}
\begin{remark}[No conflict with the previously derived momentum support]
\label{p2a:rem:v3222-octahedral-role-firewall}
The existing $m_p=1$, $\mathcal H_p\cong\mathcal H_1$, finite
Hurwitz evaluation/support and the constitutive G0, GP1--GP3
claims are retained with their exact prior scope.  An ambient
normalizer, a finite sample of a fifteen-dimensional carrier,
and an occupied physical particle source are different types.
Neither $48$ elements nor the equality of an abstract representation
dimension licenses a new particle state or a physical mass.  The
conditional physical source port must be tested downstream without
inserting a new GP or silently replacing the prior $2T$ constructor.
\end{remark}


\begin{remark}[Finite typed registry for the rooted charged-action handoff]
\label{p2a:rem:v3225-charged-origin-registry}
For the occurring rooted triality/Hurwitz parent, the finite registry of
Definition~\ref{p2a:def:v187-parent-saturation} must distinguish the
following already named outputs and their physical type boundaries:
(i) signed $D_4$ triality and the primitive odd source on its declared
rooted branch; (ii) the $C_3$ family-cycle ports and negative-Hopf
signed-incidence relations; (iii) the hierarchy--family transport of an
\emph{already generated} directed arrow; (iv) the conditional weighted
mass/incidence assembly; and (v) the same-action fermion-jet consumer of
an \emph{already occurring} Dirac member.  Items (ii)--(iii) are not
by themselves an $\mathsf{SrcVal}_{\rm odd}$ occurrence certificate.
In particular, the physical charged hierarchy-arrow origin certificate
$\mathsf{Cert}_{\rm hier}$ of
Definition~\ref{p5:def:v228-hierarchy-origin} is not discharged by
parent saturation unless an independent, licensed, result-independent
constructor has an occurring input and the required, owner-resolved,
noncentral charged arrows among its frozen outputs.  The charged
assembly of Theorem~\ref{p5:thm:v397-charged-preterminal-srcval-realization}
retains that hypothesis.  This registry records a missing typed action
edge; it does not add a fourth GP, grant a new constructor, or turn a
geometric relation into a physical source value.
For observables depending only on the left return $RR^\dagger$, the
right-unitary representative $R\mapsto RU$ may be quotiented only after
actual action membership and common marked-readout conditions are
certified.  A physical right-orientation claim separately requires an
independently occurring right-sensitive source/mark.
\end{remark}

\section{Conclusion and constitutional export boundary}
\label{p2:sec:conclusion}

Starting with the strong geometrically primary assertion
\(\mathsf E_R:\exists\mathfrak M\,\operatorname{Geo}_R(\mathfrak M)\), G0
types its actual occurrence as invariantly distinguishable geometric content
without selecting a witness or candidate.  GP1
restricts whole role candidates to finite/bounded structured realization.  GP2
applies no-insertion relative to \(P_{<R}\) and the role geometry while making
each candidate internally complete and exact:
\[
 \operatorname{Pres}_{\mathfrak M}(\mathfrak M)
 =\operatorname{Gen}_R(\mathfrak M).
\]
Each candidate's compact geometric closure is a compact topological group;
when its exact operations are realized affinely and faithfully on the finite
state/test span, Haar averaging gives a finite unitary representation and
the parent is compact Lie.  A carrier-linear lift is a separate certificate.  For any separately constituted comparison
presentation, GP1-TB gives a canonical compact completion candidatewise, and
compact recurrence makes the closure of its forward isometry monoid a group
(Theorem~\ref{p2:thm:gp1-bounded-isometry-closure}), so reversal is derived
rather than inserted.  Currying each candidate comparison metric generates full
mathematical profiles that detect its geometric action kernel; operational
B-test use additionally requires a G0 test-occurrence certificate, of which
metric-curry closure is the full-profile realization.
GP2 thereby defines the candidate-internally exact-and-host-faithful admissible
model class under the constructor-relative generated closure of
Definition~\ref{p2:def:gp2-generated-closure}.  In particular,
Theorem~\ref{p2:thm:gp2-generated-reciprocal-role-closure} shows that once a
later generated reciprocal separating spectrum has a faithful autonomous
variance and no prior joint-support relation is present, GP2 itself
objectifies that reciprocal role and supplies its relation-free product with
the original role; the legacy \((\mathrm{BR}_{\rm ind})\) and
\((\mathrm{IND}_{xp})\) labels therefore record proved antecedents rather
than additional physical postulates.  Registered generated boundaries remain
optional signature items for roles that actually constitute a bounded
word/depth grammar; they are not the default source of physical content.  GP3
then compares whole complete candidates through role-faithful structural
comparisons and asks whether an irreducible least class exists before one
candidate's self-data are frozen.
Candidatewise endpoint comparison together with endpoint irreducibility gives
the unique least \(U(1)\) scalar parent-role skeleton and source-global
\(SU(2)\cong S^3\) noncommutative parent-role skeleton; the extra
\(SO(3)\) quotient relation has already been excluded by GP2.  The
scalar-succession specialization then constitutes the typed
Paper-I-A Hopf-flow provenance occurrence
\(\mathsf{HS}_t\), and
Theorem~\ref{p2:thm:transported-succession-orientation} transports the
positive primitive Reeb lattice to the autonomous scalar parent.  This
removes a downstream sign choice while leaving duration and energy units,
Hamiltonian selection, and actual evolution open.  The
Peter--Weyl/local comparison gives an independent skeleton-level replay, while
the endpoint theorem supplies the universal comparison from every admissible
primitive realization in its declared domain needed for the GP3 conclusion.
If Geometric Primacy supplies richer intrinsic directions, brackets,
relations, or tests that the endpoint does not realize, the comparison is
inadmissible and the theorem abstains rather than deleting them.  Once GP3
selects \(\mathfrak M_R\), its complete exact content is frozen through
\[
 P_R=P_{<R}\cup\operatorname{Gen}_R(\mathfrak M_R).
\]
For every finite rooted chronological presentation the independently defined
generated and co-realized structural presentations are linked by
\[
 \boxed{
 \rho_{\mathbf Q}:
 \operatorname{GeoEff}_{\Sigma,P_0}(\mathbf Q)
 \xrightarrow{\ \cong_{\rm typed}\ }
 \operatorname{PhysStr}_{P,\Sigma}(\mathbf Q).}
\]
Thus a rooted generated structural descendant requires no second premise
merely to be structurally co-realized in its native sort.  The equivalence
does not promote ambient mathematics, collapse an ungenerated isomorphism,
create a cross-role relation, select an unresolved-family member, identify a
geometric operator with the physical Hamiltonian, or actualize a preparation,
history, outcome, nonzero member, empirical value, or calibration token.
The paper still does
not claim a unique unmarked complete structured lift, because distinct
faithful carrier presentations over that skeleton need not be compared by the
primitive role geometry.

The same-owner automorphism theorem adds a further positive closure without
changing the constitutional postulates.  Every licensed intrinsic constructor
word is equivariant under unmarked automorphisms of its complete frozen input
presentation, and a scalar/action descendant with trivial target action is
therefore invariant.  For a constituted continuous automorphism family this
gives \(dS(R_\xi)=0\) wherever the generated first variation exists.  Thus the
native action-invariance antecedent used by Noether is derived from the frozen
same-owner presentation and GP2 exact host faithfulness; only its ordinary
variational consequences remain to be taken.  Transformations that change a
generated source, boundary datum, coupling, orientation, or other retained
mark are outside this theorem.  More generally, theorem outputs are separated
into internal closures, whose physical hypotheses are already generated, and
analytic certifications, whose regularity/existence hypotheses must be proved
independently before any output enters a dependency frontier.

Nontriviality, connectedness, branch type, OR1, three-dimensional Lie
regularity, roundness, and center sensitivity are therefore theorem
consequences of the applied closure.  The numerical radius remains outside
the selection claim.  The primitive metric need not equal the
canonical negative-Killing metric used downstream.  Under the optional (FD)+(LR) hypotheses on any retained faithful realization
over the least skeleton, the two distances are
globally bi-Lipschitz after an unfixed positive rescaling, without becoming
equal.  Paper~II still separates generic finite-level mathematics from physical
selection.  The initial momentum carrier is assigned a direct
constructor-relative host derivation in II-B, with \(m_p=1\) read off only
after the GP2-generated type-support \(\mathcal G^{\rm car}_{\Sigma_p}
\cong\mathcal H_1\) is established.  An assertion about the entire
GP3-irreducible physical carrier additionally requires the
whole-structured comparison certificate of
Theorem~\ref{p2b:thm:momentum-carrier-rigidity}.  On the position side, current
Paper~II-C makes the actual residual occurrence the finite generated frontier
history \((\mathrm{HIST}_x^{\rm fin})\), derives
\(m_x=m_x^{\rm res}(F_x)\) from each occurring frontier, and derives
monotonicity from retained inclusion; a nonconstant history still requires
new generated higher-harmonic content and a physical-time relation.  Raw
ambient coordinate data and absolute normalizations remain
completion-relative, while the normalized protected even algebra on
\(\varrho_{\rm sp}>0\) is completion-independent and is represented by
\(\Cl(0,3)\).  Recovery regimes and stronger local/global comparisons remain
subject to their separately generated or registered conditions.

The exact finite converse certificates of
Theorems~\ref{p2a:thm:v267-heat-bregman-converse},
\ref{p2a:thm:v267-one-frequency-tagged} and
\ref{p2a:thm:v267-rooted-one-heat} permit complete observable
identification under weaker operator-reconstruction requirements.
Their hypotheses remain same-owner action inputs; scalar spectral data
alone do not prove physical rooted alignment.

Theorem~\ref{p2a:thm:generated-exact-solution-moduli} packages the resulting
framework-wide exhaustion rule.  A prior generated typed diagram determines
its complete exact-solution moduli, not a preferred member by default.  A
unique GP3-least class may close the role; otherwise the unresolved family is
retained while every licensed readout constant on that family descends as
common generated content.  Independent role problems remain a relation-free
product on the factored admissibility domain. An independently generated
common-value equation adds its fibre-product incidence while all other joint
conditions remain.  The theorem changes neither the G0 occurrence
boundary nor the rule that a particular history, state, process member, or
calibration token requires its own declared output or occurrence witness.

For finite point resolution, the same exactness discipline proves a sharp
domain boundary.  Generic \(\varepsilon\)-net support does not imply group
closure, the Hurwitz Gram pattern, optimal 24-point packing, or an exceptional
binary-group role.  Paper~I-B exceptional geometry derives the unmarked \(2T\)
core only after either an occurring 24-point optimal kissing support or one
common group-closed least-exceptional role has independently been established.
Paper~II-F, Section~\ref{p2a:app:generic-net-exceptional-obstruction} proves a sharper
representative-selector obstruction: an unmarked oriented round \(S^3\)
admits no nonempty finite isometry-natural embedded subset, and an unmarked
\(SU(2)\) admits no 24-point automorphism-natural embedded subset.  This
obstruction does not exclude a natural orbit or family of 24-cell supports,
or a symmetry-invariant functional or process whose unique minimizer orbit is
\([2T]\).  Licensing and generating such a role or mechanism is a separate
antecedent.  Thus the remaining exceptional branch is entered by an actual
exceptional-support occurrence or by one of the exact independently generated
orbit-level antecedents of
Corollary~\ref{p2a:cor:two-conditional-2t-discharge-routes}, not by an
embedded representative selected from homogeneous data alone.





\printbibliography

\end{document}
